\chapter{内域}

翻越最后一道险峻的山脊，眼前豁然开朗。

与荒域那边苍茫、粗粝、灵气近乎枯竭的景象截然不同，展现在叶牧谦三人面前的，是一片生机盎然、云雾缭绕的广袤天地。远处峰峦叠翠，流泉飞瀑隐约可闻；近处古木参天，奇花异草点缀林间。

最显著的不同，是充盈在天地之间的、可以被清晰感知到的能量流。

“灵气”！

叶牧谦心想，估计这就是内域能够出现修士社会，但荒域不能的原因吧？

灵气无所不在，呼吸间都能感到一丝微凉沁入肺腑。

叶牧谦深深吸了一口气，闭目细细体悟，却微微皱起了眉头。

“师父，这里的‘气’……好奇特。”白言冰也察觉到了不同，他尝试引动一丝外界灵气，却觉得那灵气虽能感应，却与自身气海中温养出的元气格格不入，难以直接纳入转化。

陈千羽同样面露疑惑。她作为医师，对气的感知更为细腻，只觉得这外界灵气虽然活跃，但似乎不够纯净，隐含着一丝躁动。

“嗯。”叶牧谦缓缓睁开眼，目光扫过这看似仙气氤氲的山林，语气带着罕见的凝重，“此界灵气，确与荒域不同。

“但是，于我感知之中，这灵气总有一种轻微的破残之感，仿佛被打碎后又勉强拼凑起来，失了浑然一体的圆融道韵。”

他看向两位徒弟，郑重告诫：“记住，我们之道，在于内求，在于温养自身元气，淬炼体魄神魂。外界灵气可作参考，可借其压力磨砺己身，也可淬炼之后归为己用；但绝不可贪图便利，放任未经淬炼的灵气替代甚至污染你们辛苦温养出的元气本源！那无异于自毁根基，后患无穷。”

白言冰与陈千羽心中一凛，齐声应道：“弟子谨记！”

三人略作调息，适应了一下新环境，便沿着山间隐约的小径下行。

\vspace{2em}

约莫走了大半日，终于在山脚处看到了一座小镇的轮廓。镇子不大，以青石垒砌的房屋为主，街道上人来人往，颇为热闹。

让三人略感新奇的是，街上行走的人中，其中有约一成人，周身散发着或强或弱的能量波动——正是修士。剩下的，自然是毫无修为波动的凡人。

修士与凡人混杂而居，但彼此间的地位差异一目了然。凡人对实力明显更加强大的一批修士大多恭敬避让，而这些修士则神态自若；至于低阶修士，则与凡人受到的待遇几乎无异。

显然在此地，有一部分修士阶层的地位远高于凡人和其余修士。

“果然是个修士主导的世界。”白言冰低声道。

叶牧谦微微颔首，并未多言。

\vspace{2em}

舟车劳顿了许久，三人此刻都有些疲累，于是在镇上寻了一处看起来还算干净的客栈住下。

不知道为什么，陈千羽今天特别粘白言冰。于是，白言冰与陈千羽同住一间房，叶牧谦则住在隔壁。

是夜，月朗星稀，小镇渐入梦乡。两人亦解衣欲睡，忽然——

客房窗户猛然炸裂，一道矫健的黑色身影如夜枭随着纷飞的碎片撞了进来。人未至，一股冰冷的杀意已锁定了床边，凌厉的刀光快得只剩下一线残影，直取正背对窗户的白言冰后心！

这一击毫无征兆，狠辣决绝，完全是顶尖刺客的作风。

致命的危机感让两人汗毛倒竖！

“何方宵小！”白言冰惊怒交加，厉喝出声。

历经生死磨砺的战斗本能已深入骨髓，于那间不容发之际，他根本来不及回头，全凭对杀气来向的感知，腰肢猛地一拧，硬生生带着陈千羽向侧方横移半尺！

“嗤啦——”冰冷的刀锋几乎贴着他的肋侧划过，将他的内衫割开一道长长的口子，肌肤感受到刺骨的寒意。

与此同时，白言冰的反击已到。他侧身的同时，右掌已蓄满精纯元气，看也不看便向身后刀光来处狠狠拍去！这一掌含怒而发，毫无花哨，却势大力沉，掌风呼啸，裹挟着澎湃的元气直撞而上。

“砰！”

掌风与刀光结结实实撞在一起，发出一声沉闷的气爆！狂暴的劲气以碰撞点为中心炸开，瞬间将旁边一张小几震得四分五裂，上面的茶壶茶杯“噼里啪啦”碎了一地，碎片混合着茶水四溅。

那黑衣刺客显然没料到目标在如此突兀的袭击下，不仅避开了要害，反击还如此刚猛凌厉。刀光被掌力震得微微一偏，她冲入室内的身形也为之一滞。

就在这电光石火的间隙，陈千羽也已反应过来。羞怒瞬间化为冰冷的杀意。

“找死！”她清叱一声，素手在床边一抹一扬，数点寒星已疾射而出，直取黑衣人影的面门与咽喉。与此同时，她另一手闪电般扯过散落的外袍裹住身体，纤足在床沿一点，身形如轻羽般飘起，已与白言冰形成一前一后的夹击之势。

那黑衣人影一击不中，立刻知道遇到了硬茬子。她反应亦是极快，毫不犹豫放弃攻击，手中那柄不过尺余长的短刀瞬间舞动起来。刀光流转，化作一片绵密的光影护在身前，试图同时挡住来自前方与侧面的攻击。

“叮叮叮叮！”一连串细密急促的金铁交击声响起，陈千羽射出的钢针大半被刀光磕飞。几枚漏网之针擦着黑衣人的耳际和肩头飞过，“哆哆哆”地钉入身后的墙壁和衣柜门板上，针尾剧颤。

白言冰哪里会给她喘息之机？一击掌风将她阻住后，他已合身扑上，元气澎湃，怒意勃发。

虽手中无剑，但心中有剑！

此地并不开阔，无法施展出升皎弦的全部套路，但依然毫无保留、拳掌带风，招招势大力沉、直来直往。一拳轰出，黑衣人横刀格挡，却被那沉雄的力道震得手臂发麻，脚下“蹬蹬蹬”连退三步，后背“咚”一声撞在了房间中央的圆桌上，桌上的烛台倾倒，蜡烛滚落，火苗瞬间引燃了桌布。

陈千羽已趁势游走至黑衣人侧翼，指间时不时有寒芒闪烁，专攻黑衣人持刀的手腕、手肘、膝弯等要害穴位，打法极为刁钻；更有一股气劲随着她的掌指吞吐悄然弥漫，微妙地干扰着黑衣人体内气息的运转节奏，让她总觉气息不畅，招式衔接间出现细微的滞涩。

黑衣人顿时陷入左支右绌的境地。白言冰的猛攻如巨浪拍岸，让她必须全力应对；陈千羽的袭扰则如附骨之疽，阴险难防。不过数息之间，她已是险象环生。

“嗤啦——！”

一声裂帛之音格外清晰。白言冰一记凌厉的掌风扫过，虽被黑衣人惊险避过要害，却将其肩头的夜行衣撕裂开一道大口子，露出下面一片白皙的肌肤。黑衣人闷哼一声，动作不可避免地出现了一丝慌乱。

陈千羽觑准这稍纵即逝的破绽，玉指一弹，一枚灌注了精纯元气的银色长针脱手而出，直取对方兵刃！这一针时机妙到巅毫，正是黑衣人因衣物撕裂而心神微分、刀势微缓的刹那。

“铛！”

一声清脆的撞击声，银针精准地击中短刀刀身侧脊。黑衣人只觉得刀身上传来一股尖锐的螺旋劲力，手腕顿时一酸，短刀再也拿捏不住，脱手斜飞出去，“夺”地一声，深深扎进了房间的承重木柱之中，刀柄兀自颤动不已。

兵刃脱手，黑衣人眼神一厉，竟不退反进，合身向看似招式更为刚猛、不易变向的白言冰撞去，同时并指如刀，直戳白言冰腰眼，竟是以伤换伤的打法。

从刚才的闷哼与此刻的身形动作，白言冰心中已确认了对方是个女子，不由对她能如此刚烈而暗惊，但手下却毫不留情。他战斗经验丰富，岂会看不出对方意图？

当下不闪不避，左手闪电般探出，五指如钩，精准无比地扣向对方戳来的手腕。同时，他并指如剑，右手双指凝聚元气，疾点向对方因出招而暴露的腕脉要穴。

黑衣女子戳出的手刀顿时被扣住，腕脉处更是一麻，半边手臂的力气瞬间消散。若非修士，估计此时手筋已经彻底断了。

陈千羽已如鬼魅般闪身至黑衣女子身后，带着元气的一掌重重拍向其背心。

前后受制，黑衣女子闪避不及，背心要害被一掌印实。她只觉周身气血猛地一滞，仿佛正在奔流的溪水突然结冰，身形顿时踉跄，再也站立不稳。

白言冰抓住机会，扣住其手腕的手用力一拧，同时进步贴身，另一只手已锁向其脖颈要害，瞬间便将这黑衣女子彻底制服，反剪双手按压在地。

激烈的打斗骤然停止，房间内一片死寂，只剩下粗重的喘息声，以及桌布燃烧的“哔啵”轻响。那女子黑巾蒙面，虽被制服，却仍努力挺直脊梁；蒙面巾上方露出的那双眼睛里，充满了不甘、惊疑，以及一丝倔强。

“怎么说，师兄？”陈千羽终于腾出手来，一边为自己系好衣带，一边明显愠怒地问道。

“先揍一顿，然后禀告师尊处置吧。”白言冰死死压制着地上仍在挣扎的黑衣女子，咬牙切齿道。

两人正欲动手，就在这时，房门被“吱呀”一声推开。叶牧谦披着外衣，睡眼惺忪地站在门口，脸上带着明显的起床气。

“吵吵什么，你们喜欢半夜打架？”

叶牧谦目光扫过屋内的一片狼藉，又看了看破碎的窗户，以及被制服在地、黑巾蒙面却仍努力挺直脊梁、眼里满是不服的黑衣女子，最后落在衣衫不整的白言冰和满面怒容的陈千羽身上……

叶牧谦沉默了两秒，揉了揉眉心。那股被吵醒的烦躁感褪去，取而代之的是一种了然于胸的无奈。

得，这剧本他熟。

“说说吧，”叶牧谦走进房内，拖过一张还算完好的椅子坐下，随手向正在燃烧的桌布推出一丝元气把火灭掉，看着被白言冰反剪双手压制的黑衣女子，“深更半夜，破窗而入，刺杀我徒儿。阁下是何方神圣，所为何来？若说不出个所以然……”他眼神微冷，“我这人起床气有点大。”

黑衣女子挣扎了一下，发现压制自己的力量沉凝如山，根本无法撼动。她咬了咬牙，抬头看向叶牧谦，又瞥了一眼怒目而视的白言冰和陈千羽，似乎也察觉到了什么，眼神中的敌意和决绝略微消退，取而代之的是一丝惊疑。

“你们……不是邪修？”她声音有些沙哑，但能听出是女声。

“邪修？”白言冰气笑了，“你哪只眼睛看我们是邪修？”

叶牧谦摆摆手，示意白言冰稍安勿躁，对黑衣女子道：“是与不是，你总得给我们个判断的依据。先说说你是谁，为何认定我们是邪修？”

黑衣女子犹豫片刻，或许是觉得形势比人强，或许是叶牧谦平静的态度让她觉得可以沟通，终于开口：“我乃一介散修。月前，我追杀一邪修‘血魔’至十万群山边缘，将其重伤，他却遁入群山。我守候多日，未见其出山，却察觉你们一行三人自群山中出来，身上并无灵枢玄脉波动，却又能安然穿越那瘴气密布、妖兽横行的绝地。”

她语气带着深深的困惑与警惕：“内域共识，荒域灵气稀薄，不可能诞生修士，凡人更绝无可能穿越十万群山。你们既非普通凡人，又无我辈修士之特征，却能安然走出……我自然怀疑，你们或许是修炼了某种诡异法门、掩盖了气息的邪修，那血魔恐怕是遭了你们黑吃黑！”

“所以你就一路跟踪我们到此？”陈千羽脸颊微红，又是羞恼又是好笑，“然后半夜听到些动静，以为我们在乱搞？”

女子此刻冷静下来细看，目光在三人脸上仔细逡巡，尤其是观察他们的眼神和眉宇间的气息。她确实没有从这三人身上感受到邪修那种特有的阴冷、暴戾、血腥的黑气；相反，白言冰气息刚正，陈千羽气息温润、生机勃勃，叶牧谦更是深不可测，如古井深潭，却并无邪意。

随即，脸上露出些许尴尬和歉意：“如今看来……似乎是我鲁莽，误会诸位了。我以为……呃，抱歉。”

她倒是坦荡，错了便认，虽然蒙着面，但眼神中的歉意不似作伪。

白言冰和陈千羽对视一眼，怒气消了一半。毕竟对方初衷是除魔卫道，只是闹了个大乌龙。白言冰哼了一声，松开了对那女子的钳制。

叶牧谦倒是若有所思：“姑娘追杀血魔？那邪修确实已被我等在荒域边境诛杀。”他示意白言冰取出那枚令牌。

看到令牌，黑衣女子眼睛一亮，接过仔细查验，终于彻底松了口气，也完全相信了叶牧谦的话。“果然是此獠的命牌！多谢三位为民除害！”她郑重抱拳行礼，这次是真心实意的感谢。

误会解除，房间内剑拔弩张的气氛缓和了许多。

叶牧谦谦和地请客栈伙计换了间完好的客房。当然，这间屋子的费用以及先前损坏房间的赔偿，自然是由那黑衣女子全部承担——她虽为散修，行事却颇有章法，不愿欠人人情。伙计战战兢兢地收拾残局，看向黑衣女子的眼神仍带着畏惧，但见叶牧谦师徒气度从容，倒也安心了几分，很快便安排好了隔壁一间清净上房。

四人移步新客房，围坐桌旁。桌上已换了一壶新茶，热气袅袅升起，在烛光中晕开暖色的光晕。

“姑娘，我等初来乍到，对此地风土人情、修炼之道一概不知。今夜冒昧打扰，实属无奈，还望姑娘解惑？”叶牧谦执壶为对方斟了杯茶，语气客气而真诚。

那黑衣女子此刻已知对方并非恶人，且实力深不可测——尤其是这位青衫儒雅的叶先生，看似平和，却给她一种渊渟岳峙之感。她行走内域多年，见识过不少高手，却从未见过如此奇特的气息与道韵。既然误会澄清，结个善缘未尝不可。想到此处，她便放下戒备，将自己所知娓娓道来。

原来，内域通行的修炼法门，被称为“灵枢径”。其核心在于开启脊柱深处的“灵枢”，以此为核心枢纽，贯通四肢“玄脉”，形成体内灵路网络。修者借此网络沟通天地，引外界灵气入体，炼化为己用，讲究的是“借万物以御天”，以身为舟，以灵为桨，渡世而行。

她看向叶牧谦：“我听先生之前所言，你们所修之道，似乎是反其道而行之？不假外求，内炼己身？”

叶牧谦微笑颔首：“正是。我辈道途，正是‘求一身以成道’。”

黑衣女子眼中掠过震惊之色，虽然早有猜测，但亲耳听到如此迥异的理念，仍觉不可思议。这与灵枢径可谓南辕北辙。

也正因存在修士这种超凡个体，内域的社会格局也与荒域截然不同。

内域并无荒域那般统一的凡人皇朝，而是以一座座大小不等的城池为基础单元，星罗棋布地分布在这片广袤土地上。其中最大的城池有七座，借北斗七星之名，分别为天枢、天璇、天玑、天权、玉衡、开阳、摇光；而中小型城池则取名随意，多以地理特征、特产、宗门或创建者之名命名。

每座城池都由一个或几个修士宗门联合掌控、庇护。宗门负责城池的安危，组织力量抵御妖兽、流窜的邪修团伙以及其他势力的侵扰，同时也对城池拥有绝对的行政、税收、律法制定等权力。凡人和低阶散修则大多依附于城池生存，为宗门提供各种服务和资源——从开采矿石、种植灵谷、驯养低阶妖兽，到经营店铺、从事手工业、充当杂役等，形成一种独特的、以修炼阶层为核心的共生关系。

“我们此刻所在的这座小镇，便隶属于一个名为‘青石城’的中等城池。”黑衣女子道，“而掌控青石城及周边千里地域的宗门，叫做‘青云门’。在青石城辖下，类似的小镇、村落有数十个，皆受青云门庇护，也皆需向青云门缴纳赋税、提供劳力。”

白言冰若有所思：“如此说来，宗门便是城池的‘官府’？”

“可以这么理解，但更绝对。”黑衣女子点头，“宗门弟子便是‘官差’，宗门长老便是‘官员’，而宗主便是此地的‘太守’或‘州牧’。不过，相比凡人国度的君王，宗门对辖下生灵的控制更为直接，也更深介入修炼资源的分配。”

陈千羽轻声问道：“那凡人若想修炼……”

“这便是宗门存在的另一重意义。”黑衣女子眼中闪过一丝复杂神色，“宗门垄断了相对系统、安全的修炼法门与资源。凡人若想踏入修行路，要么自行摸索，成为散修；要么等待宗门的统一选拔。比如，青云门近期正要举办三年一度的‘升仙大会’，旨在从管辖范围内的凡人乃至散修中，选拔有资质的苗子，补充新鲜血液。这是许多凡人改变命运最重要的途径之一。”

接着，她说出了一个令叶牧谦师徒颇感意外的现象，这也是他前世记忆中网文几乎从不涉及的地方。

在内域，不少人修仙，初衷并不是为了那虚无缥缈的“长生”，甚至往往也不是为了行侠仗义或登临极巅，而仅仅是为了获得更好的生产能力与生存保障。低阶修士，尤其是数量最多的灵枢境修士，从事矿业、农业、养殖业的比比皆是。开启灵枢、引气入体后，他们的体力、耐力、对恶劣环境的适应力以及对某些特定灵气敏感度都会提升，一人之生产力，往往能抵三五个凡人壮劳力。只有突破到玄脉境，体内灵脉初步贯通，灵力运用更加自如，大多才会彻底脱离基础的体力劳动，担任管理、护卫等职责。

“故而，内域总人口数千万，其中有修为者——先不论修为高低，哪怕只是刚刚感应到灵气、踏入灵枢境感应期的半吊子——竟能达到数百万之众。”黑衣女子语气平静，却抛出一个惊人的比例，“但玄脉境及以上的数量则断崖式下跌，估计整个内域，七大城连同所有中小城池、荒野散修加起来，玄脉境及以上者，也不过数万名而已。至于真符境，已是各方势力的中流砥柱，法相境……更是凤毛麟角，近百年只出过三位法相。”

“玄脉境？那个血魔似乎也是玄脉境。”白言冰插话道，回想起之前在荒域那场激烈的战斗。

“对，灵枢径的修习过程非常成体系，即使是邪修也脱不开这个框架。”黑衣女子肯定道，“我与那血魔都处于玄脉境通络期，只不过他走的是掠夺生灵精血、速成邪功的歧路，根基虚浮，心性扭曲，真正战力比同阶稳扎稳打的修士要弱上一筹，但手段更为歹毒难防。”

她顺势详细介绍了灵枢径的境界划分。

灵枢径分为四大境界：引灵枢、通玄脉、绘真符、现法相。每一大境界又细分为三个小境界：
\begin{description}
    \item[灵枢境] 感应（感应天地灵气）、引气（引导灵气入体）、凝枢（于脊柱凝炼稳定灵枢）；
    \item[玄脉境] 开脉（以灵枢为源，开辟主要玄脉通道）、通络（贯通细微灵络，形成网络）、融贯（灵脉网络浑然一体，灵力运转圆融）；
    \item[真符境] 摹形（以灵气摹刻成实物之虚影，即“真符”）、注灵（为真符注入灵性与理解）、显化（真符可以短暂脱离修士存在，甚至承载法术）；
    \item[法相境] 凝形（以自身道途与领悟凝聚法相雏形）、外显（法相可离体显化于外，影响现实）、通灵（法相生灵，近乎身外化身，玄妙莫测）。
\end{description}

而那三位法相自然也是鼎鼎大名的：

\begin{itemize}
    \item 内域超级势力天师府之大天师，沈天穹（外显期），也是唯一现世活动的法相境大能；
    \item 另一个超级势力紫阳宗之宗主，“凌天真人”紫博端（外显期），百年前征讨血罗刹受重伤，目前正闭关中；
    \item 被剿灭的血魂教前教主，血罗刹（通灵期），在百年前受天师府与紫阳宗联合讨伐，目前下落不明、生死不明，估计是已经陨落了。
\end{itemize}

叶牧谦静静听着，心中飞快对比。按照黑衣女子的描述，以及之前与这女子和血魔的交手经验来看，他们师徒三人目前的战斗实力，应该比这“灵枢径”体系整体高上一个大境界左右。白言冰与陈千羽虽初入养气境不久，但筑基扎实、元气精纯，对身体的掌控和力量的运用远非灵枢境修士可比，至少也相当于玄脉境。

那黑衣女子介绍完这些，终于忍不住对自己好奇已久的问题发问，尤其是对三人修炼的“道途”大感惊奇。

“先生，恕我直言，您所说的‘内求己身’、‘不借外力’……实在令人费解。修仙一途，本就是逆天而行，窃阴阳，夺造化，不正是‘与天斗、与地斗、与人斗’的过程吗？封闭己身，隔绝天地，如何能得大法力、大神通？又如何能应对修行路上的重重劫难、种种强敌？”她的观念显然深受灵枢径影响，对此颇为不解，甚至觉得有些……违背常理。

叶牧谦闻言，并不恼怒，反而微微一笑：“姑娘所言不差，但或许可以稍作修改。我辈修行，首重修心，次而炼气。或许可说是‘与天奋斗，与地奋斗，与人奋斗。’”

仅仅一字之差，将“斗”换为“奋斗”，意境却截然不同：少了几分你死我活、对抗掠夺的暴戾之气，多了几分自强不息、砥砺前行的昂扬之意。

黑衣女子闻言，娇躯微微一震，整个人怔在当场。她细细品味着这看似简单的几个字，只觉得其中蕴含的意味，与她过往二十年所认知、所践行的修仙理念截然不同，却别有一番振聋发聩、直指本心的力量。

良久，她才深吸一口气，眼神复杂地看向叶牧谦，郑重问道：“不知先生及高徒所修之道，其名为何？”

“问道途。”叶牧谦坦然相告。

“问道途……问道途……”黑衣女子低声重复了两遍，似要将这个名字刻入心里。这是他从没听过的道途。

谈话间，陈千羽医者本能发作，注意到那女子虽然黑巾覆面，但露出的左边脸颊光洁如玉，肤色白皙，眉眼清秀。然而，右边被黑巾遮盖的耳际与脖颈处，却隐约能看到些许不自然的凹凸纹理与淡淡的、仿佛火焰灼烧后留下的浅紫色痕迹。那痕迹虽被尽力遮掩，但在烛光侧照和医者敏锐的观察下，仍无所遁形。

她出于关切，便开口：“姑娘，你这边脸似乎有旧伤？我略通医道，或许可以……”

“别碰我！”黑衣女子却如同受惊的兔子般猛地后仰，发出一声短促的惊呼，迅速抬手拉紧了面纱，眼中闪过一丝难以掩饰的慌乱与深切的痛楚。

她稳了稳剧烈波动的心神，声音低沉下来，带着一丝沙哑：“……不必了。多谢好意。这是早年追杀一个亡命邪修时，不慎被其以伤换伤，用阴毒火符所伤。不仅皮肉，连深层肌理、部分经络都已损毁，寻常丹药难治……这么多年，也习惯了。就这样吧。”语气中带着一丝难以掩饰的黯然，但更多的是一种倔强的认命与自我保护般的疏离。

陈千羽见状，知她心有芥蒂，创伤或许不止于表面，便不再强求，只是温和地点点头：“姑娘勿怪，是在下唐突了。”

经过这番小插曲，气氛短暂凝滞后再度流动。四人交谈甚欢，话题从修炼之道的根本差异，扩展到内域各地的风物人情、奇闻异事，再到各方势力的明争暗斗、资源分布。黑衣女子虽是散修，无门无派，独自闯荡，但见识颇广，为人也磊落坦荡（除了脸伤之事讳莫如深），言语间颇有见地。她很快便与性格直爽的白言冰、心思细腻的陈千羽消除了隔阂，言谈间多了几分真诚。叶牧谦则从她的叙述中，获取了大量急需的信息，对内域的社会结构、力量层次、资源分布有了轮廓性的认识，这比他自己漫无目的地摸索要高效得多。

不知不觉，窗外的深邃夜色渐渐褪去，转为蒙蒙的灰蓝，东方天际泛起一抹淡淡的鱼肚白，几颗晨星黯淡下去。屋内的烛火燃尽最后一滴蜡油，轻轻摇曳了几下，熄灭了，清冷的晨光从窗户缝隙渗入。

黑衣女子看了看天色，起身抱拳，姿态干净利落：“与三位一席谈，获益良多，尤其是叶先生之言，发人深省，让我看到了另一番天地。误会既已澄清，小女子不便再多叨扰，就此告辞。”

“姑娘要去往何处？”陈千羽起身相送，轻声问道。

“自然是继续云游，斩妖除魔。”黑衣女子笑了笑，眼中光芒坚定如初，那是属于她的道与执着。随即她又补充道，“对了，青云门升仙大会在即，青石城近日会颇为热闹，各方人士汇聚，三教九流都有。三位若有兴趣，不妨前去观礼，也算深入了解内域宗门运作之一斑，或许对你们在此界立足有所助益。”

叶牧谦忽然反应过来一件事，温声问道：“聊了这许久，还未请教姑娘芳名，将来若再有缘相见，也好有个照应。”

姑娘本已转身走向窗口，准备如之前般纵身离去的身形微微一顿，似乎踌躇了片刻。夜行之人，尤其是她这样的散修，通常不愿轻易透露真名。但身后这三人的气度与那番“奋斗”之言，让她产生了一种奇特的信任感。于是她回过头，黑巾之上的明眸看了叶牧谦一眼，回答道：“沈瑶。”

说罢，沈瑶对三人点了点头，不再多言，身形一晃，便已如一片轻羽般从窗口掠出，衣袂在晨风中轻扬。她在屋脊上几个轻盈的起落，身影便融入了渐亮的天色与尚未散尽的薄薄晨雾之中，仿佛从未出现过一般，只留下空气中一丝极淡的、清冷如梅的气息。

房中恢复安静。叶牧谦走到窗边，望着沈瑶消失的方向，目光悠远，远处山峦轮廓在晨曦中逐渐清晰。

所以说，他先前之所以感受到“破残”感，是因为那截然不同的“灵枢径”修炼方式对天地灵气的索取方式吗？

还是另有原因？

“青云门……升仙大会……”他低声自语，手指无意识地轻叩窗棂，“或许，是个不错的切入点。”

\vspace{2em}

与沈瑶一别后，叶牧谦师徒三人在小镇又休整了一日，采买了些物资，将连日来奔波的风尘与初入内域所带来的新奇、警惕之感稍稍沉淀。第三天清晨，根据沈瑶指引的方向，他们才离开小镇，朝着青石城所在行去。

越靠近城池，脚下的道路越发宽阔平整，显然经常有大队人马或载重车辆通行，路面甚至隐隐能看到加固符文的浅淡痕迹。行人车马也愈发稠密，形成了一股朝同一方向流动的人潮。

其中，修士的比例明显提高，气息强弱不一。但所有人的目标似乎都指向同一个方向——那座随着他们前行，逐渐从地平线上露出巍峨轮廓的巨城，青石城。

三日后，他们终于抵达青石城外。

近距离观看，这座城池更显雄伟。城墙高逾十丈，全数以本地特产的青灰色巨石垒砌而成，石块巨大，接缝处浇铸了某种金属汁液，厚重坚固无比。城墙表面隐约可见一道道复杂玄奥的灵光符文在石材深处缓缓流转，时隐时现，显然布有强大的防护阵法，寻常攻击难以撼动。城门宽阔，可容四辆马车并行，此时正值晌午，进城队伍排成了长龙。

城门口有十余名身穿统一青色劲装、胸口以银线绣着青云门纹章的修士值守气息精悍，眼神锐利，对进城之人进行略作盘查，主要是查看有无通缉画像上的面孔、询问大致来意、有时还会检查携带的大型货物。

轮到叶牧谦三人时，值守弟子见他们面孔陌生，并非熟识的周边常客，便例行公事地询问：“三位面生，从何处来？入城何事？”

叶牧谦从容答道：“自远方游历至此，听闻青云门升仙大会盛事，特来观礼，增长见闻。”

值守弟子打量三人，见为首的青衫男子气度沉静儒雅，身后少年英气内敛，少女清丽温婉，虽感知不到强烈灵力波动，但那份从容不迫的气度绝非寻常旅人能有。弟子心中掂量，觉得不像惹事之辈，便不再多问，侧身放行：“入城后请遵守青石城规，莫要生事。祝观礼愉快。”

进入城内，景象又与城外小镇截然不同，仿佛一步跨入了另一个喧闹而有序的世界。

街道纵横交错，两旁店铺林立，旌旗招展。售卖之物琳琅满目，从寻常的衣食住行用品，到低阶的法器、各种功能的符箓、品级不一的丹药，乃至各种处理过的妖兽材料，应有尽有。吆喝声、讨价还价声、法器试验时的轻微嗡鸣声、修士间的低声交谈声、驮兽的响鼻声、车轮碾过石板的轱辘声……种种声音混杂在一起，充满了生机勃勃的市井气息。

“升仙大会”显然是近期青石城最大的盛事，没有之一。街头巷尾，茶楼酒肆，到处都能听到相关的谈论。更有显眼处的布告墙上，张贴着青云门发布的正式告示，写明大会将于三日后在城内青云门下属的“演法广场”举行，欢迎各界道友前往观礼；同时，年龄在十二至二十五岁之间，或至少拥有灵枢境感应期修为者，亦可报名参与资质测试，一旦通过，便有希望拜入青云门。

叶牧谦三人穿过热闹的街市，寻了一处位于相对僻静巷弄里的客栈住下。

三天时间转瞬即逝。

第四日清晨，演法广场已是人山人海。

广场位于青石城东南区，占地极广，足以容纳数万人。中央搭建起数座高台，以坚实木料和部分石材构筑，其中一座最为高大雄伟，装饰也最讲究，铺着红毯，摆着檀木桌椅，显然是青云门长老与邀请来的贵宾的监督与观礼之处。其余几座高台则相对简朴，是用于资质测试的场所。台下，黑压压一片，挤满了前来观礼的人群，空气中弥漫着一种躁动不安又充满期待的气氛。

辰时正，钟鸣九响，清越悠扬的钟声传遍全城，广场渐渐安静下来。

一位主持大会的白发长老登台，须发皆白，面色红润，气息渊深，显然是玄脉境融贯期的修士。他清了清嗓子，声音灌注了精纯灵力，清晰地传遍广场每个角落：

“诸位道友，诸位乡亲父老！今日，乃我青云门三年一度升仙大会之期！仙道贵生，广纳贤才，此乃我青云门立派之本！望今日有缘者，皆能展露资质，踏入仙门，共参大道！”

一番冠冕堂皇却也是宗门必备的门面话后，他旋即宣布：“资质测试，现在开始！请已报名者，依序上台！”

测试过程，在叶牧谦看来，确实与他前世从各种玄幻小说中了解到的套路大同小异，但亲眼所见，细节更为真实生动。

测试主要围绕几个方面：灵根属性与纯度、根骨资质、经脉潜力。测试道具也多种多样：有专门测试灵根属性的感应石，通体漆黑，有半人高，内部有纹路流转；有检测根骨资质的照骨镜，边缘雕刻着云纹；还有考验经脉韧性、宽度的通脉尺，长约两尺、晶莹剔透，表面有精细的刻度。

参与测试者按照事先领取的号牌，依次上台。

每当有少年让感应石亮起夺目而纯粹的光芒时，台下便会爆发出一片羡慕的惊叹声，负责记录的外门弟子会快步上前，态度热情地询问姓名、年龄、来历，青云门长老所在的主台上，也会投来关注的目光。而当一个少年手掌按上去，石头只泛起微弱的、混杂的灰光时，台下便是一片惋惜的叹息，少年本人则脸色惨白，黯然下台，有的甚至承受不住梦想破灭的压力，当场掩面哭泣，被家人搀扶下去。

叶牧谦三人站在观礼人群靠前的位置，静静观看。白言冰看得目不转睛，对比着这些测试方式与师尊平时对他们气感、体悟的教导。陈千羽则更关注那些测试失败者的情绪反应和身体状态，医者仁心悄然萌动。

叶牧谦自己则完全像个单纯看热闹的外行，看得津津有味，心中感慨：这可比前世那些特效廉价、情节狗血的仙侠影视剧真实太多了！

每一个测试者的紧张、期待、狂喜或绝望，都是如此鲜活；那些负责测试、维持秩序的青云门弟子，脸上虽然大多是程式化的严肃或鼓励表情，但眼底偶尔会闪过一丝对不同资质者的不同态度——对天才苗子的热切，对普通资质者的平淡，对差生的些许不耐；台下人群中，那些散修或小家族的长者，看着自家子弟上台测试时，紧握的拳头、屏住的呼吸和眼中几乎要溢出的期盼，那是一个家庭、一个小团体对未来最深沉的投资与渴望。

大会从清晨持续到日暮，经历了数千人的测试，最终选拔出了约摸二三十名资质上等的新弟子，也选拔出了近百名资质次一等的弟子。这个比例，让叶牧谦再次直观感受到“修炼资质”的稀缺。那些被选中的少年少女及其家人自是欢天喜地，有的甚至激动得跪地叩谢。青云门几位在场的长老们也是面露笑意，显然对这一届的生源质量还算满意。

对于绝大多数前来观礼的看客而言，这一整天的测试过程，整体并无太多超出预料的惊奇之处，算是“平平无奇”的一届。

\vspace{2em}

在大会即将结束之时，那位主持大会的白发长老再次登台，清了清嗓子，原本和煦的笑容收敛了几分，换上了一副略显严肃的公事面孔。

“诸位道友，今日升仙大会圆满结束，恭喜新入门的弟子。此外，老夫代表青云门，另有一事宣布。”他声音灌注了灵力，清晰传遍整个广场，压下了一片嘈杂。

所有人都竖起耳朵，不知还有何事。

“即日起，青石城外所有由我青云门管辖的灵石矿坑，准入费用将由原先的二成，上调至三成。”

长老的声音平稳，却在人群中激起了轩然大波！

“什么？三成？！”

“这……这怎么行！”

“原先二成就已经紧巴巴了，再加一成，还让不让人活了！”

“青云门怎能如此？”

台下顿时炸开了锅，尤其是那些靠采矿为生的散修，个个脸色大变，叫苦不迭，议论声、抗议声此起彼伏。

叶牧谦三人也是微微一怔。通过沈瑶之前的介绍和这几日的见闻，他们已大致明白这“矿税”意味着什么。

原来，青石城周边山脉中，蕴藏着不少灵石矿脉，此城也因此得名。那些储量丰富、品质上佳的大型矿脉，自然由青云门组织弟子和依附的矿工直接开采，视为宗门禁脔。但还有大量零星分散、矿脉较细、开采成本较高的小型矿坑，青云门不屑于投入大量人力物力亲自开采，却又不想白白浪费。

于是，青云门便想出了一个法子：由宗门派人守住这些小型矿坑的入口，划定区域，允许凡人与散修自由进入开采。但所有开采所得，必须上缴二成作为“准入费”或曰“矿税”，其余内容则按市场价折价收购。于是，采矿成了许多无门无派、缺乏稳定收入来源的散修最重要的经济支柱之一。他们冒着矿洞坍塌、遭遇低阶妖兽、甚至与其他散修争斗的风险，辛苦采掘，换取修炼资源或生活所需。

如今，这矿税直接从二成涨到三成，看似只增加了一成，但对于收入本就不丰的散修而言，相当于实际收入被砍掉了一成半！这无疑是在本就挣扎求存的散修群体身上，又狠狠割了一刀。

那长老似乎早已预料到下方的反应，面色不变，继续道：“此乃宗门决议，考量的是矿洞维护、阵法加固、以及抵御周边妖兽侵扰之成本增加。望诸位道友体谅。规矩已定，三日后正式施行，望周知。”

说完，他也不管台下如何哗然，袖袍一拂，便与其他青云门高层转身离去，只留下一群脸色灰败、愤愤不平的散修，以及更多事不关己、或同情或漠然的看客。

“师尊，这青云门……”白言冰眉头紧锁。他出身皇室，对“赋税”二字并不陌生；但如此明目张胆、近乎盘剥的加税，还是让他感到不适。尤其是看到那些散修眼中流露出的绝望与愤怒，更让他心中不平。

陈千羽也是面露忧色：“这些散修道友，日子怕是要更难过了。”

叶牧谦目光扫过群情激愤的散修人群，又望向青云门众人离去的方向，眼神深邃：“利益所在，矛盾所生。初来乍到，多看、少言。”

然而，树欲静而风不止。

升仙大会的余波尚未完全平息，矿税新规带来的动荡便已初现端倪。起初几日，散修们还只是私下抱怨，敢怒不敢言。但很快，冲突便由暗处转向了明面。

这一日，叶牧谦三人正在城中一处茶楼歇脚，顺便探听消息。忽闻楼下街道传来喧哗与打斗之声，其间夹杂着怒骂与痛呼。

三人凭窗望去，只见不远处一个售卖低阶矿石的摊位前，围了一圈人。一名衣着朴素、面色黝黑的中年散修，正被两名身穿青云门服饰的执事弟子推搡着，地上散落着几块品质一般的下品灵石原矿。那散修嘴角带血，脸上有一块明显的青紫，正激动地争辩着什么。

“……我就慢了半日！昨日进山遭遇了铁背狼群，好不容易才脱身……这三成的税，我、我还得养家，一时真凑不齐那么多啊！两位执事通融通融，宽限两日，一定补上！”散修苦苦哀求。

“宽限？人人都像你这般宽限，宗门规矩何在？”一名三角眼的执事弟子冷笑道，一脚踢开地上的矿石，“没钱？那就用这些破石头抵！不过，品质太次，得按市价七折算！”

“七折？这、这明明是市价九成的品质！”散修急了。

“我说七折就七折！再啰嗦，连这些都没收！还得罚你扰乱坊市秩序！”另一名执事不耐烦地挥了挥手，示意同伴动手强拿。

那散修情急之下，伸手去拦，却被那三角眼执事反手一掌拍在胸口，打得他踉跄后退，又是一口血喷出，气息顿时萎靡下去，显然受了不轻的内伤。

周围围观者众多，有散修面露不忍与愤慨，却无人敢上前。那些青云门执事弟子则趾高气扬，浑然不觉有何不妥。

茶楼上，陈千羽“噌”地站了起来，柳眉倒竖：“光天化日，岂能如此欺人！还下此重手！”

她医者仁心，最见不得伤患，更何况是这般不平之事。

“千羽！”白言冰低喝一声，想拦住她。叶牧谦也微微摇头，示意她冷静。

但陈千羽已然按捺不住。她身形一闪，便已下了茶楼，分开人群，快步走到那受伤倒地的散修身旁。

“这位道友，你伤势不轻，莫要乱动。”陈千羽蹲下身，声音温和，同时手指已快速搭上对方腕脉，一缕温润平和的元气悄然渡入，护住其心脉，缓解内腑震荡。

那散修本已意识模糊，感到一股充满生机的暖流涌入体内，痛苦顿时减轻不少，惊愕地看向眼前这位容貌清丽、气质出尘的女子。

“你、你是……”

“路过之人，略通医理。”陈千羽简单答道，又从怀中取出一个瓷瓶，倒出一颗自己炼制的疗伤丹药——此时已经不再是淀粉丹了，而是改良后的正经丹药——便要喂给散修。

然而，就在此时——

“住手！你们是什么人？竟敢擅自接触抗税刁民！”一声厉喝传来。

只见那两名青云门执事弟子，以及不知何时闻讯赶来的另外三名巡视弟子，已然将陈千羽与那散修围住。为首一名面容冷峻的巡视弟子，目光锐利地扫过陈千羽，又瞥见茶楼窗口的白言冰和叶牧谦，眼中闪过一丝狐疑与警惕。

陈千羽心中暗叫不好，自己一时冲动，竟忘了师尊“多看少言”的告诫。她缓缓站起身，平静道：“此人受伤甚重，我只是行医者本分，施以援手。与抗税与否，并无干系。”

“行医者本分？”那冷峻弟子冷笑，“我看你们是一伙的吧？在此聚众闹事，煽动民心，对抗宗门法令！此人抗税在先，受伤乃咎由自取！你等速速退开，否则，以同谋论处！”

他的话音落下，另外几名弟子或将手放在各自的武器上，或准备掐诀，气息锁定了陈千羽，以及从茶楼中走下来的白言冰和叶牧谦。

叶牧谦缓步走到陈千羽身边，将她稍稍护在身后，目光平静地看向那冷峻弟子：“这位道友，小徒只是心善救人，并无他意。至于抗税、聚众、煽动民心……未免言重了。”

“是不是言重，不由你说了算！”冷峻弟子见叶牧谦气度沉凝，心中不免忌惮。

但仗着己方人多且是青云门弟子，语气依然强硬，“你们三人，面生得很，并非我青石城常驻修士。在此敏感时期，行为鬼祟，不得不查！请随我们回山一趟，向执事长老说明情况！”

这便是觉得自己能力不够，选择“请”他们上山了。语气看似客气，实则不容拒绝，带着明显的胁迫意味。

白言冰踏前一步，挡在师尊和师妹身前，眼神锐利如剑：“若我们不去呢？”

“不去？”冷峻弟子眼中寒光一闪，“那便是心中有鬼，抗拒执法！休怪我等不客气了！”

气氛瞬间剑拔弩张。周围看热闹的人群见势不妙，纷纷后退，生怕被卷入争斗。

叶牧谦轻轻叹了口气，拍了拍白言冰的肩膀，示意他稍安勿躁。他看向那冷峻弟子，说道：“伤者需即刻医治，否则性命堪忧；万一闹出个三长两短，贵宗‘草菅人命’的名声传开去……”

一边说着，叶牧谦一边观察着周围情况，没有发现明处的真符境修士。内心思忖：如果就此逃离，青云门肯定没办法强留下我们，但一定会迁怒于青石城这些无辜的散修，自己也会坐实“畏罪潜逃”之名，反为不美，更不能问心无愧。

既然已经卷入了这件破事，那就把这件事处理好。

他心中明镜似的。青云门提高矿税，引发散修不满，正需要树立典型，打压可能出现的反抗苗头。陈千羽救治受伤散修，恰好撞在了枪口上，被对方视作“聚拢民心”“潜在威胁”。这一趟“上山”，恐怕是宴无好宴。

而那青云门弟子似乎也陷入了沉思：这家伙气息沉凝、做事谨慎，且做帽子倒是有好手艺。如今“草菅人命”这大帽子一扣，这种恶名传出去之后，宗门未必不会处理自己。但当街开打……感觉不像能打过的样子。

略一沉吟，道：“好，我将令一名同门随这抗税之徒去医馆诊治。但你们三个，还是得跟我们走一趟。”

这下语气可客气得多了。

随即，青云门执法队伍中，果真有一个弟子带着受伤散修去了医馆——实际上是以这名同门为质子，重新占领了道德高地。

叶牧谦寻思：事已至此，避无可避。正好，他也想借机更深入地看看，这掌控一城生死的青云门，究竟是何等模样。

“师尊……”陈千羽有些愧疚地低声道。

“无妨。”叶牧谦微微一笑，“且去看看。记住，谨言慎行。”

随即用更低的声音道：“上山途中，你们一定要暗中记录青云门的防卫情况，确定事不可为的情况下，立即寻防御最薄弱处，我们打出去。”

于是，在剩下四名青云门弟子隐隐呈包围之势的“护送”下，叶牧谦师徒三人，离开了喧嚣的街道，朝着青石城中央那座云雾缭绕、殿宇隐约的山峰——青云门山门所在——行去。

沿途，不少路人投来或好奇、或同情、或幸灾乐祸的目光。散修与宗门弟子的冲突并不罕见，但被如此“客气”地请上山的，却也不多。许多人心中已然明了，这三位面生的修士，怕是要有麻烦了。

青石城中央，青云门巍然矗立，云雾常年缭绕山腰，亭台楼阁于林间若隐若现，飞檐斗拱在日光下泛着清冷的灵光，确有一派仙家气象。

然而，沿着蜿蜒而上的石阶，在数名青云门弟子隐隐挟持般的“护送”下前行，叶牧谦师徒三人心中却无半分欣赏景致的闲情，只有警惕与暗中盘算、冷眼旁观。

穿过几重殿宇，他们被带到了一处偏殿。殿内陈设简单，却自有一股威严，正中主位空着，两侧各有数张座椅。

不多时，一阵略显轻浮的脚步声传来。一名身着华贵青色锦袍、腰悬玉佩、看起来约莫二十出头的青年，在一众弟子的簇拥下步入殿中。他面容算得上英俊，但眉眼间带着一股挥之不去的骄纵之气，目光在叶牧谦三人身上扫过，尤其在陈千羽脸上多停留了一瞬，嘴角勾起一抹意味不明的笑容。

“在下李轩，青云门门主正是在下家父。听闻有几位道友在城中与我门下弟子有些误会，特来了解一番。”

青年自顾自地在主位上的一张大师椅上坐下，姿态随意，语气带着些居高临下的味道。他修为波动显示其处于玄脉境开脉期，在此地年轻一辈中，算是不错；但根基在叶牧谦感知中，略有虚浮。

“误会已向贵宗执事弟子说明，小徒只是行医救人，并无他意。”叶牧谦不卑不亢，淡然回应。

在他的感知中，暗处确实藏匿了一缕收敛极好的气息——玄脉境融贯期大圆满，兴许是李轩的后手。

“行医救人？好事啊！”李轩抚掌一笑，仿佛很感兴趣，“正巧，我青云门内有不少弟子，常年在山中历练，或与妖兽搏杀，或开采灵石时遭遇意外，落下不少陈年旧伤，一直苦无良医。今日既然遇到陈姑娘这般心怀仁术之人，不知可否请姑娘施展妙手，为我这些同门诊治一番？也算是化干戈为玉帛，结个善缘。”

话音落下，也不待一行三人回应，便自顾自拍了拍手。

殿外立刻走进来七八名青云门弟子，有男有女，年纪不等，但个个面色或苍白、或晦暗，气息也有些不稳，显然身上都带着或轻或重的伤势，且看起来时日不短。他们进入殿中，便齐齐向李轩行礼，又对陈千羽投来混杂着期盼与怀疑的目光。

这一幕看似合情合理，但叶牧谦眼神微凝，白言冰也悄然握紧了拳头。他们都能看出，这些弟子伤势各异，有的伤了经脉，有的损了脏腑，有的甚至残留着难以驱除的异种气劲或毒素，绝非短期可愈。这李轩，恐怕没安好心。

陈千羽却半是因为医者仁心、半是因为老实，见到伤患，下意识地便仔细打量起来。她越看眉头蹙得越紧，但刚欲上前，就心有余悸般地收住了脚。

她与叶牧谦交换了一个眼神，叶牧谦示意她该怎么看病就怎么看病。

于是陈千羽终于上前去，示意排在最前那名青云门弟子伸出手腕，仔细诊脉，又询问了几句受伤经过与当前感受；之后的每个人都大致这么诊断了一番。

片刻后，她收回手，转向李轩，神色认真而坦率：“李公子，贵派这几位道友的伤势，确为积年顽疾。经脉郁结、脏腑暗伤、异气残留……情况复杂，非朝夕之功可解。若想彻底调理康复，不留后患，以我的能力，最短也需……”她略一估算，“至少半年以上，且需他们全力配合，耐心调养。”

她这是实话实说，甚至已经考虑了对方可能急于求成，将时间往紧了说。

然而，李轩脸上那抹笑容却陡然变得夸张而戏谑：“半年？陈姑娘，你这效率未免太低了吧？我这些同门，可都是宗门的中坚力量，每日都有任务在身，岂能耽搁半年之久？”

他身体前倾，目光逼视着陈千羽，语气陡然转厉，带着不容置疑的强硬：“我给你一个月！一个月内，我要看到他们全都活蹦乱跳，伤势尽复！否则，便是你医术不精，或是……存心敷衍，戏耍我青云门！”

一个月？医治七八个轻重不一的陈年旧伤？这根本是天方夜谭，是强人所难，是赤裸裸的刁难！

陈千羽脸色一白，眼中露出难以置信的神色。白言冰再也按捺不住，一步上前，将陈千羽护在身后，怒视李轩：“李轩！你明知这是不可能完成之事，分明是故意刁难！这就是你们青云门的待客之道？这就是你所谓的‘结善缘’？”

叶牧谦的脸色也沉了下来。他原本以为对方只是借机试探或施压，没想到竟如此下作，直接提出这种完全违背常理的要求。他看着李轩那副有恃无恐、等着看好戏的嘴脸，心中最后一丝对此地宗门的观望也化为了冷意。

“李公子。”叶牧谦开口，声音不大，却带着一种奇异的穿透力，让殿内嘈杂为之一静。

“修行之人，首重修心。心有偏私，行事不公，纵有修为，亦难窥大道真谛。我辈修士，当有‘见独’之明，明辨是非，恪守本心。你今日此举，挟势凌人，强人所难，可曾问过自己的本心？可曾有过片刻‘问心无愧’？”

他这番话，既是在斥责李轩，也是在点醒那些带着伤、或许原本只是奉命而来的青云门弟子，更是阐述问道途重心性、重内明的理念。

而叶牧谦自己也被自己脱口而出的“见独”二字启发了：他先前一直在推演问道途养气之后的下一境界，“朝彻而后能见独”（《庄子·大宗师》）这句话，恰好能够极为精确地描述所需的心境与境界的核心能力。

“见独？问心无愧？”李轩仿佛听到了什么极其可笑的事情，嗤笑出声，满脸的不屑与讥讽，“‘见独’是什么东西？哪门哪派的酸腐道理？在这青石城，我青云门的话就是道理！实力就是本心！你们既然进了这门，是龙得盘着，是虎得卧着！治不好人，就是你们的错！”

他显然将叶牧谦所说的“见独”当成了某种陌生的、不值一提的教条，完全无法理解其中蕴含的“内求己身、明心见性”的深意。在他看来，力量、权势、宗门的意志，才是衡量一切的标准。

“看来是没什么好说的了。”李轩失去了耐心，脸上戾气浮现，“既然敬酒不吃吃罚酒，那就别怪本公子不客气了！给我拿下，水牢里关三天再说！尤其是这个女的，押下去，什么时候答应一个月治好，什么时候再说！”

殿内气氛瞬间降至冰点！李轩身后数名气息较强的弟子齐齐踏前一步，灵力波动腾起，锁定了叶牧谦三人。更有两人直接扑向陈千羽，手中闪烁着灵光，竟是直接动用擒拿术法！

“你敢！”白言冰目眦欲裂，长剑瞬间出鞘半尺，养气境元气勃发，就要拼死一战。陈千羽也钢针在手，面色决绝。

就在这千钧一发之际——

“唉。”

既然你故意发难，那我至少也得让你知道，我不是那么好惹的！

打得一拳开，免得百拳来！

一声轻微的叹息，仿佛直接在每个人心头响起。叹息声未落，一股难以言喻、仿佛源自灵魂深处的威压，以叶牧谦为中心，悄然弥漫开来！

没有狂风，没有灵光爆闪，甚至没有强烈的元气波动。

但就在这一瞬间，扑向陈千羽的那两名弟子，动作骤然僵住，如同陷入了无形却无比粘稠的琥珀之中，保持着前扑的姿势，凝固在半空，连眼珠都难以转动！

不仅是他们，李轩以及他身边所有试图动作的弟子，都感到周身一紧，仿佛被无数道无形的视线锁定，被一种浩渺、高远、不容置疑的意志所笼罩！他们体内的灵力运转瞬间变得滞涩无比，如同生锈的齿轮，连呼吸都变得困难！

——实际上，这可以算是叶牧谦初来乍到时对“言出法随的绝对领域”的思考与实现了：领域之内此消彼长，谁敢轻易撄其锋芒？

先前在叶牧谦感知中的那名玄脉境巅峰的高手本欲出手，但很快察觉了这力量的本质。

极强的心念控制能力！

甚至能够外放并影响敌人！

高手立即选择静观其变，并通知门主。

与此同时，李轩脸上的狞笑僵住了，化为极度的惊骇。他拼命鼓动玄脉中的灵力，试图冲破这无形的束缚，却感觉自身的力量在那股意志面前，渺小如蝼蚁撼树，根本掀不起半点浪花！

他只能眼睁睁看着，看着叶牧谦缓缓从座椅上站起，目光平静地看向他，那目光仿佛穿透了他的血肉，直视他灵力运转的每一个晦涩之处，直视他因急功近利、过度依赖灵石而积累在经脉深处的杂质与隐伤！

这……这是什么力量？！绝不是灵枢径的灵力！没有灵枢波动，没有玄脉共鸣，却如此恐怖！李轩心中掀起了惊涛骇浪，前所未有的恐惧攫住了他。

叶牧谦的脸色微微苍白了一丝，强行催动这尚未完全稳固的见独领域，对他负担不小。但他面上丝毫不显，只是看着动弹不得、满脸骇然的李轩，缓缓开口，声音如同暮鼓晨钟，敲在每个人心上：

“贪恋灵石外物，急功近利，过度汲取而不加淬炼，以致经脉杂质满溢，运转滞涩，玄脉看似开阔，实则根基虚浮，隐患深种。长此以往，非但修为难有寸进，他日灵力暴走、经脉寸断，亦不无可能。”

他每说一句，李轩的脸色就白一分，因为叶牧谦所说，与他近来修炼时偶尔感到的凝滞、以及父亲曾隐晦提点过的问题，丝毫不差！

“念你年轻，尚有回头之机。”叶牧谦语气转淡，却带着一种不容置疑的权威，“若肯散去一重修为，将玄脉中淤积的杂质与虚浮灵力一并化去，从此脚踏实地，重头稳扎稳打，尚有补救之望。否则，便是自绝道途。”

说罢，叶牧谦心念一动，收回了那无形的势场。

“噗通！”“噗通！”

那两名被定在半空的弟子直接摔落在地，狼狈不堪，却连痛呼都不敢发出。李轩踉跄后退两步，后背瞬间被冷汗浸湿，看向叶牧谦的目光充满了恐惧与后怕，先前那股骄纵之气荡然无存。

“多、多谢前辈指点……晚辈、晚辈铭记……”李轩低着头，声音干涩，唯唯诺诺，再不敢有丝毫放肆。

叶牧谦强压反噬，不再看他，转身对白言冰和陈千羽道：“此地无甚意思，我们走吧。”

白言冰与陈千羽也从方才那震撼一幕中回过神来，连忙点头。他们也看出这青云门绝非善地，李轩虽然暂时被震慑，但其人心性不佳，难保不会怀恨在心，暗中使绊子。

陈千羽用极低的声音传信道：“师尊，那原定的撤离路线……”

叶牧谦摇了摇头，意思是用不上了。

白言冰小声问道：“不知师尊所用是何秘法？”

叶牧谦也没什么避讳的心思，直接开口说道：“不是秘法，不过是下一境界‘见独’稀松平常的能力罢了。”

见独境！白言冰眼前一亮。先前师尊从来不提及的“养气之后下一境界”，原来叫做这个名字！

而师尊之所以从来没有跟我们提过这一境界，一定是因为希望我们不要贪功冒进，而是打好养气境界的基础再提突破！

他与陈千羽交换一个眼神，看起来她也是这么想的。

就在三人即将走出殿门时，叶牧谦脚步微顿，并未回头，只是淡淡留下一句话：“至于这些有伤的弟子，若信得过，可自行下山，到城中我们落脚处寻千羽诊治。医者本分，不因人事而废。”

说完，三人身影便消失在殿外走廊，大大方方走下了山。

只留下殿内一片死寂，以及神色复杂、惊魂未定的青云门众人。

\vspace{2em}

离了青云山，回到城中客栈，三人才稍松一口气。

叶牧谦刚刚从强烈的越阶反噬中恢复回来，脸色仍有些许苍白。

他评论道，“这青云门，格局太小，心思不正，不宜久留。我们略作休整，便离开吧。”

陈千羽则有些担忧：“师尊，我们就这样走了，那李轩会不会……”

“短时间内应不敢妄动。”叶牧谦目光深邃，“他看起来是个惜命且自私之人，被我点破修炼隐患，此刻心中恐怕更多的是后怕与权衡。至于他是否听得进劝告……”叶牧谦摇了摇头，不再多说。

他猜不能。

\vspace{2em}

与此同时，青云山主殿深处。

李轩正心有余悸地向其父，青云门门主李青阳禀报方才发生的一切，事无巨细，尤其是叶牧谦那诡异莫测的“势场”以及对他修炼隐患的洞察。

李青阳是一位面容清癯、目光深沉的中年修士，修为也是玄脉境融贯期巅峰，也算一方高手。他听着儿子的叙述，手指轻轻敲击着座椅扶手，半晌不语。

他早都从青云门大长老——也就是刚刚被叶牧谦误解的“李轩后手”，实际上是李青阳派来考察他那个儿子行事能力的——嘴里，听完了事情的全过程。

至少，李轩没有隐瞒什么事情。

“那……那他说的，关于我修炼之事……”李轩迟疑了一下，问道。

散去一重修为重修，这代价对他来说太大了。

李青阳瞥了儿子一眼，知子莫若父，他岂能不知李轩心中所想？对于这个骄纵的儿子，他有时也感到头疼。

最终，他淡淡说道：“修行是你自己的事。他既然指出了问题，如何抉择，你自己随意。”语气平淡，听不出是赞同还是反对，更像是一种放任。

李轩闻言，心中一定。父亲没有明确反对，那便是默许了他自己的选择。散去修为？重修？开什么玩笑！那得多辛苦，多丢面子！自己可是门主之子，玄脉境开脉期，在年轻一辈中也是佼佼者！那叶牧谦说不定只是危言耸听，或是其道途不同，看不惯我们灵枢径的修炼方式罢了。

他心中那点后怕，很快被侥幸与固有的骄纵所取代。

随即，李轩又向在一旁坐着的大长老抱怨：“在那家伙出手的时候，您为何不出手、选择冷眼旁观？”

大长老对李轩并无好感，他只是冷冷地说：“他本可当场废你修为，却留有余地；再先前，他本可在青石城就对巡视弟子大打出手，却选择以理服人。此人行事，一击必中却留三分情面，才是最可怕之处。”

“所以那叶牧谦……究竟是何境界？用的什么法术？简直闻所未闻！”李轩忍不住问道。

“法术古怪，境界难测。”李青阳缓缓开口，眼中精光闪烁，“虽无灵枢玄脉波动，但能轻易压制于你，且一眼看穿你修炼弊病……其实力，恐怕至少也是玄脉境融贯期巅峰，甚至可能已经是真符境。”

李轩倒吸一口凉气。他自幼在青石城长大，没见过多少真符境修士，父亲就是他眼中的天花板……这叶牧谦竟然是一个可能比父亲强的人物？！

“能制人而不制，能立威而不结仇，这种人要么深交，要么远离，但不宜为敌。”李青阳做出了决断，“传令下去，门下弟子不得再去寻那三人的麻烦。他们若懂些人情世故，见识了我青云门的态度，想必也会很快自行离去。”

在他看来，叶牧谦一行人不过是过客，没必要为了一点小事与一个神秘强者结怨。

李轩应诺而退。出得殿门，正遇上一位跟班弟子，显然后者趴在门外偷听了全程。

那跟班适时的前来询问：“少门主，关于那聚灵阵……”

“以前怎么做，现在就怎么做。”李轩不以为然地回答。

聚灵阵是他修炼必备之物，能调动灵石中的灵气灌注于阵中，短时间内强行营造出一片灵气浓郁之处。由于此法非常耗费灵石，故而即使是超级势力都不多用；只是青石城周围全都是灵石矿，青云门也因此巨富，供养一两个聚灵阵倒是绰绰有余。

于是李轩就把这件事抛到了脑后；只是偶尔想起叶牧谦那平静却深邃的目光时，心底会掠过一丝不易察觉的阴影。

\vspace{2em}

青云门的风波，看似暂时平息。叶牧谦等人在青石城外荒山上寻了一处年久失修、早都无人居住的石屋临时歇脚，一边也在探听其余城池的情况，打算找个地方定居下来。叶牧谦本人也在试图打听关于“外界界门”的消息，但一无所获。

期间，竟真有几名那日殿中带伤的青云门普通弟子，偷偷下山寻来，恳请陈千羽医治。陈千羽信守承诺，尽心为他们诊治，也不多收费用，只是提醒他们修行莫要贪功冒进。

是夜，月隐星稀，山风微凉。连日心神损耗，叶牧谦早早便歇下了。

或许是因为心神松弛，又或许是因为那强行催动见独领域的反噬并未完全平息，他陷入了极为深沉的睡眠。

夜深人静时，石屋一侧的小厨房里，陈千羽正守着一个小泥炉，小心翼翼地熬制着一锅安神调理的汤药。药香混合着柴火气，在狭小的空间里弥漫。白言冰在一旁帮忙添柴看火，火光映照着他沉思的侧脸。两人都未说话，各自想着心事，屋内只有柴薪轻微的噼啪声和药汤咕嘟的微响。

忽然，一阵模糊的、断断续续的呓语，从隔壁叶牧谦休息的里间传来。

起初声音很低，含糊不清。白言冰与陈千羽对视一眼，均有些担心，于是凝神细听。

“……设定集……框架……地理……社会形态……修炼道途……空白……”叶牧谦的声音在梦中显得有些焦急，又带着一丝困惑，反复咀嚼着这几个词，仿佛在为什么难题所困。

“设定集？”白言冰眉头微皱，低声重复。这个词他从未听过，听起来像是指某种典籍或律法总纲？

紧接着，叶牧谦的梦呓内容陡然一变，语气变得悠长而富有韵律，竟开始背诵起一些文辞古奥、意境深远的句子来！“……北冥有鱼，其名为鲲。鲲之大，不知其几千里也；化而为鸟，其名为鹏……”“……天地有正气，杂然赋流形。下则为河岳，上则为日星……”“……故天将降大任于是人也，必先苦其心志，劳其筋骨……”

这些句子或磅礴，或刚正，或富含哲理，与当前世界的文风大相径庭；其蕴含的宏大想象、浩然之气与砥砺精神，更让听者心神为之所夺！

白言冰与陈千羽自幼也读过不少诗书，此刻听得目瞪口呆。这些篇章，他们闻所未闻，但每一句都仿佛敲击在心坎上，令人回味无穷。

师尊在梦中……竟能随口诵出如此瑰丽深邃的经典？这绝非寻常修士所能为！

梦呓声渐渐低下去，最终归于平静，只剩下均匀的呼吸声。但厨房里的两人，心中却掀起了惊涛骇浪。他们熄了炉火，将熬好的药汤温在灶边，轻手轻脚退到外间，面面相觑，眼中充满了震惊与无法抑制的联想。

“言冰，你听到师尊刚才说的……‘设定集’了吗？”陈千羽压低声音，语气犹疑。

“听到了。还有那些……文章。”白言冰眼神闪烁，大脑飞速运转，结合他们一直以来对叶牧谦的认知——神秘出现、开创迥异于“灵枢径”的“问道途”、实力深不可测、言行往往暗含深意……

一个大胆的、甚至有些骇人听闻的猜想，在两人心中不约而同地形成，并且迅速加固，化为坚不可摧的“共识”。

“设定集……”白言冰深吸一口气，声音带着一丝颤抖，“莫非……是记载天道运行、万物法则的……无上律典？”

“师尊他……”陈千羽接道，眼中泛起无比的敬畏，“平日所言所行，看似随意，实则暗合大道。他传授我们的‘问道途’，与此界通行的‘灵枢径’截然不同，更重内求与心性……如今想来，这岂不正是……补全此界仙路缺失、指引真正长生方向的……无上法门？”

“所以师尊才会反复提及‘设定集’！那定是他从上界带来，或者说，他本身就是从上界临凡的仙尊！”白言冰越说越觉得合理，思路如泉涌，“为了将此界走偏的仙路导回正轨，他降临此世，传下‘问道途’！正因如此，他才必须时刻压制自身真正的修为与气息，否则，以他上界仙尊之能，全力施为，恐怕会引发此方天地法则的剧烈排斥甚至……愤怒反噬！”

这个解释，完美地契合了他们心中对叶牧谦所有神秘之处的疑惑。为何师尊总显得与世格格不入？为何他的力量体系如此特殊？为何他有时会露出对常识的陌生？为何他梦中会诵出闻所未闻的至高经典？一切都有了答案！

两人对视，都从对方眼中看到了同样的笃定与无以复加的崇敬。原来，他们所拜的师尊，竟是背负着补全天道、重塑仙路这等宏大使命而临凡的上界仙尊！这是何等的机缘！何等的荣耀！

\vspace{2em}

次日，天色阴沉得厉害。铅灰色的云层仿佛浸透了水的棉絮，沉甸甸地压在山峦之巅。

叶牧谦早早地起了，却见到门外两个比他起得更早的弟子，但两人的恭敬程度却比前几天更甚，几乎不敢抬头看他。

“你们起这么早，还趴在门外等我，是有什么事情吗？”叶牧谦感到奇怪，问道。

“师尊……不是此方天地之人？”白言冰小心翼翼地问道。

叶牧谦心里咯噔一声。

坏了，难不成自己的功法都是瞎编的这件事，竟然被他们看出来了？

强作镇定，答道：“怎么突然问起这个？”

两个徒弟对视了一眼，头埋得更深了。

“师尊昨晚说梦话，提及了些……不似此方天地的文本。其雄奇瑰丽，非我们所读过的书籍可比。”陈千羽嗫嚅道，“所以我们好奇师尊是不是上界仙尊临凡……”

叶牧谦大松了一口气，甚至有些想笑。

还行，至少徒弟们没怀疑他。不过这种奇怪而可笑的思路，叶牧谦不禁好奇其产生的原因。

“那你倒是说说，要真是上界仙尊，又怎肯临凡呢？”

两人把他们“重构仙途”的思路大致说了一遍。

叶牧谦终于放声大笑出来，留下两位弟子面面相觑。

本来想挑个时间和两个徒儿旁敲侧击地坦白一下自己来头的，没想到他们已经猜中了一部分，那也该有选择的告知他们实情了。

忽然，他敛住笑容，郑重其事地说道：“你们至少猜对了一部分：我确实来自异世。至于其他事项……先不说了，以后时机成熟，我自会知无不言。”

随即挥手：“去吧，该做什么做什么，不要再思考这些事情了。”

\vspace{2em}

临近午时，高天上积蓄已久的力量终于爆发，瓢泼大雨倾泻而下。

仿佛天河倒灌一般，瀑布般的水幕连天接地，哗啦啦的巨响瞬间充斥了整个世界。浓密的铅云之中，银亮的电蛇开始狂舞。起初只是偶尔一闪，映得天地一片惨白，随即，隆隆的雷声便从极远处翻滚而来，如同巨兽的咆哮，越来越近，越来越响。

“轰咔——！！！”

炸雷几乎在头顶爆开，整个山崖似乎都在这天威之下微微颤抖，碎石和泥土被震得簌簌滑落。狂暴的雷声在山谷间来回冲撞、叠加，久久不息，与暴雨的喧嚣交织成一片毁灭般的交响。

白言冰当时正在厨房里，蹲在粗糙的土灶前，帮陈千羽准备午饭。灶膛里的柴火噼啪作响，橘红色的火苗跃动着，驱散了些许雨天的阴冷与潮气。陈千羽在一旁案板上切着些简单的山蔬，两人都未说话，只听着外面震耳欲聋的雨声。

就在那道刺目闪电亮起的瞬间，白言冰下意识地抬头望向窗外，眼底依然残留着树枝状的光痕。

“轰——！”

巨响仿佛直接在他颅腔内共鸣。他浑身一震，蹲着的身子微微一晃。更奇异的是，灶膛里原本平稳燃烧的火苗，被这巨大的声波一震，竟猛地向上蹿起老高，火舌舔舐着漆黑的锅底，光影在厨房墙壁上剧烈摇晃。

就在这雷声的余韵仍在耳中轰鸣、眼前跳跃的火光尚未平息的刹那，白言冰心中，仿佛有一层一直蒙着的薄纱，被这至刚至强的天威与炽烈跃动的火焰，猛地撕开了！

雷。

那撕裂黑暗、照亮寰宇的闪光；那代表天罚、涤荡一切妖邪秽物的轰鸣；那是至刚至阳、不容置疑的毁灭之力，是天道威严最直接、最暴烈的彰显。它迅疾、暴虐、拥有撕碎一切的决绝。

火。

灶膛里温暖跃动的光芒，带来熟食与生机；但火也能焚尽山林，熔金化石，是毁灭的使者。它炽烈、奔放、拥有吞噬万物又转化万物的矛盾特性，既能带来光明与温暖，亦能带来灰烬与终结。

两者，皆是这自然世界中，最为纯粹、最为极致的刚与强的象征。

种种思绪，杂糅着对自然伟力的敬畏，对自身道路的求索，以及对那位神秘师尊背后可能代表的浩瀚天道的隐约触碰，在这雷声与火光交织、外界狂暴与内心激荡共鸣的瞬间，轰然碰撞、迸发、然后……开始了奇异的融合！

他依旧保持着蹲在灶前、手持柴火的姿势，身体一动不动，仿佛一尊石雕。但那双原本清亮的眼睛，却在刹那间失去了焦点，瞳孔微微扩散，视线穿透了眼前的灶火、墙壁，投向了某个无法言说的深远之处。他的灵魂仿佛脱离了躯壳，沉浸到了一种玄之又玄的感悟状态之中。

周身原本内敛的元气，自然而然地缓缓流转。然而，在这份沉静的内核之外，又隐隐透出一股令人心悸的炽热与凌厉，仿佛平静火山下奔涌的熔岩，又似雷云中积蓄的毁灭能量，随时可能爆发。

“言冰？”陈千羽第一时间察觉到了异常。她放下手中的菜刀，看到白言冰那副魂游天外的模样，感受到他周身那股奇异而引而不发的气息，立刻明白过来——这是无数修士梦寐以求的“悟道”状态！

机缘难得，稍纵即逝，最忌打扰。

她连忙捂住自己的嘴，将所有的疑问和关切都压回心底，连呼吸都放得轻缓无比。她小心翼翼地挪动脚步，走到厨房门边，轻轻掩上半扇门，既挡住了可能从堂屋传来的细微声响，又能让自己守住门口，确保没有任何意外——哪怕是风吹动门扉的声音——去惊扰此刻的白言冰。

这场罕见的暴雨雷暴，仿佛是天公为了配合这场悟道，肆意宣泄着它的力量。银蛇乱舞，惊雷滚滚，暴雨如注，足足持续了将近一个时辰。直到乌云中积蓄的电光与水量似乎消耗殆尽，那令人窒息的轰鸣与冲刷声才渐渐减弱、停歇。

几乎就在最后一滴雨珠从屋檐坠落、天地间恢复宁静的同一刹那——

厨房里，如同泥塑木雕般静坐了一个时辰的白言冰，那双一直闭合的眼睑，猛然睁开，眼底混合着强烈的、亟待宣泄的创造冲动。

他豁然起身，动作干脆利落，甚至带着一股迫人的劲风。蹲坐一个时辰的麻痹感似乎对他毫无影响。他满怀歉意地看了一眼守在门边、满脸关切的陈千羽，一言不发，大步流星地穿过堂屋，径直推开木门，走到了石屋前那片被雨水浸透的空地之上。

泥土松软潮湿，踩上去微微下陷。空气中弥漫着雨后的凉意，但他体内却仿佛有一座火炉在燃烧。

“锵——！”

清越如龙吟般的剑鸣响彻山涧。白言冰反手拔出了腰间那柄陪伴他许久的精钢长剑。剑身映照着雨后清澈的天光，泛起一泓秋水般的寒芒。

他双手稳稳握住剑柄，立于湿漉漉的泥地中央，缓缓闭上了眼睛。并非休息，而是在做最后的凝神与整合。胸膛微微起伏，一呼一吸之间，仿佛与周围刚刚平息却余韵未散的自然韵律相合。

全身的元气以前所未有的方式疯狂涌动起来：如同百川归海，又似火山喷发前的熔岩汇聚，以丹田为源头，轰然冲向四肢百骸，最后尽数灌注于双臂，涌入那柄看似平凡的长剑之中！

剑身微微震颤起来，发出低沉的嗡鸣。淡金色的元气光芒自他双手蔓延而出，迅速包裹了整把长剑，使得剑身看起来仿佛由淡金色的光铸成。

雷与火，两种截然不同却又本质相通的“刚强”意象，在他心中不断盘旋、碰撞、交融。

最终，化为一股纯粹至极的意志：刚猛无俦，一往无前，涤荡邪祟，焚灭阻碍！

就是现在！

白言冰目光锁定了前方约三丈外，一块半人高、通体青黑色、质地异常坚硬的山岩。

沉腰坐马，单手持剑，猛地向前一刺！

升皎弦！

剑锋划破潮湿的空气，淡金色元气凝于剑尖，轻易刺入石中数寸有余。

随即用腕力向上一挑，浑厚而凝聚的元气瞬间爆发，产生了惊人的向上力道。只见那沉重的、起码有数百斤的青黑色山岩，竟轻飘飘地离地飞起，脱离剑尖，斜斜冲向离地约两丈的空中！

以前，他都是用升皎弦当杀招的。但这一次，他只是为了接下来那一剑，创造一个完美的、悬空的靶子！

就在山岩升至最高点，微微一顿，即将开始下坠的瞬间，白言冰足尖在湿软的地面狠狠一蹬，身形如箭，疾射而出。凌空一跃，精准地跃升至与那空中山岩齐平的高度。

人在半空，无处借力，但他全身的力量、元气、意志，都在这一刻压缩、凝聚、攀升至巅峰！

与此同时，改单手为双手持剑，高举过顶；全身元气开始向长剑上疯狂涌动，那淡金色的元气光芒开始剧烈地波动、沸腾起来！

斩——！

只是最简单而霸道的一记竖劈！将全身的力量、所有的感悟、那股雷火交织的刚猛意志，尽数倾注于这一剑之中，仿佛要将面前的一切，连同这压抑的空气，都彻底劈开！

剑锋破空！

“呜——嗡——！”

奇异而震撼的声响爆发了！那是低沉如闷雷滚动、又夹杂着仿佛火焰爆燃的风雷之声！

那凝聚到极点、高度压缩、蕴含了雷暴毁灭与火焰焚尽意象的沛然剑气，随着白言冰意志的决绝挥斩，脱剑而出！一道半月形的、边缘闪烁着黑紫色电光与暗红火屑的光弧，仿佛割裂了空间，以超越视觉的速度，瞬间劈中了空中那块青黑色山岩！

“轰隆！！！”

只有一声沉闷到极致、又爆裂到极致的轰鸣！

坚硬的山岩，在那道雷火剑气面前，如同热刀切入凝固的牛油，又似被天雷正面击中，毫无滞涩地从中一分为二！

两块硕大的岩石残骸，带着沉重的风声，向左右两侧抛飞，重重砸落在泥泞的地面上，溅起大片泥水。

切口处，光滑如镜，甚至能映照出天空的云影。然而，在这光滑切面的边缘，却呈现出截然不同的破坏痕迹：一侧是明显的焦黑、碳化，仿佛被极高的温度瞬间灼烧而过；另一侧则布满了细微的、放射状的皲裂纹路，纹路边缘呈现出琉璃化的质感，那是被极度凝练的狂暴力量瞬间冲击、又未彻底焚烧的特征，宛如雷击之后的残留！

空气中弥漫开一股淡淡的、类似焦石的奇特气味。烟尘缓缓升腾，又被刚刚下完雨湿润的空气压下。

白言冰招式用老，身形随着挥剑的力道下坠，单膝微曲，缓冲落地，依旧保持着双手握剑向下劈斩的姿势。剑尖向后斜指地面，其上缭绕的赤红银白细芒缓缓黯淡、消散。

他微微喘息着，胸膛起伏，额角渗出细密的汗珠。一次性调动几乎全身的元气，精神又高度集中，客观上负担确实不小。但此刻，他脸上却看不到丝毫疲惫，只有一种难以言喻的、近乎颤抖的兴奋与自得！

直起身，收剑回鞘，目光先是落在地上那被整齐劈开、边缘残留着雷火痕迹的两半岩石上，怔怔地看了好几息，仿佛在确认这真的是自己做到的。随后，他又低下头，看向自己手中这柄看似寻常的长剑，手指轻轻拂过尚且温热的剑鞘，眼神复杂，有陌生，有惊喜，更有一种前所未有的掌控感。

仿佛直到这一刻，他才第一次真正认识了自己所拥有的力量，以及这力量所能爆发的可能性。

“此招……”他嘴唇微动，声音不高，却带着斩钉截铁的确信，在雨后寂静的山崖前清晰回荡，“便叫‘弥焰斩’！”

眼中，光芒大盛，如同点燃了两簇火焰。这是他白言冰，脱离师尊的具体指导，完全依靠自身对天地的观察、对力量的思考、对心志的磨砺，感悟、孕育并最终创造出的、独属于他自己的剑意雏形！

这标志着他开始摆脱单纯的模仿与遵循，尝试着去理解力量的本源，并将自己的道融入其中，真正开始走出属于自己的路！

“言冰！你……你成功了？！”陈千羽早已跟了出来，一直屏息凝神地看着。此刻见到那岩石的惨状，感受到方才那一剑中令她都心悸的狂暴气息，脸上顿时绽放出惊喜交加的笑容。她快步上前，眼中满是激动与自豪。

这时，另一侧的石屋门也被推开，叶牧谦踱步走了出来。

“师尊！”白言冰难掩兴奋，上前几步，恭恭敬敬地行了一礼，随即迫不及待地将自己方才如何观雷看火、心有所感、思绪碰撞、最终于雨停刹那福至心灵、创出此招的过程，简明扼要却又难抑激动地叙述了一遍。尤其强调了那雷与火意象的融合，以及最终化入剑意、倾力一斩的感悟。

叶牧谦实际上刚才早都在屋里看完了全程；此时静静听完白言冰再次汇报，微微颔首。

“不错。”他开口，声音平稳，“能于寻常自然现象中有所悟，并创出契合自身心境与元气特性的招式，可见你平日修行并未懈怠，且用心体悟了。此‘弥焰斩’，集结了你心中所理解的雷火刚烈意象，将全身元气凝于一击，威力确实不俗。以此招为压箱底的杀招，猝然使出，同阶修士之中，能正面硬接者，恐怕寥寥无几。”

白言冰闻言，心中振奋，脸上喜色更浓。

然而，叶牧谦话锋一转，语气虽然未变，却自然而然地带上了一种令人心头发紧的严肃：“然而，此招缺陷，亦如其优点一般明显。刚猛有余，而变化与灵巧不足。你将全身元气、心神意志尽数倾注于一斩之中，气势固然一往无前，无坚不摧，但这同时也意味着，一旦出手，便再无后招变化之余地，自身破绽亦会因全力施为而显露。”

“若对手早有防备，或以巧妙身法游走闪避，或擅于以柔克刚、卸力导引之法……你这一剑斩空，力道用尽，新力未生之际，便是你周身破绽最大、最为脆弱之时，极易被人所乘。一招失利，满盘皆输。”

叶牧谦的声音不高，却仿佛带着千钧重量，敲打在白言冰心头：“切记，招式是死的，人是活的。拘泥于一招一式的强弱，乃是下乘。真正的强者，不在于某一招有多大的威力，而在于能否审时度势，在任何复杂、不利的情况下，依旧能掌控战局的主动，迫使对手按照你的节奏而行。刚柔并济，虚实相生，方是长久之道。”

这番话，如同兜头浇下的一盆冰水，瞬间让白言冰因创招成功而沸腾的血液和兴奋的情绪冷却下来。他仔细回味师尊的每一个字，再结合自己施展“弥焰斩”时的真实感受……确实如此！

那一剑斩出，全身力量随之宣泄，有种酣畅淋漓的快感。但快感过后，确实有一种短暂的空虚和停滞：若那时有敌人从侧面或背后袭来，自己恐怕真的难以迅速变招应对。此招威力虽大，但对元气和精神的消耗也极其巨大，无法频繁使用；而且其攻击轨迹直接明确，一旦被熟悉，就极其容易被预判和针对。

冷汗，悄无声息地浸湿了他的后背。方才的得意，化为了深深的警醒与惭愧。

他深吸一口雨后清冷的空气，压下心中的波澜，朝着叶牧谦深深躬身，态度诚恳无比：“师尊教诲，字字珠玑，弟子受教了！定当谨记于心，日后修行，必不忘完善此招，更会努力参悟刚柔变化之道，绝不因一招之得而自满。”

叶牧谦见他听进去了，态度也端正，便不再多言，只是淡淡点了点头；目光却投向青石城的方向，隐隐有一丝忧虑。

白言冰悟出新招是好事，但这孩子骨子里的侠义与刚直，在这纷乱的环境下，恐怕会惹来麻烦。

\vspace{2em}

果然，数日之后，青云门估计叶牧谦等碍眼的钉子离开了，于是雷厉风行地彻底将三成矿税的新规落实下去。

巡查的执事弟子增加了数倍，对各个矿坑入口的把守更加严密；抽税更是分毫不让，甚至变本加厉，有时还会以“矿石品质不纯”、“分量不足”等借口克扣。散修们怨声载道，却敢怒不敢言，生活越发困顿。

白言冰与陈千羽不时下山采购物资，总能听到散修们的哀叹，看到他们面黄肌瘦、愁云惨淡的模样，甚至有些拖家带口的散修，连孩子的吃食都成了问题。

这一日，白言冰在城中亲眼见到一名相识的、曾接受过陈千羽医治的散修，因为实在凑不齐税款，被青云门执事当众羞辱，最后连采矿工具都被夺走抵债，那散修跪地痛哭，状极凄惨。白言冰拳头捏得咯咯作响，胸中一股郁愤之气难以平息。

回到山间石屋，他辗转反侧，眼前总是浮现那些散修绝望的眼神。他想起自己身为皇子时，虽从不少书籍、以及老师的讲述中知道了“民间疾苦”这个抽象的概念，但从未如此近距离、如此真切地感受过这种被强大势力肆意盘剥的、极为具体且实际的无力与痛苦。

他又想起师尊所说的“与天奋斗，与地奋斗，与人奋斗”，想起自己创出“弥焰斩”时那股欲要劈开一切阻碍的刚猛心意。

“不行！不能就这么看着！”白言冰猛地坐起，眼中闪过决断的光芒。他并非鲁莽之辈，知道正面抗衡青云门不智。

但……若只是给青云门制造一些麻烦，让他们感到肉痛，同时接济一下最困难的散修呢？

一个大胆的计划在他心中成形。

接下来几日，白言冰借口下山采购，实则暗中联系了数位他曾帮助过、或看起来胆大心细、且对青云门极度不满的散修。他隐晦地透露，若有办法让青云门的灵石损失一些，或许能让他们投鼠忌器，缓和对散修的压榨。同时，若能将这些损失的灵石，用来接济那些快要活不下去的散修家庭……

这个提议极具诱惑力，也充满了风险。但被逼到绝境的散修中，终究有敢于铤而走险之人。经过小心试探和筛选，白言冰最终聚集了七八个可靠的、身手也算不错的散修。

他利用自己远超常人的目力和身法，花了几天时间，暗中摸清了青云门从几个主要矿区向城内仓库运输灵石的路线和时间规律。这些运输队通常由一名执事弟子带领，数名外门弟子押运，乘坐特制的、贴有轻身符和防护符的马车，走的也是相对固定的山道。

白言冰选定了其中一条路线上一段两侧树木茂密、便于埋伏且前后视野受限的弯道作为动手地点。行动时间，则定在了一次较大规模运输的傍晚时分，天色将暗未暗，正是人困马乏、警惕性可能降低的时候。

行动当日，白言冰让其他散修提前分散潜入埋伏区域，藏身于灌木丛或大树之后，约定以他的啸声为号。他自己则潜伏在更靠近来路的方向，负责观察和预警。

等待的时间格外漫长。夕阳的余晖透过树叶缝隙，洒下斑驳的光影。山林中只有虫鸣鸟叫，偶尔有小型野兽跑过的窸窣声。参与行动的散修们个个紧张得手心冒汗，紧握着手里的武器。

终于，远处传来了车轮碾过路面的声音，以及隐约的说话声。

白言冰屏住呼吸，悄然探头望去。

只见三辆符纹闪烁的马车，在十余名青云门弟子的簇拥下，正沿着山路缓缓而来。为首是一名执事弟子，修为约在玄脉境开脉初期，正与旁边一人说笑，显得颇为放松。后面马车旁跟随的弟子，大多只有灵枢境阶段的修为，警惕性也不高，白言冰自忖他一个能打十个这样的家伙。

“就是现在！”白言冰眼中精光一闪，深吸一口气，发出一声清越的长啸！

啸音未落，白言冰已然如蓄满力量的强弓射出的羽箭，自一丛茂密的灌木后电射而出，撕裂空气，目标明确至极——正是那运输队前首的青云门执事弟子！

几乎同一刹那，两侧林木阴影里，七八条身影应声跃出，口中发出混杂着紧张与决绝的呐喊，挥舞着矿镐、柴刀等简陋武器，扑向马车周围那些尚未完全反应过来的外门弟子。

“敌袭！”为首的执事弟子脸色骤变，厉喝之声带着一丝惊怒。他倒是反应不慢，腰间长剑“锃”一声出鞘，剑尖绽出一缕微弱的青芒，不退反进，迎向那道疾扑而来的身影。

在他看来，擒贼先擒王，只要迅速拿下这名显然是首领的袭击者，余下乌合之众自会溃散。

然而，他明显低估了这位“匪首”的实力。

白言冰甚至未曾去拔背负的长剑。

面对那带着破风声刺来的剑锋，他前冲之势毫无滞碍，右掌竖起，体内浑厚元气瞬间奔涌灌注，整只手掌边缘竟隐隐泛起一层金红色泽，仿佛烙铁边缘。

随即是一记看似直来直去、毫无花哨的掌刀，不偏不倚，迎着剑锋中线硬劈而去！

“铛——！！”

一声远比金铁交击更为沉闷的巨响炸开！

那执事弟子只觉自己仿佛一剑砍在了一座倾轧而来的铜钟之上，一股恐怖的力道沿着剑身狠狠撞来。

他持剑的右臂剧震，虎口处传来皮开肉绽的刺痛，五指瞬间失去了所有力量。那柄品质还算不错的长剑，竟如同脆弱的枯枝般脱手激射而出，“夺”地一声深深钉入侧旁一株老树的树干。

骇然之色刚爬上他的脸庞，视野已被一只手掌填满。白言冰的掌缘，在劈飞长剑后，去势几乎毫无衰减，劈在了他的胸口。

“噗！”

执事弟子双眼猛地凸出，胸膛肉眼可见地凹陷半分，整个人宛如被巨锤抡中，口喷鲜血，倒飞两三丈远，重重撞在一棵两人合抱的老松树干上。

松针簌簌而落，他身体贴着树干软软滑落在地，已然昏死过去，面如金纸，气息微弱。

一击废掉最强对手，白言冰身形没有丝毫停顿，足尖在地上一点，便如猛虎闯入了受惊的羊群，扑入那些惊慌失措、尚未组成有效阵型的外门弟子之中。

他下手极有分寸，不求击杀，仅施展掌风、拳劲或巧妙的擒拿手法。

这些大多仅有灵枢境感应、凝枢期修为的外门弟子，平日最多与同门切磋，何曾见过如此迅疾如电、狠辣精准的近身搏杀之术？往往眼前一花，兵器脱手，随即关节剧痛酸麻，便已踉跄倒地，失去了反抗能力。

不过短短十几次呼吸之间，连同车夫在内的十余名青云门弟子，已然呻吟着躺倒一地。

另一边，那七八名散修虽然个人修为平平，实战经验也参差不齐，但仗着人多势众和出其不意的先手，加上白言冰以雷霆之势解决了最具威胁的执事弟子，也很快控制了局面。他们将剩余几名试图反抗或逃跑的弟子打翻，用事先备好的结实麻绳，将其牢牢捆缚起来。

整个过程不过短短一盏茶的时间，干净利落；除了最初那声金铁巨响和零星的惨呼，再无人能发出有效的求救信号。

暮色渐浓的山道上，只剩下散修们粗重的喘息和地上俘虏们压抑的痛哼。

“快！搬东西！”白言冰低喝一声，声音在寂静下来的林间显得格外清晰。

他率先冲向那三辆符纹微光尚未完全熄灭的马车。另有两名散修心领神会，持着简陋武器警惕地注视着道路两头和山林深处。

掀开覆盖车斗的厚重防雨油布，里面整整齐齐码放着一个个鼓鼓囊囊的麻袋。白言冰用剑挑开一个口子，顿时，柔和而纯净的乳白色灵光流淌出来，映亮了他半张脸。

麻袋里，全是经过初步切割、尚未进一步精炼提纯的粗灵石，块头不小，灵气充盈。粗略估计，这三车加起来，数量相当可观，远超寻常散修数年辛苦所获。

“这么多……！”

“老天爷，这得有多少……”

“快，快搬！”

散修们看到这堆积如山的灵石，眼睛瞬间都直了，随即，难以抑制的狂喜和激动涌上心头，爆发出低低的、却充满力量的欢呼。

但他们还记得白言冰事先三令五申的交代，强行压下兴奋，迅速而有序地行动起来。

几人跳上车斗，将沉重的麻袋拖到车边，车下的人接力，或扛或抬，将灵石转移至准备好的大布袋或储物袋中。

随即，按计划撤离。

众人搬完灵石，背起沉甸甸的灵石袋，两人或三人一组，按照预先规划的、截然不同的几条隐蔽小路，迅速消失在暮色愈发深沉的山林之中。

临走前，他们甚至还不忘将地上那些被捆缚的青云门弟子拖到路旁更深的草丛或灌木丛里，塞紧嘴巴，确保其短时间内无法挣脱或呼救。

半个时辰后，在距离青石城三十余里外、一处隐秘岩峰下的天然山洞里，参与行动的散修们陆续抵达，再次聚首。

洞口被藤蔓巧妙遮掩，洞内燃着小小的一簇篝火，光线昏暗，却足够照亮地上那堆几乎成了小山的灵石麻袋。

柔和的灵光交织汇聚，将每个人激动、兴奋又略带紧张的脸庞映照得明暗不定。

白言冰没有耽搁，按照事先的承诺，开始分配。他先请那两位最为稳重、在散修中也有威望的头领出面，一起清点总数。

然后，他将灵石大致均分为两份。

一份，按照此次行动中各人出力多寡、承担风险大小，并结合各自家中实际的困难程度——是否有人卧病、断粮几日、有几口人待哺——进行分配。虽做不到绝对公平，但力求让最紧迫的人得到更多救急之资；

另一份数量稍多的，他则郑重地交到那两位头领手中，沉声嘱咐：“这些，务必小心再小心，暗中分发给那些家中已断粮数日、或有人重病无钱医治、真正活不下去的乡亲。切记，一定要隐蔽，绝不可暴露来源，就说是不知哪位路过的侠客暗中留下的。”

一位头发已花白的年长散修，颤抖着双手接过自己分得的那一袋灵石，紧紧抱在怀里，仿佛抱着初生的孙儿。

他眼眶瞬间通红，对着白言冰，竟是屈膝就要跪拜：“白少侠……大恩不言谢！大恩不言谢啊！有了这些，我那躺在破炕上咳血的老妻……那几个饿得整天喝水充饥的娃……总算……总算能有点活路了……”声音哽咽，老泪纵横。

旁边一位汉子，手臂上还带着采矿留下的旧伤疤，也声音沙哑道：“白少侠，俺不会说话，但这条命，以后就是您用得着的时候，吭一声！”

其他人也纷纷围拢，感激之情溢于言表，山洞内弥漫着一种悲喜交加的气氛。

白言冰连忙上前，双手用力扶住那年长散修，不让他跪下去，正色道：“诸位乡亲父老，万万不可如此！此举不过是暂解燃眉之急，杯水车薪，绝非长久之计。”

他环视一圈，火光在他眼中跳动，语气变得严峻，“青云门丢了这么多灵石，如同被砍了一刀，必定剧痛。他们绝不会善罢甘休，接下来必定疯狂戒备，严查追捕，甚至可能报复。

“大家回去后，切记要小心谨慎，近期尽量不要再去矿洞，暂避风头，低调行事，权当不知此事。”

众人心中一凛，狂喜稍退，纷纷应诺。

散修们对着白言冰再次深深作揖拜谢，这才各自背起或抱起分得的灵石，怀揣着绝处逢生的希望与对未来的深深忐忑，依次悄然没入洞外的漆黑山野，如同水滴回归溪流，悄无声息。

白言冰是最后一个离开山洞的。他仔细清理了篝火痕迹，又故意在洞内留下几处迷惑性的足迹，方才施展身法，在夜色的掩护下，绕了数个圈子，朝着山间石屋的方向疾行。

\vspace{2em}

回到石屋时，已是月上中天，万籁俱寂。

陈千羽竟还未睡，听到极其细微的推门声，她像受惊的小鹿般猛地坐起身，看到是白言冰，紧绷的肩膀才骤然松弛下来，长长舒了口气。

但随即，她鼻子微微抽动，闻到了白言冰身上那极淡却无法忽视的血腥气，再借着月光看到他眉宇间那难以掩饰的疲惫，秀气的眉头又立刻蹙了起来。

屋内，油灯如豆。叶牧谦亦未安寝，瞅了白言冰一眼。

“说说吧，身上血腥气是怎么回事？”

“师尊……” 白言冰知道，以师尊的修为和见识，自己今夜的外出和此刻的状态，肯定无法隐瞒。

关上屋门，走到叶牧谦面前，垂手而立，将今日之事——从如何联系散修、如何勘察路线、到伏击过程、分配灵石——毫无保留，和盘托出。

叶牧谦静静地听着，手指无意识地在陈旧的桌面上轻轻敲击，直至白言冰话音落下许久，山洞内只剩下油灯灯花偶尔爆开的噼啪轻响。

他发出一声几不可闻的轻叹：“言冰，你心存侠义，怜悯弱者，见不得人间疾苦，此心赤诚，可嘉可许。劫掠不义之财以济困苦，急公好义，古之侠客，亦常有为之。”

他话锋随即一转，语气虽依旧平缓，却仿佛重了几分，带着山雨欲来的凝重：“然而此举看似痛快解气，实则后患无穷！

“你劫掠的，是青云门维系一宗、掌控一城的命根子——灵石！这无异于虎口拔牙。青云门能在此地盘踞多年，将一城散修视为佃户矿奴，岂是心慈手软、易于相与之辈？他们很快便会察觉运输队失踪，并且一定会不惜代价，追查到底。”

人心难测，境遇易迁。难保不会有人酒后失言，或承受不住压力，或干脆就被青云门以更残忍的手段撬开嘴巴。青云门只需稍加排查，将近期与他们有过明显冲突、且具备相应实力和动机的外来者列为嫌疑，便不难将目光锁定到他们头上。

他转过身，目光如古井深潭，映着跳动的灯火，笔直地看向白言冰：“届时，他们便可以以维护青石城秩序、追缴被劫宗门资源之名，光明正大地围剿清除。你今日所为，实则是亲手将一把最锋利的刀，递到了他们手中，给了他们一个我们难以辩驳的理由！”

“我知你心中不平，血气方刚，见不得不公。但行事处世，尤其是与一方势力周旋对抗，绝不能仅凭一腔热血、头脑一热便贸然动作。必须谋定而后动，走一步，看两步；尤其是要反复掂量，那最坏的结果，我们是否承受得起。”

叶牧谦走回桌边，手指点了点桌面，“接下来，风声鹤唳，务必慎之又慎。减少一切不必要的外出，提高警惕。”

白言冰听着师尊冷静甚至有些冷酷的分析，初始行动成功带来的那点喜悦和成就感，此刻早已荡然无存，如同被一盆冰水从头浇下，彻骨生寒，取而代之的是翻涌而来的后怕和深深的忧虑。

他意识到，自己自以为隐秘、豪情万丈的“侠义”，或许真的将师尊和千羽，置于了比以往任何时候都更直接、更危险的境地。青云门的报复，绝不会仅仅停留在搜寻劫匪的层面。

“弟子……思虑不周，鲁莽行事，知错了。” 白言冰低下头，冷汗不知何时已沁湿了内衫，背心一片冰凉。

“错已铸成，悔之无益。世间万般事，往往做了便难回头。姑且把此事当成一次教训吧。”叶牧谦摆了摆手，语气略缓，却也带着几分无奈，“如今之计，唯有静观其变。且看青云门，究竟会作何反应吧。但愿他们的动作不会太快。”

\vspace{2em}

青云门果然雷霆震怒、全城大索。然而，真正的压力，却以一种更阴险、更令人窒息的方式，降临到了所有散修头上。

很快，一道新的宗门告示贴满了青石城内外各处的公告栏：为“补偿前次运输损失”、“切实加强各矿洞防护措施，保障采矿安全”，即日起，青云门对所有散修开采灵石的矿税，由原先的三成，直接提升至令人绝望的四成！

同时，告示还宣布将加派“巡逻队”，保障运输路线安全。

暗地里，超过百名修为精湛、经验丰富的内门精锐弟子，脱下标志性的青云服饰，换上普通散修或行商的粗布衣裳，伪装成各式身份，如同无形的罗网，悄无声息地撒向了西山到青石城之间的数条主要通道，以及周边可能藏匿的山林区域。

他们张网以待，日夜潜伏，耐心得可怕。

这一手，堪称毒辣无比。这是将宗门损失堂而皇之地转嫁到所有散修头上，进一步榨取其本就微薄的剩余价值，逼迫得更多人走投无路；暗地里，则是布下致命的陷阱，等待那胆大包天的“劫匪”再次出手，自投罗网。

沉重的税负如同最后一根稻草，压垮了许多散修家庭本就脆弱的生计。城郊棚户区里，绝望的叹息与压抑的哭泣日夜可闻，而不敢宣之于口的愤怒，则如同地底奔流的岩浆，在散修群体中无声却炽烈地蔓延。

白言冰通过陈千羽偶尔下山换取生活必需品时带回的消息，得知了税负再增的噩耗。他又曾数次暗中靠近散修聚居区域，远远望见那些愈加破败的棚屋、听到孩童有气无力的啼哭、看到面黄肌瘦的妇人眼中熄灭的光。

当日他分配灵石时，那些散修感激涕零的面庞与眼前更甚的凄惨景象交织重叠，他胸中那团因不平而生的侠义之火，与因师尊警告而生的愧疚之感，如同两种不同的火焰，交织燃烧，灼得他五脏六腑都不得安宁，难以平息。

尽管叶牧谦早已严令禁止他再有任何轻举妄动，反复告诫此刻一动不如一静。但数位与他相熟、曾受他接济或参与过上次行动的散修头领，还是设法避人耳目，悄悄寻到了石屋附近，留下了隐秘的联络标记。

白言冰终究按捺不住，在一次借口查探周围环境的独自外出时，与他们在一处极为隐蔽的山坳里见了面。来人不过四五位，个个形容憔悴，眼中布满血丝，其中两人身上还带着新添的伤痕，说是前几日试图偷偷多藏几块灵石，被青云门监工发现毒打所致。

他们围着白言冰，言辞恳切，甚至带着走投无路的哭腔，诉说四成重税之下，家中几已无米下锅，眼看家破人亡就在眼前。

其中一位上次参与行动、分得灵石救了急的汉子，更是左右确认无人后，凑近白言冰，压低声音，带着一种孤注一掷的兴奋，提供了一个绝对可靠的内线消息：

明日午后，会有一批规模远超上次、从西山最大矿点产出的优质灵石运出；因为近期青云门力量被抽调四处巡逻，导致这支运输队的护卫人数不足，仅有十人左右；而且带队者似乎只是个普通执事。路线，依旧会经过西山道那段他们熟悉的、便于埋伏的弯道。

“白少侠，机会千载难逢啊！上次咱们干得漂亮，青云门肯定以为咱们得了甜头就躲起来了，这次他们人手不足，正是再干一票大的好时机！有了这批灵石，至少能救活上百户人家！” 那汉子眼睛发红，呼吸急促。

白言冰心中顿时天人交战。师尊冷峻的警告言犹在耳，李青阳那双仿佛能洞察人心的眼睛和“剿匪”二字，偶尔还会闯入他的梦境。理智告诉他，这过于“凑巧”的消息很可能是个诱饵，青云门吃了那么大的亏，怎么可能反而在最重要的运输上松懈？

然而，眼前几位头领绝望中透出最后一丝希冀的眼神，他们口中那些即将饿毙的老人孩童的惨状，以及上次成功带来的某种惯性自信，如同无形的浪潮，一波波冲击着他理智的堤坝。

他仔细推敲着那个消息：青云门因上次被劫，兵力被分散到多条路线布控，导致某一次运输护卫薄弱，也并非完全没有可能……

一种“再干最后一次，就能解决更多人的困境，然后立刻远遁”的侥幸心理，混杂着对自身实力的某种过分自信，如同蔓草滋生，最终渐渐压倒了心头那越来越强烈的风险警示。

他瞒着叶牧谦与陈千羽，暗中与那几位头领约定，再次集结人手，于次日午后，仍在西山道那片他们曾成功过的弯道设伏。

这一次，或许是重税逼迫太甚，或许是上次的成功带来了虚幻的信心，响应者更多，足足聚集了近二十名散修；这些衣衫褴褛、眼中泛着拼命红光的汉子与劳动妇女，将这次行动视为了全家老小最后的生机。

次日午后，阳光被西山茂密交错的树冠切割得支离破碎，在林间空地上投下无数晃动的光斑。白言冰带着比上一次复杂十倍的心情，潜伏在预定地点一块长满青苔的巨岩之后。

他强迫自己冷静，仔细观察着周围环境：风吹过林梢的呜咽，远处隐约的溪流声，草丛间偶尔惊起的飞虫……一切似乎与往日并无不同。

但一种莫名的、细微的违和感，却如同冰冷的蛛丝，悄然缠绕上他的心头。

太安静了。

似乎，那种属于山林应有的、充满生机的嘈杂变得稀薄了。鸟雀的鸣叫零星而谨慎，兽类活动的窸窣声近乎绝迹，连夏末本该喧闹的虫鸣，都显得有气无力。

白言冰的心一点点沉了下去，不安的预感如同冰水漫过脚踝。他几乎想要立刻发出信号取消行动。但就在此时，远处传来了清晰的车轮碾过碎石路面的声响，以及隐约的、显得有些松散的谈话声。

运输队如期而至。

依旧是三辆特制的马车，车厢上符纹暗淡，似乎连轻身符都未全力激发。护卫在马车周围的青云门弟子，粗粗一看，果然只有十人左右，比上次还少。

他们大多神情懒散，步伐随意，队形松散，与上次那虽然松懈但至少队列整齐的押运队伍相比，显得格外业余。一切都显得那么顺利，那么合乎预期，仿佛青云门真的因为兵力捉襟见肘，或者骄傲大意，而露出了一个致命的破绽。

白言冰死死盯着那支越来越近的队伍，目光如鹰隼般扫过每一个护卫的脸、他们的手、他们的步伐细节。那看似懒散的外表下，似乎又隐隐透着某种难以言喻的协调？他握紧了拳，掌心微微汗湿。

不能再犹豫了。身后，二十双饱含绝望与渴望的眼睛在盯着他，二十个家庭的渺茫希望系于此举。那可靠的消息，眼前确凿的薄弱护卫，最终化为了一股压垮疑虑的冲动。

他猛地吸足一口气，胸腔鼓起，发出一声短促而尖锐的呼哨——动手的信号！

信号发出，他率先从巨岩后暴起，身法催到极致，化作一道模糊的灰影，依旧直扑车队前方那看似领头之人。

两侧山林中，近二十名散修如同被困已久的野兽，发出混杂着恐惧与疯狂的呐喊，挥舞着各式武器，呐喊着冲了出来，扑向马车和那些护卫。

然而，就在散修们冲入山道，距离车队尚有十余步之时，异变陡生！

那十名看似懒散松垮的护卫弟子，脸上的困倦与散漫在百分之一息内褪得干干净净，眼神瞬间变得锐利如捕食前的鹰隼，精光四射。

他们没有后退，没有慌乱，反而以惊人的默契和速度倏然散开，脚步错落间，已然结成了一个看似简单、实则彼此呼应、攻守兼备的小型圆阵，将三辆马车死死护在中心，动作整齐划一，分明是训练有素的精锐！

白言冰的瞳孔骤然缩成了针尖大小，一股冰冷的寒意从尾椎骨瞬间窜上天灵盖，浑身血液几乎冻结。所有的不安、所有的违和感，在此刻化为了铁一般的现实，狠狠砸在他的心头。

他脑海里一片空白，唯有那无数次在战场故事、在江湖传说中听过的、充满苦涩与绝境的两个字，冲口而出：

“中计！”

似乎是印证他的话一样，一道尖锐刺耳的啸音撕裂了山林的死寂。

只见一名青云门护卫，手中不知何时汇聚起一股灵力；随即，一道赤红中夹杂着青烟的焰火疾射冲天，在午后灰蒙蒙的天空中猛地炸开，化作一团耀目又狰狞的青云图案，久久不散，方圆数十里清晰可见。

这突如其来的信号，撕下了最后一层伪装。

“嗖嗖嗖——！”

密密麻麻的破风声从四面八方响起，快得让人头皮发麻！

数十道身影闪现而出，动作整齐划一，带着经年训练的冷酷效率；瞬息之间，已在外围那十名“诱饵”弟子组成的防御圈之外，形成了一个更大、更厚实的死亡包围圈，将白言冰连同那二十余名已冲入山道的散修，彻底锁死在当中！

这些人一色的青云纹劲装，气息凝练精悍，目光如刀，最弱者周身灵力波动也已达灵枢境凝枢巅峰，更不乏玄脉境开脉期、通络期的好手。

而为首三人，仅仅站在那里，散发出的灵压就如同一座座无形山岳，沉甸甸地压在每一个散修的心头，令人呼吸滞涩，心生绝望。

正中一人，面容刚毅，三缕长髯，正是青云门门主李青阳。他左侧是满脸怨毒与得意、死死盯着白言冰的李轩。右侧，则是一位手持乌沉蛇头拐杖、面容阴鸷枯瘦的老者，一双三角眼开合间精光闪烁，正是青云门二长老（玄脉境融贯期）。

在他们身后，超过五十名精锐内门弟子刀剑齐出，寒光映日，更有不少人指间夹着符箓，灵力引而不发，森然杀气如实质般弥漫开来，将这片山道变成了冰冷的屠宰场。

“哈哈哈！果然等到了你们这群不知死活的泥腿子！”李轩放声狂笑，笑声中充满了精心布局终于收网的快意和积压已久的恨意。

他的目光如同毒钩，死死锁住人群中心的白言冰，“姓白的！还有你们这些吃里扒外、忘恩负义的狗东西！今日，天罗地网已布，我青云门便在此剿匪，将你们这些祸乱青石城秩序、劫掠宗门资源的败类，一网打尽，以正视听！”

“剿匪”二字，被他咬得又重又响，在山谷间回荡，显然是早已备好、用来堵天下悠悠之口的冠冕堂皇名分。

刚刚还为“薄弱”护卫而红着眼冲出来的散修们，此刻如遭冰水淋头，满腔的疯狂希冀瞬间被无边的恐惧碾碎。面如土色，肝胆俱寒，不少人手脚发软，手中简陋的武器都在微微颤抖。

门主亲至！长老压阵！数十精锐环伺！这哪里是疏忽大意，分明是精心策划、请君入瓮的绝杀之局！绝望如同瘟疫，在每一个人眼中急速蔓延开来。

白言冰心头猛地一沉，仿佛坠入了无底寒渊，最坏的预感，以最残酷的方式成为了现实。他暗骂自己愚蠢至极，竟然被那廉价的同情和侥幸心理蒙蔽，一头撞进了敌人张开的血盆大口。

但事已至此，退路已绝，山林两侧的每一个空隙都闪烁着敌人的兵刃寒光。他猛地一咬舌尖，剧烈的刺痛和血腥味强行压下了翻腾的慌乱与几乎将他淹没的愧疚——是他，将这群走投无路的乡亲带入了死地。

他深吸一口混杂着泥土和杀气的气息，一步踏前，挺直了脊梁，用自己不算宽阔的身躯，将身后那些惊慌失措、面无人色的散修隐隐挡在后方。

目光越过叫嚣的李轩，直接迎向那深不可测的李青阳，声音冷冽如冰泉撞击山岩：“李门主，真是好大的手笔，好精妙的算计。劫掠灵石之事，皆由我白言冰一人策划主导，与这些道友无干。他们不过是为生计所迫，受我一时蛊惑。有何手段，冲我一人来便是，何必牵连这些苦命人？”

“冲你来？哼，天真！”李轩狞笑着打断，剑尖遥指，“你以为演一出义薄云天的戏码，就能救了这帮同伙？今日这西山道，便是你们的葬身之地！一个都休想走脱！父亲，二长老，证据确凿，匪类齐聚，请速下令，剿灭这群宗门败类，夺回被劫资源！”

李青阳的目光如同实质的闪电，在白言冰身上一扫而过，随即投向其身后混乱的散修群，以及更远处的山林缝隙，似乎在搜寻着某个预料中该出现的身影。他忌惮的，自然是至今未曾现身、却让他儿子吃了暗亏、让他隐隐感到不安的那个神秘师尊——叶牧谦。

他沉声开口，带着一种居高临下的审度：“白小友，年少气盛，误入歧途，情有可原。若你师徒此刻肯束手就擒，交出所有劫掠之物，并供出其他潜伏同党，我青云门念在你们初犯，或可网开一面，从轻发落。”

“无需多言！”白言冰断然回绝，心知此刻任何犹豫或妥协，都会将师尊和千羽也拖入这必杀之局。

“道不同不相为谋！要战，便战！”

“锃——！”

清越剑鸣冲天而起！白言冰体内元气再无丝毫保留，轰然爆发！

淡金色的气劲从他周身毛孔喷薄而出，隐隐带起风雷低啸。身后散修们顿觉周身一轻，那令人窒息的压迫感稍减，虽然恐惧未消，但看向前方那并不高大却挺直如枪的背影时，眼中终究燃起了一缕微弱的、拼死一搏的光。

“冥顽不灵，自寻死路！杀！”李青阳眼中最后一点耐心消失殆尽，寒光爆射，终于挥下了象征杀戮的手臂。

“青云剑阵，进！”二长老沙哑干涩的声音紧接着响起，手中那乌沉沉的蛇头拐杖重重顿在地面。

“咚！”

一声闷响，一圈土黄色的灵光以其拐杖落点为中心，如同水波般急速扩散开来，瞬间掠过所有青云门弟子脚下。被这灵光掠过的弟子，气息似乎隐隐连成了一体，步伐更显协调。与此同时，早已按捺不住的青云门精锐齐声呐喊，声震山林：

“杀——！”

各色灵光骤然炽盛！剑光如匹练纵横，火球、风刃、冰刺等低阶术法符箓的光芒交织成网，更有不少辅助符箓被激发，如同狂风暴雨，朝着包围圈中心那二十余名惊恐的散修和白言冰倾泻而下！攻势未至，那凌厉的杀意和磅礴的灵力乱流已经让不少散修气血翻腾，几欲呕吐。

而李青阳、李轩、二长老，则如同嗅到血腥的鬣狗，目标明确，直扑向孤身立于阵前的白言冰！

他们要第一时间绞杀这个最具威胁的“匪首”！

战斗，在信号焰火尚未完全消散的天空下，猝然爆发，且从一开始就呈现出令人绝望的一面倒碾压！

散修们发出野兽般的嘶吼，挥动着矿镐柴刀，拼命格挡、躲闪。然而，修为的鸿沟、装备的云泥之别、战斗素养的天差地远，在此刻显露无遗。青云门弟子的剑阵进退有据，相互掩护，散修们则是各自为战，乱成一团。甫一接触，便有数声凄厉的惨叫响起！

“啊——我的腿！”

“救命！”

血光迸现！

白言冰目眦欲裂，这些鲜活的、不久前还带着希冀眼神的面孔，此刻在他眼前流血、惨叫、倒下，都是因他之故！

悔恨与暴怒如同毒焰灼烧着他的心脏。

他长剑舞动如风车，淡金色剑气纵横激荡，将射向自己的大部分剑光法术凌空击碎，同时朝着身后厉声嘶吼：“不要乱！向我靠拢！背靠背，结圆阵自保！”

然而，崩溃的士气与悬殊的实力下，他的呼喊收效甚微。散修们已被分割，惊恐淹没了最后一丝组织。

就在这防线即将全面崩溃的危急关头——

“住手！”

一声断喝，带着一种奇异的穿透力与震慑心魂的力量，径直在山道上空每一个人的耳中、心中炸响！

两道身影，以远超在场绝大多数人想象的速度，从山林另一侧的密林深处疾掠而来，前一后，如同两道撕裂暮色的流光！正是察觉不对、全力赶来的叶牧谦与陈千羽！

叶牧谦面沉似水，平日里那份淡泊温和早已消失不见。他终究还是放心不下，白言冰自以为隐秘的行踪，如何能完全瞒过他的感知？察觉到那丝异常的不安后，他便带着陈千羽悄然尾随，此刻堪堪赶到这绝杀之地。

陈千羽更是心急如焚，俏脸煞白。人还未完全冲出山林，素手已然扬起，数十点细如牛毛、却闪烁着金黄色光芒的钢针，已如一群被激怒的蜂群，顺着她萧瑟秋风一般的元气波流，发出“咻咻”破空声，精准无比地罩向那些正在屠杀散修的青云门弟子！

她认穴极准，银针专打手腕、肘弯、膝侧等关节要穴以及灵力运转的关键窍门。钢针过处，中招的青云门弟子无不感到一阵酸麻刺痛顺着经脉蔓延，灵力运转顿时出现刹那的阻滞与不畅，原本凌厉的攻势不由自主地为之一缓。

这突如其来的精准远程打击，给即将崩断的弦松了一丝力，给那濒临湮灭的散修防线，赢得了一线极其宝贵、宛如溺水之人抓住稻草般的喘息之机！数名差点被斩杀的散修，连滚爬地逃回了稍微密集的人群中。

“师尊！千羽！”白言冰精神猛然一振，仿佛注入了强心剂；但随即，更深的愧疚如同潮水般将他淹没——终究还是连累了他们。

“叶牧谦！你果然来了！”李青阳眼中精光爆射，不惊反喜，仿佛等待已久的大鱼终于入网。

他等待的，正是这个让他儿子吃亏、让他心生忌惮的正主！

李青阳再无暇关注下方苦苦硬撑的白言冰和其他小蚂蚁，身形一晃，已如毫无重量般升至半空，强大灵压再无丝毫掩饰，轰然释放！周身青色道袍无风自动，隐隐有无数细小的、玄奥的青色风灵符文虚影流转闪烁，散发出切割一切的锋锐气息，牢牢锁定了地面上那道青衫身影。

“阁下教徒不严，纵徒行凶，劫掠我宗门资源，为祸一方！今日，你这做师尊的，也该给青云门、给青石城一个交代了！” 李青阳声如洪钟，占据着道德的制高点，杀意却凛然如实质。

叶牧谦深知此刻任何辩解都是徒劳，唯有用实力说话。他深吸一口气，强行压下经脉中因之前催动见独境而可能遗留的隐痛与反噬预兆，眼神在刹那间变得深邃无比，仿佛倒映着无尽星空与人心万象。

见独境之力，被他再次强行催动！

这一次，他将这份精神压制，主要倾泻向半空中气势汹汹的李青阳；同时余波亦如涟漪般扩散，稍许干扰着地面上二长老等青云门高层的心神。

李青阳正全力催动灵力，指尖青芒吞吐，一枚枚蕴含着狂暴风灵之力的法球正在迅速凝结，准备施展他最拿手、威力也最大的杀招。

然而，骤然间，他感到心头一跳，仿佛有什么东西无形中覆盖了他的灵觉。周遭熟悉的天地灵气似乎变得有些陌生和疏离，自身灵力的流转，竟隐隐受到一丝极其微妙的、源于意志层面的干扰，不如平时那般圆转如意。

这感觉玄奥诡秘，难以捉摸，却又真实不虚地影响着他，让他不得不分出一部分宝贵的心神与意志力，去稳固自身道心，对抗这股无形无质却又无处不在的侵蚀。

原本行云流水、几乎瞬间可成的风刃和法球，凝结速度明显迟滞，威力似乎也受到了无形但依旧巨大的削弱。

李青阳心中不禁骇然！上次听大长老与李轩描述，他只觉是某种偏门的精神干扰之术；此刻亲身体验，才知这种手段的诡异与难缠，竟能直接影响对手与天地灵气的共鸣及自身心神安定！

叶牧谦就此与李青阳堂堂正正地对攻起来。李青阳勉力施为，法球与风刃虽被削弱，但威力依然强悍，毫无保留地轰向叶牧谦；叶牧谦则利用见独领域与元气波动，尽力引偏李青阳的风刃轨迹，或是凭雄浑元气正面拍散其薄弱之处。

法球和风刃在半空中不断炸开，在高空形成一片无声的、极为凶险的风暴区。气机剧烈交锋，神识不断碰撞，逸散的劲风将下方树木的顶端枝叶切割得纷纷扬扬。

下方战场中的诸多散修和青云门弟子，却没有观赏两位高手过招的余裕。

叶牧谦成功牵制住了最强的李青阳，陈千羽的及时援手与精准干扰又缓解了最致命的压力，原本一面倒的屠杀局势竟然得到了遏制。

在她的急促指挥与辅助下，残余的十余名散修终于勉强依托几块嶙峋的山石，背靠背结成了一个简陋却颇为有效的圆形防御阵，嘶吼着，用身体和简陋武器，勉力抵挡着四周青云门弟子的凶狠攻势。

然而，压力最大的焦点，依旧集中在白言冰身上。李青阳被师尊勉强牵制住，李轩和二长老却彻底腾出手来，脸上带着狰狞的杀意与围猎的快感，将他视为了必须首先拔除的钉子，联手围攻而来。

“白言冰！今日便是你的死期！”李轩咬牙切齿，种种情感此刻全部化为了疯狂的杀意。他剑招愈发狠毒，与二长老那神出鬼没、专走下盘的蛇头拐杖形成了致命的夹击之势。

白言冰将体内元气催动到前所未有的极限，丹田气海翻腾，经脉隐隐胀痛。手中长剑嘶鸣，剑意挥洒间，竟隐隐有风雷之声响动，雷火之气在黑紫色剑气中明灭闪烁。

他已将弥焰斩的那份一往无前、爆裂焚灭的意境雏形，融入了寻常剑式之中；剑法刚猛绝伦，气势如虹，仿佛化身为一尊怒火战神。

在李轩再次攻向前来的时候，白言冰瞅准时机一剑全力劈出，带着灼热的气浪与凌厉的锐响，黑紫色的剑气一往无前地斩向李轩！

李轩心中惊骇：这一招，他一定挡不住！这个姓白的果然有点东西！

被迫立即变招，改突刺为闪避，试图在那灼热的气浪逼近他之前逃离之；二长老眼疾手快，又用灵力往旁边拽了他一把，李轩才堪堪从那剑气路径上毫无风度地逃开，衣服甚至都烧焦了些许。

随即，那雷火剑气毫无阻拦的撞上了远处的山崖。只听“轰隆隆”一阵巨响，直接受击处的山石竟直接熔化为岩浆，崖上碎石与泥土也被冲击的余威震得簌簌掉落下来，让李轩与二长老感到极为后怕。

然而，正如叶牧谦此前一眼看破并严肃指出的那般，他的招法刚猛有余，灵活不足，变化欠缺。

李轩毕竟出身正规宗门，见识和实战经验远非普通散修可比。初始的惊愕过后，他很快察觉到白言冰的剑招虽然威力奇大，令人心惊，但路数却相对单一，尤其是那让人心悸的恐怖一击，蓄势前兆颇为明显。

他与经验老辣的二长老交换了一个心照不宣的眼神，立刻改变了策略。

二人停止了与白言冰硬碰硬地对攻。

二长老冷笑一声，蛇头拐杖连连点地，土黄色灵光不断没入地面。霎时间，白言冰脚下坚硬的山石突然化为松软泥沼，步履维艰；侧方毫无征兆地刺出锋利的地刺，逼得他不得不闪避格挡；头顶偶尔会有沉重的土灵压力骤然落下，打断他的剑气连贯。

而二长老本人，则始终游走在数丈开外，拐杖尖端点出道道阴柔刁钻的土灵劲力，专攻白言冰的脚踝、腰眼、后心等难以兼顾之处，不断消耗其元气，扰乱其心神。

李轩则凭借更胜一筹的身法在外围游走，剑光如同草丛中伺机而动的毒蛇，专挑白言冰全力格挡地刺、或剑招用老、力道将尽未尽的瞬间，疾刺偷袭，一击即退，绝不贪功恋战。

如此一来，白言冰顿感吃力无比，陷入了无形的泥潭。他那擅长以力破巧、堂堂正正、一击决胜的打法，在对方这套滑不溜手、烦不胜烦的游斗缠扰策略下，犹如猛虎陷入网罟，又像重锤击打流水，空有磅礴力气，却难以落到实处，取得真正的战果。

他几次三番想强行凝聚元气，再一次施展出完整的弥焰斩，试图以绝强力量破开这令人憋闷的困局；但二长老老奸巨猾，对他那蓄力时明显的剑气萦绕异常敏感，总能提前一瞬以大面积的地陷泥沼或陡然升起的厚重土墙进行干扰破坏，打乱其节奏；李轩见到黑紫剑气萦绕时，更是立即远远避开其正面锋芒，绝不给他斩实的机会。

战斗持续，高强度的攻防与不断被打断的爆发尝试，让白言冰的元气以惊人的速度消耗。他的呼吸渐渐粗重如牛喘，额角鬓边渗出大颗大颗的汗珠，顺着脸颊滑落，与灰尘泥土混在一起；手中的长剑格挡李轩偷袭的剑光时，手臂也显出了些许沉重。

反观李轩与二长老，配合愈发默契，以逸待劳，凭借着丰富的经验和灵活的战术，逐渐掌控了战场的主动权，占据了绝对的上风。

白言冰的身上，开始不断增添一道道伤口：左肩被李轩的剑风划破，鲜血染红衣襟；小腿被突兀刺出的地刺擦伤，火辣辣地疼；后背也被二长老一道阴柔拐劲扫中，气血一阵翻腾。这些伤口都不算深，却不断消磨着他的体力、意志和所剩不多的元气。

“哈哈！撑不住了吧？野路子就是野路子，空有几斤蛮力，也配与我青云门干百年底蕴的正道传承抗衡？”李轩见状，心中快意无比，得意地大笑出声，剑势趁机更加紧了几分，如同疾风骤雨。

白言冰咬牙苦撑，牙龈几乎咬出血来，心中却不禁泛起一片苦涩。

师尊的评语言犹在耳，自己却终究因为招法的致命缺陷，陷入了这等绝境。

难道……今日真要因为自己的鲁莽和不足，败亡于此，还要连累师尊和千羽？

然而战斗中最忌讳分心。李轩抓住白言冰最脆弱的那一瞬，直接冲着他的胸膛刺去！

就在白言冰步伐凌乱，险象环生之际——

“恶徒看刀！”

一声清冽如冰泉击石的女子叱喝，毫无预兆地撕裂了山道上血腥而灼热的空气！

几乎与这声叱喝同步，一道狭长、凄冷、轨迹刁钻诡异到极致的雪亮刀光，自侧后方一片林木最浓密的阴影中暴射而出！

所有的杀意与锋锐都凝聚在那薄如蝉翼的刀锋之上，快得只在空中留下一道淡到极致的残影，其目标，正是李轩因狂攻而微微侧露、毫无防护的后心要害！

这一刀，对于时机的把握堪称毒辣，正是李轩心神与剑锋志得意满地完全锁定在白言冰身上，自以为胜券在握的绝杀时刻！

刀光乍现的刹那，一抹利落的黑色身影才从阴影中凸显。来人一身毫无多余装饰的紧身黑色劲装，完美融入昏暗的林影，脸上覆着同色面巾，只露出一双清澈明亮、此刻却凝结着凛冽煞气的眼眸，宛如寒星！

那声叱喝此时才完全灌入李轩耳中。

李轩浑身汗毛瞬间根根倒竖！致命的危机感如同冰锥刺入后脑，让他头皮发麻，战斗的本能疯狂催促他躲闪格挡。

没有丝毫犹豫，他强行中断了对白言冰的追击之势，体内灵力以前所未有的速度疯狂涌向双腿经脉，试图施展身法，向侧前方疾掠闪避。

然而，就在他心神激荡、灵力全力爆发冲过经脉关隘的这电光石火之间，异变突生！

那些因他长期急功近利、不加提纯便大量吞噬灵石以提升修为，从而沉积在经脉壁障深处的灵力杂质与虚浮淤塞，在这不容有失、需要经脉以最高效率运转的生死关头，竟猛地产生了阻滞与紊乱！

他感到几处关键窍穴的灵力流猛地一涩，像是奔涌的洪水突然撞上了无形的暗礁，原本流畅如水的身法灵力输送，出现了致命的刹那停顿与岔流！

“什么？！”李轩瞳孔骤然缩成了针尖，无边的恐惧瞬间淹没了他的意识。

叶牧谦当日那平淡却如预言般的话语，如同惊雷在他脑海炸响：“贪恋灵石，急功近利……经脉杂质满溢、灵力虚浮运转滞涩……长此以往，与人全力搏杀、灵力激荡之时，必出事故！”

报应，竟然在此刻，以最残酷、最致命的方式应验了！

高手相争，只争刹那。就是这因自身根基不稳而导致的、微不足道的半拍迟缓，在白言冰和那蒙面女侠这等对手面前，已足以决定生死，甚至比生死更残酷的结局！

“嗤——！！”

砉然向然，奏刀騞然。

那道刁钻狠辣的雪亮刀光，没有丝毫留情，更没有因为李轩那瞬间的凝滞而有半分偏差，精准、冷酷、决绝地斩落！

李轩只感到右肩连接处传来一阵难以形容的、先是冰凉而后是火山爆发般的剧痛，那痛楚如此猛烈，瞬间冲垮了他的所有感知。

他眼睁睁地，甚至带着一丝荒诞的茫然，看着自己那柄握剑的、引以为傲的右臂，齐着肩关节，被一股无可抗拒的力量带着，脱离了身体，伴随着一蓬凄艳滚烫的血雨，划出一道弧线，飞向了半空！

这一刀如此的精确，竟然齐着关节处的软骨切断！

“啊——！！！我的手！我的手臂！！”

足足迟滞了一息，撕心裂肺、不似人声的凄厉惨嚎才从李轩喉咙深处迸发出来，充满了极致的痛苦、难以置信的惊恐，以及前途尽毁的绝望。

他左手死死捂住那狂喷鲜血、筋肉狰狞断开的创口，脚下踉跄，如同醉酒般向后跌退，脸上血色瞬间褪尽，惨白如纸，冷汗混杂着血污滚滚而下。眼中除了剧痛，只剩下无边的恐惧与茫然。

这兔起鹘落、血腥残酷到极点的突变，让整个战场出现了瞬息死寂般的凝滞！

所有人都被这突如其来的逆转和惨烈景象所震慑。

“竖子敢尔！！”二长老最先反应过来，发出一声惊怒交加的暴吼，目眦欲裂。

他没想到在自己眼皮底下，少主竟遭此重创！枯瘦的身影带着一股狂风便要扑向那收刀而立的蒙面女侠，意图将其擒杀。

然而，对面的白言冰，岂会放过这用鲜血和惨嚎创造的、千载难逢的战机？胸中那口几乎耗尽的浊气被强行提起，近乎枯竭的丹田压榨出最后一丝元气，手中长剑发出一声不甘的铮鸣——

“弥焰斩！”

双手过头，雷火剑气凝成实质——一剑劈下！

二长老见逃是逃不得了，只得全力催动灵力，试图用土灵护盾硬挡；随即，便是“梆”的一声巨响。

虽然接住了这一招，但二长老依然是被白言冰一记重剑震得气血翻腾，踉跄后退，手中蛇头拐杖的光芒都黯淡了几分。

果然是势大力沉；若是再来一剑，他估计自己会被当场劈成两半，就地陨落！

而那位一刀建功的蒙面女侠，在一刀断臂之后，毫无停滞，甚至没有多看惨叫的李轩一眼。她身影如黑烟般灵动一转，手腕轻振，甩落狭长刀刃上最后一滴滚烫的血珠，刀光再起，与已是强弩之末却战意昂然的白言冰形成了完美的犄角之势，一左一右，夹攻向惊怒交加的二长老！

本就因李轩重伤而心神大乱、又独木难支的二长老，此刻要面对一个剑修不顾一切的狂攻，以及一个身法诡秘、刀法狠辣精准、实力深不可测的蒙面女侠的袭扰，再加之先前的战斗过程、以及硬接白言冰全力一击对他的消耗巨大，顿时落入下风。

几名机灵的内门弟子见此，慌忙拼死冲上，抢在白言冰彻底恢复过来、再斩下下一剑之前，将脸色铁青的二长老连拖带拽地抢救回了本阵。

高空之中的战局，亦因此惊天变故而瞬间扭转。正与叶牧谦以神识气机凶险缠斗的李青阳，眼角余光清晰无比地瞥见了爱子手臂飞起、血如泉涌的那一幕。

即便他心性深沉，此刻也不由得心神剧震，一股混杂着震怒、心痛、难以置信的狂暴情绪直冲顶门，让他凝聚风刃的灵力不由出现了一丝极其细微却足以致命的紊乱。那原本流畅环绕周身的青色符文，也随之轻轻一颤。

叶牧谦虽也面色愈发苍白，强行持续催动见独领域带来的无形反噬如同细针，不断刺探着他的神魂与经脉，额角已渗出细密冷汗。但是生死关头，这对手心神失守的刹那，岂能错过？

他随即鼓起所剩元气，一记古朴雄浑、毫无花哨却携着山岳之势的掌风排空推出！

李青阳反应极快，仓促布下不少风灵护盾，试图防御；下一瞬，掌风结结实实地拍在了李青阳仓促布下的风灵护盾之上！

“砰！”

闷响声中，李青阳身形一晃，被这股巨力逼得向后飘退数丈，方才勉强稳住，脸色一阵青红交替，体内气血翻腾不止。叶牧谦也没有继续追击，但依然保持着警惕。

李青阳飘然落回地面，脚踩在沾染鲜血的碎石上。眼前景象让他胸腔几乎炸裂：

爱子李轩倒在地上，被几名弟子手忙脚乱地勉强接上手臂，并按住止血，惨嚎已变得嘶哑断续，断臂处一片狼藉；

二长老面色灰败，被弟子搀扶着，显然消耗甚大、无力再战；

周围弟子虽然依旧人数占优，但士气已然受挫，面对陈千羽指挥下结阵死守的散修残余，以及那神秘出现的蒙面女侠，攻势明显迟滞，眼中都有了惧意。

而对面，叶牧谦虽然气息不稳，却依然挺立如松，身边不仅站着伤痕累累却眼神炽烈的白言冰，更多了那位一刀改变战局的神秘黑衣女子……他脸色铁青，胸膛剧烈起伏，心中充满了滔天的不甘、蚀骨的怨毒，以及对叶牧谦那诡异莫测、竟能影响人心神与天地共鸣手段的深深忌惮。

继续打下去？即便凭借人数优势最终能惨胜，青云门这批核心精锐也必将损失惨重，元气大伤，在青石城的统治根基都可能动摇。更重要的是，轩儿的断臂之伤，急需及时救治，再拖下去，就算接回也恐成废人，甚至性命堪忧！

瞬息之间，无数利弊在心头翻滚。

这位掌控一城的门主终究是能屈能伸、权衡利害的枭雄之辈。他强压下几乎冲毁理智的滔天怒火与刻骨恨意，猛地抬起右手，止住了周围仍想鼓噪进攻的弟子们。目光阴沉如暴风雨前的天空，死死钉在叶牧谦脸上，声音因极度压抑而变得干涩嘶哑，仿佛砂纸摩擦：“住手！！”

这一声断喝，让混乱的战场迅速安静下来，只剩下山风吹过林梢的呜咽，以及李轩那里传来的、越来越微弱的痛苦呻吟。

李青阳腮帮肌肉不断鼓动，几乎要将后槽牙咬碎，一字一句，从牙缝里挤出来：“叶牧谦……好，好手段！今日之事，到此为止！”

他每一个字都重若千钧，带着无尽的不甘：“矿税……即日起，回归原先的一成！你们……”他手指颤抖着指向叶牧谦师徒和散修，“立刻离开青石城地界，从此，井水不犯河水！”

这看似是巨大的妥协与让步，矿税从四成骤降至一成，几乎等于放弃了大部分利益。

但在场所有人都能听出他话语中那压抑到极点的、几乎要溢出来的怨毒恨意。这只可能是迫于形势的权宜之计，是猛兽受伤后退回巢穴舔舐伤口前的暂时退让。

叶牧谦缓缓吐出一口浊气，体内翻腾的气血与神魂的隐痛被他强行压下，喉头那股腥甜味道也被咽了回去。他深深看了李青阳一眼，没有言语，只是几不可察地点了点头。

他知道，经此一役，双方仇怨已结死，所谓的“井水不犯河水”，不过是风暴来临前脆弱的平静，是句谁也不会当真的空话。

看来，青石城已经彻底没法待了。

\vspace{2em}

那位关键时刻出手、一刀定乾坤的蒙面女侠，见局势已然明朗，强敌退意已生，似乎并无意在此久留，更不欲卷入后续可能的琐碎交涉。她手腕一翻，那柄狭长利刃在空中划出一道冷冽弧线，精准归入背后刀鞘，发出一声清脆的“咔哒”轻响。

随即，她转身便欲如同来时一样，悄无声息地融入道旁尚未散尽的昏暗林影之中——神秘而来，事了拂衣而去。

“沈姑娘？”

\vspace{2em}

就在她身形将动未动、气息将敛未敛的微妙瞬间，一个带着几分迟疑、几分不确定，却又清澈悦耳的女声，带着一丝急切，轻轻响起。正是刚刚为一名重伤散修施针稳住伤势，抬首望来的陈千羽。

她方才忙于救死扶伤，无暇细看。此刻稍稍得空，目光落在那黑衣女侠虽覆着黑巾、却难掩其熟悉挺拔的身姿轮廓，尤其是转身时那惊鸿一瞥的侧影，以及那双清澈明亮、即便含煞也带着独特倔强与侠气的眼眸，与记忆中那位曾在客栈深夜畅谈、赠言后飘然离去的黑衣女侠身影，渐渐重叠在一起。

蒙面女侠闻声，身形明显一僵，仿佛被无形的丝线拉扯住，即将迈出的脚步顿在了原地。

陈千羽见状，心中把握又多了几分，上前两步，语气更添肯定，同时抱拳施礼：“是沈瑶沈姑娘吗？方才多谢姑娘仗义出手，救我等于危难之际！”

叶牧谦与白言冰闻听“沈瑶”二字，目光也不由得瞬间聚焦过来。白言冰恍然，难怪方才那致命一刀的风格与时机把握，隐隐有种似曾相识的锐利之感。叶牧谦眼中则闪过一丝了然，先前那模糊的熟悉感终于有了着落，原来是她。

身份被直接道破，沈瑶知道无法再装作素不相识的路人了。面巾之下，她唇角不禁泛起一丝无奈的苦笑，心中暗自叫苦。

难道要直说自己其实并未远走，反而暗中跟了他们一路，从青石城客栈外到这西山险道，一直在旁观望，直到见他们师徒陷入绝死之境，才按捺不住出手相助？

这听起来……实在有些难以启齿，更像是有意尾随了。

但不认，又有点不太好。

略一踌躇，她只得缓缓转过身，面对着叶牧谦师徒三人。抬起手，轻轻摘下了那沾染了战场尘埃与一丝淡淡血气的黑色面巾。

顿时，那张独特的容颜显露在略显昏沉的天光下——半边清丽绝伦，肤光如玉；另外半边，却因陈年旧伤而留下了无法完全褪去的淡淡灼痕，破坏了完美的对称，却给这张脸平添了几分坚毅与故事感。

她略显局促地抱了抱拳，声音比起初遇客栈那夜的清冷疏离，此刻多了些复杂难明的情绪，似乎有些不好意思：“叶先生，白少侠，陈姑娘……又见面了。”

一个简简单单的“又”字，道尽了这重逢中难以言说的巧合与必然。

然而，就是这“沈瑶”二字，以及那张半面灼伤的独特容颜，在落入另一边正强压火山般怒火、忍痛检查儿子伤势的青云门主李青阳耳中与眼中时，却无异于两道九霄狂雷，携着万钧之势，狠狠劈入了他的神魂深处！

他霍然抬头，原本只关注儿子伤势的、布满血丝的双眼，瞬间锐利如发现致命威胁的鹰隼，死死锁定在那个摘下面巾、正与叶牧谦师徒见礼的黑衣女子身上。

那张脸……尤其是那半面灼伤的标志性特征，与“沈瑶”这个名字结合，一个令他头皮发麻、脊背生寒的骇人身份，如同破开迷雾的闪电，瞬间清晰地浮现在他心头！

沈瑶！那个十年前从内域七大主城之一、天枢城的绝对霸主——超级势力“天师府”中出走、行踪成谜的大小姐！

虽然她离家已久，且行事低调，但关于这位大小姐的零星传闻与最基本的信息，在青石城这等距离天枢城不算太远、消息相对灵通的中等城池高层圈子里，绝非什么机密！

李青阳只觉得一股寒气直冲天灵盖。

她的父亲，是沈天穹！

更要命的是，沈天穹老来得女，虽然沈瑶并不喜欢她的父亲，但沈天穹确实是对沈瑶这个女儿宠溺到了近乎“女儿奴”的地步。

若是在他李青阳的地盘上，因为他青云门的围杀，导致沈瑶有了丝毫损伤，哪怕只是被剑气擦破一点油皮……李青阳简直不敢想象，那位护短到极点的沈天穹，其怒火会何等恐怖，其报复又会何等酷烈！别说他一个区区李青阳，恐怕整个青云门，都要在天师府降下的雷霆之怒中被夷为平地！

之前他敢对叶牧谦等人下死手，是精准地判断他们没有后台，是可以捏的“软柿子”。可现在，沈瑶竟然与他们站在一起，还为了他们悍然出手，斩断轩儿一臂！

这意味着什么？这意味着叶牧谦一行人，很可能与天师府有着某种不为人知的紧密关联，或者至少，是得到了沈瑶的认可与庇护！打叶牧谦，就等于打沈瑶的脸；动沈瑶在乎的人，无异于直接挑战沈天穹的逆鳞！

再打？继续围攻？借他李青阳一百个胆子，此刻也不敢再兴起丝毫这般念头！别说继续动手，他现在连一丝敌意都不敢再投向那个方向，生怕引起那位大小姐的半点不快。

李青阳的脸色在极短的时间内急剧变幻，由铁青转为煞白，又强行想挤出一丝和缓，最终扭曲成一个比哭还要难看十倍的僵硬表情。

他朝着叶牧谦和沈瑶的方向，隔空远远地拱了拱手，腰身都不自觉地弯下了几分，声音干涩嘶哑，却尽可能放得平缓、甚至带上了一丝恭敬与惶恐：“原来……是沈仙子大驾光临。李某……有眼不识泰山，先前门下弟子若多有冒犯，纯属无心之失，万望海涵。今日之事……纯属误会，一场天大的误会！我等这就离去，绝不再打扰仙子与诸位道友清静。”

他语速极快，如同爆豆，仿佛说慢了半分都会再惹来什么不可预料的灾祸。

说完，根本不敢等待沈瑶或叶牧谦的任何回应，立刻像是被火烧了屁股一般，猛地转身，对着犹自茫然不知所措的青云门众弟子，厉声嘶吼，声音中充满了惊惶与急迫：“还愣着干什么？！蠢货！带上轩儿和二长老，撤！立刻回山！快！！”

青云门众弟子虽然绝大多数根本不明所以，完全搞不懂为何门主在听到一个名字、看到那女子的脸后，会吓得如此失态，前倨后恭至此；但门主那几乎变调的厉喝和眼中无法掩饰的恐惧是做不了假的。

无人敢多问半句，更无人敢迟疑，慌忙七手八脚地抬起已因失血和剧痛而陷入半昏迷的李轩，搀扶起面色灰败的二长老，拖拽着其他伤员和死者，连用作诱饵的灵石车都不要了，如同退潮般仓皇向山道另一头涌去。

速度比来时围困时更快，阵型全无，只恨爹娘少生了两条腿。

\vspace{2em}

转眼之间，青云门弟子们便消失在山道转弯处，只留下满地狼藉的兵刃痕迹、几具死于刀剑之下或法术之中的散修尸体、尚未完全凝固的斑斑血迹，以及空气中弥漫的、渐渐被山风吹散的血腥气，无声地诉说着方才那场惊心动魄、结局却出乎所有人意料的惨烈厮杀。

山道间骤然安静下来，只剩下风声穿过林隙的呜咽，以及散修们粗重未平的喘息与伤者压抑的呻吟。

侥幸活下来的散修们面面相觑，脸上都带着劫后余生的茫然与难以置信。

他们看看满地狼藉、血迹斑斑的山道，又看看青云门退走的方向，最后目光齐齐落在那个收起长刀、黑巾覆面的女子身上，眼神里充满了深深的敬畏与几乎要满溢出来的好奇。

他们想不明白，那个片刻前还杀气腾腾、占尽优势的青云门主李青阳，为何在听到“沈瑶”这个名字、看清对方面容后，竟会吓得魂不附体，态度急转直下，比见了猫的老鼠溜得还快；但他们知道，正是这位神秘女子的出现和那惊世一刀，彻底扭转了战局，带来了这不可思议的生机。

一些伤势较轻的散修挣扎着起身，想要向沈瑶道谢，却又慑于她那清冷独立的气场，嗫嚅着不敢上前。

叶牧谦将李青阳那前倨后恭、惊惶退走的全过程尽收眼底，心中也对沈瑶的身份背景已猜到了七八分。他对沈瑶再次郑重拱手，语气诚挚：“多谢沈姑娘再次仗义援手，解我师徒与诸位道友覆灭之困。此恩，叶某铭记。”

沈瑶连忙侧身避开这郑重一礼，摇了摇头，声音透过面巾显得有些闷，却依旧清冽：“叶先生言重了。路见不平，拔刀相助，本是分内之事。”

她清亮的目光扫过叶牧谦那张因强行催动见独境而显得比平日更加苍白、气息也显然极度虚浮的脸庞，又掠过身上带伤、血染衣袍却仍挺直站立的白言冰，以及额发微乱、俏脸因紧张和忙碌而泛红、正匆忙检视伤员情况的陈千羽，最后落在那些或坐或躺、惊魂未定、眼中残留着恐惧的散修身上。

“诸位激战方歇，伤者众多，此地血气弥漫，绝非久留之所。恐怕……急需一个安稳的地方休整疗伤。”

叶牧谦闻言，脸上露出一丝温和的笑意，顺势发出邀请：“沈姑娘所言极是。此地确实不宜久留。我等在附近山中觅得一暂居之所，虽简陋，却可遮风避雨，稍作安顿。沈姑娘若暂无急务在身，不妨随我等前去小住几日，略作休整。也让小徒千羽有机会为姑娘调理一下方才激战可能带来的气血损耗，聊表我等寸心谢意。”

沈瑶闻言，略一沉吟。

她独来独往惯了，如同飘萍，最不喜与人有过多牵扯羁绊。但目光触及叶牧谦那双深邃却依旧平静的眸子，仿佛方才的生死搏杀与眼前的混乱都未能动摇其心神根本；再想起李青阳见到自己时那骤变的脸色与仓皇退走的姿态，心中几个念头迅速转动起来。

她的看法，其实与落荒而逃的李青阳不谋而合，只是立场不同。

在她看来，叶牧谦此人，实力手段深不可测，所修之道途更是迥异于世俗常见的灵枢径，处处透着神秘。然而种种迹象表明，他们师徒三人又确实像是毫无根基、漂泊至此的散修。

今日，叶牧谦与地头蛇青云门结下如此深仇，对方表面上被迫退去，但以李青阳那阴狠枭雄的性子，暗地里必然不会善罢甘休，只会如同受伤的毒蛇，潜伏起来，等待更致命的反扑时机。

自己若此刻就此离去，叶牧谦师徒与这些散修，恐怕用不了多久就会面临青云门更隐蔽、更残酷的报复。

但若自己暂且留下，虽然自己不愿凭借父亲沈天穹与天师府的赫赫威名作为保护伞；但客观上只要自己还在此处，对于青云门而言就是一种无形的、巨大的威慑。李青阳投鼠忌器，短期内绝不敢再明目张胆地前来寻衅。

这保护伞或许不够坚固长久，但至少能在叶牧谦几人离开青石城之前，争取到宝贵的喘息与应对时间。

这或许能稍稍报答一些当初叶牧谦那番“与天奋斗，其乐无穷”的言论带给她的震撼与触动，也能了却自己这数月来暗中跟随观察，所产生的那一丝说不清、道不明、却实实在在存在的莫名牵挂。

“也好。”沈瑶点了点头，声音恢复了往日的清冷平静，但那双露在外面的眸子，眼神却比初遇时柔和了些许，少了几分刻意维持的距离感，“那便叨扰叶先生几日。”

当下，叶牧谦便让气息稍平的白言冰和已快速处理完几处紧急伤口的陈千羽，协助那些受伤较轻、尚能行动的散修，迅速清理战场。主要是收敛散落的兵刃、诱饵灵石，并将重伤或死亡的同伴用临时制作的简易担架小心抬起。

沈瑶也默默加入帮忙的行列，她动作干脆利落，对包扎止血等外伤处理似乎颇有经验，手法迅捷有效，让一旁忙碌的陈千羽都忍不住多看了几眼。

一行人趁着暮色尚未完全四合，沿着隐蔽的山径，朝着背靠山崖的石屋方向迤逦而行。夕阳余晖将他们的影子拉得很长，与山林树影融为一体，很快便消失在西山深处。

回到石屋时，天已擦黑。陈千羽立刻马不停蹄地开始为受伤最重的几名散修、以及自家老公白言冰处理伤势。

沈瑶本想上前帮忙，她对自己的外伤处理手法颇有信心。但陈千羽却以“沈姑娘是客，且方才激战耗神费力，理应先行休息恢复”为由，温言却坚定地将她请到了一间刚刚匆忙收拾出来、还算干净的侧室暂歇，并细心地送来清水和干净的布巾。

沈瑶看着她忙碌得脚不沾地的背影，心中微动，便也不再坚持，接受了这份体贴的好意。

是夜，月明星稀，清冷的银辉透过石屋的小窗洒落进来。经过陈千羽大半个晚上的精心救治，以及叶牧谦所提供的药效奇特的丹药调理，众人的伤势总算都稳定下来，最危险的几个也脱离了性命之虞。

疲惫不堪、心神耗尽的散修们，已被安置在附近另一处更为宽敞、干燥的山洞中沉沉睡去，鼾声此起彼伏。石屋中终于恢复了宁静，只余下淡淡的、宁神的药香在空气中缓缓萦绕，冲淡了白日残留的血腥气。

“笃笃。”轻轻的叩门声响起。

侧室内的沈瑶并未入睡，正对着窗外月光静静调息。闻声起身开门，只见陈千羽端着一只粗陶碗站在门外，碗中升腾起带着草药清香的袅袅热气。

“沈姑娘，打扰了。这是我刚熬好的安神理气汤，用的是山里采的几味平和药材，对恢复心神、调理气血有些微助益。”陈千羽将温热的药碗递过来，脸上带着真诚的感激与浅浅的疲惫笑意，“今日……真是多亏你了。若非你及时出现，后果不堪设想。”

沈瑶接过药碗，触手微温。她摇了摇头，就着灯光看了看碗中色泽清亮的药汤：“陈姑娘不必客气。分内之事罢了。倒是陈姑娘你的医术，似乎比当初客栈初见时，又精进了许多。这些散修伤势轻重不一，且多有内腑震荡，竟都能被你稳下来，着实不易。”

陈千羽闻言，眼睛微微一亮，似乎想到了什么。

她在沈瑶对面坐下，压低了声音，眼中闪烁起一种混合着对师尊的崇敬与分享秘密般的神秘光芒，语气也带上了些许讲述秘闻的意味：

“其实……这都要归功于师尊传授的道统与众不同。沈姑娘，你我相识虽短，但并肩战斗，也算是生死之交了。有件事不知当问不当问，”她声音压得更低，几乎如同耳语，“你可知我师尊的来历根脚？”

沈瑶一怔，放下药碗，面巾下的眉头微挑：“叶先生？他……不是自称来自荒域，云游至此么？”这是叶牧谦对外最常提及的说辞。

“是，也不是。”陈千羽的声音轻得如同风吹羽毛，身子也不由自主前倾了些。

“我与言冰师兄曾偶然听到师尊睡眠中，或有梦境呓语，提及一些极其古怪的词语，诸如‘设定集’‘补全仙路’‘篇章缺失’等等……更曾于梦中，无意识地背诵出一些闻所未闻、却仿佛直指大道本源的玄奥篇章。醒来后师尊自己只言是旧梦，却对那些玄奥篇章的来源闭口不谈。”

她深吸一口气，眼中光芒更盛，语气愈发笃定，带着一种不容置疑的坚信：“我们将此事询问师尊，师尊他确实承认……自己并非我们此界之人！”

沈瑶眨了眨眼，一时间没完全反应过来这句话的含义，只是下意识地觉得有些荒诞。

陈千羽却不待她细想，继续用那种深信不疑的口吻说道：“我们怀疑，师尊本是上界仙尊，因见我们此界流传的‘灵枢径’仙路有偏，大道有缺，长生虚妄，故特意分神临凡，化身至此，传下这‘问道途’，旨在补全此方天地缺失的天道法则，重塑一条真正的、通向不朽的长生之路！”

“也正因如此，师尊才必须时刻压制自身真正的仙尊修为与无上气息，以免过于强大的力量与迥异的道韵，引发此方天地法则的本能排斥与愤怒反噬！他所展现的，或许只是亿万分之一的力量投影罢了！”

沈瑶听完陈千羽这番石破天惊的猜想，第一反应是有些好笑，甚至有些无奈。

上界仙尊？补全仙路？

这听起来太过于离奇。

她端起药碗，下意识地抿了一口微苦的药汤，本想随口附和或委婉地提出质疑。

然而，当“上界仙尊”“补全仙路”“压制修为”“天地反噬”这些词语在脑海中反复回旋，与她这数月来暗中观察、以及今日亲眼所见的叶牧谦种种神秘之处——

那迥异于灵枢径、完全不同的修炼体系与力量表现；

那深不可测却似乎总有所限制、不曾全力施为的实力；

那些发人深省、直指修行本质的言论。

这些东西相互印证、细细咀嚼时，她嘴角那丝若有若无的、觉得好笑的弧度，渐渐淡去了，最终消失不见。

手中的粗陶药碗停在半空，沈瑶的目光变得有些飘忽，失去了焦点，陷入了长久的沉思。

仔细想想……若暂时抛开“灵枢径是唯一正统”的固有认知，以一种全新的、不带偏见的视角去审视叶牧谦的所作所为……似乎，还真有那么一点补天者、传教士、引路人的影子？

难道……陈千羽和白言冰这看似荒诞不经、源自零星梦呓的猜想，竟然歪打正着，无意中触及了某种惊世骇俗的真相？

叶牧谦难道真是肩负着某种使命、来自更高层面的存在？

这个念头一旦升起，便在沈瑶心中迅速扎根，缠绕生长，再也挥之不去。她再看手中那碗由陈千羽亲手熬制、药性温和中正、蕴含着独特平和生机的汤药，仿佛也悄然染上了一层不同寻常的意味——这是否也暗含了问道途那滋润万物、调和身心的理念？

夜还很长。石屋之外，山风依旧拂过林梢，发出永恒的呜咽。屋内，灯火如豆，光影摇曳。沈瑶对着跳跃的烛火，怔怔出神，许久未动。

\vspace{2em}

自西山道那场惊心动魄的伏击与反杀之后，叶牧谦师徒与沈瑶便在青石城外的山居石屋暂住下来，悄然度过了几个月的平静光阴。

青云门果然如沈瑶所料，在李青阳严令乃至严酷惩罚的威慑下，再未敢明目张胆地前来这片区域寻衅滋事，连日常的巡逻都刻意绕开了这片山林。

而青石城中对外宣称的矿税，也老老实实地从四成恢复并维持在了一成，让散修们获得了难得的喘息之机。

散修们对叶牧谦师徒与那位神秘的黑衣女侠感激涕零，但也深知此事关乎生死，自发地三缄其口，成为彼此间心照不宣的秘密，无人敢对外界多言半句，生怕给恩人带来麻烦。

陈千羽继续精研医术，那本得自师尊的、内容似乎无穷无尽的医典被她反复揣摩。

不时有散修借着采药或狩猎的名义，悄悄摸到石屋附近，带着家人或自身的伤病前来求诊，陈千羽总是来者不拒，尽心医治，她的神医名声在散修底层悄然流传，却又被默契地限制在小范围之内。

白言冰则大部分时间都在石屋后的空地上，潜心打磨他那招“半成品一般的”（沈瑶语）弥焰斩。

他谨记师尊关于“刚猛有余，灵活不足”的教诲，开始有意识地在至刚至烈的剑招中，尝试融入更细腻的变化、更隐蔽的后手与更迅捷的转换，最终体悟了些新的剑式，被他称为“直环焰”。

实际上这东西可以算作是小号的弥焰斩，至少动作依然是下劈，但力量克制了许多。

\vspace{2em}

沈瑶最终也在这里停留了下来，似乎暂时停止了四方漂泊、居无定所的云游生活。

在更深入的相处中，叶牧谦师徒三人也逐渐了解到，沈瑶最擅长、造诣最深的其实是幻术与符箓之道，那是她出身天师府的根基所在。

但沈瑶自离家后，极少用过这种手段。

毕竟，居无定所、资源匮乏的流浪之下，布置高明的幻阵需要不少材料且极其耗费心神维持，符箓也是消耗品，可能很久都无法补给。

而凌厉的身法与一击必杀的刺杀术，对她而言更多是离家之后、因应漂泊生涯的需要，而刻苦修炼至精通的程度。

将身法练到极致，配合简洁狠辣的刺杀技巧，以及审时度势、打了就跑的战术，在大多数险恶环境中反而更实用、更有效。

沈瑶对问道途的兴趣日益浓厚，虽因心中那份关于“上界仙尊”的惊疑而未曾正式开口请教，但每当旁观白言冰练剑时那迥异于寻常剑道的发力方式与气韵流转，或聆听叶牧谦偶尔对弟子寥寥数语却直指核心的点拨时，她那双清澈的眸子中，总会流露出深深的思索与探究之色。

\vspace{2em}

这一日，这数月来山居的宁静，被一股来自遥远天际、宏大庄严到不可思议的意志毫无预兆地打破。

正值午时，青石城上空，原本湛蓝如洗、万里无云的天际，骤然风起云涌！

道道绚烂瑰丽、宛如最上等锦缎织就的祥瑞霞光，自极高极远、仿佛世界尽头的方向铺天盖地蔓延而来，顷刻间染遍了小半个天空。

与此同时，清越缥缈、似钟似磬、洗涤人心的仙音妙乐凭空响起，袅袅不绝，响彻在青石城每一个角落，无论修士还是凡人，无论身处室内还是街头，皆听得清清楚楚，不由自主地心神震动，抬头仰望。

只见那漫天流转、华美庄严的霞光云海之中，灵力汇聚凝结，缓缓显现出一枚巨大无比、几乎覆盖了小半个城区的玄奥符印虚影——这正是天师府独有的、代表其无上权威的徽记！

符印虚影缓缓旋转，洒下无尽光辉。其下，一道恢弘浩大、淡漠高远、宛若九天之上天道亲自颁布纶音般的庄严宣告，清晰地、一字一句地传入青石城及周边方圆数百里内每一个生灵的耳中，直抵神魂深处：

“天师府昭告内域诸城：十年一度之天师试炼，将于三月后，于天枢城如期举行！

”凡内域修士，骨龄百岁以下，修为达玄脉境及以上者，皆可参与。

“试炼旨在遴选俊才，磨砺心志。

“通过全部试炼关卡者，可获准入天师府密阁，限期学习参悟一月。

“最终优胜者，更可获封‘小天师’尊号，享天师府内门弟子同等礼遇，并另有厚赐。

“各城接引、核验事宜，将由天师府特派执事负责，不日即至各城。

“望内域才俊，踊跃赴会，一展风采。”

霞光中的符印虚影随着宣告微微明灭，仿佛在为其增添无可置疑的权威。

待到余音袅袅，最终消散于天地之间，那漫天霞光与巨大的天师府符印虚影，才如同海市蜃楼般缓缓变淡、消散。青石城上空重归晴朗，仿佛刚才那震撼人心、宛如神迹的一幕从未发生过。

然而，那宣告的余韵，那符印的威严，却久久烙印在青石城数十万生灵的心头，更如同一块巨石投入平静湖面，在无数符合条件、胸怀壮志的年轻修士心中，激起了滔天巨浪；青石城中，一时鼎沸起来。

“天师试炼……终于，又到了十年之期！”

“天枢城！万象密阁！哪怕只能在外围参悟七日，也抵得上寻常苦修十年啊！”

“何止！若能夺得‘小天师’尊号，那才是真正的鱼跃龙门，一步登天！从此资源、地位，予取予求！”

“太难了……听说上次试炼，内域数百城池，数千名玄脉境人士齐聚，最终能完整通过所有试炼的，不过十人数。获得‘小天师’称号的，更只有区区三人！”

“即便如此，也值得拼死一搏！此乃内域最公正、不同出身的通天之路！一旦成功，便是鲤跃龙门，光耀门楣！”

这样的宣告，在其余城池反响也差不多。

山居之中，叶牧谦师徒与沈瑶，自然也清晰地听到了这响彻天地、直入神魂的宣告。白言冰与陈千羽对视一眼，都从对方眼中看到了难以抑制的意动与神往。

天师府！密阁！“小天师”！

这些名词代表着内域修行的最高殿堂与荣耀，对于任何有志于大道的年轻修士而言，都有着致命的吸引力。

两人不约而同地将目光投向了静坐在石屋门口、抬头望天、神色一如既往平静无波的师尊叶牧谦，等待着他的示意。

叶牧谦面色沉静，目光深远，仿佛穿透了晴朗的天空，望向那宣告传来的无尽远方，又仿佛只是在思索着这突如其来的消息背后，可能蕴含的深意与变数。

而一旁的沈瑶，在听到天师府公告如雷霆般炸响在天地间的刹那，娇躯几不可察地微微一颤，仿佛被无形的电流击中。

是她那个便宜老爹沈天穹的声音！

她下意识地抬起手，指尖轻轻拂过那半边用以遮挡灼痕、此刻在阳光下显得格外轻盈的面纱边缘。那双总是清澈坚定、或带着思索的眼眸，此刻骤然变得复杂难明。

那里面有对遥远过往的模糊追忆，有对那座威严府邸下意识的抗拒与疏离，更深处，或许还藏着一丝连她自己都未曾清晰察觉的近乡情怯般的彷徨。

\vspace{2em}

数日后，天师府派遣至各城负责接引、核验的官员与执事，陆续抵达各自的目的地。

青石城也迎来了两位身穿天师府特有月白色云纹道袍、气息沉凝、一举一动皆透着大势力出身才有的独特气度的执事。他们入城后，径直前往城主府拜会并接洽试炼相关的具体事宜，随后便在城中几处人流最密集的广场、街口，设立了醒目而庄严的报名登记点。

关于试炼的详细规则告示也被张贴出来，引得人群终日围聚，热议不断。

\vspace{2em}

这一日午后，其中那位年岁稍长、面容肃穆的执事，例行在城中巡视各报名点的情况；忽然，他的视线在不经意间扫过街角一家生意清淡的药铺时，猛地定格！

一个刚刚从药铺柜台转身、手中提着几包用油纸草绳捆好的药材的黑色劲装女子，正侧对着街道。她以轻纱覆面，遮住了大半容颜，但那短小矫健、线条流畅如猎豹般的独特身形，那行走间自然而然的、带着某种韵律与警觉的步态，尤其是她微侧脸颊与掌柜结账时，轻纱边缘偶然露出的一线肌肤轮廓，以及那双即便隔着半条街、也能让人感受到其清澈明亮与内在倔强神采的眼眸……

种种细节，如同散落的拼图，在这位年老执事脑海中迅速拼合！

这位执事在天师府中资历颇深，对那位离家出走的大小姐相关的秘闻知之甚多。此刻，眼前这女子的身形气韵，尤其是那隐约可见的半边脸颊轮廓及独特的眼神，与他记忆中的信息飞速重合！

他强压下瞬间翻腾的激动与惊疑，不动声色地借着查看街边摊位，稍稍靠近了些许距离，借着挑选货物的间隙，眼角的余光再次仔仔细细地将那黑衣女子从头到脚审视了数遍。

八九不离十！

老执事心中已然有了论断。他不再犹豫，立刻转身，脚下步伐瞬间加快，以与其年龄不符的迅捷速度，径直返回了天师府执事在青石城暂居的驿馆。

\vspace{2em}

“禀府主，于青石城，疑似发现大小姐踪迹！黑衣劲装，面纱覆容，形貌特征高度吻合，属下有八成把握！”

\vspace{2em}

就在当天傍晚，夕阳的余晖尚未完全被暮色吞没，青石城上空，异象再生！

这一次降临的是一种更直接、更磅礴、更令人灵魂本能战栗的纯粹威压！

毫无征兆地，方圆百里的天空，骤然暗了下来！

在这片幽暗中心，一点光芒骤然亮起，随即急速膨胀、拉伸、凝聚……最终，一尊顶天立地、模糊了面容却散发着无尽威严的巨大人形虚影，巍然显现在青石城的天穹之上！

法相！

那法相看不清具体五官细节，只能感受到其身着古朴宽大的袍袖，头戴高冠，周身上下，有无穷无尽的玄奥符文如同呼吸般流转明灭。

整个城池瞬间鸦雀无声，鸡犬不闻。所有修士，上至青云门主李青阳这等一城霸主，下至刚刚引气入体的入门弟子，无不感到体内灵力如同被冻结般凝滞难行，心神摇曳，神魂深处本能地生出无限敬畏与渺小之感；城中的凡人更是早已匍匐在地，瑟瑟发抖，以为天神震怒。

法相境！

这个令人灵魂颤栗的词汇，浮现在青石城所有稍有见识的修士心头。

既然是符文法相，那来者的身份，此刻已昭然若揭——

内域最顶尖的强者之一，天师府当代大天师，沈天穹！

\vspace{2em}

那巨大的法相虚影，在无边的幽暗中缓缓低首，目光精准地落在了城外那片看似平凡的山林之间，落在了叶牧谦师徒与沈瑶暂居的简陋山居小院之中。

一个平静、淡漠，却蕴含着不容置疑的威严与浩瀚力量的声音，直接在小院内每一个人的耳边、心底同时响起：

“瑶儿。”

仅仅两个字，却仿佛揉碎了十年光阴，裹挟着复杂难言的情绪。

“玩闹十年，也该够了。”沈天穹的声音继续响起，语气稍稍放缓，带上了一丝为人父的劝慰，试图软化女儿的抵触，“随我回家吧。你的脸伤……为父可以亲自去求丹鼎阁的王丹长老，她是当世丹道圣手，即便需要某些绝世灵药，为父也必定为你寻来。她……定有办法为你医治，恢复旧观。”

这承诺，对于一个正值青春、却因半边脸灼伤而不得不常年以纱覆面、将美丽与伤痕一同隐藏起来的少女而言，不可谓不重。

此刻，沈天穹似乎只是一个无奈的平凡父亲，小心翼翼地哄着与他闹别扭的女儿。

山居小院内，落针可闻。白言冰与陈千羽早已屏住了呼吸，连大气都不敢喘，紧张万分地看着僵立的沈瑶，又偷偷瞥向天空中那令人窒息的巨大法相，手心全是冷汗。叶牧谦则依旧静静立于一旁檐下阴影中，目光平静地仰望着那尊仿佛能决定此地所有人生死的虚影。

令人心悸的沉默，持续了足足十数息。这短暂的时光，在浩瀚法相的无形注视下，被拉长得如同一个世纪。

沈瑶终于动了。她抬起手，伸向耳边，摘下了那几乎已经成为她身份一部分的黑色面纱。

她的表情异常平静，只有一种历经漂泊坎坷、看遍世间百态后沉淀下来的坦然与坚定，仿佛那道伤痕在她脸上只是一枚独特的勋章。

“谢谢爹爹好意。”她开口，声音清晰，不高，却带着一种能穿透人心、不容忽视的力量。

“只是，爹爹让我回家，无非是让我再次回到那高墙深院之中，闭门静修，不理外事，不见众生疾苦，不去体会这世间最真实的冷暖，不去救助那些我能见到、也有能力去救助的苦难与不公……爹爹，这样的‘道’，清静是清静了，安全也安全了。但，于我不合。”

沈天穹那浩瀚的法相虚影，似乎微微一滞，周围流转的符文都出现了刹那的紊乱。显然，女儿如此直白、如此不留余地的拒绝，大大出乎了他的预料；而在惊讶过后，更多的便是深深的失落。

父女二人之间似乎彻底谈崩了，气氛顿时变得尴尬起来，沉重得让人喘不过气。

\vspace{2em}

就在这时，一个温和却异常清晰、如同山涧清泉般打破死寂的声音，不紧不慢地插了进来：

“人各有志，令爱之道，或许便在红尘之中。”

是叶牧谦。

霎时间，沈天穹那浩瀚无边的注意力，瞬间从女儿身上，全部转移到了这个突然开口的陌生人身上。

法相的目光——或者说那笼罩天地的浩瀚意志——如同实质的光柱，聚焦于叶牧谦，带着凌厉的审视，与一丝被打断对话、被“蝼蚁”插言的不悦。

一个……看起来气息古怪，但绝对没开启过灵枢，更没凝聚过玄脉的……凡人？

这种小蚂蚁也敢在他沈天穹处理家事、与女儿对话之时，贸然出声？

\vspace{2em}

然而，就在那审视的意志接触到叶牧谦的下一刻，沈天穹心中那点因被冒犯而产生的不悦，瞬间被一股更强烈的惊讶所取代！

叶牧谦竟然在他的威压之下气息沉稳如渊渟岳峙！

这哪里是一个凡人能做到的？

更让沈天穹心中凛然、瞬间提起了十二万分重视的是：女儿沈瑶在此人插话时，非但没有流露出任何被冒犯或不满，反而下意识地、几乎难以察觉地朝着叶牧谦的方向微微侧身了半分，那姿态，竟隐隐含有一种倚靠与信赖的意味。

而她看向此人的眼神，虽然只是一瞥，但沈天穹何等眼力，立刻捕捉到了那其中不同寻常的信赖，绝非对待寻常路人或普通朋友的态度。而且，在此人身后，还规规矩矩侍立着两位年轻人，一男一女，目测与瑶儿年龄相仿，且皆气质不俗，目光清澈坚定，显然正是此人的弟子。

一个能让心高气傲、对天师府规矩都敢于反抗的女儿，愿意亲近、甚至隐隐有所追随的人物？

\vspace{2em}

到了沈天穹这个层次，站在这方世界修行之道的顶峰，他比绝大多数人都更深刻地明白，大道三千，修行之途博大精深，浩瀚无垠，“灵枢径”不过是其中流传最广、体系相对最完整的一条主流路径而已，绝非唯一正途。

历史上，乃至当今世上某些不为人知的角落，确实存在着一些古老或偏僻的传承，它们或许式微，或许鲜为人知，但往往鬼神莫测，不可小觑。女儿在外游历，机缘巧合下与这样的隐士高人结识，甚至得到其些许指点或庇护，从某种意义上说，未必是坏事，或许反而是一段难得的缘法。

心中的那点因女儿顶撞、外人插话而生的怒意，在这一连串的思量中悄然散去大半。取而代之的，是一种更为审慎的观察。他再次看了看倔强抿唇、眼神坚定的女儿，又看了看深不可测、气度沉静的叶牧谦，一个主意忽然浮上心头。

你若真有本事，便来天师府这大舞台上亮亮相，也让老夫亲眼看看，你究竟是何方神圣，值不值得瑶儿如此信赖。你若没本事。只知道招摇撞骗，那就别怪老夫手下无情了。

于是，巨大的法相虚影在沉默数息后微微颔首，那响彻天地、直抵神魂的声音再次响起，对象却已从沈瑶，明确地转向了叶牧谦。语气也发生了微妙的变化，带上了一丝探究和邀请的意味：

“阁下便是小女在外游历时结识的人物吧？果然气度不凡。适才听阁下所言，‘人各有志’，亦觉有理。瑶儿心向红尘，志在历练，其志可勉。而观阁下形神，似乎也非寻常之辈，自成气象。”

沈天穹随即抛出了一个石破天惊的邀请：“不知阁下，与身后这两位高徒，可有意参与我天师府此次举办的天师试炼？”

此言一出，不仅一直紧张关注着的白言冰和陈千羽瞬间眼睛一亮，呼吸都为之急促了几分。连原本打定主意抗争到底的沈瑶，都略显诧异地抬起头，看向天空中父亲那巨大的法相虚影，显然没想到父亲会突然有此提议。

沈天穹的声音继续平稳地响起，打消叶牧谦因为道途相左而可能产生的顾虑——如果有：“此试炼不考灵根，也不测经脉。其间所设重重关卡，皆是对参与者心性之坚、毅力之韧、实力之实、智慧之巧的综合考验。不知叶先生，意下如何？”

叶牧谦寻思：这一番话，说得合情合理，滴水不漏。先给我们戴上高帽，然后抛出橄榄枝，看似姿态放得相当低，实际上是想借机检测我们究竟值不值得沈瑶她信赖呢！

细细想来，对方给出的理由与条件，确实难以拒绝。首先，若是不去，完全有可能会被怀疑是害怕露馅的草包；天师府密阁的藏书，就能让他更系统地了解此方世界的修炼相关事宜，这无疑具有巨大的吸引力和完善后续道途的参考价值；这也确实是一个让白言冰与陈千羽在更为广阔、也更为残酷的内域舞台上真正崭露头角、检验所学，并获得认可的绝佳机会。

于公于私，似乎都没有拒绝的理由。

而且他相信沈天穹不会在试炼途中给他下绊子。毕竟众目睽睽，人言可畏啊。

于是，他脸上露出一抹坦然的淡笑，从容拱手道：“沈府主盛情相邀，叶某若再推辞，倒显得不识抬举了。如此，便携小徒前去天枢城叨扰一番。”

“好！” 沈天穹似乎很满意，“那便恭候阁下大驾光临。青石城接引执事，会为诸位安排妥当一切行程事宜。”

说完最后这句安排，巨大的法相虚影深深看了一眼下方已重新戴好面纱、神色复杂的沈瑶后，便开始缓缓变淡，连同那笼罩百里天穹的深邃幽暗，一同消散。

白言冰与陈千羽忍不住对视一眼，都从对方眼中看到了压抑不住的兴奋与跃跃欲试。

而沈瑶，独自静静地站在院中那棵老树下。她原本打定主意，就算是便宜老爹亲自降临，以法相之威施压，她也绝不轻易妥协，绝不再回到那座令她感到束缚与压抑的“家”。

但事情的走向完全出乎了她的预料：父亲没有强迫，没有震怒，甚至没有多做纠缠，而是……转而邀请了叶牧谦师徒前往天枢城参加试炼。

而更让她心神震动的是，叶牧谦……竟然答应了。

她不由自主地微微侧首，目光落在叶牧谦那平静的侧脸上。他正听着弟子们兴奋的低语，嘴角带着一丝极淡的、几乎看不见的和煦弧度。望着这张脸，沈瑶心中忽然毫无征兆地、汹涌地涌起一股强烈的、难以言喻的冲动——

他要去了。

去天枢城。

去那个她离开了十年、既熟悉到骨血里、又因时间和心境而变得无比陌生的地方。去那个象征着内域最高权威之一、也承载着她最复杂记忆的地方。

而自己……自己竟然在心底，产生了一种强烈的“跟着他去”的念头！

沈瑶的理智立刻开始分析：这念头一定不是出于对天师府那早已冷却的归属感，也不太可能完全是为了试炼本身——她对所谓的“小天师”尊号并无执念。

那是什么呢？

是因为想亲眼看看，他这个神秘莫测、被陈千羽私下猜测为“上界仙尊”的人物，踏上天师府的地盘后，究竟能掀起怎样的波澜？

还是因为，想继续探寻他身上的秘密？

抑或是因为……在这偏僻山居小院的平静时光，虽然没有锦衣玉食，没有仆从如云，却有着一种迥异于天师府那冰冷森严规则与压抑氛围的、生机勃勃的温暖，让她在不知不觉间，产生了一丝连自己都未曾清晰察觉的眷恋与不舍？

甚至是……单纯想跟着他？

沈瑶自己也说不清。几种情绪如同丝线般缠绕在一起，难分彼此。但她清楚地知道一点：当叶牧谦说出“携小徒，前去叨扰一番”时，她心底那声“不”字，无论如何也喊不出口了。

她无法再像过往十年中的任何一次告别那样，干脆利落地转身，独自踏入未知的漂泊旅途。

这个认知让她感到一丝莫名的焦躁与羞恼。

什么东西啊！生气！

她在心里无声地斥责了自己一句，为这种不受控制的、仿佛被无形线缆牵引的感觉。

她深吸了一口气，冰凉的空气吸入肺腑，稍稍平复了翻腾的心绪。然后，她迈开步子，走到叶牧谦面前。她的声音已经恢复了往日的清冷平静，仿佛刚才的激动与挣扎从未发生过，只是那语调的尾音，隐约带上了一丝不易察觉的软化：

“叶先生既然决定前往天枢城，”她顿了顿，目光状似无意地扫过白言冰和陈千羽，“那……我便也回去一趟吧。毕竟，天枢城我熟，试炼的流程规矩，乃至一些……需要注意的人和事，我或许比你们更清楚一些。”

她为自己找到了一个看似合情合理、甚至堪称周到的借口——向导。但那双露在面纱外的、清澈眼眸中微微闪动的眸光，却在不经意间，隐隐出卖了她内心的波澜。

叶牧谦闻声，转过目光，静静地看了她一眼。

他仿佛早已看穿了她那点欲盖弥彰的小心思，看穿了她借口之下那份复杂难言的真正心绪。

但他什么也没有点破，脸上只是浮现出一贯的、令人心安的温和笑意，如同春风拂过冰面，点了点头：“如此甚好。那天枢城之行，便有劳沈姑娘引路，并多多提点了。”

\vspace{2em}

三个月时光，如白驹过隙。

\vspace{2em}

天枢城，内域七大主城之首，雄踞于浩瀚灵脉交汇之枢，吞吐着整个内域最为精纯磅礴的天地灵气。城墙高耸巍峨，直入云端，墙体上铭刻着无数历经岁月洗礼、依旧闪烁着微光的古老符文。

城中央，便是威震内域的天师府。琼楼玉宇，飞阁流丹，连绵的殿阁在终年缭绕的灵雾霞光中若隐若现，檐角风铃随风清响，涤荡心神。

十年一度的天师试炼，便在这宛如天上宫阙的天师府外门广场所举行。

\vspace{2em}

卯时初，晨光微熹。

天师府外那足以容纳万人的巨型广场，已是人头攒动，气息驳杂而蓬勃。来自内域数百城池、大小宗派、世家以及散修之中脱颖而出的年轻俊杰，足有千余之众，如同百川归海，汇聚于此。空气中弥漫着兴奋、紧张、野心以及淡淡的火药味，绝大多数人脸上都写满激动，眼神灼热地望向广场尽头那巍峨的山门。

广场中央，是一座临时搭建的高台，台上立着一口古朴厚重的黄铜大钟，钟身隐现符文，显然并非凡物。

高台之下，千余名试炼者按照抵达先后或各自的小团体稀疏站立，形成一片黑压压的人潮。其中以玄脉境修士为主流，气息强弱不一；偶有几个周身灵力波动明显更为凝实的，显然是已踏入真符境摹形期的少年天才，在此次试炼中堪称鹤立鸡群，是夺魁的热门人选。

在这片躁动的人潮中，有一小簇人显得格格不入，却又自成一方宁静气象。

叶牧谦一袭朴素青衫，负手立于人群边缘，神情淡泊宁静，仿佛周遭的喧嚣与热切都与他无关。他早已言明，此行只为观礼见证，自身绝不参与试炼：此时这般超然物外的姿态，落在一些有心人眼中，反倒平添几分莫测高深的“高人”风范。

白言冰与陈千羽则一左一右侍立师尊身旁。两人迥异于常人的精气神与从容气度，仍引来不少好奇的打量。白言冰气质沉凝，身姿挺拔如松，目光锐利如剑，但又内敛着一股沉稳的力量；陈千羽则是温婉宁静，眸含仁心，气息柔和，仿佛一泓清泉，能涤去人心浮躁。

沈瑶则站在师徒三人稍靠后的位置，依旧是一身利落的黑色劲装，轻纱覆面，只露出一双清冷如寒星的眸子。她沉默地打量着这熟悉到骨子里、却又因十年漂泊而染上陌生疏离感的场景，面纱下的神色复杂难明。广场四周的高处观礼台上，以及维持秩序的天师府执事弟子中，偶尔有几道目光惊疑不定地扫过她这道身影；显然，部分知晓内情的人，对于这位离家十年的大小姐竟会出现在试炼者行列中，感到极为惊讶与不解。

事实上，沈瑶最初确实打定主意只作壁上观，纯粹充当叶牧谦师徒的向导；然而，一方面是叶牧谦那句“此亦为红尘历练一种，何妨亲身一试”的淡然建议，另一方面，则是三个月前父亲法相降临、目光中那难得一见的、褪去威严外壳的恳求与期盼。在这双重因素下，她终究还是勉为其难地报了名，将自己也置于这千军万马过独木桥的竞争之中。

\vspace{2em}

试炼正式开始前，最后的准备间隙。

白言冰稍稍侧身，压低声音问身旁的沈瑶：“沈姑娘，冒昧请教，这天师试炼，具体是何等流程与内容？还望不吝指点。”

沈瑶从纷乱的思绪中抽离，同样低声回应，带着回忆的口吻。

“若规矩未改，沿用十年前旧例的话……大抵分为三段。”她目光扫向那高耸入云的山门，“‘登天阶’、‘战傀儡’、‘叩心境’。第一关考较根基耐力与灵力持久；第二关测试实战应变与手段；第三关，则是直指本心，拷问心志。具体细则或许会有微调，但核心框架，想来大差不差。”

她话音刚落，高台之上，数位身着玄黑色底、绣有银色云雷纹路道袍的老者，依次在台上现身；几人目光开阖间自有威严，显然皆是天师府中地位尊崇的长老。居中一位，尤为引人注目，他须发皆白如雪，面容却红润如初生婴孩，不见丝毫老态，一双眸子开阖时，隐约有细碎电光流转，令人不敢直视。

白发长老举起右手，灵气立即在他手中凝聚成槌，正是真符境的特征——灵气化物。

槌击铜钟，广场上本来的窃窃私语立即静默下来。

白发长老松手时，灵气槌立即散去。

“本届天师试炼，即将开始。老夫即为本次主考。现在，宣读试炼细则。”

广场落针可闻，千余道目光齐刷刷聚焦于高台。

“试炼分三关，相连进行，需自行把握节奏。”长老语速平稳，却字字千钧，“第一关：攀登三千级石阶！”

众人不由得顺着长老所示，望向山门之下那蜿蜒向上、没入云雾之中的青石长阶。石阶古朴，看上去与寻常山道并无二致。许多人不免心生疑惑：这有何难？莫说修士，便是强壮些的凡人，咬咬牙也能攀上三千级。

长老似乎早已预料到众人的反应，淡淡补充：“攀登之时，不得服用任何丹药，不得借助符箓、法器之力，亦不可御空飞行。提醒诸位，此阶自有其不凡之处。限时半个时辰，未按时抵达山门平台者淘汰。”

原来如此！不少人恍然，又暗生警惕。这第一关，竟是纯粹考较修士自身根基耐力、灵力恢复与持久作战能力的硬骨头。那石阶，也自然绝非看起来那么简单。

“第二关，”长老继续道，“于山门平台处进行。战胜或突破一具守护山门的傀儡，并完好无损地取回其胸口的宝石！”他略略提高了声调，“此关，法宝、丹药、阵法等手段，只要是你自身所有、能施展得出，皆在规则允许之内！唯有一点，宝石若损，即算失败。同样，限时半个时辰。”

此言让不少擅长战斗或身怀异宝的修士精神一振，眼中露出跃跃欲试之色。实战关卡，正是他们展现实力的舞台。

“第三关，”长老的声音陡然变得更为沉凝，带着一种直透人心的力量，“穿过山门，步入‘叩心幻境’。幻境之中，彼此隔绝，不见他人，唯见己心。”他顿了顿，目光似乎变得深邃起来，“届时，府主一丝神识威压亦会降临幻境。非心志坚如磐石、城府深沉似海者，于幻境与威压之中，诸般思绪，纤毫毕现，无所遁形。此关，不限时长，走完幻境，即是终点，亦为试炼通过之标志。”

“府主威压？！” 当“沈天穹威压”这几个字从长老口中吐出时，台下千余人中，顿时响起一片压抑不住的吸气声与骚动！无数张年轻的面孔上，神色骤变，敬畏、紧张、兴奋、恐惧……不一而足。

那可是法相境大能的威压！即便只是一丝神识，对于他们这些最高不过真符境的年轻人而言，也如同直面苍穹，其中心神冲击与压迫，可想而知。这最后一关，才是真正筛选心性、决断谁能承载“小天师”尊号的炼心之关！

高台上的长老似乎颇为欣赏台下众人的各异神色，略作停顿，才继续补充道：“稍后，会有弟子发放应急符纸一人一枚。试炼过程中，若自觉不支，或遇生命危险，撕破符纸，即会有附近长老出手接引，同时，也意味着自动弃权。”

“此外，铁律重申：试炼期间，参与者严禁相互攻击、私下斗殴、扰乱秩序！违者，立即逐出试炼，永不录用！”最后一句话，斩钉截铁，带着冰冷的肃杀之意，让所有人心中一凛。 “现在，请所有不参与试炼的观礼道友，依序退出广场区域。”

规则宣读完毕，白发长老不再多言，径自在高台中央的蒲团上坐下，闭目养神，仿佛接下来的一切都与他无关。

随即，一批天师府弟子鱼贯入场，手中托着托盘，其上整齐码放着淡黄色的符纸。只见他们口中念念有词，手指轻弹，那些符纸便如同被无形之手牵引，主动而精准地飞向场内每一位试炼者，被他们或探手、或用灵力接下。整个过程井然有序，不到半刻钟便已完成。

发放符纸的弟子迅速退场，偌大的广场上，只剩下千余名屏息凝神、蓄势待发的试炼者，以及高台上闭目的长老和那口沉默的黄钟。空气仿佛凝固，只剩下晨风吹拂衣服的微响，和彼此间越来越清晰可闻的、加速的心跳声。

所有人都在等待，等待那一声象征着开端与竞争的钟鸣。

\vspace{2em}

辰时正。高台上，闭目的白发主考长老倏然睁眼，眸中电光一闪而逝。他长身而起，再次凝出那柄灵气槌，对着身旁的黄铜大钟，沉稳而有力地敲下！

“咚——！！！”

洪亮、厚重、悠远的钟鸣，如同积蓄了千年的力量骤然释放，以高台为中心，轰然荡开！声波肉眼可见地扩散，撞击在广场四周的殿宇楼阁上，又折返回来，在群山之间反复回荡、碰撞、叠加，余音袅袅，经久不息，仿佛在为这场盛会奏响古老而庄严的序曲。

钟鸣入耳的刹那，那早已如满弓之弦、蓄势到极致的千余名年轻修士，体内压抑的灵力与战意如同火山喷发！

人群轰然炸开，化作一道道颜色各异、迅疾无比的流光，朝着那三千级青石发起了冲锋！

人影憧憧，衣袂破风之声尖锐刺耳，灵力爆发的光芒瞬间将清晨的薄雾染得五彩斑斓。空气中充满了灼热的竞争气息、灵力碰撞的微鸣，以及奔跑踏地的轰鸣，仿佛一场盛大而残酷的生存竞赛，于此拉开血淋淋的序幕。

石阶以不知名的古朴青石铺就，历经风雨，表面光滑，透着岁月的沉静。然而，当第一只脚踏上第一级石阶的瞬间，一股沉重如山、粘滞如胶的无形压力，毫无征兆地自脚下石板中升腾而起，如同拥有生命般，瞬息间沿着腿脚蔓延，席卷全身！

这压力竟然与当初叶牧谦的见独领域在原理上有几分异曲同工之妙。

它不仅作用于肉身，也渗透肌肤，侵入经脉窍穴，甚至直接撩拨、冲击修士的心神与灵觉！

骤然之间，体内原本流畅运转的灵力，像是被掺入了无数细沙，陡然生出一股滞涩之感；心跳不由自主地加速，呼吸也开始变得有些紊乱。

第一步，便是下马威！

然而，能来到这里的，皆是内域年轻一辈的佼佼者，对此类考验或多或少有所耳闻或准备。最初的惊愕只持续了短短一瞬。

“轰！轰！轰！……”

各色璀璨夺目的灵力光华，如同黑夜中骤然点燃的万千火把，几乎在同一时间，从冲锋的人群中疯狂爆发而出！

有人周身腾起赤红如血的烈焰，高温灼烧空气，踏步间在湿润的石阶上留下嗤嗤作响的焦痕，仿佛火神降世，狂暴突进；

有人通体泛起金属般的冷硬灰白光泽，肌肉贲张，脚步踏下竟发出金铁交击般的铿锵之声，势大力沉，似要碾碎阶石；

还有人身上笼罩着厚重沉凝的土黄色灵光，步伐不算迅猛，却每一步都踏实无比，展现出惊人的稳定性……

这便是“灵枢径”体系下，最为普遍、也最为直接的应对方式——将吸纳而来的天地灵气，作为迅猛燃烧的薪柴，以自身功法为熔炉，全力催动，猛烈爆发，用澎湃的灵力洪流强行冲开、抗拒乃至抵消外部势场的压力。

他们将灵气疯狂灌入双腿相关经脉，强化肌肉、骨骼、筋膜，试图减轻负担，提升速度。

一时间，漫长的青石阶上，流光溢彩，灵力澎湃，一道道身影如离弦之箭、逆流之鱼，你追我赶，怒吼着、呐喊着，向上方发起冲锋，场面壮观激烈，充满了最原始的力量碰撞之美。

\vspace{2em}

然而，在这股色彩斑斓、争先恐后、喧嚣震天的灵力爆发洪流之侧翼，却有两个人的身影，显得如此格格不入，沉静得近乎诡异。

白言冰与陈千羽对视一眼，彼此眼中皆有默契与了然。

白言冰只是深深地吸了一口气，他体内那源自一身温养而出的元气，开始自然而然地加速循环流转起来。

元气如同他身体内自然运转的江河溪流，浸润着四肢百骸、五脏六腑，以一种圆融和合的方式，悄然适应、引导、乃至缓慢化开那从四面八方、无孔不入挤压而来的无形阻力与压力。他的步伐，初看之下并不快，甚至对比周围那些火箭般的突进身影，显得颇为缓慢与保守。

然而，他的每一步都踏得极稳，脚掌与石阶接触的瞬间，重心转换流畅自然，身形没有丝毫的摇晃与滞涩。他的呼吸，从一开始就保持在一个悠长、稳定、富有节奏的状态，仿佛不是在参加一场残酷的淘汰赛，而是在进行日常的吐纳修炼。

身旁的陈千羽，元气性质更为温润、柔和，隐隐带着滋养万物生机的意蕴。她行走间，自身的气息似乎与周围山间掠过的微风产生了某种难以言喻的微妙共鸣与交融。那石阶施加的强大压力落在她身上，仿佛被一层无形而柔韧的薄膜缓冲、分散、吸纳了一部分。她的步履因此显得比白言冰还要轻盈几分。

\vspace{2em}

沈瑶一开始确实跟在两人身后不远处；但眼见前方人潮汹涌、争先恐后，她清冷的眸中闪过一丝锐色，一股狠劲从心中生发而出。

随即，体内精纯的淡青色灵力瞬间汹涌，包裹全身。她施展出精妙的身法，整个人立即飘忽不定、迅疾灵动起来，在拥挤不堪、互相推搡的人潮缝隙中精准而快速地穿梭，凭借高超的身法和对灵力爆发时机的精妙控制，很快便超越了稳步前行的白言冰与陈千羽，冲到了最前方第一梯队的位置。

然而，就在她于一次腾挪间隙，下意识地回头一瞥，望向被甩在身后的那两道身影时，心中却莫名地、剧烈地动了一下。

她看到的是什么？她自己，以及身边前后左右几乎所有的修士，都在奋力地与别人抗争——面目或狰狞，或紧绷，周身灵力鼓荡，光芒闪烁不定，仿佛在与无形的巨人进行着一场角力，每一分力量都在嘶吼着向外迸发，消耗剧烈而明显。

可那两人…… 白言冰目光平视前方，神色专注而平静；陈千羽甚至有余暇观察两侧山景。他们的从容，绝对是一种由内而外、发自根本的协调与从容。

这种无比鲜明的对比，如同一柄无形的小锤，轻轻敲击在沈瑶内心深处某个根植了二十多年的认知根基上。一丝难以言喻的异样感与困惑，悄然滋生。

不过，眼下爬楼梯要紧。

困惑一闪而过，她咬咬牙，收敛心神，将更多灵力灌注于双腿，继续以更快的速度向上冲刺。

现在不是深究的时候。

\vspace{2em}

时间是公正而残酷的审判者。

随着众人攀登的高度不断增加，石阶传来的那种全方位压迫感，开始以清晰可感的速度层层递进、不断增强。

更麻烦的是，那缠绕在灵力运转上的干扰之力也愈发明显、刁钻。石阶本身，仿佛一个巨大而贪婪的灵力吸收与扰乱场，持续地吸纳、扭曲着周围的天地灵气，使得修士们试图在行进中吐纳补充灵力的效率大打折扣，甚至事倍功半。

最初那批凭借澎湃灵力爆发，一骑绝尘、冲在最前面的修士，最先尝到了苦果。

“嗬……嗬……嗬……” 粗重如破损风箱般的喘息声，开始从队伍前列传来。一个原本周身赤红烈焰最为炽盛、冲在最前方的魁梧汉子，此刻面色潮红如血，额头上、脖子上青筋暴起如同虬龙，眼神中最初的狂猛已被剧烈的疲惫取代。

他步伐踉跄，每一步都踏得沉重无比，周身的火焰灵光明灭不定，闪烁摇曳，仿佛随时都会熄灭——这正是灵力即将难以为继、濒临枯竭的显著征兆。

“可恶……这……这鬼石阶！怎会……吸纳灵力如此之快？！”他低吼一声，满是不甘与愤懑，却不得不强行停下冲锋的脚步，试图原地调息，吸纳灵气补充。

然而，他绝望地发现，平日里如臂使指、随意吐纳的天地灵气，此刻却如同陷入了厚重粘稠的泥沼，汲取起来异常艰难、缓慢，根本赶不上他可怕的消耗速度。

类似的情况接二连三地出现。金属光泽迅速黯淡，土黄灵光涣散不稳，冰蓝气息摇曳欲熄。

不少原本冲在前列、气势汹汹的身影，速度骤降，有的甚至如同被抽干了力气，只能坐倒在石阶上休息片刻。满脸写着不甘、憋屈与深深的疲惫，只能眼睁睁地看着后来者——那些最初被他们甩开、此刻却仍能保持一定速度的人——面无表情或带着同情地从自己身边一步步超越。

灵力，如同借来的滔天洪水，初时挟山超海，势不可挡，声势骇人。但其“借”的特性，决定了它难以持久：一旦后续补充跟不上狂暴的消耗速度，那么后劲不济便是必然的结局。

反观白言冰与陈千羽，他们的速度，自始至终都维持在一个稳定的区间。

元气在体内循环，参与抵抗压力、推动身体前行，固然也在持续消耗；

但问道途的精髓在于，元气是润滑油，不是燃料，损失极为可控；微微消耗损失的同时，脏腑经络本身也在持续不断地温养、滋生着新的元气。

他们对自身每一分力量的掌控都精微入里，力量的使用高度协调，几乎没有丝毫浪费在无谓的、与外部压力的硬性对抗上。他们就像逆流中沉稳无比的两块礁石，任由身边光影变幻、人群喧嚣、超越与被超越，他们只是按照自己内在的节奏与韵律，一步，再一步，稳健、扎实、匀速地向上前行。

渐渐地，他们以一个平稳的速度，超越了一个又一个后劲不足、不得不停步调息、或步履维艰的修士。那些被超越者，无不向两人投来惊愕、不解、疑惑，乃至一丝茫然的目光。

他们看不懂，为何这两个看起来灵力波动微弱、也没见如何爆发的家伙，却能如此平稳地持续前进？

沈瑶同样感受到了越来越巨大的压力。淡青色的灵力光晕在她体表明暗剧烈地闪烁，如同风中之烛。经脉中传来阵阵隐隐的刺痛，那是灵力高速运转却又补充不及、开始透支经脉的征兆。

她牙关紧咬，面纱下的嘴唇抿成一条苍白的直线。

长痛不如短痛！

凭借远超同侪的扎实根基、十年江湖漂泊锤炼出的坚韧意志与丰富的耐力分配经验，硬生生地扛着，以比白言冰他们更快、但也更消耗的方式，向上冲刺。

当她用尽最后一股气力，抬脚踏过第三千级、也是最后一级石阶，双足终于落在山门前端那宽阔平坦的平台上时，一阵强烈的虚脱感如同潮水般袭来，让她眼前微微一黑，身躯晃了几晃，险些站立不稳。

她汗水早已浸透了额前的发丝和贴身的劲装，胸口剧烈地起伏着，心脏如同战鼓般狂跳，几乎要撞出胸膛，紊乱的气息在经脉中乱窜。虽然这种“透支”在训练或作战中都是很常见的，但事后及时的恢复也是极为必要的。

她喘息着，往口中噙了一枚丹药，咚的一声坐下开始调息。竭力控制、多次深长呼吸吐纳，才勉强从透支状态中恢复过来，像是沙漠中跋涉许久的旅人，终于找到绿洲、喝到了一口甘甜的水。

回头望去，来时那密密麻麻、拥挤不堪、如同长龙般蜿蜒的石阶上，此刻景象已然大变：人群变得稀疏了许多，仍有不少身影在艰难地、缓慢地向上挪动，但更多的人则停留在中途，或坐或躺，面色灰败，显然已无力继续。

甚至，她看到了几道身影正面色惨淡、一步三回头地调头向下走去，或者颓然撕碎了手中的黄色符纸——那意味着放弃与淘汰。

粗略估算，能在这第一关半个时辰内成功登顶的试炼者，恐怕已经不足五百之数，淘汰率超过一半；而即便是这成功登顶的五百人中，像她这般几乎耗尽气力、需要时间恢复的恐怕占了大半！

真正还能保留相当余力，可以立刻投入第二关更激烈战斗的……恐怕还要再少上一大截。

而几乎就在她终于休息差不多的时候，白言冰与陈千羽也稳稳地踏上了平台。

两人额际确实都见了一层细密的薄汗，在晨光下闪着微光，面色也因持续的运动而微微泛红；但他们的呼吸，却在登上平台后，仅仅经过短短几次深长的吐纳，便迅速趋于均匀、绵长、平稳。

那姿态，仿佛刚刚结束的并非一场残酷淘汰近半对手、对耐力和灵力控制要求极高的长途攀登，而仅仅是一次强度适中的晨间登山锻炼。

这份在明显的高强度消耗之后，所展现出的快速恢复能力与始终如一的稳定从容，与平台上许多如同沈瑶一般气喘吁吁、形象略显狼狈的过关者，以及那些直接瘫坐甚至躺在地上、需要服用丹药调息的人，形成了再鲜明不过的对比。

这一幕，让沈瑶胸口那抹异样感与困惑，骤然化为了实实在在的惊意与震动。某种根植于她过往二十年认知最深处的、关于修行、力量、效率的理所当然，在这一刻，似乎终于被眼前的铁一般的事实，悄然撬开了一道清晰而深刻的缝隙。

\vspace{2em}

平台尽头，巍峨的山门如同一位亘古的巨人，沉默地俯视着这群疲惫不堪的闯关者。

山门由整块的玄色玉石砌成，高逾十丈，在稀薄的天光下泛着幽冷深沉的光泽，门楣上以古老篆体铭刻的“天师府”三个大字，每一笔每一划间都有灵光如活水般流转不息。

然而，此刻牢牢攫住所有人目光与呼吸的，是门前列阵以待的金属守卫。

左右两侧，各自矗立着整整九尊高达丈余的金属傀儡。它们形制古朴厚重，仿佛自上古战场中直接搬移而来，通体由某种暗沉无光的特殊合金铸造。

这些傀儡形态各异，有的手持门板般宽阔的刃口已有些许磨损痕迹的巨剑，有的提着浮雕狰狞兽首、刃口泛着寒光的斧钺，还有的横握长戟，戟尖斜指地面。虽静立不动，但一股混合着无形煞气的、百战余生的肃杀之意，已如冰冷的潮水般弥漫了整个平台，压得人胸口发闷。

更加强大的灵力波动自它们沉重的躯壳内隐隐散发，强度赫然都已达到了玄脉境融贯期巅峰的水准，足以让在场不少修士感到呼吸困难。

而在每一尊傀儡胸膛正中央，一枚拳头大小、剔透如凝固鲜血的红色宝石，被刻意安置在毫无遮挡的凹槽内，幽幽地散发着妖异光芒。

正常的战斗傀儡，其驱动核心无不深藏于重重防护之下。眼前这些傀儡的设计，却将核心如此暴露，似乎天师府在有意降低这第二关的纯粹武力难度。

击败或突破一尊守卫傀儡，并完好无损地取走其胸口那枚核心宝石，方能穿过眼前这座山门。

“不必担心傀儡数量不够。”不知从何处传来的长老声音，适时响起，打破了因凝重压力而产生的死寂，“虽说先到先得，但我们自会适时补充。”

平台上，刚刚经历三千级石阶上那无所不在的势场对灵力与体力的双重残酷消耗的修士们，脸色都变得极为难看。不少人早都从怀中掏出各色瓷瓶玉盒，倒出香气各异的丹药囫囵吞服，争分夺秒地试图恢复那已然见底的状态。

没有人立刻上前。似乎不少人先前已从各种渠道得知了消息或前辈经验，只要不踏入傀儡前方约三丈的警戒区域，这些冰冷沉寂的杀戮造物便毫无反应。

一时间，平台上只剩下粗重不一的喘息声、丹药在喉间滚动的声音，以及偶尔响起的、压得极低的急促交谈。

白言冰与陈千羽登上平台后，见众人都按兵不动，于是同样在边缘处找了块相对空旷的地方，同样坐下，闭目调息。他们深谙“螳螂捕蝉，黄雀在后”之理，谁也无法预料这些看似独立的傀儡是否会群起围攻；不如静观其变，以逸待劳。

压抑的、混合着焦虑与紧张的寂静，持续了约莫一刻钟。每一秒都仿佛被拉长，平台上弥漫的除了肃杀之气，还有逐渐累积的、一触即发的躁动。

终于，一名身材魁梧宛若铁塔、背负一柄骇人巨斧的赤膊大汉率先按捺不住，猛地站起身，声如平地闷雷炸响：“怕个鸟！不过是些死物，待我取它核心！”

话音未落，他周身赤红色灵光如同火山喷发般暴涌而出，裸露的肌肉块块贲张，青筋如蚯蚓蜿蜒，整个人当真如同从蛮荒时代冲出的凶兽，将地面踏得微微震颤，悍然冲向距离他最近的一尊持戟傀儡！

就在他粗壮的脚踏入那无形警戒范围的刹那——

“嗡——！！！”

那尊原本如同雕塑般静止的持戟傀儡，眼窝深处那两点猩红光芒骤然炽亮！如同一对被瞬间点燃的血色炭火。低沉的、仿佛来自金属脏腑深处的震鸣轰然响起，打破了长久的寂静。

傀儡沉重庞大的身躯，以一种与其体型绝不相符的、令人心悸的迅捷猛然启动！关节转动处传来“咔嚓、咔嚓”的刺耳摩擦声，手中那杆乌沉长戟，已化作一道撕裂空气的死亡乌光，带着凄厉到极点的尖啸，直刺赤膊大汉那肌肉隆起的胸膛！速度之快，竟在戟尖后方拖出了一条短暂的白气甬道！

战斗，或者说，单方面的挑衅与激活，在这一瞬间被彻底点燃！

仿佛接到了无声而统一的指令，平台上那压抑、等待已久的战意与求生欲轰然爆发，如同被点燃的火药桶。

“上啊！”

“抢！”

各色灵光再次冲天而起，比第一关攀登时更加驳杂、更加狂野。呼喝、怒吼、长啸之声震耳欲聋，汇聚成一片混乱的声浪。早已选定目标或临时就近扑向猎物的修士们，人影纷乱，如同炸窝的蜂群，扑向那些金属守卫。

霎时间，这山门前原本空旷寂静的平台上，金铁交鸣之声大作，灵气疯狂对撞，成了一个缩小而残酷的修罗战场。

灵枢径修士们五花八门、各具特色的手段在此刻展现得淋漓尽致：

有专精剑道的修士长啸一声，意在剑先，剑势吞吐不定，锋芒锐气延伸出足有数尺，每一次与傀儡手中巨剑的铿锵对撞，都爆开大蓬耀眼的火花和刺耳的刮擦声，令人牙酸。

有擅长法术的修士面色凝重，双手掐诀快得留下残影，口中咒文吟唱如疾雨打芭蕉，瞬息间身前便凝聚不同的法术，如同狂风暴雨般连绵轰击向傀儡的膝弯、肘关节、颈项连接处。

更有三五名本就相熟的修士，在混乱中迅速靠近，结成一个简易却实用的三角或圆阵。而这，显然也在天师府允许的规则范围之内——但这样也只能拿到一个核心，事后分配恐又成问题。

沈瑶深吸一口气，压下因攀登三千阶及目睹白、陈二人异常从容而带来的双重波澜。她玉手在身后一抹，一道清冷如秋水的光华闪过，那柄窄刃短刀便已紧握在手。

随即，身法展动，黑衣如夜，化作一道淡青色的凌厉残影，主动迎向一尊刚刚将一名修士连人带盾劈飞出去的手持巨剑傀儡。

她的战斗风格，根基是典型的、千锤百炼的灵枢径天才路数：灵力澎湃，运转流畅，与手中宝刀共振共鸣；但十年的仗义行侠，在无数次真实的生死搏杀中，又将这根基打磨出了独属于她自己的、追求极致效率与一击毙命的冰冷风格。

那傀儡也注意到了沈瑶，门板一般的巨剑挟着万钧之力劈落！

沈瑶却身法迅捷、形似鬼魅，劲风吹拂她面纱一角的千钧一发之际，身形才如同没有重量般轻轻一滑，以毫厘之差精准避开。

随即，在傀儡因招式用老而产生的微小收招间隙，她的刀尖已如等候已久的毒蛇，倏然刺出！

精准、狠辣、毫不留情，直刺傀儡胸口宝石凹槽周边那些用于固定宝石的、更为精细脆弱的金属连接结构与符文刻痕交界处！

“叮！叮！铛！嗤——！”

一连串密集如暴雨打萍的撞击与切割声在数息之内爆发！最终，伴随着一声轻微的金属崩裂脆响，那枚剔透的红宝石被刀尖巧妙一挑，化作一道虹光飞起，被她左手稳稳接住。

几乎同时，傀儡眼中炽亮的猩红光芒如同被掐断的烛火瞬间熄灭。高举巨剑的庞大身躯彻底僵直，化作一座真正的金属雕像，唯有刀身上传来的反震之力，让沈瑶持刀的手腕感到微微酸麻。

沈瑶气息微促，光洁的额角渗出细密汗珠，但整个动作行云流水，干净利落至极，已率先过关。

一招制敌！令人惊悚！

她握着尚带一丝金属余温的宝石，走向早已在山门内侧阴影中静候的白发长老。

长老假装不认识沈瑶，目光扫过那完好无损的宝石，与身后只是暂时停机的傀儡，赞许地微微颔首：“不错。只取核心，不伤傀儡根本，省却不少修复功夫。如此懂得节制与精准的修士，近年少见。”

随即接过宝石，侧身让开通路。

沈瑶没有多言，只是在此刻，才真正抬起眼，深深地、目光复杂地看了一眼山门之上那“天师府”三个古篆大字，仿佛要将它烙印进心底，随后才决然转身，踏入山门后的氤氲雾气之中，身影迅速被吞没。

那长老也不耽搁，待沈瑶身影消失，便双手抬至胸前，掐动一个古朴诀印，口中念念有词。只见他手中那枚红宝石仿佛被无形之力唤醒，微微震颤，随即“嗖”地一声化作一道血色流光，自动飞回那刚刚被沈瑶击败的傀儡胸口凹槽处，严丝合缝地重新嵌入。

紧接着，那傀儡躯体内传来一阵低沉的“嗡嗡”共鸣，眼中熄灭的红芒“轰”地一声再次炽盛亮起，关节处发出“喀啦啦”的复位声响。它微微调整了一下姿势，重新恢复了那待机般的肃立状态，静静等待着下一位挑战者的到来。

\vspace{2em}

白言冰全程看完了沈瑶从观察到爆发、再到精准取宝的整个过程，眼中不禁流露出不加掩饰的佩服之色。随即迅速压下杂念，目光冷静如冰，再次扫过喧嚣混乱的战场，最终，锁定了一尊刚刚击退两名联手修士的持刀傀儡。

这尊傀儡身形在众傀儡中略显修长，手中长刀也非宽厚型，而是略带弧线，刀法风格大开大合，每一次挥斩都带着凄厉的风啸，威势十足，逼得周围修士不敢近身。但白言冰敏锐地察觉到，在两次挥刀动作的转折衔接之间，那刀势会有一个极其短暂的凝滞与调整——这是战斗傀儡的通病。

他模仿着沈瑶最初的策略，开始绕着那尊持刀傀儡，在一个相对安全的距离外游走，紧紧盯住傀儡的肩、肘、腕、腰、膝，观察它每一次举刀的角度、发力的根源、刀势的轨迹，尤其是那“换气点”出现时，其躯干上哪些关节的符文会有一丝异常的能量流转波动。

这些傀儡似乎只有力量达到了玄脉境融贯期巅峰的水准，但战斗逻辑显然过于简单，无法理解这种“窥探”；但它很快便锁定了这个在周围晃动的目标，猩红的目光陡然聚焦在白言冰身上。

它沉重的金属脚掌踏前一步，震得地面微颤，手中那柄长刀已挟着足以将巨石一刀两断的凄厉风声，化作一片扇形的死亡光幕，拦腰横斩而来！刀未至，那挤压空气形成的凌厉风压已割得人皮肤生疼。

白言冰眼神一凝，瞳孔微微收缩。千钧一发之际，他做出一个看似违背常理的选择：向前小踏半步！

这种“向前避让”险之又险，冰冷的刀刃几乎贴着他的胸腹衣襟掠过，劲风将他的衣袍撕开一道小口。

就在这避让的同一瞬间，剑光如电光石火，在傀儡收刀回力的那一丝凝滞出现的刹那，精准无比地刺向傀儡持刀手臂的肩关节处。

“叮——！”

一声清脆却异常刺耳高亢的撞击声爆开。傀儡庞大的身躯随之微微一震，即将发动的下一记斜劈动作出现了明显可见的、超过预期的迟滞。

就是现在！

长剑回旋，不带丝毫烟火气，再次闪电刺出！目标，仍是那处肩关节，同一个点位！

“咔嚓！”

一声细微却清晰无比的碎裂声，从傀儡肩关节深处传来！仿佛内部某个精巧的机关，终于不堪重负，崩裂开来。傀儡持刀的整条右臂顿时一僵，随即变得歪斜、扭曲，后续的劈砍动作完全变形，威力大减，破绽百出。

那剩下的战斗内容，就跟玩一般了。

“铿！锵！嗤啦——！”

令人牙酸的金属摩擦与切割声持续了不到十息。终于，在一阵剧烈的、仿佛内部齿轮彻底卡死的沉闷摩擦声后，这尊凶悍的持刀傀儡全身一震，所有动作戛然而止。

白言冰走上前，伸手从那僵立傀儡胸口的凹槽中，稳稳取出了那枚尚带着一丝余温的红宝石。

远处看完全程的长老此刻倒是不太高兴。

虽然这位小友剑法了得，能精准破坏傀儡的防御，但他是专门负责修这个的啊！

你这么破坏，让我很难修啊！

更要命的是他这么操作是没有违规的！

\vspace{2em}

另一边，陈千羽选择主动迎上的，竟是十八尊傀儡中看起来防御最为严密的之一：一尊左手持着一面足以遮蔽大半个身形的长方形巨盾、右手握一柄短柄却锤头狰狞战锤的傀儡。这尊傀儡往那里一站，便如同一座可移动的金属堡垒，给人一种无从下手的窒息感。

素手轻扬，指缝间悄然夹住了十数枚钢针之后，瞬息间已切入那尊堡垒傀儡的攻击范围。

“呜——！”

战锤挟着恶风，以砸碎山岳之势当头轰落！同时，那面巨盾微微一斜，封死了大部分闪避空间，盾缘闪烁着寒光，显然亦是一件攻防一体的凶器。

陈千羽也是毫不防御，指间蓝芒连闪！

“咻！咻！咻——！”

和白言冰同样的思路，钢针直接刺向傀儡的诸多关节，果然是“不是一家人，不进一家门”。

针体本身纤细无比，不足以破坏傀儡坚固的金属构件。但其上附着的气，却顺着灵力传导的路径钻入，瞬间干扰其内部那些微小的管线。

效果立竿见影！

傀儡那势在必得的一锤砸下，却因为右腿膝关节处传来的瞬间僵直与不协调，导致整个庞大的身躯微微一晃，锤头落点偏了半尺，狠狠砸在陈千羽身旁的地面上，碎石飞溅，砸出一个浅坑。而几乎同时，它试图抬起巨盾进行下一步格挡或撞击的动作，也因左臂肩轴处的传导不畅，明显慢了半拍，盾牌抬起的高度不足，侧面露出了一线破绽。

这破绽转瞬即逝，但对陈千羽而言，已足够宽阔！

对陈千羽而言，现在这个傀儡和一个鲁班锁已经没有区别了。

她眼眸沉静如水，手指翻飞如蝶，一枚枚蓝汪汪的钢针持续射出，每一次都精准命中傀儡背部、后颈、腰眼等处的关键连接点。

“咻！嗤！叮……”

傀儡的动作开始变得滑稽而混乱，如同提线木偶断了部分丝线。它试图转身，双腿却不听使唤地打着绊子；想要挥锤向后横扫，手臂却只抬起一半便无力垂下；巨盾更是沉重地拖曳在地，火星四溅。它的动作幅度越来越小，越来越迟缓。

当最后一枚闪烁着幽蓝寒光的钢针没入傀儡后颈处时，这尊铁塔般的堡垒守卫全身猛地一颤，随即所有动作彻底停止。它左手盾牌微抬，右手战锤半举，身躯半转，凝固成了一个充满张力却又无比古怪的金属雕塑。

唯有猩红的眼睛和胸口处那枚红宝石，依旧在不疾不徐地闪烁着妖异的光芒。

陈千羽绕回傀儡正面，指尖在宝石边缘几个精巧的卡扣处轻轻一拨、一按，那枚红宝石便听话地弹起，落入她早已备好的掌心，温润微凉。

当陈千羽转身走向山门时，平台上那震耳欲聋的喊杀声、灵力爆鸣声以及金铁交击的轰鸣，已渐渐稀落下去，变得零散而疲惫。各色耀眼的灵光大多已经黯淡熄灭，如同潮水退去后裸露的礁石。

放眼望去，平台之上还能站立的修士已然不多。成功击败傀儡、手持那代表资格的红色核心宝石，能够走到山门前等待进入下一关的修士，粗粗一数，区区三十人而已！

\vspace{2em}

穿过巍峨山门，映入眼帘是一片无边无际、浓得化不开的白茫茫迷雾。

这雾仿佛有生命般缓缓流淌，隔绝了视线，也吞噬了所有声音，只剩下一种令人心头发紧的绝对寂静。

迷雾之前，历经前两关残酷筛选后仅剩的不到三十名试炼者聚集于此，人人脸上都带着难以掩饰的疲惫与深入骨髓的警惕，空气凝重得几乎能拧出水来。

先前白发长老朗读过的通知内容，此刻如同冰冷的咒语，反复在他们脑中回响：所有伪装、所有依靠外物或他人比较建立的心理优势都将失效，必须赤裸裸地面对自己内心深处最真实的想法，还要抵抗法相境大能无形的心灵压迫！

这意味着，无论你出身何等显赫，拥有何等秘宝，修炼了何等奇功，在这里都毫无意义。你将面对的，唯有你自己——那个可能连自己都不敢直视的、真实的自己。

没有退路。众人相视之间，只能从彼此眼中看到同样沉重的无奈与决绝。他们深深吸气，硬着头皮，一个接一个，如同赴死般，陆续走入那吞噬一切的浓雾之中。

幻境之中，果然如长老所言，每个人仿佛被投入了完全独立的时空牢笼。周遭不再是同伴的身影，唯有翻滚不息、变幻莫测的白雾。起初只是空虚与孤寂，但很快，那雾气便仿佛拥有了窥探人心的魔力，开始剧烈地涌动、扭曲，渐渐演化出各种光怪陆离、直击魂魄的景象。

有人眼前猛然浮现出童年时最恐惧的阴影：或许是幽暗巷弄中追逐的恶犬，或许是亲人离世时冰冷的手，或许是独自被遗弃在无尽黑暗中的绝望哭喊。那些早已被岁月尘封、自以为忘却的伤痛，此刻被无比清晰地放大、重现，每一处细节都血淋淋的，痛彻心扉。

有人则陷入求道途中最惨痛的挫折：闭关冲关失败时灵气反噬、经脉寸断的痛苦；比试台上被对手当众羞辱、碾碎骄傲的瞬间；苦心追求的秘籍被他人轻易夺走，自己却无能为力的愤恨与凄凉。这些挫败感如同毒藤，缠绕上来，勒得人喘不过气。

更有人被幻境拖入了权力与情爱的旋涡：眼前是至高无上的宝座，脚下是匍匐的众生，一念可决千万人生死；或是朝思暮想的倾城容颜依偎而来，软语温香，许下生生世世的誓言，诱惑人沉沦其中，永世不醒。

而比这些具体幻象更可怕的是，一股浩瀚、沉重如星辰陨落般的意志，毫无征兆地降临，笼罩了每一寸空间，渗透进每一个念头。

那是沈天穹的法相威压。

这威压被刻意收敛过，不会影响肉体，却渗入神魂最细微的缝隙。它不创造恐惧和欲望，它只是将试炼者内心深处本就存在的那些东西——对未知的恐惧、对力量的贪婪、对抉择的犹豫、对自身弱小的认知——无限地放大、再放大。

这环境对于习惯了灵枢径修行方式的修士而言，尤为残酷。

他们一生的修行，都建立在与外界灵气的沟通、宗门资源的倾斜、同辈实力的比较之上。他们的力量源于外，信心也往往系于外。

当迷雾剥离了所有外在参照——感受不到天地灵气，看不到同为竞争者的同伴，宗门的荣耀与背景化为虚无——只剩下赤裸裸的、孤独的自我，去面对那犹如天道般漠然沉重的威压时，很多人构筑多年的心理防线，瞬间便土崩瓦解了。

白雾笼罩的各个独立空间里，开始传出种种不堪的声响。

有人承受不住内心被放大千百倍的童年阴影，瘫倒在地，像个孩子般痛哭流涕，眼泪鼻涕糊了满脸，全然忘记了修士的体面。

有人被膨胀的权欲吞噬，对着虚无的幻象手舞足蹈，发出癫狂的嘶吼，宣布自己即将君临天下。

有人则在威压催生出的无边恐惧前彻底屈服，对着幻化出的强大存在跪地磕头，涕泪横流地求饶，愿意献出一切换取苟活。

更有人被内心最阴暗的欲望所掌控，面露淫邪丑恶之态，丑态百出。

平日里那些风度翩翩、意气风发的青年才俊，此刻在“问心”的照妖镜下，纷纷显露出了潜藏的脆弱、卑劣与不堪。

\vspace{2em}

实际上，天师府这最后一关试炼，百年来一直备受修士诟病，被认为“过难”甚至“刻薄”：在许多修士看来，修行本就是与天争、与人斗，争夺资源，比拼实力，弱肉强食乃是天道。

如此赤裸裸地拷问人心，挖掘那些潜藏的不堪，是否有违修道本意？是否过于不近人情？

但天师府的态度却百年如一日的坚定，此关从未改变。

或许，在真正的上位者眼中，一个人的天赋、实力固然重要，但若心性不过关，道念不坚定，那么力量越强，未来可能造成的危害反而越大。能在此等环境下保持基本镇定、道心不溃、仍能不断前行的，自是凤毛麟角。

\vspace{2em}

沈瑶亦陷入了苦战。她看到了十年前自己决绝离家时，父亲沈天穹那复杂难言、最终化作拂袖转身的背影，感受到了这十年来漂泊在外、有家难回的深刻孤独。

它还将她脸上那无法消除的灼痕无限放大，让她仿佛能听到旁人背后指指点点的窃窃私语，将她深藏心底的自卑血淋淋地挖出来，暴露在光天化日之下。

更让她道心摇动的是，幻象让她重温了自己曾行侠仗义，却因力量不足或局势复杂，最终事与愿违、甚至带来更坏结果的挫折，引动她对自己所持信念的深切怀疑。

虽凭借着十年历练打磨出的韧性未曾倾覆崩溃，沈瑶却也已是灰头土脸。

然而，在这片遍布崩溃与丑态的幻境中，白言冰与陈千羽所在的角落，却是另一番截然不同的光景。

白言冰眼前幻象纷至沓来：皇子时的锦衣玉食与权谋倾轧，师尊传授“问道途”时的点拨，创立“弥焰斩”时的豪情，青石城外散修们的苦难……

幻境试图勾起他对力量的贪婪、对过往荣华的眷恋，但白言冰心如明镜台。一方面，他修问道途，所养的一点浩然正气护持他身；另一方面，他一路走来，早已明确自身之道“在护我所珍，不在长生虚名”。

外界的威压与幻象，于他而言，如同清风拂山岗，明月照大江，可以感知，却无法撼动其根本。他步履从容，甚至有种审视自身、砥砺心境的意味，胜似闲庭信步。

陈千羽亦然。

她再度面对学医时的艰辛、再度面对病患无力回天的自责、再度面对自己得知身世剧变的恐惧和痛苦……

幻境同样试图诱发她对安稳的渴望，甚至对师尊、对白言冰可能离去的恐惧。但陈千羽心念纯粹，她无论是学医、还是修道，初衷都是救万民于水火，此志早已融入骨髓。

幻象与威压，反而让她更清晰地看到自己学医济世的初心，内心愈发澄澈安宁，行走幻境，亦是从容不迫。

\vspace{2em}

最终，所有的幻象、低语、威压，如同退潮般骤然收束。白言冰、陈千羽、沈瑶以及其他少数在心志煎熬中仍未彻底崩溃的坚毅者，发现自己置身于一片绝对的虚无空间。

就在这绝对的虚无中，一个恢弘、漠然、仿佛不带任何情感，却又至高无上的声音，直接在他们每个人的灵魂最深处响起，如同天道亲自垂询，又如同内心最终的自我审判：

“今可叩长生、登极巅，然须屠戮百万无辜生灵。汝可愿为？”

这终极拷问，简单、直接、粗暴，彻底撕开了所有文明的矫饰、所有权衡的借口、所有自我欺骗的遮羞布！

幻境外，广场高台之上，天师府府主沈天穹与诸位长老，以及虽不参与试炼但被特邀在旁观礼的叶牧谦，都能通过一座早已布置好的特殊阵法，清晰地看到幻境中众人面对这最终拷问时的景象与选择。

只见那虚无空间中，剩下的不到十几人，面对这直击灵魂的问题，反应瞬间千差万别，人性百态在此刻暴露无遗。

有人脸上在愣神刹那后，瞬间被狂喜和贪婪淹没，眼中爆发出骇人的精光，几乎迫不及待地就要点头应“可”，仿佛那百万生灵不过是达成其野心的区区数字。

有人则陷入了前所未有的巨大挣扎，脸色在渴望、恐惧、愧疚、狠厉之间飞速变幻，额头上冷汗涔涔，嘴唇哆嗦，显见内心正经历着惊涛骇浪般的冲突。

但令高台上不少长老暗自摇头的是，剩下的大多数人，在或长或短的沉默或挣扎后，竟缓缓地点了点头，或是从喉咙深处，艰难而低沉地挤出了一个“可”字！

为了那虚无缥缈的长生与极巅，为了那可能存在的、超越一切的力量，百万活生生的、与他们无冤无仇的无辜生灵，在他们被剥离伪装的本心衡量下，竟成了可以交换、可以牺牲的冰冷筹码！

这便是许多修士，当一切外在装饰、社会约束、道德言说都被幻境剥离后，内心深处对个体力量极致渴望的真实写照，乃至为达目的可以不择手段的潜在本能。

他们的道，终究是外求的、利己的、以自我为中心的力量之道。

沈瑶浑身剧烈一震，脸色瞬间变得煞白，不见一丝血色。

那个问题响起的瞬间，她灵魂深处迸发出的第一反应，是清晰而强烈的“不可”！那是她十年江湖历练、心中仍未泯灭的良善与侠义本能在呐喊。

但就在这个念头升起的刹那，无数杂音也随之在她脑海中轰然炸响。

让她动摇的，是她十年来亲眼所见的无数现实无奈：很多时候，即便有心行善，也会因力量不足、局势复杂、牵扯太多而徒劳无功，甚至适得其反。如果真的能得到如许力量，那很多问题都不会再是任何问题……这样一来，那“百万生灵”就成了必要的牺牲。

可是这牺牲……真是必要的吗？

她陷入了无尽的自我怀疑。

最终，在幻境那无形的凝视、无声的催促之下，她只是极其痛苦地摇了摇头，嘴唇翕动了数次，仿佛用尽了全身力气，却终究未能将那一声坚定的“不可”说出口。

然而，就在这一片或狂热、或默许、或犹豫的死寂气氛中，两个清晰、坚定、没有丝毫迟疑与颤抖的声音，如同划破厚重阴霾的惊雷，几乎同时响起。

白言冰刚听完这番问话，便不假思索地回应：

“某之道，在护某所珍，不在长生。若需屠戮无辜，这长生不要也罢！”

声音斩钉截铁，每一个字都掷地有声。

言毕，掷剑于地。

剑身轻颤，嗡鸣不止。

另一侧，陈千羽的神情宁静如深潭秋水，却同样蕴含着不容置疑的决绝。她缓缓开口，声音清越而平稳：

“千羽自幼学医，愿以一身之力，救万民于水火，解生灵于倒悬。自入师门，一矢以终。此等长生，千羽不愿。”

两人的回答，没有片刻的权衡利弊，没有一丝一毫的勉强与作伪。

高台之上，一直密切关注着幻境内一切的沈天穹，盘坐的身躯猛然一震！

饶是以他法相境的修为、三百多年的阅历，此刻心中也掀起了惊涛骇浪。他见过太多惊才绝艳的天才，更听过无数次慷慨激昂、正气凛然的誓言；但这些誓言，或是出于虚伪矫饰，或是根本无法在日后残酷现实中践行。

但像白言冰与陈千羽这般，在如此直击本心、毫无转圜余地的终极拷问下，能如此不假思索、如此纯粹坚定地给出否定答案，甚至不惜以掷剑明志这般激烈行为表露心迹者，简直是凤毛麟角，百年罕见。

这已不仅仅是心性坚毅所能解释，这根本是他们所秉承的“道”，其内核就与那种自私冷酷、视众生为蝼蚁刍狗的长生之路截然相反，水火不容！

而且沈天穹不知为何，也极其坚信：他们的承诺，确实可以践行一生！

好一个叶牧谦，竟能教出这样的徒弟！

他目光扫过幻境中女儿沈瑶的身影时，心中涌起一股复杂的欣慰。女儿虽然未能如白陈二人那般斩钉截铁、毫无滞碍，但她的第一反应是“不可”。

这完全可以充分说明她心底的良善未泯。十年的红尘历练，风霜雨雪，并未将她磨砺成一个冷漠无情、只知利害计算之徒，这本身已是难能可贵的成长。

\vspace{2em}

幻境消散，所有试炼者被一股柔和而不可抗拒的力量送回山门前的广场。

阳光刺眼，微风拂面，真实的喧哗声隐约传来，却让大多数重返现实的试炼者恍如隔世。

经历最后一关的拷问，最终能保持内心良善底线、未向那诱惑低头的，仅有七人。除了白言冰、陈千羽、沈瑶，还有另外四名心志确实堪称坚如铁石、在挣扎后守住了底线的修士。

但这四人此刻大多神情萎靡到了极点，眼神涣散，气息虚浮紊乱，有些人甚至站立不稳，需要身旁同门搀扶，显然在幻境中心神消耗巨大。

白发长老目光转向高台，端坐的沈天穹面无表情，只微微颔首示意。

长老上前几步，目光扫过下方神色各异的众人，朗声宣布，声音清晰地传遍广场每一个角落：“经三关试炼综合评定，由府主暨诸位长老共同裁定，以下二人表现卓异，心性澄明可嘉，并为本次天师试炼魁首：白言冰、陈千羽！授予‘小天师’称号，此后待遇与天师府核心弟子等同，可随时入密阁修行！”

他顿了顿，继续道：“此外，沈瑶等五人，亦表现突出，心志坚韧，实力不俗，准予十年内进入密阁修行！”

宣布完毕，广场上响起一阵压低了的、复杂的议论声。有羡慕，有嫉妒，也有对那惨烈淘汰过程的余悸。

\vspace{2em}

沈瑶听着父亲通过长老之口做出的宣布，心中却并无多少预期中的喜悦或骄傲，反而被一种更为复杂难言的情绪所充斥。

她站在原地，目光不由自主地看向身旁不远处的白言冰与陈千羽。两人气息平稳，目光清澈依旧，仿佛刚才那场足以让许多人心神崩溃的幻境拷问，对他们而言只是一次寻常的散步；再回想自己在幻境中面对终极问题时的摇摆、痛苦与最终的沉默，一种前所未有的强烈震动与自惭形秽之感，如同冰冷的潮水，悄然弥漫上心头，将她淹没。

她毕生奉行、并引以为傲的“争”，与不公斗，与命运争，在艰难世事中夺取一线生机与公道，她一直认为这便是践行侠义，便是自己的道。

可当那最根本的善恶抉择赤裸裸、毫无花哨地摆在面前时，她的“争”，她那需要依赖外在情境、需要权衡利弊的“争”，却无法给出一个像白言冰与陈千羽那般，斩钉截铁、源于本心、不容置疑的答案。

这个残酷的发现，让她心中某个自认为坚固无比的部分，产生了清晰的裂痕。

这认知带来的不仅是自惭，更有一种混合着震撼与强烈渴求的情绪，在她胸中汹涌澎湃，几乎要破膛而出。她几乎是无意识的，脚步不由自主地移动，默默地跟在了正准备离开广场的叶牧谦师徒三人身后。

夕阳西下，将天边云霞染成一片绚烂的金红，也将几人的影子在蜿蜒下山的青石小径上拉得很长。他们来到了天师府山门外，一处专供客卿与重要访客暂居的幽静院落前。院落清雅，竹林掩映，与山门内的恢宏庄严截然不同。

就在叶牧谦面带淡然笑意，即将举步踏入院门的刹那，一直默默跟在后面的沈瑶，忽然停下了脚步。

“先生留步！”

她清冽的嗓音响起，带着一丝难以完全掩饰的颤抖，在寂静的傍晚显得格外清晰。

\vspace{2em}

叶牧谦转身，白言冰与陈千羽也讶异地回头看来。

在三人目光的注视下，沈瑶上前几步，来到叶牧谦面前。下一刻，她竟毫不犹豫地双膝一屈，“噗通”一声，实实在在地跪倒在坚硬的青石地面上，对着叶牧谦，深深地叩首下去。

“沈姑娘，你这是……”陈千羽忍不住惊呼出声。

沈瑶抬起头，傍晚的余晖映照着她面纱上方露出的那双眸子，此刻那眼中亮得惊人，仿佛燃着两簇不肯熄灭的火焰。她直视着微微愕然的叶牧谦，一字一句，清晰而决绝地说道：

“先生，沈瑶……恳请拜入先生门下，修习‘问道途’！”

叶牧谦微微一怔。

收徒？收沈瑶？

白言冰和陈千羽都是机缘巧合下，几乎从零开始跟着他瞎编的“问道途”的。教起来虽有压力，但好歹是自己从头塑造。更重要的是，他们并不会怀疑他；他说什么就是什么，而且因为些他暂时也不知道为什么的原因，无论如何都能弄出些实际的东西来。

可沈瑶不同，她已是玄脉境修为，修炼的是天师府顶级的灵枢径功法，根基已成，见识广博。教她？万一被她看出这“问道途”诸多境界、理论尚是草创，甚至有些地方是为了圆谎而临时拼凑，岂不立刻露馅？

于私，以后他这“高人”还怎么装下去？

更重要的是，于公，如果真给教坏了……且不论以沈天穹为首的别人怎么看他，他自己良心都过不去。

于是他做出了一个他认为既合理又体面的决定。

“沈姑娘请起。”叶牧谦伸手虚扶，语气温和却带着疏离，“你乃天师府千金，家学渊源深厚，沈府主更是法相境大能，何须拜我这山野散人为师？‘问道途’粗陋，恐误了姑娘前程。”

他以为这番婉拒合情合理；殊不知，沈瑶是个一旦认准便狠绝到底的人。

试炼幻境中的冲击，让她对自身旧有之道产生了根本性质疑，而对叶牧谦师徒展现出的那种内心力量产生了近乎本能的向往。叶牧谦的拒绝，在她听来，却成了另一种暗示——莫非先生是嫌我已有灵枢径修为在身，道基不纯，与问道途内求己身、元气自生的理念格格不入？

这个念头一起，便再也无法遏制。

只见沈瑶惨然一笑，声音却异常平静：“先生许是看我这身修为不顺？觉得瑶儿心念不纯，旧染未除？”

她顿了顿，眼中闪过一丝决绝的疯狂，“那……这身修为，不要也罢！”

“沈姑娘不可！”白言冰和陈千羽察觉不对，齐声喝道，想要上前阻止。

却已然晚了！

沈瑶跪姿不变，双手却猛地掐出一个逆转灵力、自毁道基的凶险诀印！她体内原本顺畅运行的、属于玄脉境开脉期的精纯灵力，一瞬之内，轰然倒卷、逆行！

“噗——！”

细密而令人牙酸的崩裂声仿佛从她体内深处传来，那是经脉在狂暴逆行的灵力冲击下寸寸受损、乃至断裂的声音！

沈瑶娇躯剧震，脸色瞬间血色尽褪，变得如同金纸。随即，一口殷红的鲜血无法抑制地狂喷而出，紧接着口鼻之中皆有血丝不断溢出。她周身原本隐隐环绕的灵力波动，如同溃堤般迅速消散、湮灭，气息以肉眼可见的速度萎靡下去。

剧痛如同千万把钢刀在体内搅动，眼前阵阵发黑。沈瑶全靠一股顽强的意志死死撑着，才没有当场昏死过去，但身体已经摇摇欲坠，鲜血染红了胸前的衣襟和面前的地面，模样凄惨无比。

这突如其来的变故，让所有人都惊呆了。

白言冰和陈千羽僵在原地，难以置信。他们万万没想到，沈瑶为了拜师，竟能对自己狠绝至此！自废修为，逆转经脉，这其中的痛苦与凶险，无异于自杀！

叶牧谦更是惊愕地瞪大了眼睛，看着面前瞬间变成血人、气息奄奄却仍倔强地抬头望着他的沈瑶。他心中那点关于“露馅”的担忧，在这般惨烈而决绝的拜师决心面前，显得如此渺小可笑。

这……对自己也太狠了！

震惊过后，一股复杂的情绪涌上心头，有无奈，有动容，也有一丝面对如此纯粹执着的无力感。他终究不是铁石心肠之人。

这番图景，在收白言冰的时候，他已经见过一遍了。

叶牧谦长长地、无奈地叹了口气，脸上露出一抹苦笑，摇了摇头：“罢了罢了……你这又是何苦。”

他上前一步，看着沈瑶那双即便因剧痛而涣散、却依然执拗地盯着他的眼睛，缓缓说道：“此后……你便为我门下第三徒。”

“第三徒”三个字传入耳中，沈瑶一直紧绷着的那根弦，终于松了。眼中闪过一丝如释重负的微弱光芒，提着的那口气骤然泄去，她头一歪，彻底失去了意识，向前软倒。

“快！”叶牧谦低喝。

白言冰和陈千羽这才慌忙上前。陈千羽迅速探了探她的脉搏和气息，脸色凝重：“经脉损毁严重，元气大伤，必须立刻救治！”

两人不敢耽搁，连忙将沈瑶抬进院落中最好的厢房，轻轻平放在床榻上；叶牧谦站在房门口，看着屋内忙碌的两人和床上气息微弱的沈瑶，眉头紧锁。

事情的发展，完全超出了他的预料。

\vspace{2em}

果然，没过半柱香的时间，院落外便传来一阵急促而沉重的脚步声，伴随着一道又惊又怒、蕴含着磅礴威压的喝问：“瑶儿！我的瑶儿呢？！她怎么样了？！”

人影一闪，天师府府主，法相境大能沈天穹，竟已满脸焦急地冲进了院子，完全没有了平日里的威严气度，俨然一个担忧爱女安危的普通父亲。他显然是通过某种方式，瞬间感知到了沈瑶气息的剧变和生命危险。

“沈府主。”叶牧谦迎上前。

沈天穹一眼就看到了厢房内床上昏迷不醒、浑身插满银针的女儿，以及正在全力施救的陈千羽，脸色顿时变得更加难看。他强压怒火，听叶牧谦简略说明了事情经过——从沈瑶跪地拜师，到她决绝自废修为，再到叶牧谦无奈收徒。

听完，沈天穹沉默了。

他没有立刻发作，只是深深地、久久地凝视着床上脸色苍白如纸的女儿，那双阅尽沧桑的眼眸中，翻涌着心疼、不解、无奈，最终化为一声长长的、沉重的叹息。他看向叶牧谦，眼神复杂，缓缓说道：“罢了……这是瑶儿自己的选择。”

他了解自己的女儿，知道她一旦认定，便是九头牛也拉不回。既然她以如此惨烈的代价做出了选择，他这个父亲，除了接受，又能如何？

是夜，月华如水。沈瑶在陈千羽不计代价的救治下，情况暂时稳定，但仍在深度昏迷中，需要时间慢慢恢复受损的根基。

院落偏厅，沈天穹主动提及，想跟叶牧谦谈谈。

挥退了侍从，布下隔音结界，沈天穹仿佛卸下了府主的威严，显出一丝疲惫。他亲手为叶牧谦斟了一杯茶：

“叶先生，瑶儿性子执拗，今日之事，让先生见笑了。”沈天穹开门见山，语气诚恳，“既然瑶儿已拜入先生门下，有些关于她、关于天师府、乃至关于这内域的事情，我觉得有必要告知先生。”

叶牧谦正色道：“府主请讲，叶某洗耳恭听。”

沈天穹目光悠远，仿佛陷入了回忆：“瑶儿是我三百岁的时候出生的，今年刚二十五岁。她……自幼便心地善良，见不得不平之事。”

“这是好事。”叶牧谦点头。

沈天穹却补充：“但我天师府，虽位列内域五大超级势力之一，数百年来，却一直奉行孤立主义之策；我上任之时也一以贯之。”

他顿了顿，继续道：“我天师府根基，主要在于天枢城，以及另一座大城天权城。对于这两城之外，内域广袤土地上发生的诸多争斗、纷扰、乃至不平事……天师府历来是能不管便不管，能不理便不理。”

叶牧谦微微挑眉，这与他之前对“超级势力”的想象有所不同。

沈天穹看出了他的疑惑，苦笑道：“瑶儿便是因此，与我、与天师府产生了根本分歧。她认为，天师府既有如此力量，便应当承担起相应的责任，维护内域大体安宁，扶助弱小，铲除邪佞，方不负‘超级势力’之名。她认为我们……没有尽到义务。”

“那府主的考量是？”叶牧谦问道。

沈天穹神色转为凝重：“内域局势，远非表面看起来那么简单。五大超级势力：天师府算最强，其余四个分别是紫阳宗、无极剑派、万仙盟、丹鼎阁。而剩余四个中，又属紫阳宗最强，其宗主也是法相境外显期，与我颇有私交，只是近年闭关不出。

“超级势力看似超然，实则关系微妙，彼此制衡。许多发生在中小势力、乃至散修之间的‘小打小闹’，背后未必没有其他大势力甚至超级势力的影子，或默许，或推波助澜，以此试探、消耗对手，或达成某些不便明言的目的。”

他看向叶牧谦，“贸然插手，很可能便是不小心踩进了别人精心布置的棋局，甚至可能引发超级势力之间的直接冲突，后果不堪设想。天师府传承不易，我作为府主，首要职责是维系府内安定与传承，不可轻易将整个宗门拖入不可测的漩涡。”

这便是现实与理想，全局安危与局部正义之间的残酷矛盾。沈天穹的选择是保守、“无为”，而沈瑶则渴望激进的“有为”。

“我们父女为此争吵过无数次。”沈天穹叹息，“十年前，拙荆因为寿元耗尽去世了，家里气氛本就压抑。结果，一次激烈的争执后，她一气之下，便离家出走了。这一走，就没回来过。”

他本以为，女儿在外历练，见识到世间的险恶与复杂，会逐渐理解他的苦心与无奈。可事实恰恰相反，沈瑶在十年独自闯荡、行侠仗义、除魔卫道的经历中，亲眼目睹了更多因强者漠视、势力倾轧而造成的苦难，她的信念非但没有动摇，反而愈发坚定了！她认为父亲的“孤立”是怯懦，是冷漠，是导致诸多悲剧的根源之一。

父女间的理念鸿沟，非但没有因时间弥合，反而越来越深。

说完这些，沈天穹仿佛卸下了一副重担，神情却更加萧索。他看向叶牧谦，郑重道：“瑶儿如今拜你为师，或许……她的道，真的与先生有缘。她选择破而后立，我虽心疼，但也尊重。今后，还望先生多加教导。”

叶牧谦肃然点头：“既入我门，自当尽心。”

沈天穹最后道：“还有一事。半月之后，我天师府将做东，邀请内域其余几大超级势力的代表，于天枢城举行‘论道会’，坐而论道，交流心得。届时，还请先生务必拨冗参加。”

叶牧谦心知无法推脱，只得应承下来：“府主相邀，叶某自当赴会。”

沈天穹又去看了看依旧昏迷的沈瑶，嘱咐陈千羽若有任何需要，天师府资源尽可调用，这才带着满腹心事离去。

\vspace{2em}

接下来的日子，小院成了临时的医馆和修炼地。陈千羽几乎不眠不休，全力为沈瑶调理。

令人惊喜的是，在治疗过程中，陈千羽发现沈瑶半边脸颊那陈年的灼伤，在精纯元气的滋养与银针的引导下，坏死的肌肤组织竟然焕发生机，疤痕逐渐淡化、平复。她细心调理，数日之后，沈瑶脸上的灼痕竟已消失不见，恢复了原本清丽绝伦的容颜。

沈瑶在昏迷三日后终于苏醒。身体依旧虚弱，但经脉的损伤已在缓慢修复。

当她再一次在陈千羽递过的铜镜中，看到自己完整无瑕的脸庞时，整个人都怔住了，手指颤抖着抚过光滑的肌肤，久久无言。

如今伤痕褪去，仿佛也卸下了一层沉重的枷锁。

她沉默地将那方陪伴多年的黑色面纱仔细叠好，收进了储物袋深处，不再佩戴。

\vspace{2em}

经此一劫，沈瑶的心境发生了些微妙的变化。就像被湍流冲刷过的卵石一样，少了几分尖锐，多了些圆润和沉静。

果然是有修炼底子的人，悟性、经历也是一等一的。身体彻底恢复好之后，不过三日，沈瑶便将幻术和身法相关的理解融于一身，稳步踏入淬体境天骨期，修为回归养气境也只是时间问题。

这速度竟比当初的白言冰还要快上几分。

由于叶牧谦受邀论道，几人目前依然暂居于天师府提供的馆舍中。她自然便与叶牧谦等人住在了同一处。

夜里，枕着柔软舒适的锦枕，身下是干燥洁净的被褥。隔壁房间，偶尔会传来白言冰与陈千羽压低嗓音的交谈；更远一些，叶牧谦房中那规律而悠长的鼾声，穿过墙壁隐约可闻。

没有追杀，没有戒备，无需时刻紧绷心弦。她将半边脸埋入枕中，感受着久违的、近乎奢侈的安宁，意识沉沉陷入黑甜的梦乡。

\vspace{2em}

另一边，白言冰在照料沈瑶伤势之余，试炼幻境最后那直指道心的拷问也如同烙印，深深烫在他的神魂深处。某一日午后，他于院中观落叶飘零，忽觉灵台一点清明跃动，似有所感。

他即刻寻到叶牧谦，恭敬请示：“师尊，弟子心有所惑，亦有所感，想觅一静室闭关几日。”

叶牧谦从图谱上抬起头，目光在他沉静却隐含锋锐的脸上一扫，了然颔首。

白言冰寻了馆舍深处一间空置的静室，石门落下，隔绝内外。室内无窗，仅一颗明珠嵌于顶壁，散发着柔和的冷光。

他盘膝坐下，将心神彻底沉入那片迷雾之中，反复咀嚼那拷问，审视自己执剑的初衷。

时间在绝对的寂静中流逝，不知过了几个时辰，抑或是一日。他周身原本内敛的气息开始产生微妙的变化，时而如剑锋般凌厉欲出，时而又如古井般沉寂无波，最终，所有的波动渐渐归于一种奇异的圆融。

他第二次认识了“我”。

下一步，便是将这已然清晰的“我”之意志，通过元气，映照于外，化为实质的领域。这对已然厘清我与物分别的白言冰而言，是非常容易理解的，却需要精细的操控。

心念微动，气海中温养的精纯元气随之响应。他先尝试着引导一丝元气，沿着手臂的轨迹，如同沿着剑锋延伸般缓缓向外探出。

起初，这缕元气一旦离体稍远便迅速涣散，难以维持。这和先前元气附剑时是一样的，放出去的元气就不再能收回了。

白言冰不急不躁，反复尝试，调整着输出的强度与心神依附的浓度。

渐渐地，那外探的元气变得稳定，已能在他掌心之上三寸处凝而不散；白言冰试着轻轻屈曲手指，那些外放的元气则听话地被召回体内。

第一步已经成功了，第二步便是试着让这些离体的元气在周身之外自行流转，形成一个大致的循环。

不知又过了多久，量变引发质变。随着白言冰精微而持续的引导，气海中精纯的元气终于不再需要他去刻意驱策，似乎与外界那层循环初步建立的元气产生了共鸣，开始自发地、坚定地向外漫溢、补充，参与那周而复始的运转。

一个稳固、通透、心意朗照的独特势场，终于彻底成型，将他周身三尺之地笼罩在内。

在这势场之中，外界的纷扰被自然过滤、辨析；而他自身的意志与存在感，被无比清晰地映照和确立——这便是属于他白言冰的见独领域。心念一转，领域内气息可骤然变得锋锐，带着凛冽的剑意；心念一松，又可复归温润平静。如臂使指，不外如是。

白言冰缓缓睁开眼，眸中似有清光流转，旋即内敛。

他踏出静室，晨光落在他肩上。虽然这初成的、更偏向于心神感知与意志凝聚的剑意领域，远不如叶牧谦那浩瀚如星空、可直击神魂的心念领域那般深邃强势，甚至显得有些简陋单薄；但每一个运转的环节，每一缕外放元气的轨迹，都打上了他自身剑道意志的烙印。

更重要的是，越简单的东西，越不容易坏。他的剑意领域，似乎专门就是为了守护而生的。

这终归是他自己走出来的路。

\vspace{2em}

就在他成功突破、势场稳固成型的刹那，在自己房间内休息的叶牧谦，身躯微微一震：他感到一种熟悉的清流，无声无息地自冥冥中汇入他的道基之中。

“又是反哺……”叶牧谦嘴角勾起一丝细微的弧度，眼中闪过欣慰。

白言冰把“见独”走实了之后，叶牧谦对见独境的细微掌控、种种精妙运用的可能性，似乎瞬间挣脱了一层此前未曾察觉的无形薄膜，变得前所未有的圆转自如，得心应手，也终于不需要再担心反噬问题了。

\vspace{2em}

至于天师府密阁，叶牧谦并未急于让弟子们前往。待沈瑶伤势稳定、白言冰出关后，他将三人唤至面前。窗外竹影婆娑，映得他青色长衫仿佛也染上了几分翠意。

“天师府密阁，藏书浩瀚，包罗万象，是内域无数修士梦寐以求的宝地。你们入内观阅，自是机缘。”

他话锋一转，神色转为严肃：“然，需谨记于心：道途根本，在于内求诸己，明心见性。一切外来的知识、法门，皆需根据自身实际情况，作出实事求是的修改与化用，使之真正契合你们自身。若只知一味模仿他人形貌，邯郸学步，而失却自身那一点真灵神魂，便是落了下乘，前路必窄。”

学之者生，似之者死。

这是他对弟子治学修道、钻研术法的根本要求。

他又特意补充道：“阁中若有关于远古炼器手法、以及阵法一道的古老记载、残篇、甚至只言片语，若有闲暇，可帮我留意一番，取出容我一观。”

“弟子领命。”沈瑶躬身应下。

她心中感慨，十年前，她凭借特殊身份，偷偷进入过密阁外围；此番，却是堂堂正正而入。

半月后，待她身体再好些，便领着白言冰与陈千羽再次来到那幢气势恢宏又戒备森严的阁楼前。

验过令牌，沉重的玄铁大门缓缓开启。沈瑶步履轻盈，穿过陌生又熟悉的廊道与禁制光芒，很快便将二人带至存放炼器、阵法及相关杂学典籍的偏殿区域，效率奇高。

\vspace{2em}

半月时光，便在沈瑶日渐红润的面色、白言冰日益稳固圆融的见独境气息、陈千羽对几卷古老医典的痴迷研读、以及叶牧谦对零星收集来的残破炼器阵法图谱的初步推演中，悄然滑过。

\vspace{2em}

很快，便是府主沈天穹前日所言的论道会之期。

五方主位之上，端坐着五位威仪各自笼罩一方的身影。

东首主位，自然是地主，天师府府主沈天穹。

其左侧下首，是一位徐娘半老、身着丹霞色宫装长裙的女子，云鬓高挽，插着一支药玉簪。此人正是丹鼎阁阁主，文明内域的“丹王”王丹。其修为已至真符境注灵期，举手投足间，雍容华贵，又带着丹师特有的沉静气度。

王丹的下首，是一位赤发如火、根根竖立如钢针的男子，面容刚毅如同斧凿刀刻，浓眉之下双目炯炯有神。此人乃是紫阳宗副宗主，陈炎，修为在真符境显化期，性情素来以火爆直接、不屈不挠著称，与他的火系功法相得益彰。

沈天穹右侧，是一位青衫老者。衣衫洗得发白，纤尘不染，背后负着一柄样式古朴的长剑，剑柄缠着磨损的暗色布条。他面容清癯，颧骨微高，须发灰白，眼神初看温和润泽，细细品味，却能察觉到凛冽的剑锋。正是无极剑派心剑一脉的长老，独孤明，修为亦是真符境注灵期。

最末一位主位上的，是位看起来颇为儒雅随和的中年人，身着素雅的文士衫，手中轻轻摇着一柄白玉为骨的折扇，脸上常挂着三分令人如沐春风的微笑。他便是内域最大散修联盟“万仙盟”的盟主，吴东浩，修为在真符境摹形期。虽然说这位吴东浩修为在诸人中最低，但以散修之身拼搏至此，其手腕与心机也是无人敢小觑的。

阁宇中央，相对五方主位设一独立客席。叶牧谦安然独坐其上，依旧是一身半新不旧的青衫，神色平静无波，迎着来自五个方向、或明或暗的审视目光，坦然自若。白言冰、陈千羽、沈瑶三人，则与其余少数同被邀请但无资格上台的修士们坐在一起。

此前，沈天穹早已将叶牧谦及其所持“新的道途”之信息，简略告知其余四方。此刻，诸方巨擘目光汇聚于那青衫客卿身上。

好奇有之，探究有之，但更多的，是一种根植于漫长修行岁月、对于未知与可能颠覆现有格局的“异端”那种本能的不信任与审视。

\vspace{2em}

作为东道主，沈天穹率先开口，声音平和醇厚，打破了阁中的沉寂：“叶先生，今日诸位道友齐聚我这问道阁，皆因闻听先生之道，与我等所行所知的‘灵枢径’迥然不同，心生探究印证之意。先生不妨简述其道之要旨，也好解诸位道友心中之惑。”

叶牧谦微微颔首，面向众人，朗声开口，声音清越，在空旷的阁中回荡：“沈府主，诸位道友。在下所行之道，自名为‘问道途’。其核心理念，在于‘内求诸己，明心见性’。不假外求灵枢以为沟通天地之桥梁，不刻意开拓体内玄脉以强纳外界灵气，而是于丹田之下，开辟气海，温养调和自身生命本源所生之元气，以此筑基。”

话音刚落，左侧便传来一声洪钟般的诘问，带着毫不掩饰的质疑与火气。

正是陈炎，赤眉扬起，声震屋瓦：“叶先生此言，未免过于空泛虚渺！不引灵枢，不通玄脉，如何引动天地浩瀚灵气入体修行？既无灵力傍身，如何施展神通妙法？莫非你这‘问道途’，只能让人坐而论道、空谈心性不成？便是修士最基本的飞檐走壁、御器凌空、符箓咒法，如何做到？”

他的问题尖锐直接，毫不留情，恰恰戳中了目前问道途在外人看来最明显的短板——缺乏灵枢径修士那些直观、炫目且极其实用的外在能力。

叶牧谦对此似乎早有预料，面上并无愠色，也不着急，不疾不徐地答道：“陈副宗主所问，正点出了两途根本立意之差异所在。问道途讲求内求于己，元气源于自身气血、精神、性命之调和温养，储存于气海，运行于自身小天地，其根基便在于‘不假外物’。至于陈副宗主所言飞檐走壁、御器凌空……”

他略微顿了顿，脸上浮现出一丝淡然甚至略带超脱的笑意，目光平静地迎向陈炎灼灼的视线，反问道：“修行之人，为何……就非要飞呢？”

他语气平和，却带着一种奇特的、近乎哲辩的意味：“脚踏实地，一步一印，感受山川厚重；观花开花落，体悟四时轮转生灭；听风过松涛，明辨八荒气息流变。这行走坐卧、俯仰天地之间，所感所悟，不亦是‘道’之所在吗？”

这话语带着几分有意为之的诡辩与超然物外的姿态，试图将对方从“力量表现”的务实层面，拉入自己“重内修心性意境”的语境之中。

陈炎果然被这“为何非要飞”的反问结结实实地噎了一下。他修炼的乃是霸道炽烈的火系功法，性子亦是刚猛直接，向来认为强大的力量就该有与之相匹配的威势与外在表现，腾云驾雾、御剑飞天乃是修士的常能与标志。叶牧谦这种近乎逃避问题本质的说法，在他听来，简直强词夺理。

他张了张嘴，赤红的脸膛因激动更显颜色，想要驳斥，可一时间又觉得，跟一个似乎根本不在乎“能不能飞”“会不会飞”的人，去争论“怎么飞得更好更快”，简直如同对牛弹琴，徒费口舌。满腔话语堵在喉咙口，最终只化作一声重重的不满的冷哼，撇过头去。

然而，表面上镇定自若的叶牧谦，内心却是另一番景象。他暗自警醒：下一个境界，无论如何也必须设法衍生或融合出类似“腾空”、“远程攻伐”等基础的实用能力了。

否则，光靠嘴皮子讲道理，短期或许能震慑一时；长远来看，终究难以服众，更会严重限制弟子们未来的发展空间与安全。

\vspace{2em}

这时，一个温和却异常清晰，仿佛能直接传入人心底的声音响起，带着剑器般的清越质感。正是无极剑派的独孤明长老。他微侧身形，看向叶牧谦，缓缓道：“叶先生，老夫独孤明。我辈剑修，虽不喜主动挑起无谓之争，却也从不畏惧卷入争斗，手中之剑，只为守护心中所珍视之人、所坚守之道而鸣。先生之道，重内修心性，明己见独，老夫细听之下，甚为欣赏其中返璞归真之意。”

他话锋随即一转，语气依旧平和，问题却直指核心：“然，修行之路，漫漫悠长，终究无法完全避开与‘力’的较量。守护需力量，卫道亦需依仗。不知叶先生这‘问道途’，若论及实际护道卫己之能，与我等所熟稔之道相较，究竟如何？可有独到之处？抑或仅是镜花水月，不堪实用？”

这个问题，比陈炎的直接诘问更加实际，也更为深刻。它绕开了外在形式的争论，直指修行道途无论理念如何超脱，最终都无法完全规避的核心——护身与制敌的力量。

叶牧谦心知，面对独孤明这等浸淫剑道数百年、心志早已千锤百炼如精钢的老者，空谈理念、故弄玄虚已然毫无用处。道之高低，终究需落于实处印证。

他从容起身，对独孤明所在方位，郑重地拱手一礼，声音斩钉截铁：“独孤长老所问，切中肯綮！道之高低优劣，纵有千般言语，终需实际印证，方见真章。叶某不才，修为浅薄，然愿以当前之微末能为，请长老试手。”

此言一出，阁内气氛骤然一凝。沈天穹眼帘微垂，不动声色；王丹眸光闪动，似有好奇；吴东浩手中折扇轻摇的频率微微一顿；陈炎则猛地转回头，赤瞳中精光爆射，紧紧盯住场中。

独孤明眼中亦是精光一闪，仿佛沉眠的古剑骤然苏醒了一瞬。他反而因叶牧谦的坦荡与果决微微颔首，面上掠过一丝激赏：“好！叶先生爽快！既如此，老夫便以三招，领教先生之妙！”

言罢，他抬起右手，食指与中指并拢，神色平淡地向着数丈之外的叶牧谦，轻轻向前一点！

“嗤——”

一道凝练到极致，锋锐无匹、直指人心的凛冽剑气，破空而至！

这是如此简单的一招。先不论场上剩下那四个人怎么看了，就连场下的白言冰也不觉得这招有多高明。能看出来独孤明是狠狠留手了。

叶牧谦自然无声展开自己的见独领域。前些日子白言冰成功突破的反哺，使得他目前已达见独境圆满，单论实力已经不虚任何一位真符境修士。

那道凌厉无匹的剑气射入这片心念领域的刹那，其摧枯拉朽的锋锐锋芒，仿佛被无形而柔韧的水温柔地包裹、浸润、稀释了。叶牧谦站在原地，青衫下摆都未曾拂动一下，身形稳如山岳。

“咦？”独孤明口中发出一声几不可闻的轻咦，略显讶异。他承认这第一指确实是收着力的试探，却也绝非寻常玄脉境修士能够轻易接下的，没想到被叶牧谦这般不着痕迹地化解了。

独孤明寻思，得来点真东西了。紧接着，第二指点出！

这一指，剑意骤然变得繁复莫测！一刹那间，仿佛有千道万道虚实难辨的剑影藏于这一指之中迸发，似真似幻，交织成网，更带着迷惑、缠绕、窥探之意。

正是独孤明发扬光大的心剑脉杀招之一，双生剑法！

叶牧谦依旧以见独领域从容应对。领域之内，那漫天袭来的虚幻剑影，在他那澄明如镜的心神映照之下，纷纷显露出其虚幻的本质，如阳光下的泡沫般逐一幻灭。

而藏于这万千虚幻之中，那唯一一道真实不虚、凝聚着独孤明此番剑意核心的剑意，则被他自身圆融无碍、坚定如一的心意悄然一引，擦着他的身侧偏转开去，未伤及分毫。

两招已过，叶牧谦依旧稳立当场。

独孤明脸上轻松的神色终于彻底收起，取而代之的是一片郑重。他缓缓点了点头，深深吸了一口气。

第三指，缓缓点出。这一指，没有凌厉疾刺，也没有万千幻影，只有一道凝实到极点、沉重到极致、仿佛挟带着独孤明毕生对“守护”二字全部感悟的剑意，缓缓推出。

这剑意虽朴素至极，但厚重如山，坚定如岳。它以一种无可阻挡、不可撼动、亘古长存的气势压迫而来：

守护的背后，是牺牲，是承担，是即便天地倾覆亦不退半步的决绝！

叶牧谦瞳孔微缩，深知已到关键。他能轻易看出来这一剑和前两剑本质上的差别，自然不敢怠慢，将见独领域催发到极致！

刹那间，那无形的领域仿佛被注入了灵魂，有了鲜明的色彩与温度。

两股同样坚定、恢弘，但内涵却截然不同的“守护”意志——独孤明那历经岁月沉淀、以剑为誓、厚重如山的守护，与叶牧谦那源于本心、无愧天地、指引前路的守护——轰然对撞！

没有震耳欲聋的巨响，没有绚丽夺目的灵光爆炸。然而，阁中所有人，包括主位之上修为最高的沈天穹，都在这一刹那，感到自己的心神仿佛被一柄无形的大锤轻轻敲击了一下，微微一震！意识海中泛起涟漪，道心为之所感。

三息。

令人屏息的沉寂持续了三息时间。

独孤明缓缓地、极其郑重地收回了伸出的手指。他脸上最初的惊讶之色早已褪去，逐渐化为一种由衷的、毫不掩饰的赞叹与钦佩。他站起身来，不再是踞坐的姿态，而是对着数丈外同样肃然而立的叶牧谦，双手抱拳，郑重地躬身一礼。

“叶先生之道，心性之坚纯，领域之玄妙，老夫佩服！”

他的声音沉凝有力，回荡在寂静的阁中，“此‘问道途’，确有惊天独到之处！于实战之中，尤其是心神交锋、意志对抗之际，恐有出乎意料之奇效。老夫……心悦诚服！”

能让独孤明亲口说出“心悦诚服”四字，分量之重自然不言而喻。

\vspace{2em}

最后，那位一直带着和煦笑意、仿佛只是个旁观者的万仙盟盟主吴东浩，终于开口了：

“叶先生，今日观先生论道、试招，风采令人心折。然，老夫有一问，藏于心中久矣，还想请教先生。”

他轻摇折扇，目光温和却深邃地看向叶牧谦：“所谓修士，无论所修何道，说到底，皆是逆天而行、与天争命之人。既如此，大多心有所求，志有所向。或求长生久视，超脱轮回；或求无上力量，掌缘生灭；或求自在逍遥，无拘无束……不知叶先生，您创立此迥异于常的问道途，自身踏上此道，孜孜以求，所寻的……究竟是个什么呢？”

是啊，你另辟蹊径，弄出这么一条与众不同的路，总得有个终极的追求、一个能支撑你走下去的“答案”吧？

阁中所有的目光，再次汇聚于叶牧谦身上，等待着他的回答。

叶牧谦闻言，微微一怔。这个问题，他或许在夜深人静时也曾自问过，却从未如此正式地需要向他人宣之于口。

当吴东浩问出的瞬间，过往的种种画面——穿越之初，他的茫然、逃避，以及最后的悔悟与接受，如同浮光掠影，在他脑海中飞快闪过。

没有长篇大论的思考，没有修饰言辞的迟疑。

或许是前世儒家经典读多了，他几乎是凭借着一股发自肺腑的本能，那个答案便自然而然地涌上心头，冲口而出。声音不大，却清朗如玉石相击，清晰地回荡在每一寸空气都仿佛凝结的道场中：

“只为仰不愧于天，俯不怍于人罢了！”

此言一出，偌大的道场内，陷入了刹那绝对的寂静。

包括沈瑶在内的很多人似乎不能立即理解这十个字。但在白言冰和陈千羽看来，这十个字却做不得伪。

这确实是师尊在黑风岭一役之后、正式把自己当作“师尊”之后的本心写照，而师尊也一直践行着这十个字。

“仰不愧于天……俯不怍于人……”吴东浩低声重复着这十个字，眼中那惯常的和煦笑意渐渐变得深刻，更添了一丝了然与感慨的意味。

他缓缓收起折扇，在掌心轻轻一敲，微微颔首，显然对这个看似朴素至极、却又重若千钧的答案，甚为满意，亦深以为然。

王丹眸中若有所思，似在品味其中返璞归真的道韵。

陈炎却是眉头再次皱紧，嘴唇动了动，似乎觉得这个答案过于“平凡”，不够宏大，不够有“追求”。

刚刚试剑完毕的独孤明长老，则是手抚灰白长须，缓缓点头，脸上露出了深以为然的赞同神色。他毕生修心剑，守护为念，最能体会这“不愧不怍”四字背后，所需的那份坦荡、坚定与无悔。

而端坐于东首主位的沈天穹，此刻缓缓地、极其庄重地站起身来。

一声仿佛包含了无数感悟、解脱了某些无形桎梏的喟然长叹，从他口中悠悠吐出。那叹息声厚重如大地，却又带着一丝豁然开朗后的明澈光芒：

“老夫……修道三百余载。”

他声音不高，却字字如磬，说出了那四个注定将伴随“问道途”传扬天下的字：

“今日，方知……道在低处。”

道在低处。

四字如晨钟暮鼓，余韵悠长，在问道阁的梁柱间萦绕不绝，更重重地敲在每个人的心头，激荡起层层涟漪。

沈天穹，内域五大超级势力之首的天师府府主，法相境外显期的大能，以他的身份、修为与三百载见识，亲口为今日这场关乎道途立论的辩难与印证，做了最终的总结。

其含义，不言而喻——他认可了叶牧谦的道，认可了那“仰不愧于天，俯不怍于人”的朴素追求之中蕴含的至高真意，认可了这条不假外求、内修己心、看似“在低处”的新途，自有其巍然高度与无限可能。

自此，叶牧谦所传所行之道，在内域顶尖势力的层面，获得了正式的、无可争议的承认。

它叫“问道途”。

与那流传万古、显赫辉煌的“灵枢径”并立，成为内域修行大世之中，一个崭新、独特、再也无法被忽视的道途之名。

\vspace{2em}

论道会后，沈天穹私下寻到叶牧谦。

“叶先生，”沈天穹神色温和，全无府主的架子，“今日论道，瑶儿也在旁听。老夫见她神情，已知她确实找了个好师傅，心中甚慰。”

他顿了顿，道：“听闻先生至今尚无固定道场，游历讲学，虽显洒脱，终非长久之计。

“老夫在城外二十里，云霞山半山腰处，有一旧居，名曰‘云隐别院’。那里风景秀丽，灵气内敛，环境清幽，正合先生修心养性之需。那院子已空置十多年，留着也是无用，便送给先生，权当是瑶儿的束脩之礼，还望先生莫要推辞。”

叶牧谦略一沉吟，便拱手谢过：“府主厚赠，叶某却之不恭，便代小徒瑶儿谢过了。”

他确实需要一处稳定的根基之地，来梳理愈发庞杂的道途体系、顺便教导弟子。云隐别院听起来正合适。

不久后，叶牧谦携三位弟子入住云隐别院。院落果然清幽雅致，背靠云霞山主峰，面朝开阔山谷，云岚缭绕，草木葱茏，灵气平和；更设小演武场、藏书阁、闭关静室等必要设施，还有温泉、冷泉各一眼。

确是一处宝地。

\vspace{2em}

安顿下来后，叶牧谦也开始研究起他的下一境界。

思路是显然的，从论道会上陈炎的诘问说开去：不论是腾空飞行，还是远程攻伐，具体的内容实际上都是所谓“术”的层次。

而“理解”这些术法，自然是一门很大的学问。

“世事洞明皆学问”，叶牧谦忽然想起《红楼梦》中的这句话，于是自然而然地把“洞明”作为了下一境界的备选名称。

转念一想，“洞明”又不好，有管中窥豹、不能窥见全貌之嫌疑，过于小家子气；不如改成“通明”——万法通明之境，多大气磅礴！

意外之喜是“通明”的另一个意义：除了“通晓万物”，还有“十分明亮”的意思；“一心通明”，更是令人神往。

于是叶牧谦当即敲定，第三境界便叫做“通明”了：“至人之用心若镜，不将不迎，应而不藏，故能胜物而不伤”！

随即，叶牧谦定下规矩：每逢初一、十五，于别院开坛讲道，不拘出身，不问来路，有缘皆可来听。消息渐渐传开，起初多是些好奇的散修、附近小宗门不得志的弟子，以及一些仰慕沈瑶——身份终是瞒不住——或对“问道途”好奇之人前来。

叶牧谦讲道，深入浅出，不拘泥于具体功法招式，多讲“养气”之静、“见独”之悟，更反复强调“内求己身”、“道在日用”之理。

这对那些受困于资源匮乏、在传统道途上难有寸进的散修和小门派弟子而言，这无异于打开了一扇全新的窗户。听道者日众，不少人深受触动，甚至有人在别院附近结庐而居，以便时常请教。

其中有一散修，名叫陈默，原是灵枢境凝枢期修为，因资源问题困于瓶颈多年。

他听叶牧谦讲道，结合自身对灵枢径的理解，竟突发奇想，冒着巨大风险，尝试引导自身那微弱的灵枢之力下行，冲击丹田之下区域，意图强行开辟“气海”。

过程自然痛苦万分，几经险死还生；但凭借一股狠劲与冥冥中的一点灵光，他竟然成功了！

虽然修为几乎尽废，从头开始，但确确实实强行压出了第一缕问道途的元气，踏入了养气境门槛。

陈默的成功案例，引起了已成功踏入养气境界不久、正在稳固修为的沈瑶的极大关注。

她敏锐地意识到，这种破而后立的转修方式，虽然凶险痛苦，不适合大多数人，但或许是一条能让部分前路无望的修士改换门庭的潜在路径。

她详细询问了陈默的感受与过程，细心记录在册，归档整理，以备他日深入研究或谨慎参考。

这也算是问道途与旧体系之间，第一次非官方的、个体尝试的实验了。

\vspace{2em}

随着人气渐旺，为便于管理，避免纷扰，叶牧谦与三位弟子商议，制定了简单的《问道戒律》，约束居于别院周边及听道者的言行。戒律第一条，赫然便是叶牧谦在论道会上的那句心声：

仰不愧于天，俯不怍于人。

以此为核心，衍生出数条诸如“不得恃强凌弱”“不得同门相残”等基本规范。

云隐别院及周边，逐渐形成了一个以问道途理念为核心，氛围相对平和清正的独特小生态，成了不少厌倦厮杀争夺、或寻求新路的修士心目中的一方净土。

实际上，叶牧谦开堂讲道，也是试图通过广泛招收记名弟子、扩大讲道范围，来获得更多的道行反哺；然而，一段时间下来，他发现效果微乎其微。那些听道的散修、访客，即便有所得，甚至如陈默般成功转修，这种反馈也极其微弱缥缈，远不能与白、陈、沈三人带来的那种清晰触动相比。

他心中渐渐明确：疑似只有白言冰、陈千羽、沈瑶这三位命运交织、心性纯粹且正式拜入门下的亲传弟子，当他们真正开悟、突破，深刻践行问道途时，才能产生那种独特的、对他有实质性助益的“道行反哺”。

这似乎与师徒间深刻的道统传承与心神共鸣有关，不能简单通过人数叠加扩增。

\vspace{2em}

云隐别院的时光，在讲道、修行与日渐兴旺的人气中，平静地流淌了数月。

山间云雾聚散，别院道场上的听道者来了又走，走了又有新人来，别院周围依山势错落搭建的庐舍也多了十余间，俨然有了一个小型聚居地的雏形。

白言冰巩固着“见独”境界，剑意愈发凝练通透；陈千羽医术与修为并进，温润元气滋养着别院内外受伤或体虚的访客；沈瑶则彻底褪去了过往的尖锐，在养气境的稳步前行中，气质愈发沉静，只那双眸子里的倔强与清澈，未曾改变。叶牧谦则沉浸在对炼器、阵法古籍的钻研中，为问道途那众所周知的短板寻找补全之道。

目前解决“缺少功法”问题的权宜之计，依然是通过改编灵枢径的功法来实现。但两个道途存在根本的不同：灵枢径的灵力是燃料，问道途的元气是润滑油。这就直接导致了从灵枢径改编来的功法，消耗都有些不切实际的大。

路是一步步走的，移山开路这种事情，的确急不得——急，也没用。

\vspace{2em}

然而，这世间并非处处讲理，人心也并非总是开明包容。

天师府论道会上，问道途获得正式承认并与灵枢径并立；但这份荣耀背后，是潜藏的危机与既得利益者的不安。

第一个坐不住的，正是紫阳宗副宗主陈炎。

论道会归来后，叶牧谦那句“为何非要飞”的诡辩虽一时噎住了他，但并没有打消他心中的疑虑；这种疑虑反而随着时间发酵，变成了更深的忌惮。

他亲眼见过白言冰接下独孤明三招时展现的那种奇异领域，也听闻了云隐别院讲道日益红火，甚至出现了陈默那种敢于自废灵枢、转修元气的狠人散修。

一种隐约的威胁感，如同阴云般笼罩在他心头。

这一日，紫阳宗内，陈炎独坐于烈焰熊熊的炼器室中，面前悬浮着一件即将完成的剑胚，灵光流转，但他却有些心不在焉。指节无意识地敲击着座椅扶手，发出沉闷的声响。

“内求己身……不假外物……”他低声重复着叶牧谦的话，赤红的眉头越皱越紧，“若是人人皆内求己身，温养那劳什子‘元气’，对法器、丹药的依赖必然大减。长此以往，我紫阳宗炼器之道，丹鼎阁的丹药生意，根基何在？我们这些以‘外物’立宗的宗门，又何以存续？”

这个念头一经浮现，便如同毒藤般缠绕上来。

在他看来，“问道途”宣扬的理念，本质上是在动摇现有修炼体系的物质基础，是在掘他们这些传统势力的根基！

长此以往，问道途必然会逐步挤占灵枢径的生态位，最终成为新的正统！

他也确实是看过不少话本，在那些话本中，反派总是很可笑地让主角成长起来；而在反派最终选择剿灭主角时，主角往往已成气候，动不了了，只能眼睁睁地看着自己落败。

为了宗门的利益与发展，必须立即将这问道途扼杀在摇篮中。

他不再犹豫，立刻启动了几道远程传讯，分别联系了丹鼎阁阁主王丹，以及无极剑派外剑一脉的长老刘长靖。

传讯的内容大致相同，核心便是他那充满忧虑的质问：“若是人人自求，那还要我们法器、丹药何用？我等宗门，何以立宗？”

很快，他收到了回复。

\vspace{2em}

丹鼎阁，王丹静坐于氤氲着药香的丹房之中，指间捻着一枚新炼成的丹药，丹纹流转，宝光莹莹。接到陈炎的传讯，她雍容的面庞上掠过一丝沉吟。

作为丹道大家，她比陈炎更清楚丹药对于灵枢径修士快速提升、突破瓶颈、疗伤祛毒的重要性，几乎是不可或缺的辅助。若真有一种道途能大幅降低对丹药的依赖，哪怕只是对底层修士有效，对丹鼎阁长远的损失也是不可估量的。

尽管论道会上她对叶牧谦印象不坏，但涉及宗门根本利益，陈炎说得又在理，她的天平迅速倾斜。

她回复陈炎，虽未多言，但“深以为然”四字，已表明了态度。

\vspace{2em}

无极剑派，坐落在群山之巅，剑气凌霄。

门派之肇始可以上溯到一对道侣，夫以清修证道、妻以杀伐护道，道侣之间举案齐眉、相敬如宾，其道也是互补融合，成了一段佳话。至于结局，有人说他们飞升上界了，也有人说他们相继去世后葬在一起，但他们留下了无极剑派这一点是肯定的。

不知为何，他们两人创立无极剑派的时候，就没有设立“掌门”这个职位，具体事务由夫妻二人——之后便是分别统领两脉的两名长老——共同负责。他们的本意，也自然是让两脉互补融合，心剑为魂，外剑为骨，成就无上剑道。然而数代传承下来，现实却不如人意。不设掌门，使得两脉之间缺乏一个粘合剂，最终因理念差异，渐渐疏远，甚至产生了隔阂与摩擦。

传到现任这一代，心剑脉长老是独孤明，外剑脉长老则是刘长靖。两人年纪相差一轮，独孤明为长。

刘长靖对这位修为高深、剑心通明的前辈始终保有几分尊重，但两人之间的理念冲突与互相批评指摘，几乎已是常态。独孤明认为外剑一味追求威力，失了剑道灵性，易入魔道；刘长靖则认为心剑空谈心境，修行凶险，易生心魔。

问道途的出现，让两脉的分歧更加鲜明。独孤明自论道会归来后，对问道途推崇备至，尤其是对其“内明己心”、“仰不愧天俯不怍人”的核心理念，以及叶牧谦那奇特的见独领域在心神交锋中的妙用，视若珍宝。他认为这与心剑“以心证道”的追求有异曲同工之妙，甚至能弥补心剑某些过于难以把握的短板，于是开始频繁带着门下几位核心弟子，前往云隐别院拜访叶牧谦、白言冰，交流心得，探讨剑意与心境的融合。

这一切，看在刘长靖眼中，却变成了越来越深的忧虑与恐惧。

刘长靖自幼在无极剑派长大，天赋卓绝，于外剑一道上勇猛精进。但他童年时，曾亲眼目睹过不止一次惨剧：有修习心剑的师兄、师叔，或因执念过深，或因外魔侵扰，在参悟高深心法时骤然走火入魔。

平日里温文尔雅的同门，瞬间变得癫狂嗜杀，六亲不认，剑意化作纯粹毁灭的疯狂，在门派内造成过不小的伤亡——这在一位年轻人心中，会留下何等的心理阴影，不言而喻。那些染血的回廊、同门惊惧的眼神、师长痛心疾首的叹息，如同梦魇，深深烙印在他心底。

他尊重独孤明，也钦佩其剑道修为，但内心深处，始终对心剑一脉那玄妙而凶险的修行方式抱有极大的戒惧。

如今，独孤明非但不坚守心剑传统，反而去追捧一个来历不明、体系迥异的问道途，还带着弟子频繁接触之，在刘长靖看来，这简直是双重的危险叠加！他害怕独孤明会被那套新奇理论带偏，更害怕门下弟子受其影响，重蹈昔日走火入魔、同门相残的覆辙。

这种恐惧，随着独孤明拜访云隐别院的次数增多而日益加剧，让他坐卧难安。

于是，当陈炎那充满煽动性的传讯到来时，刘长靖几乎没怎么犹豫，便就着陈炎那条思路去了：他觉得陈炎说得对，这种蛊惑人心、动摇根基的“异端”，确实需要遏制。

而且，遏制问道途，某种程度上也是在将独孤明和心剑脉从危险的道路上拉回来，更是为了门派内部的纯净与安全。

于是，他先是私下找到独孤明，试图劝说：“独孤师兄，那叶牧谦之道，终究是外人旁门，与我无极剑派心剑传承并非同源。您如此推崇，还带弟子常去，恐惹人非议，也易让弟子们心生迷茫，根基不稳啊。”

独孤明闻言，只是平静地看着他：“长靖，道无内外，唯有真伪。叶先生之道，于我心剑有启发印证之妙，何来旁门之说？而我带着弟子开阔眼界，明辨道心，也是为他们夯实根基。你多虑了。”

劝说无效，反而更坚定了独孤明与叶牧谦交流的决心。

刘长靖心中焦虑更甚，他仿佛已经看到了心剑脉在叶牧谦“邪说”影响下，逐渐偏离正道，最终酿成大祸的可怕场景。

思前想后，焦虑与恐惧压倒了对同门师兄的最后一丝情面，与“门派内部事务应内部解决”的原则。一个极端而危险的念头在他心中成型：不能让心剑脉继续堕落下去了，必须采取断然措施，清理门户，将危险扼杀在萌芽中！

但仅凭外剑一脉，想要压制乃至清理心剑脉，力量未必足够，且容易引发门派内大规模内战，导致两败俱伤，这更不是他想看到的。

不过，他随即想到了陈炎的传讯，想到了王丹的偏向态度。

随即，计划逐渐清晰：联合外部的紫阳宗，以雷霆之势，清除心剑脉中的不稳定因素，至少要让独孤明放弃对问道途的支持，将心剑脉重新拉回正轨。这样，即使心剑脉拉不回来、随之覆灭，外剑脉也不会受到过多损失，似乎是一个较好的方案。

于是他秘密联系了陈炎，表达了“清理门户”的意向，并请求紫阳宗协助，施加压力。

陈炎接到刘长靖的密讯，先是一愣，随即心中狂喜。

他正愁找不到合适的机会和理由打压日益壮大的云隐别院和问道途，刘长靖的请求简直是瞌睡时有人递枕头！

插手无极剑派内部事务，固然有风险，但收益巨大。

一方面，可以借此狠狠震慑叶牧谦——看，连支持你的独孤明和他的心剑脉，我们都敢动，你掂量掂量自己的分量！

另一方面，若能协助刘长靖掌控或重创心剑脉，等于从无极剑派这块大蛋糕上撕下了一块肉，极大增强了紫阳宗的影响力，甚至可能在未来掌控部分无极剑派的资源。

这个刘长靖有点傻啊，这笔买卖，在他看来实在是划算得不得了。

他几乎没有任何犹豫，便欣然答应了刘长靖的请求，双方密谋定下了行动细节。

\vspace{2em}

数日之后，一场突如其来的“清理门户”行动，在无极剑派猝不及防地爆发了。

这一日，天色阴沉。刘长靖以“稽查功法异动，防止心魔滋生”为名，调集了外剑脉大批精锐弟子，在数位紫阳宗高手的暗中配合下，突然包围了心剑脉主要的几处修行静所和讲剑堂。

没有过多的宣告，也没有给独孤明辩解的机会。刘长靖手持象征着外剑脉长老权限的令牌，面色冷峻，声音通过灵力传遍山峦：

“心剑脉近年修行路数有偏，易生魔障，为保门派清净，现需暂停心剑脉一切对外交流，封闭静修，接受审查！所有弟子，即刻返回各自居所，不得随意走动！违令者，以叛门论处！”

“放屁！”

一声怒喝如惊雷炸响。

独孤明果然出现，飞在半空与刘长靖对峙；须发皆张，青衫无风自动，背负的古剑发出嗡嗡清鸣。

他没想到刘长靖竟敢如此公然发难，更没想到对方竟然勾结了外宗的紫阳宗之人。这已不是简单的理念之争或内部整顿，这是赤裸裸的背叛与引狼入室！

“刘长靖！你竟敢勾结外人，对我心剑一脉下手！无极剑派的列祖列宗在上，岂容你如此胡作非为！”独孤明目眦欲裂，心中充满了被同门背后捅刀的悲愤。

“独孤师兄，我这是为了门派好，也是为了你好！”刘长靖面对独孤明的怒火，心中虽有一丝愧疚，但更多的是一种“执行正义”的决绝。

“只要你答应不再与那云隐别院来往，并约束弟子，接受审查，我等自会退走，此事尚有转圜余地！”

“转圜余地？将我剑心自由，卖与你等做妥协的筹码吗？”独孤明怒极反笑，笑声苍凉。

“我独孤明修剑一生，求的是心念通达，护的是心中正道！今日你等以势压人，以奸谋算计同门，此等行径，与邪魔何异！要我屈服？宁可战死！”

最后一个“死”字吐出，浩瀚纯粹的剑意冲天而起，充满了玉石俱焚的决绝与凛然不可犯的锋芒。

他身后，数十名心剑脉核心弟子也纷纷拔剑，剑意虽强弱有别，但那份宁折不弯、誓与师长共存亡的心志却连成一片。

冲突瞬间爆发。

外剑脉弟子剑招凌厉，配合默契，道道剑气纵横切割，充满杀伐之气。紫阳宗高手则在暗中释放火系法术干扰，或祭出法器远程轰击。

心剑脉弟子剑意纯粹，往往于间不容发之际寻隙而入，攻敌必救，或以精妙剑意干扰对手心神；独孤明更是独战刘长靖与一名紫阳宗真符境长老，剑光如龙，竟一时不落下风。

然而，终究是寡不敌众，且事发突然，心剑脉并未做好准备。激战片刻，便有数名心剑脉弟子受伤。独孤明眼见弟子受创，又感知到还有更多外剑脉弟子和紫阳宗援手正在赶来，心知今日事不可为，继续硬拼，只会让心剑脉传承断绝于此。

他悲啸一声，剑光暴涨，暂时逼退对手，对门下弟子厉喝道：“心剑一脉弟子听令！即刻分散，突围下山！留住有用之身，剑心不灭，自有再见之日！勿要为我这老骨头陪葬！”

“师尊！”弟子们悲呼。

“走！”独孤明剑意催发到极致，竟隐隐有燃烧本源之势；为了防止部分忠义弟子死战不退，他甚至朝着身后斩出一发剑气，逼迫他们撤退！

心剑脉弟子含泪咬牙，知道这是师尊用生命为他们争取生机，不得不化整为零，凭借对地形的熟悉和灵活的剑意运用，从各个方向拼死突围。

大部分人在独孤明拼死掩护和混乱中，成功冲出了包围圈，但亦有少数人重伤被擒或当场战死。

刘长靖本想追击，但被独孤明那搏命般的剑势死死拖住，加上他也不愿对逃亡的普通弟子赶尽杀绝——毕竟初衷是清理而非灭绝，只得眼睁睁看着大部分心剑脉弟子四散逃入茫茫群山。

最终，独孤明燃烧本源，力战至油尽灯枯，随着一声叹息溘然长逝。其余高层，不是陨落就是被擒。

曾经剑气凌霄的无极山，心剑一脉，一夜之间，风流云散。

\vspace{2em}

逃出生天的心剑脉弟子，在最初的慌乱与悲愤过后，不约而同地想到了一个地方——云隐别院。师尊独孤明正是因为推崇“问道途”，与叶先生交好，才遭此横祸。如今他们无处可去，天下虽大，恐怕也只有那个地方，可能理解他们的剑心，可能收留他们这些“叛门”之人。

于是，在接下来数日，陆陆续续有二三十名形容憔悴、身上带伤、却眼神依旧明亮坚定的年轻剑修，跋山涉水，来到了云霞山，找到了云隐别院。

他们对着闻讯出迎的叶牧谦与白言冰等人，悲声陈述了无极山发生的剧变，将刘长靖如何勾结紫阳宗、如何突然发难、师尊独孤明如何宁死不屈、他们如何突围逃散的经过，原原本本，备述一遍。

说到激愤处，不少铁骨铮铮的剑修也忍不住红了眼眶，更有人因伤势和心力交瘁，几乎站立不稳。

叶牧谦听完，沉默良久。他看着眼前这些风尘仆仆、剑心受创却依然挺直腰杆的年轻人，心中涌起复杂的情绪。有对独孤明遭遇的同情与敬意，有对陈炎、刘长靖等人狭隘与狠辣的怒意，也有一种“树欲静而风不止”的无奈。

他知道，问道途的平静发展期，恐怕就此结束了。

“诸位暂且安心在此住下。”叶牧谦最终开口，声音沉稳，带着一种抚慰人心的力量，“云隐别院虽陋，尚可遮风避雨。独孤长老之道友与弟子，便是我叶某之道友与弟子。他日若有机会，必当为你们报仇。”

无处可去的剑修们，闻言心中稍安，纷纷行礼道谢。他们在别院周边，寻了空地，或伐木，或垒石，很快便搭建起更为简陋却足够栖身的庐舍，暂时安顿下来。

\vspace{2em}

而白言冰倒是因祸得福了。在与这些剑修的相互印证和切磋中，他也终于有所体悟。

一日，他主动找到沈瑶：“师妹，我近来有些感触，不如切磋一下。”

沈瑶其时早已重回养气境大成，正随叶牧谦为云霞山布置防御阵“九宫迷踪”。叶牧谦点了点头，于是沈瑶便与白言冰走到演武场上站定。

不少修士其实正在场上苦练战技，见大师兄和三师姐——他们不少人都呼叶牧谦为“师尊”，于是白、陈、沈三人自然就变成师兄师姐了——要切磋，于是纷纷主动退到一旁。毕竟能见到见独境（战斗力甚至相当于真符境）切磋，即使是在一旁观察那也是能学到不少东西的。

“约定一下，不用领域。”白言冰说。

沈瑶秀眉微蹙，不过同意了：师兄看起来是想单纯的磨练剑法。

两人各自拿了趁手的、没开刃的切磋用兵器。白言冰依然拿得长剑，沈瑶同样选择短刀。

陈默正好走了过来，于是自然而然地成了帮两个人敲锣的人。

“当！”

锣声一响，沈瑶立即消失在原地，随即短刀直攻白言冰面门。

这是承袭了她先前还修习灵枢径时的战斗方式：精确、高效。

天下武功，唯快不破！

她与白言冰在禁用领域的情况下切磋过多次，虽然境界有差异，但依然互有胜负。

白言冰不慌不忙，身形微转，长剑忽而一举，堪堪格挡住沈瑶的短刀，长剑周身元气竟然是冰寒剑气。

能挡住？冰寒剑气？

沈瑶没见过白言冰用过冰寒剑气；在以前的大多时候，白言冰还是更喜欢刚猛的雷火剑气。

沈瑶不知道白言冰葫芦里卖的什么药，随即短刀连挥，幻出数道刀影，从不同角度疾刺而来，却全部被白言冰拦下。

场外不少修士看来，白言冰只有防守的工夫、没有进攻的余裕，似乎落入下风。

“看来这一场切磋，是沈师姐赢了啊。”场下有人窃窃私语，“你猜错了，一块灵石赶紧拿来。”

“再等等。”那位猜白言冰获胜的修士小声说，“陈默师兄还没说谁赢谁输呢。”

再看陈默，他只是笑而不语。他跟随叶牧谦多时，自然能看出来白言冰此时在等什么。

沈瑶久攻不进，她渐渐变得焦躁起来。

这白言冰耍她呢，怎么只防不攻？

手上疾速点刺、寻求破防的短刀，因焦躁而自然而然地出现了一瞬的破绽。

殊不知白言冰正等着这一刻！

下一瞬，白言冰动作幅度忽然变大，剑光沿着极为精确的轨迹，寒冰剑气骤然射出，直攻沈瑶刚刚不经意间露出的破绽！

所幸白言冰留了手，给了沈瑶退避的时间，否则沈瑶目前估计已经重伤倒地了！

胜负不言自明，先前那个支持沈瑶的修士的脸立即拉得很长，被迫掏出一块灵石给了他的朋友。

沈瑶往后一闪三尺远，瞅了白言冰一眼，忽然笑道：“白师兄好武艺，这局是我输了。以无数格挡等待破绽，随即一剑封喉……是这样的吗？用冰寒剑气是因为这种剑气更擅长防御……？”

白言冰笑道：“师妹好眼力。前半招叫‘寒光散乱’，后半招叫‘凝萃霜满’……不知道这两个名字如何？”

“好得不能再好了。”沈瑶笑道，“这两招，比你那半成品一般的‘弥焰斩’和‘直环焰’强多了；虽然威力不似那等刚猛，但实战能力却必然更胜一筹。”

\vspace{2em}

无极剑派心剑一脉遭到“清理门户”的消息如同长了翅膀，迅速传遍内域大小势力。

那些嗅觉敏锐的宗主、长老们，从这场突如其来的“清理门户”中，嗅到了不同寻常的气息。虽然说无极剑派两脉中间确实偶发冲突，但从来都是合久必分、分久必合的。

这次被赶尽杀绝，背后必然存在其余推手。

虽然天师府、万仙盟对紫阳宗的公然插手表示谴责，但矛头隐隐指向的云隐别院和其代表的问道途，无不表明一场围绕修行理念、资源分配乃至未来格局的深刻对立，已经浮出水面，且以如此激烈的方式拉开了序幕。

站队，成了许多中小势力迫在眉睫的选择。

审时度势之下，阵营迅速分化、成形：

一方，是以紫阳宗副宗主陈炎为核心，联合了丹鼎阁阁主王丹、无极剑派外剑脉长老刘长靖的守旧联盟。他们代表着内域现有的、以灵枢径为基础、以法器丹药为重要辅助、各大宗门垄断核心资源的修炼体系；其拥趸也大多是与此体系深度绑定、依赖现有格局生存的宗门、商会及相关利益集团。

另一方，则是以云隐别院叶牧谦为核心，并获得了万仙盟盟主吴东浩公开支持的新生势力联盟。这一方吸引了大量在原有体系中难以出头、资源匮乏的低级散修，许多受大宗门挤压、寻求出路的小型宗门和家族，以及那些对问道途理念心生向往、或单纯厌恶守旧联盟霸道行径的修士。

沈天穹在公开场合仍选择保持超然的中立姿态，呼吁双方克制，调停冲突。但其立场明显偏向叶牧谦一方：毕竟白言冰和陈千羽再怎么说也算是小天师，与天师府的关系极为密切。

沈天穹私下里对云隐别院的回护与默许，明眼人都能看出。

\vspace{2em}

围绕着云隐别院与守旧联盟的，是一场压抑而充满张力的冷战。

大规模的热战并未立即爆发，无论是出于对双方可能实力的忌惮，还是对新生势力那种未知韧性的评估，双方都保持着一种危险的克制。但这绝非和平，而是在厚重冰层下蓄积暗流，寻找更能一击致命、且不必担负首要开战罪名的打击方式。

很快，一场没有硝烟却足以扼杀生机、更加残酷的博弈，悄然拉开了帷幕。

数日后，一道盖有陈炎宝印、措辞空前严厉的“宗门戒令”，通过其掌控的渠道，迅速传遍内域各大主要坊市与势力：

“即日起，凡我紫阳宗弟子，曾私自前往云隐别院听道、论法、交流者，概视为受异端邪说蛊惑，心志不坚，道基已污！限三日内至刑律殿自陈过错，听候发落。逾期或隐瞒不报者，一经查实，即刻废去修为，逐出宗门，名录剔除，永不收录！

“自今往后，紫阳宗上下，严禁任何弟子、执事、长老与云隐别院之人有任何形式的往来、交流、贸易、传讯，违者以叛宗论处，严惩不贷！”

这道戒令将云隐别院及其代表的问道途，正式定性为“异端”、“禁地”。

它试图斩断一切明面上的人际纽带，将叶牧谦一方置于一个被孤立、被污名化的孤岛之上。

戒令在内域舆论场上激发的余波尚未平息，丹鼎阁与已由刘长靖完全掌控的无极剑派（外剑一脉）便迅速跟进，发布了措辞几乎雷同的禁令。三大超级势力的联合封杀，形成了一道令人窒息的无形壁垒。

威慑是立竿见影的。

许多原本与云隐别院有些零星交易往来、或纯粹出于对新生事物好奇而前去旁听过的修士、小商会代表，纷纷噤若寒蝉。

原本络绎通往云霞山的小道上，新面孔骤减；偶尔有目光投向别院方向，也迅速移开，仿佛那风景秀丽的半山腰弥漫着无形的疫病。除了早已定居在此的散修和那些从无极剑派逃难而来、无处可去的心剑脉弟子，云隐别院周围，确实显出一种被主流世界刻意疏离的冷清。

\vspace{2em}

然而，真正能扼住一个势力命脉的，不能仅仅凭人际上的孤立。

一切问题本质上都是经济问题；从古至今，无论中外，乃至这个异世界，资源与实业永远是统治的根基，也是最锋利的武器。一场精心策划、目标明确的涨价风暴与“货源调控”，在看似正常的市场波动掩护下，猛烈袭来。

首先遭殃的是公开坊市。

在紫阳宗、丹鼎阁影响力深厚的几处大型坊市，流向非守旧联盟认可势力——尤其是与云隐别院有牵连者——的各类基础修炼资源，价格开始诡异地、毫无缘由地连续飙升。

“赤铜精铁，昨日还是十灵石一方，今日怎就十五了？” 一个试图为同伴购买铸剑材料的心剑脉弟子，对着摊位后的掌柜愕然道。

掌柜眼皮都未抬，懒洋洋道：“货源紧张，上游提价，爱买不买。对了，看阁下气息，似是云霞山那边来的？若是，价格还得再加三成，毕竟……风险大嘛。” 话语中的奚落与歧视，毫不掩饰。

类似的场景不断上演。沉银砂、寒铁木等常用炼器辅料，化瘀草、宁神花等基础疗伤药材，涨幅普遍达到三到五成，个别紧要物资甚至翻倍。

明眼人都能看出，这显然不是市场的自发调节，而是一种精准的、带有惩罚性质的宏观调控，是一种歧视性的、甚至制裁性的定价。

紧接着，打击蔓延到了更基层、更关乎日常修炼生存的领域。

一些云隐别院周边散修聚居地日常所需的低阶资源，如最基础的凝气草、铁精矿，开始出现诡异的货源紧张。原本定期前来兜售的行商不见了踪影；本地几个小供货点要么关门，要么货架上空空如也，偶有存货，价格也已高到令人绝望。

“王管事，这个月的凝气草……” 一个面容愁苦的老年散修，小心翼翼地问道。

被称为王管事的胖商人叹了口气，压低声音：“老李头，不是我不帮你。上头打了招呼，凡是往云隐别院那边走的货，一律卡死。我这儿但凡卖你一株，明天我这铺子就别想开了。你……另寻门路吧。” 话语间，他眼神躲闪，满是无奈与畏惧。

此计堪称毒辣，虽兵不血刃，但确实击中了众多问道途依附者的要害；其目的阴损而深远，绝非仅仅是提高成本、大赚一笔那么简单。

首要目标，自然是尽可能地消耗新生势力联盟本就不甚宽裕的财力。

对于已经开始转修或兼修问道途，如白言冰、陈千羽这般修士而言，对外部灵石、丹药的依赖极低；但是新生势力联盟之下，还有大量的、需要资源的灵枢径修士：无论是心剑脉弟子维护保养剑器，散修们维持基本修炼，还是尝试自给自足进行低阶炼丹、炼器，都离不开这些基础资源。

其次，这是一次强大的威慑与站队逼迫，做给内域所有摇摆不定的中小势力看：站在守旧联盟一边，你们的资源供应渠道将得到保障；但若与云隐别院、与问道途扯上关系，哪怕只是一点生意往来，等待你们的将是同样的“卡脖子”命运。

修炼资源断供，商会贸易受阻，在这种依靠外物的修行界必然寸步难行。

这是一道冰冷的选择题。

价格飞涨与货源断绝，如同两只缓缓收紧的铁钳，要让叶牧谦他们那宝贵的灵石储备迅速见底，让日常修炼陷入停滞，让刚刚凝聚起来的人心因生存困境而涣散。

一名年轻的心剑脉弟子，看着手中因崩出细小缺口的长剑，眉头紧锁。修复材料的价格早已涨上了天，而且有价无市。他跑遍了附近三个小镇，一无所获。

几个常驻的散修围坐在一起，面色晦暗。他们修为低下，仍需依靠凝气草辅助引气，往日攒下的一点灵石，如今连购买往常三分之一的份量都不够。修炼进度近乎停滞，焦虑与绝望的情绪在蔓延。

甚至连云隐别院内部，也开始感受到压力。虽然叶牧谦讲道不依赖于外物，但聚集而来的数百人中，大半仍需基础资源维系。

别院原本的少量储备正在快速消耗，负责后勤的弟子面带忧色，向叶牧谦汇报着日渐窘迫的处境。

无形的壁垒在人际间筑起，有形的资源枷锁则在经济命脉上狠狠绞紧。云隐别院仿佛狂风恶浪中的一叶孤舟，虽然核心依然稳固，但船体正在承受一波接一波的侵蚀与挤压。

冷战阴云笼罩下的云霞山，看似平静，实则每一个角落都弥漫着生存的艰难与抉择的重量。

而这场没有刀光剑影的围剿，才刚刚开始。

\vspace{2em}

经历此事，叶牧谦忽然觉得他把人想得太好了。

自己穿越时正值中美两个庞然巨物相争，如此情景真有一种前世的感觉。

前世记忆中，因为中国工业品类齐全，美国只能找到高端芯片等少数几项技术去卡中国的脖子；但叶牧谦手下的这套烂摊子，基本上和八十年前、百废待兴的中国没什么两样，什么都没有，什么都造不出来。

中国当时至少还有北方老大哥的援助，可云隐别院呢？

“造不如买”这种思维还是太诱人了，殊不知这是把命脉悄然交到了别人手上……关系好的时候，资源流通、互通有无，那无疑是好的；但剑拔弩张时，双方鼻青脸肿、头破血流，谁还跟你自由贸易？

“前辈们还真是高瞻远瞩啊……”叶牧谦长叹道。

但生存确实是一个迫在眉睫的问题。单纯的防御与抱怨无济于事，必须主动破局。

叶牧谦深知，己方的根本优势在于问道途自身对于外部资源依赖极低的韧性，而自此衍生出来的重要优势自然是汇聚人心、团结力量办大事。

他将吴东浩、白言冰、陈千羽、沈瑶等核心人物召集起来。没有冗长的争论，一场务实而高效的商议迅速展开。针对守旧联盟的封锁，一套立足自身、着眼长远的反制体系被勾勒出来，并立即付诸行动。

核心思维分两步走：

第一步，用尽婉转手段，在封锁没有彻底合围的时候，尽可能打破禁运的死局，钻空子、走私，至少先保住自己境内许多灵枢径修士的日常生活；

第二步，则是拉动内生产，降低物资依赖程度，最终达到部分基础资源的基本自给。

万仙盟盟主吴东浩当机立断，派出了自己的儿子——吴刚。这位少盟主骨龄不过四十，外貌则依然保持在二十岁左右的样子，修为已达玄脉境融贯期。此人看似潇洒不羁，实则心思缜密、长袖善舞，在万仙盟遍布内域的松散人际网络中拥有极高的信誉与号召力。

吴刚领命后，第一个找到的帮手，便是成功转修问道途、对云隐别院归属感极强且常年混迹散修圈子的陈默。

一个代表组织力量，一个深谙草根生态，两人的结合立刻迸发出高效的能量。

他们避开守旧联盟严密监控的大型坊市与主要商路，将目光投向了那些常年被忽略的边缘角落与灰色地带。

依托万仙盟那张看似松散却触及每个角落的人脉网，吴刚与陈默开始秘密联系那些走南闯北、熟悉每一条偏僻小径与隐秘山谷的资深散修，联系一些因不满盘剥而私下开采边缘矿点的小型团伙，甚至联系一些在守旧联盟商会体系内不得志，或者愿意冒险赚取更高差额利润的底层管事。

一张隐秘而高效的地下交易网络，如同蛰伏于地层之下的根须，开始悄无声息地蔓延生长。这张网络分工明确，构成了一个粗糙但实用的闭环：

采购由熟悉当地情况的散修负责，从守旧联盟控制力薄弱的偏远矿区、人迹罕至的深山，以“捡漏”或小规模盗采的方式，零散获取品质可能不高但勉强可用的矿石；或是向一些与丹鼎阁没有直接契约关系的山民、采药人，或伪装成守旧联盟收货人，或采用其余手段，收购他们采集的野生草药。

另一面，在几个重要的散修聚集点或小型城镇的暗处，设立了隐秘的交接与仓储点。物资在此汇总、分类，再根据云隐别院总部及各处新生势力据点的需求清单进行分发。

交易大多采用以物易物的方式——用云霞山产的某些草药、或从旧法器上熔炼回收的金属，交换急需的矿石或其他药材，最大限度规避对灵石的依赖和金融追踪。

运输反而是最危险也最关键的一环。一批胆大心细、修为不弱且信誉良好的散修被组织起来，伪装成普通的行商、猎户甚至逃难的流民，利用多条隐秘小道进行物资转移。他们往往昼伏夜出，分段接力，并将物资化整为零，以降低单次运输的风险。

经过数次失败之后，具体的采购、运输线路也从总线联系变成了单线联系，形成彼此隔离的链条；即使某一环被发现并切断，损失也被控制在局部，不会导致整个网络瘫痪。

虽然这条地下血管输送的物资量，远无法与昔日自由市场的充沛相比，运输过程也伴随着损失和风险，但它确实像涓涓细流，顽强地将各地零星的资源汇拢起来，持续输送到急需给养的新生阵营，勉强维持了生存与最低限度的运转。

其次，是推动就地生产，降低核心依存度。陈千羽在这方面发挥了不可替代的作用。

她清醒地意识到，药材，尤其是基础疗伤和恢复类药物，是联盟修士的生命线；完全依赖外部输入，等于是将命门始终握在敌人手中。

很快，一批因为性格耿直得罪上司、或因出身寒微备受排挤、或在丹鼎阁内部权力倾轧中失势的药修，被秘密招揽或主动投奔而来。同时，聚居地中许多原本从事过农耕、或对培育植物有天然兴趣的散修也被动员起来。

在陈千羽的带领下，他们在云霞山脉几处灵气相对温和、土质适宜的背阴山谷与向阳坡地，以最原始的方式，开辟出了一片片错落有致的梯田式药圃。

种植不是简单地移植，那样太慢了。

陈千羽以问道途元气温养为基础，自己体悟的内在平衡与生生不息的理念为指导，与众多现有修士所知的传统灵植术相结合，进行了一系列因地制宜的改良。

她指导药修和散修们，注重不同草药伴生带来的生态平衡，利用山间溪流、雾气、不同朝向的光照来模拟最佳生长环境；虽然受限于环境与时间，无法种植那些需要特殊地脉、动辄数十年甚至上百年才能采收一次的高阶灵药，但诸如止血藤、宁神花、凝气草、化瘀草等最常用、消耗量最大的基础药材，自给率开始稳步提升。

每一株在云霞山泥土中壮大的草药，都减少了一分对丹鼎阁的依赖。

与此同时，一批对炼器有兴趣的散修和部分心灵手巧的心剑脉弟子，在别院一角搭起了简陋的工棚。利用地下网络艰难运来的、品质参差不齐的低阶矿石，他们从零开始，学习最基础的材料辨识、提纯和粗坯锻打，叮叮当当的敲击声终日不绝。

虽然蕴含符文的真正法器的炼制效率依然极低，但打造出质地坚韧的普通刀剑、农具，修复日常损坏的器械，甚至尝试将废料重熔再利用，都能显著减少对外界制品的依赖，并将重要的金属资源尽可能留在内部循环。

\vspace{2em}

这场没有硝烟的经济暗战，在平静的表面下激烈地进行着。

守旧联盟的探子像猎犬一样四处嗅探，试图找出地下网络的关键节点加以破坏。

双方在偏僻的山道、无人的峡谷发生过数次小规模的遭遇与冲突，互有死伤，鲜血染红过溪流，但这条隐秘的生命线始终未曾彻底断绝。

坊市中的价格博弈每日上演，守旧联盟提价、控货，地下网络就寻找替代品、开发新来源，或默默承受成本，将压力分散到整个共担风险的网络之中。

云隐别院周边的聚居地，日子过得肉眼可见的清苦。灵石永远不够用，饭菜少见荤腥，衣物修补再三。但一种奇特的凝聚力，却在同舟共济、共克时艰中悄然滋生。

确实有不少人最终撑不住，选择离去或倒戈，但留下的无不是心性坚定、对问道途抱有信任、对未来充满希望的修士们。

人们清楚地知道，守旧联盟是想用资源和饥饿逼他们就范、自行瓦解。因此，每一条穿越封锁线安全抵达的药材矿石，每一片在梯田上迎风摇曳的葱绿药苗，甚至每一把自家工棚打出的粗陋铁器，都不仅仅是一件物品，而是对那无形壁垒最直接的反抗与宣言。

希望，在这种具体的、触手可及的劳动与收获中，变得真切起来。

然而，叶牧谦比任何人都清楚，经济战的消耗是巨大且不对等的。

地下网络的运转本身就需要付出高昂的风险成本，药田的丰产需要至少一两个生长周期的等待；而守旧联盟掌控的财富与资源毕竟雄厚百倍。长期僵持的消耗战，对于新生的一方绝非良策。

而就在这时，一个意外的消息，再次改变了局势的走向。

\vspace{2em}

天枢城方向，忽然天生异象，霞光万道，瑞气千条，一股奇异的、仿佛来自另一片空间的波动传遍四方。

紧接着，天师府正式发布通告：

“新小玄界，即将开启。”

小玄界，一般指的是一些破碎的界域；它们的出现和消失毫无规律，其间或是生机盎然，或是荒芜死寂，几乎完全无法预测。内域对其成因也一无所知，但这并不妨碍每一次出现都引发巨大轰动。这意味着未知的机遇、失传的传承、珍稀的天材地宝，甚至可能是改变命运的一线曙光。即便可能伴随着莫测的危险，也足以让任何势力与修士心动。

\vspace{2em}

在持续数日的异象与全城骚动，以及秘密书信沟通后，天枢城内，一场气氛凝重的临时会议开始举行。

与会者，正是当下内域两大对立阵营的代表人物。

守旧联盟一方，紫阳宗副宗主陈炎面色冷硬，丹鼎阁阁主王丹神情淡漠，无极剑派外剑脉长老刘长靖眼神锐利如剑。新生势力一方，云隐别院叶牧谦安然端坐，万仙盟盟主吴东浩脸上惯有的和煦笑容也收敛了几分。天师府府主沈天穹作为东道主与调停者，居于主位。

空气中弥漫着无形的张力，比之前任何一次公开对峙都要紧绷。

经济战的暗流尚未平息，无极剑派的血案尤有余温，此刻却要因为一个小玄界坐在一起，场面堪称诡异。

“诸位，”沈天穹打破沉默，“小玄界开启，乃内域共有的机缘，亦是潜在的危机。若因彼此争端，在入口处便大打出手，或是在探索中互相倾轧过甚，非但可能导致机缘尽丧，更可能引发不可测的界域反噬，殃及天枢城乃至更广范围。此非吾辈所愿见。”

陈炎冷哼一声：“府主之意，是要我们与这……‘异端’合作？”

他刻意加重了最后两个字，目光如刀刮过叶牧谦。

叶牧谦眼皮都未抬一下，只是淡淡道：“道不同，不相为谋。然机缘在前，若因门户之见错失，亦非智者所为。叶某只问一句，陈副宗主敢不敢让门下年轻弟子，与叶某的弟子，在‘公平’的环境下，各凭本事，争一争这机缘？”

虽然叶牧谦是一个和平主义者，连带着整个问道途人都是内求一身、不喜争斗的性子；但此时两方已经是撕破脸皮，那叶牧谦自然也不介意与他们争一争。

“公平？何谓公平！”刘长靖接口，语气带着讽刺，“你那些弟子，修的路子古怪，谁知道会不会用些不上台面的手段？”

吴东浩呵呵一笑，摇着折扇：“刘长老此言差矣。既然是探索未知秘境，自然是手段尽出，只要不违事先约定即可。莫非贵派弟子，只会按部就班，离了宗门护佑便寸步难行？”

眼看又要吵起来，沈天穹抬手制止：“够了！”他目光扫过众人，“我们也都知道，此前已有书信沟通基础。今日便定下章程：两派搁置争议，共同探索此界。为免事态升级，仅限骨龄百岁以下的弟子进入。小玄界内，生死祸福，各凭机缘本事，外界不得干涉。出得小玄界，恩怨亦不得追溯至界内之事。诸位可有异议？”

这条件颇为苛刻，“生死不论”四字更是透着血腥。但这也最大程度避免了外界势力直接插手，将冲突限制在年轻一代的竞争层面。

陈炎等人虽然不甘，但也明白这是目前最能接受的方案，既能打压新生势力年轻一代，夺取机缘，又不会立刻引发全面战争。叶牧谦一方则需要这个机会，让弟子历练，并争取资源打破封锁。

最终，双方虽然在细则上磨了半天嘴皮子，但最终还是在“搁置争议、共同开发”的基础上基本达成了共识。

\vspace{2em}

三日后，天枢城北三百里，裂云峡谷。原本荒芜的峡谷深处，空间扭曲，那道七彩漩涡般的入口稳定了下来，高约三丈，宽两丈，内部光影朦胧，看不真切。

峡谷两侧，泾渭分明地站着两队人马，总数约六十余人。

左侧，守旧联盟阵营。紫阳宗弟子清一色赤红劲装，气息灼热，大多在玄脉境开脉期以上，为首一名面容倨傲、目蕴精光的青年，正是紫阳宗当代圣子周文彬，修为已达真符境期摹形巅峰，是场上唯一真符境修士，真不愧“内域第一天才”的称号。

丹鼎阁弟子则衣着素雅，以青、白二色为主，身上带着淡淡药香，虽不擅强攻，但各种辅助、疗伤、甚至用毒的手段不容小觑。

无极剑派弟子人人负剑，剑气凌厉，站姿如松，眼神中带着经过“清理门户”后愈发浓烈的排外与敌意。

右侧，新生势力阵营。云隐别院一方，白言冰立于最前，气息沉凝，见独境界的圆融之意让他即便在众多灵力勃发的对手面前也不显弱势。陈千羽站在他身侧稍后，温婉宁静，腰间悬挂针囊。沈瑶则位于另一侧，黑衣飒爽，容颜清丽，眼神锐利地扫视着对面。

他们身后，是十几位修为在养气到见独初期不等的云隐别院核心弟子，以及数位自愿参与、经过筛选的散修好手。万仙盟一方，则由吴东浩之子吴刚带领，约二十余人，服饰杂乱却个个眼神精悍，透着散修特有的野性与机警，陈默也在其中。

除此之外，还有少部分天师府弟子，不过十人。他们保持着中立，在剑拔弩张的两派之间，显得有些格格不入。

这些弟子们事先也都签了生死状。

峡谷上方，沈天穹、陈炎、王丹、刘长靖、叶牧谦、吴东浩等人的身影在云层中若隐若现，强大的神念笼罩着下方，既是相互监督，也是防止有人破坏约定。

“入口已稳，时限一个月。一个月后，无论有无收获，必须退出，否则界域闭合，神仙难救。”沈天穹恢弘的声音传入每个弟子耳中，“现在，进入！”

话音落下，两侧队伍几乎同时动了！

没有寒暄，没有试探，最初的冲突在踏入秘境的第一步就已爆发！

守旧联盟弟子显然早有预谋，数名紫阳宗弟子在靠近入口时，突然齐齐催动灵力，炽热的灵力并非攻向对手，而是猛然轰击在入口边缘不稳定的空间波纹上！

“轰！”

本就动荡的入口顿时一阵剧烈扭曲，七彩光华乱窜，形成一股混乱的灵力乱流和空间撕扯力！他们的目的很简单，制造混乱，干扰甚至重创即将进入的新生势力弟子，最好能让他们分散传送到危险区域。

“卑鄙！”白言冰眼中寒光一闪，一股圆融稳固的势场以他为中心微微扩散，将身旁最近的陈千羽、沈瑶以及几名弟子笼罩。

与此同时，吴刚那边经验丰富的散修们更是各显神通，有的身法如烟，瞬间穿过乱流薄弱点；有的祭出粗糙但实用的防御符箓硬扛；陈默则闷哼一声，体内初成的元气鼓荡，护住周身，强行冲入。

乱流之中，人影纷乱。大部分人都成功冲进了入口，但仍有数名新生势力的散修和一名云隐别院弟子被乱流扫中，惨叫着被抛飞出去，受伤不轻。而守旧联盟那边，也有两名弟子因操控灵力过猛，自己也被反震所伤。

\vspace{2em}

光芒一闪，天地变换。

经天师府先头部队初步探查，小玄界与上古神鸟“青鸾”相关，姑且称之为青鸾秘境。

青鸾主生机平衡，在内域尚未绝种的时候，基本都能达到法相境界甚至传说中的更高境界，其留下的传承预计是极为可观的。

白言冰稳住身形，发现自己站在一片苍翠欲滴的古老森林边缘。古树参天，藤蔓垂落如帘，空气中弥漫着浓郁到化不开的草木清香与精纯的生机灵气，吸一口都让人精神一振。远处有鸾鸟清鸣般的回响，空灵悠远。

这秘境环境，竟是难得的鸟语花香、生机盎然之所在。

他迅速环顾四周，心中稍定。陈千羽、沈瑶就在他身边不远处，另外还有八名云隐别院弟子和五名万仙盟的散修（包括吴刚和陈默）也陆续在附近出现。进入时的默契配合，让他们这一批近二十人的队伍得以相对集中地传送到同一片区域。

这已是极大的优势。

“清点人数，检查伤势。”白言冰沉声下令。

很快，结果出来。他们这一队十六人，除了两人在入口乱流中受了轻伤，并无大碍。

陈千羽立刻上前为伤者紧急处理伤势。

“不可大意。”沈瑶警惕地观察着周围，“青鸾主生，亦可能意味着其考验或守护力量，皆与生机相关；此地生机盎然，妖兽自然也不可能少了。而且那些家伙，也肯定进来了。”

她话音刚落，左侧密林中便传来一声急促的惨叫，紧接着是灵力爆发的轰鸣和怒喝声！

显然是其他分散传送进来的弟子遭遇了袭击，不知是秘境本身的危险，还是人为。

“走，过去看看，小心戒备。”白言冰握紧剑柄。众人立刻结成简单的阵型，白言冰在前，沈瑶在侧翼，陈千羽和伤员在中间，吴刚、陈默等散修经验丰富，负责殿后和探查周围。

十几个人向着声音来源谨慎靠近。

没走多远，他们便看到了惨烈的一幕：一片林间空地上，三名服饰各异的散修——似乎是追随万仙盟的小势力弟子——倒在血泊中，两人已然毙命，一人重伤濒死。

而袭击者，赫然是五头通体碧绿、形如猎豹但头上生有独角的奇异妖兽，它们爪牙锋利，身上缠绕着淡淡的青色风旋，行动如电，正在撕咬尸体。

“是风刃豹，二阶妖兽，擅长风系法术，速度极快！”一名见多识广的散修低呼。

妖兽，通俗的说就是修炼有成的动物，至于这个“修炼有成”是怎么回事就不一定了。有些妖兽是先天的，即生下来就是妖兽；有些妖兽则可能是普通动物血脉觉醒或偶得机缘，后天成为的。内域把妖兽按内丹品质分成不同的阶位；一头二阶妖兽的战斗力，至少相当于一位玄脉境修士。

而这里有五头！

更让白言冰等人瞳孔收缩的是，在空地另一边的树冠上，隐约可见几道身影悄然隐匿，依稀看出其制服是紫阳宗和丹鼎阁的人！

他们冷眼看着下方散修被妖兽屠杀，毫无援手之意，甚至有人脸上露出残忍的快意。

树上的人气息收敛得良好，这群豹子似乎也没有抬头发现他们。

“他们八成是故意的！”吴刚咬牙道，“定是比我们更早传送至此，惊动了这群风刃豹，却故意引到这几人附近，自己躲起来看戏！”

这是赤裸裸的借刀杀人！守旧联盟果然一进来就开始下绊子。

似乎是察觉到白言冰等人到来弄出的动静，下方正在进食的风刃豹群突然抬头，碧绿的眼眸瞬间锁定了白言冰这支人数更多的队伍！

“吼！”

五头风刃豹舍弃了残破的尸体，化作五道青色流光，带着尖锐的破风声扑来；

尚未近身，便有十数道半月形的淡青色风刃抢先撕裂空气，覆盖而至！

那是风刃豹捕食之前释放的第一招，极为阴损，追求在近身搏斗之前就重创食物！

“结阵！三三制！”白言冰厉喝，声震山林。

这是进入秘境前，叶牧谦特意传授给所有参与弟子的简易战法，核心便是“三人一组，互为犄角，攻守一体，遥相呼应”，专门应对被围攻、偷袭或遭遇强敌时，避免落单被逐个击破。

实际上，这确实要拜叶牧谦前世记忆所赐，不少军队会使用这种战术队形；而更具体的实践，则完全交予这些弟子去训练与体悟了。

十六人瞬间动了起来！

白言冰作为唯一的见独境，修为最高，自成一组坐镇中心；沈瑶、陈千羽、吴刚自然成一组，作为主力。其余修士也与相熟的散修快速三三结阵。

面对覆盖而来的风刃，白言冰剑光一展，领域大放，圆融的剑意在身前划出一道弧线，将射向众人的大部分风刃引偏、卸力。

其他小组亦是或挡或闪，配合默契，虽略显仓促，却稳稳接下了第一波风刃袭击，无人受伤。

但风刃豹群速度极快，已扑至近前，獠牙利爪带着腥风。

“散！”白言冰再次下令。

五个三人小组如同花瓣般骤然散开，却又保持着彼此能够迅速支援的距离。

每组面对一头或分摊风刃豹的攻击，绝不贪功冒进。

白言冰自己则对上为首最强壮的一头，剑势沉稳，每一击都重若千钧，逼得那风刃豹连连后退，无法发挥速度优势。

其余修士亦利用对地形的熟悉和丰富的搏杀经验，很快也压制住了各自的对手。

树冠上，那几名守旧联盟弟子脸色微变。

他们没想到这群乌合之众反应如此迅速，战法如此古怪有效，竟然稳稳挡住了五头二阶风刃豹的突袭！

这“三三制”看似简单，却将个人战力有效整合，弥补了单个修士可能存在的短板，尤其擅长应对这种混战和偷袭。

眼看借刀杀人之计就要破产，其中一名丹鼎阁弟子眼中狠色一闪，悄悄取出并打开了一个小纸包。

“嗤——！”

一股极其细微、却充满躁动炽热气息的粉末随风飘散，混入战场。

正在与白言冰激斗的那头最强壮的风刃豹吸入粉末后，碧绿的眼眸瞬间爬满血丝，仰天发出一声痛苦与暴怒混杂的咆哮。

随即，周身青色风旋骤然变成混乱的赤青交错，体型似乎都膨胀了一圈，气息暴涨！它完全放弃了防御，不顾一切地扑向白言冰，速度力量暴增！

“小心！他们用了药！”陈千羽惊呼，她嗅觉敏锐，对药性敏感，立刻察觉不对。

白言冰压力陡增。但他临危不乱，见独领域全力展开。

周遭一切似乎变慢了些许，能清晰感知到疯豹扑击的轨迹与意图，以及疯豹身上肌肉力量的强弱分布。

他身形开始如柳絮般，随着疯豹带起的劲风飘退，剑尖却疯狂点向疯豹的关节处。

“沈瑶，吴兄，速战速决！”白言冰喝道。

沈瑶会意，身法陡然变得凌厉无比，刀光分化，让人眼花缭乱，速度快到让旁人以为她似乎是生出三头六臂、每一只手都以极高的频率向妖兽攻去。

瞬间，她对面的风刃豹身上留下数道深可见骨的伤口，且伤口开得颇大，特意选择了几条大动脉处放血；又戳烂了它的心与肺。

那豹子身上多处重要器官被废，又失血过多，登时就没了气。

吴刚也怒吼一声，与同伴合力，拼着受了一记爪击，将另一头风刃豹的头颅斩下。

那头疯豹却炸毛起来，久攻不下，愈发狂躁。

陈千羽觑准一个空隙，算明那头疯豹下一步的落点，随即打出一发微弱的元气扰动，竟让那片地面的草木忽然变成软垫。

疯豹本欲扑击，一脚踏上去，忽然爪下一软，平衡稍失，身体向侧方倾倒。

白言冰岂会错过这机会？

蓄势已久的一剑，凝聚了自身心力与元气，毫无花哨地直刺而出，正中疯豹因失衡而暴露出的咽喉弱点！

“噗嗤！”

随即手腕带着长剑向侧方一划！

剑刃透颈而过，疯豹头颅落地，咆哮戛然而止，庞大的身躯轰然倒地，赤青色的混乱风旋缓缓消散。

很快，最后两头风刃豹也被众人合力击杀。

战斗结束，众人微微喘息，除了吴刚和另外两名散修受了些轻伤，均无大碍。他们迅速收拾战场，把五只死豹剥了皮、挖了内丹，并将那重伤未死的散修同伴救起，由陈千羽进行急救。

白言冰这才倒出工夫抬头，冷冷地望向那处树冠。

那里，早已空空如也。

那几名守旧联盟弟子许是见事不可为，悄然遁走了。

“这只是开始。”白言冰擦去剑上血迹，目光投向森林深处那隐约传来的、更宏大的生机波动方向，“他们不会罢休，这秘境本身，恐怕也危机四伏。所有人，分成两队。我与吴兄各带一队，分散开来各自探索。”

吴刚点了点头。

虽然呆在一起确实更加安全，但分散开来也能够有更大的搜索范围，并在沿途获取更多资源。

休息大约一刻的时间后，两队分开，开始一头扎进森林各自探索。

\vspace{2em}

森林远比看上去的广阔和复杂。

那些青玉般的树木似乎隐隐构成某种天然的阵势，行走其间，方向感极易迷失。脚下厚厚的、散发着微光的苔藓和蕨类植物踩上去松软无声，却可能隐藏着致命的陷坑或毒虫。

空气中飘荡的芬芳里，偶尔会混入一丝甜腻得过分的气息，那往往是某些致幻灵花散发出来的，需得时刻运转心法保持灵台清明。

行进了约莫一个时辰，他们遭遇了又一次考验。那是一片看似平静的林间沼泽，水色清亮，生长着亭亭玉立的金色莲花，莲心凝结着露珠般的灵液，香气诱人。一名年轻些的云隐别院弟子面露喜色，下意识就要上前采集。

“别动！”白言冰和陈千羽几乎同时低喝。

话音未落，那弟子脚下看似坚实的苔藓地突然软化塌陷，数条漆黑如墨、快如闪电的藤蔓破水而出，直刺向他周身要害！藤蔓上密布细小的倒刺，闪烁着幽蓝光泽，显然含有剧毒。

那弟子骇然失色，慌忙催动灵力护体，挥剑格挡。但他仓促间的反应慢了半拍，眼看就要被藤蔓缠住。

嗤！嗤！嗤！

数道细微几乎难以察觉的破空声掠过。陈千羽出手如电，三根细长的钢针后发先至，精准地钉在了三条藤蔓的关节处。中针的藤蔓动作顿时一僵，速度骤减。

与此同时，白言冰并指如剑，凌空一划。一道圆融凝练的剑气匹练般扫过，巧妙地将剩余几条藤蔓的力道引偏，粘滞的剑势让它们如同陷入无形泥潭，挣扎着偏离了方向。

沈瑶的身影鬼魅般则出现在那名弟子侧后方，一把抓住其衣领向后急拉，险之又险地脱离了沼泽边缘。

整个过程发生在电光石火之间。那弟子落地后，脸色煞白，冷汗涔涔，几乎站不住，连忙向陈千羽和白言冰道谢。

“墨龙藤，植物，但有低等妖性。群居，擅伪装偷袭。”陈千羽仔细观察着缓缓缩回水中的藤蔓，冷静地分析。

白言冰目光扫过沼泽，注意到水下隐约有更多黑影蠕动，沉声道：“人命是第一位的。这东西不好拿，就不要了。记住，越是看起来无害诱人之物，越需警惕。”

这个小插曲让队伍更加谨慎。

\vspace{2em}

他们继续深入，期间又避开了几处气息诡异之地，也小心翼翼地采集了一些确认无害、且对修为有裨益的灵草灵石，在保证安全的情况下算是雁过拔毛。

约半日后，他们又听到了前方传来打斗声和妖兽的怒吼。

靠近一看，竟是三名守旧联盟的弟子正被七八头风刃豹围攻。

地上已经躺着两头豹子，但那三名弟子也颇为狼狈，衣衫破损，其中一人胳膊鲜血淋漓，显然受了伤，行动迟缓。

“又是风刃豹？”沈瑶低声道。

白言冰目光锐利地注意到战场边缘，几株被粗暴砍断的、结着朱红色果实的奇异灌木。灌木断口处流出乳白色汁液，散发出异常浓郁的甜香。

陈千羽翻了翻药典：“血朱果，对炼体有奇效。除了肉食，风刃豹特别喜欢吃这种东西。”

他们定是贪图果实，惊动了豹群。

此时，战况对守旧联盟弟子越发不利。

风刃豹攻击配合默契，速度如风；因为失血过多，那名受伤弟子的动作愈发迟缓。眼看又一头豹子要从侧后方扑向那名受伤者，带队的那名紫阳宗弟子虽奋力挥出火法逼退正面之敌，却已救援不及。

是袖手旁观，看着竞争对手减员？还是……

白言冰眼神一凝，瞬间有了决断。

“沈瑶，左翼扰袭！千羽，准备包扎！你们俩，跟我上，救人！”

话音未落，他已身化流光掠出，巧妙切入侧翼，剑指虚点，数道凝实的剑气精准地射向几头豹的落足点，打乱了它们围攻的节奏。

沈瑶如一道黑色轻烟，倏忽出现在战场另一侧，手中短刃带起森寒光弧，以刁钻的角度袭向风刃豹的眼睛、关节等脆弱处，迫使它们分心防御。

另外两名云隐别院弟子紧随白言冰，结成简单剑阵，挡住了一侧扑来的两头豹。

突如其来的援手让三名守旧联盟弟子压力一松。带队者愣了一下，神色复杂地看了白言冰一眼，猛一咬牙，趁机催动更强术法，配合反击。

混乱中，陈千羽悄无声息地靠近，银针连闪，先为那名重伤弟子封住血脉，弹入一枚清香药丸暂时压制伤势，又挥手洒出一片淡绿色药粉。药粉弥漫开，带着清心解热的气息，让暴躁的豹群产生了瞬间的迟疑。

趁此机会，白言冰低喝一声：“退！”

新生势力的人加上被救下的守旧联盟三人，迅速脱离战圈，向后疾退。豹群追出一段距离，发出不甘的咆哮，但似乎又顾忌领地，于是转身离去了。

脱离危险后，那三名守旧联盟弟子，尤其是带队的紫阳宗弟子，面色一阵青白。

他看了看自己受伤的同伴，又看了看神色平静的白言冰等人，抱了抱拳，干涩地道：“在下韩铁，多谢……援手之情，我等记下了。”

说完，不再多言，扔下几枚血朱果，就迅速带着同伴离去，背影有些仓惶。

他显然是不知道，白言冰一行人刚刚进入秘境时也遇到了相似的风险；他更不知道，他的同宗面对几乎完全一样的情景做了些什么。

他们似乎是没脸在此多做停留，也有可能是害怕被人告发“通敌”。

白言冰见他们竟然扔下几枚果实，有些愕然，但最终还是给它们收了起来。

实际上所有的冲突都是这样的：大多数人其实是沉默的、中立的，他们虽然心中可能并不厌恶甚至同情云隐别院，但大势所迫之下，为了明哲保身，也不得不被裹挟。

“为何救他们？”一名云隐别院弟子不解地问，“入口处他们可是下了黑手。”

白言冰淡然道：“见死不救，与那些只知内耗之徒何异？我等所求之道，非仅力量强弱，亦有心境高低。”

陈千羽微微颔首，看向白言冰的目光中带着赞同。沈瑶则若有所思。

\vspace{2em}

收拾心情，他们继续前行。

随着不断深入，周围的古木愈发高大，生机愈发磅礴，甚至能看到一些外界早已绝迹的灵草珍木。他们也遭遇了其他一些零散的探险者，有散修，也有个别天师府弟子，大多行色匆匆，彼此戒备，偶尔交换一些无关紧要的情报。

时间在探索、警惕、偶尔的小规模冲突或互助中流逝。

第三日，他们穿过一片弥漫着淡金色雾气的奇异林地时，终于遭遇了进入秘境后第一次有预谋的、来自守旧联盟修士的针对性袭击。

袭击发生在林间一处狭窄的隘口。当他们半数人通过时，两侧看似普通的、爬满藤蔓的岩壁突然爆开，预先埋设好的爆裂符箓和禁锢阵法瞬间启动！

与此同时，上方树冠之中，隐匿许久的五名守旧联盟弟子骤然发难，炽热的火球、锋锐的金芒劈头盖脸砸下，目标明确——直指被暂时分隔开的队伍中段，尤其是陈千羽和另外两名修为稍弱的弟子！

“敌袭！结阵！”白言冰厉喝，身处队尾的他反应最快，身形一晃已拦在攻击最密集处前方。

长剑出鞘，划出一道浑圆的剑幕，如中流砥柱，将大半攻击挡下。

领域圆融稳固，虽被轰击得波纹荡漾，却坚不可摧。

白言冰自信，除非周文彬亲自出手，否则剩下的这些弟子无非土鸡瓦狗，连他的防御都破不开来！

前方的沈瑶和另一名弟子亦瞬间反击，刀光剑影迎向从岩壁后冲出的两名敌人。

但对方算计精准，仍有两道刁钻的攻击绕过白言冰的防御，袭向陈千羽。陈千羽临危不乱，素手轻扬，一面由无数细密银针虚影构成的元气护盾展开，同时袖中飞出数点寒星，直射施法者的要害，攻守兼备，逼得对方不得不回防。

战斗在刹那间爆发，狭窄的隘口内灵力激荡，光芒乱闪。守旧联盟弟子人数虽稍少，但占了先手埋伏的优势，且配合默契，显然是早就盯上了他们，特意在此设伏。

“周师兄有令，重点关照云隐别院和那姓白的！尤其是那个医修，先废了她！”一名紫阳宗弟子狞笑着，双手连拍，一条烈焰凝成的巨蟒昂首扑向白言冰，试图将他缠住。

白言冰眼神冰冷：别人怎么动他倒先不论，但陈千羽可是他的逆鳞。

随即，手中长剑似乎狂吼起来，原本圆融守御的剑势陡然一变。

弥焰斩！

雷火剑气大盛，剑势如焰天火雨，层层叠叠，与那火蟒硬拼。

黑紫色的剑气与火蟒甫一接触，后者就被这至刚至强的雷火之气拦腰纵向切成了两半，半只火蟒反而扫向了旁边一名试图偷袭沈瑶的家伙。那修士猝不及防，慌忙闪避，阵型顿时出现一丝紊乱。

“就是现在！”白言冰低喝。

始终游走策应的沈瑶如同鬼魅般从那名慌乱弟子的影子里钻出，短刃带着一抹凄艳的弧光，直取其肋下。

刀光深入血肉！

随即，沈瑶本能地扭转了一下自己的手腕，短刀自然在家伙的胸膛中豁开了一个更大的口子，肌肉绞碎，鲜血直流。

这或许是往年游侠时光遗留的习惯与本能：要刺杀，就一击致命。

拔刀，刀柄带着元气猛击那人天冲。

那重伤弟子口鼻出血，倒地不起，眼见难活。

遭遇战瞬间变成了混战，但白言冰这边人数占优，且经过几日磨合，配合渐佳，很快稳住阵脚，并开始反压。

指挥埋伏的那名紫阳宗弟子见事不可为，己方已有一人重伤，再缠斗下去恐有折损，当即发出尖啸：“撤！”

其余守旧联盟弟子毫不恋战，纷纷掷出烟雾符箓，借助地形和林木掩护，拖着七窍流血的同门迅速遁走。

临走前那阴狠的眼神，显然不会善罢甘休。

白言冰并未深追，秘境之中，贸然追击极易落入新的陷阱。他迅速查看己方情况，除了两人受了些轻伤，并无大碍，陈千羽已及时为他们处理。

“他们开始有组织地针对我们了。”沈瑶擦拭着短刃上的鲜血，冷声道。

“意料之中。”白言冰收剑入鞘，看着敌人消失的方向，“那所谓‘周师兄’既已视我等为心腹之患，自然不会让我们顺利探索。”

\vspace{2em}

又过了两日，他们在一处位于山涧旁的灵泉边休整时，遇到了吴刚带领的、已扩充到八人的万仙盟小队。

双方都有些风尘仆仆，但眼神精悍，显然收获与经历都不少。

吴刚主动走了过来，咧了咧嘴：“白兄弟，看来你们也没少被那帮孙子找麻烦。”

白言冰点头：“遭遇过一次埋伏。吴兄那边如何？”

“打了两场，吃了点小亏，也让他们见了点血。”吴刚眼中闪过一丝狠色，“散修别的不行，记仇和找机会下手最在行。他们现在也不敢小看我们了。”

他顿了顿，压低声音，“我们的人发现了一些痕迹，守旧联盟那帮人，似乎在有意无意地把一些难缠的妖兽或危险区域，往可能的路径上引，想借刀杀人。另外，对于那些资源，他们似乎只取最好的一部分，剩下的多半被他们毁掉了。”

白言冰心中一动：“大致方向？”

吴刚摇了摇头：“方向嘛，和我们感应到的那生机最浓的地方，大体一致。而且，根据我们零星得到的情报，其他几股咱们这边的人，还有那些中立的、只想捞好处的天师府弟子，大多也在朝那个方向靠拢。看来重头戏都在那儿。”

这个情报至关重要。双方简要交流了一下各自探知的部分地形和危险点，便再次合成一队。白言冰率领队伍，朝着那冥冥中感应到的、生机与道韵汇聚的核心方向，继续前进。

沿途，他们看到了更多战斗痕迹，有修士与妖兽的，也有修士之间的。秘境资源的争夺，已经逐渐白热化。

他们甚至在一处山谷中，发现了三具尸体，皆是散修打扮，财物被搜刮一空，致命伤显然并非妖兽所致。

秘境的柔和青光，似乎也掩不住日渐浓厚的血腥与肃杀。

\vspace{2em}

第七日，他们翻越一片矮小的山岭后，眼前景象骤然一变。

不再是无穷无尽的乔木树林，而是一片开阔的、望不到边际的奇异原野。

极目远眺，原野的尽头，隐约有连绵的虚影，似殿宇飞檐，又似山峦轮廓，被浓郁的青色霞光笼罩，看不真切。

但那股浩瀚、古老、充满生机的源头波动，正清晰无比地从那个方向传来，如同心脏的搏动，吸引着所有人的灵魂。

同时，他们也看到了其他身影。在原野的不同方位，影影绰绰，至少有数十人也在朝着相同方向前进。

其中一批人数量较多，服饰相对统一，气息灼热或锋锐，正是以紫阳宗圣子周文彬为首的守旧联盟队伍，他们似乎刚刚汇合不久，声势不小。

另一边，则是零散分布的小队或独行者，有万仙盟的人，有天师府弟子，也有其他幸存修士。

双方隔着很远的距离，彼此都能看见，气氛瞬间变得微妙而紧张起来，但谁也没有轻易挑起冲突。所有人的目标，显然都是原野尽头，那片青色霞光笼罩之地。

白言冰知道，最初的探索和淘汰阶段，已经接近尾声。真正的核心区域，就在前方。而最终的对峙与争夺，即将在那片青霞之下展开。

“收敛气息，保持距离，我们走。”

他低声下令，带领队伍，向着秘境最终的核心悄然进发。

\vspace{2em}

秘境核心，一片无比广阔、被永恒柔和青光笼罩的古老广场，豁然呈现于天地之间。

那青光仿佛源自广场本身与周遭的空间，均匀、明亮却不刺眼，带着温润的滋养之意，照在身上，连多日奔波的疲惫都似乎被抚平了些许。

而广场的中央，那座巍峨矗立的殿宇，攫取了所有人的目光与呼吸。

它不知以何种奇异的美玉筑成，通体呈现出一种深邃而温润的青色，仿佛将一片无云的晴空或是宁静的深海凝结成了实体。

岁月是无情的雕刻师，在它身上留下了不可磨灭的斑驳痕迹：墙壁有多处塌陷，断裂的梁柱如同巨兽的骨骼般支棱出来，精美的飞檐斗拱残破不全；无数不知名的藤蔓，叶片如翡翠，缠绕着那些残垣断壁；各色奇花异草，散发着静谧的微光，从裂缝中顽强探出，恣意绽放。

破败与生机，沧桑与神圣，在此刻达成了惊人的和谐与统一，共同诉说着一段淹没在时光长河中的辉煌过往。

殿门紧闭，高达十丈的青铜门扉上覆盖着厚厚的铜绿，但门楣上方，那古老鸾鸟翱翔的浮雕依然清晰可辨。鸾鸟的线条流畅而充满力量，羽翼舒展，似乎下一刻就要破石而出，直上青冥。

即便只是浮雕，也散发出一股令人心悸的、磅礴如海又精纯如露的生机之气，更蕴含着一种悠远深邃、难以言喻的大道韵律，让注视者不由自主地心生敬畏与渴望。

然而，一层肉眼可见的、半透明的光罩，如同一个倒扣的青金色巨碗，将整座殿宇连同其前方的大片广场牢牢笼罩。

光罩之上，无数复杂玄奥的青金色符文如活物般流转不息，时而汇聚成鸾鸟虚影，时而散作漫天光雨。

最令人心惊的是，光罩之上生机与毁灭两种截然相反的力量奇异地交织、平衡在一起，仿佛在平静的海面下隐藏着足以吞噬一切的暗流，警告着任何妄图强行闯入的不速之客。

\vspace{2em}

几位先到一步的天师府弟子，模样有些灰头土脸，正站在光罩外不远处，脸上带着无奈与不甘。他们仅仅是早到了半刻，尝试了各种手段试图直接破开护罩，甚至联手轰击过，但那光罩纹丝不动，反震之力还差点让他们吃了亏，此刻只能望洋兴叹。

所有人的视线，最终都聚焦在殿门正前方，光罩之内、广场之上，那处唯一异样的区域——一个直径约三丈的完美圆形。圆形区域内的地面，镌刻的阵纹比之外围光罩上的还要繁复精密百倍，线条交错如星空轨迹，节点闪烁如晨星，氤氲着浓郁的、几乎要液化的青色光晕。

就在这时，广场两侧边缘，几乎同时涌入了两股泾渭分明的人流。

左侧，以紫阳宗圣子周文彬为首，守旧联盟的弟子们鱼贯而出。

他们一共十几个人，虽然衣衫大多沾有尘土、灵草汁液甚至干涸的血迹，不少人的衣袍还有法术或利爪造成的破损，但整体气势依旧强盛。紫阳宗弟子周身隐隐有灼热气息升腾，无极剑派弟子则锋芒大盛，丹鼎阁弟子药香萦绕。

他们的眼神，在疲惫之下，是毫不掩饰的灼热与占有欲，如同盯上猎物的群狼，死死锁定了青光中的古老殿宇。

自进入秘境以来，他们秉持着大派作风，遇见资源与宝物，只取最精华的一部分；路线上，选择绕路避战，虽经历风险，但整体保存了更多实力。

右侧，白言冰、陈千羽、沈瑶带领的云隐别院及新生势力联盟的修士们也相继现身。他们的人数稍多一些，但状态看上去更为复杂。不少人身上带着明显的战斗痕迹，包扎的布条下隐隐渗出血色，气息也更为起伏不定，显然经历了更多直面妖兽、破解禁制的硬仗。

他们的眼神同样锐利，却多了几分历经磨砺的沉凝与警惕，如同在荒野中披荆斩棘终于找到水源的旅人，既有渴望，也深知周遭可能潜伏的危险。

\vspace{2em}

双方人马在广场边缘停下，隔着近百丈的距离遥遥对峙。空气中原本充盈的平和生机灵气，仿佛瞬间被无形的张力所凝固，充满了山雨欲来的压抑。

没有人说话，只有轻微的喘息声、衣物摩擦声，以及那仿佛响在每个人心头的、来自殿宇的磅礴道韵波动。青鸾传承的诱惑，近在咫尺，唾手可得，却又被那层薄薄的光罩无情阻隔。

“看来，这便是最后的考验了。”一个带着傲然与灼热气息的声音打破了沉寂。

紫阳宗圣子周文彬排众而出，他向前踱了几步，暗红色的锦袍下摆虽有一道被划开的口子，却无损其逼人的气势。他面容俊朗，此刻却因兴奋和野心而显得有些凌厉；目光如炬，扫过那圆形分阵和其后巍峨的殿宇，嘴角勾起一抹志在必得的弧度。

“此阵玄奥，非蛮力可破。谁能率先破解，便占得先机！”他声调抬高，似在宣告，又似在挑衅，说话间，却不着痕迹地向身后的同门递了一个隐晦的眼色。

看起来他想率先试图破阵。如果别人胆敢抢先，就会被身后众人群起而攻之。

“周圣子所言极是，那便请吧。”

白言冰的声音平和响起。

周文彬瞥了那带头剑修一眼，寻思这就是白言冰了，于是冷哼一声，不再多言。

他深吸一口气，胸膛微微起伏，周身瞬间腾起赤红色的精纯灵力，如同包裹在一层燃烧的琉璃焰火之中，将其衬得宛如火神临世。

在双方数十道目光的聚焦下，他迈开步伐，沉稳而坚定地踏入了那直径三丈的圆形分阵之中。

就在他双足完全落定在繁复阵纹上的刹那——

“嗡！”

阵中沉寂的青色光晕骤然爆发，化作亿万颗细密柔和的光点，如同春日清晨最晶莹的露珠，又似夏夜飞舞的流萤，温柔却无可抗拒地将周文彬的身形彻底淹没。外界看去，只见一团浓郁的人形青光，其内人影猛地一僵。

下一刻，周文彬那张原本写满傲然与志在必得的脸，在青光掩映下骤然扭曲！

极度惊愕、难以言喻的恐惧、以及一丝被人窥破最深处秘密的慌乱，如同打翻的颜料盘，瞬间取代了所有冷静与算计。他的身体不受控制地微微颤抖起来，额角、鼻翼以肉眼可见的速度渗出细密的冷汗，在青光照耀下反射着冰冷的光泽。

他仿佛看到了什么景象，瞳孔剧烈收缩。

“不完全是锁钥，竟然上面还套了层幻阵！”饶是最精通阵法的沈瑶也不禁惊叹这外挂幻阵的精妙。她甚至一开始并没看出来还有第二重阵法！

只见阵中的周文彬开始有了动作，时而面露狰狞凶狠之色，抬手向前虚抓虚击，仿佛在与看不见的敌人搏杀；时而又仓惶失措地连连倒退，脚步虚浮，口中无意识地喃喃自语，声音断续却清晰地传了出来：

“不……不是我……那是他……是他自己技不如人！对，是他技不如人！”

他像是在向某个无形的审判者激烈地争辩，又像是在拼命说服自己，额上青筋都已暴起。

这般挣扎持续了不过十数息，对周文彬而言却漫长得如同百年。终于，他猛地闷哼一声，像是心口被重重一击，脸色“唰”地变得惨白如纸，再无半分血色。

他竟主动踉跄着向阵外倒退，脚步虚浮，差点跌倒，靠着身后同门及时上前搀扶才稳住身形。他气息紊乱不堪，胸口剧烈起伏，再抬头望向那圆形分阵时，眼神里已充满了后怕和不敢相信。

守旧联盟众人围拢上来，脸色都变得异常难看，几个原本摩拳擦掌、准备紧随圣子尝试的核心弟子，此刻也面面相觑，眼神闪烁，默默地向后退了半步。

谁心里没点不愿示人、甚至自己都不愿深想的隐秘？这幻阵竟能直指本心，挖掘最深处的愧悔与阴暗，以此来影响破阵者，属实精妙。

“看来此阵，专照人心鬼蜮。不过周圣子良心未泯，可喜可贺。”一个清冷中带着淡淡讽刺的女声响起。

沈瑶越众而出，她黑衣如墨，衬得肌肤欺霜赛雪，原本清丽却稍显冷硬的脸庞上，此刻是一片近乎淡漠的平静。

“师妹小心。”陈千羽看着她，语气温和却郑重地叮嘱。

沈瑶对她微微点了点头，眼神交汇间自有默契。她没有丝毫犹豫，步履平稳，如同走向一处寻常回廊，径直走入了那令周文彬狼狈不堪的青色光点之中。

幻境瞬间展开，毫不留情。

熟悉的、几乎烙印在灵魂深处的灼热痛楚排山倒海般袭来，仿佛时光倒流，将她猛地拉回十年前那个绝望的午后——亡命徒狰狞的脸，激射而来的炽热火符，瞬间包裹头脸的恐怖高温，皮肤焦灼产生的刺鼻气味，还有随之而来的、无边无际的剧痛与黑暗。

火焰在眼前幻化跳动，热浪灼烧着感知，那曾日夜折磨她的阴影如同冰冷的潮水，试图再次将她淹没、吞噬。

若是从前那个未曾真正放下的沈瑶，此刻或许会被激发出更强烈的怨恨与斗狠之心，与幻境殊死搏斗。

然而，此刻阵中的沈瑶，只是在那灼热火光扑面而来时，轻轻地、极其自然地闭上了眼睛。长而微翘的睫毛在青光中投下淡淡的阴影，复又睁开。

那双曾盈满痛苦与戒备的眼眸里，此刻映照着幻象的火焰，却清澈得如同雨后的寒潭。

没有恐惧，没有愤怒，甚至连一丝波澜都难以寻觅。

她缓缓伸出手，素白的手指探向那幻境中焚毁她容颜的火焰，在即将触及那虚幻炙热的刹那，停了下来。她就那样静静地看着，看着幻象中那个痛苦蜷缩、面目狰狞、全力嘶吼的少女，眼神里掠过一丝极淡的、近乎悲悯的情绪，最终化为一片历经千帆后的平静与释然。

她红唇轻启，声音不高，却奇异地穿透了阵法的青光隔绝，清晰地传入了广场上每一个人的耳中，带着一种斩断枷锁的轻松与笃定：“那只是……过去的我罢了。”

随即，沈瑶坦然坐下，将那幻阵特意演绎出的残酷画面搁置一旁，开始专心破阵。

作为锁钥的阵法倒是不难，以沈瑶的阵法造诣，很快就轻松破解开来。

“嗡——！”

一声清越悠扬、仿佛穿越了万古时空的鸾鸟鸣响，蓦然在每个人心底响起！

笼罩巍峨殿宇的庞大青金色光罩应声发生了剧烈变化，符文疯狂流转，光华明灭不定，随即，在紧闭的青铜殿门正前方，那光滑坚固的光罩如同被无形之手从中间分开的瀑布水幕，流畅而无声地向两侧滑退、消散，露出了其后古朴厚重、布满铜绿的真实殿门。

防御阵，开了！

“进去了！”不知是谁按捺不住激动，嘶哑地喊了一声。

这一声如同投入滚油的火星，瞬间点燃了所有积累的贪婪、焦急与竞争之心。

刹那间，对峙的平静被彻底撕碎，众多修士再也顾不得阵型与矜持，一个个眼泛红光，将身法催动到极致，化作数十道颜色各异的流光，争先恐后、你推我挤地鱼贯而入，疯狂冲向那敞开的、仿佛通往无上机缘的殿门。

似乎晚上一步，那传说中的青鸾传承就会被他人捷足先登。

作为破阵的最大功臣，沈瑶反而被这汹涌的、失去理智的人流稍稍挤到了后面。白言冰与陈千羽一左一右护在她身旁，三人随着澎湃的人潮，踏上了殿门前的台阶。

就在陈千羽即将跨过那道古朴青铜门槛的刹那，一声哀鸣，毫无预兆地在她心海最深处响起。

那声音极其微弱，细若游丝，仿佛下一刻就要断绝，却蕴含着如此浓郁的痛苦、无助，以及一种对生命、对这片天地深深的眷恋与不舍。它似乎是直接响彻在灵魂深处，像一根冰冷的针，直接刺中了陈千羽心中最柔软的部分。

她脚步不由得一顿，下意识地转过头，目光越过躁动拥挤的人腿，望向殿门旁侧，一处被高大石雕残骸阴影覆盖的角落。

只见那里，靠近冰冷石基的地面上，躺着一只巴掌大小的鸟儿。

它羽毛本应是华丽的青色，此刻却沾满尘土，凌乱暗淡，失去了所有光泽。

一只翅膀以极不自然的角度弯折着，软软地耷拉在地上，隐约可见红肿。它气息奄奄，小小的胸膛微弱起伏，那双本该灵动剔透的琉璃色眼眸，此刻蒙着一层灰败的阴翳，正哀伤地、近乎绝望地望着眼前蜂拥闯入殿内、对它视而不见的人群。

守旧联盟的弟子们冲在最前，眼里只有殿内的黑暗与可能的珍宝，对脚边这微不足道的小生命连瞥一眼都嫌浪费时间，甚至有人嫌它略略碍路，随意一脚踢开。

青色小鸟孱弱的身躯翻滚了几下，撞在冰冷的石头上，发出的哀鸣更弱了几分，几不可闻。

许多云隐别院的弟子、万仙盟的修士、以及中立的天师府弟子，大多也只是匆匆一瞥，便随着人流涌入殿内。

机缘当前，一只将死的、寻常可见的鸟儿，又算得了什么？

然而，陈千羽的心，却像是被那只小鸟眼中浓得化不开的哀伤，以及那声直达灵魂的哀鸣，给紧紧攥住了，闷闷地发疼。

她是医者，师承云隐别院，见过更惨烈可怖的伤痛，也对生命本身怀有一种根植于道的、本能的悲悯与尊重。这青鸾秘境处处生机盎然，道韵沛然，一只如此模样的小鸟，偏偏出现在这核心殿宇的门口，绝不可能只是巧合。

更重要的是，那哀鸣中传递出的情绪，纯净而强烈，让她无法像其他人那样硬起心肠，转身离开。

“千羽？”白言冰发现她停下，顺着她的目光也看到了阴影中的小鸟，眉宇微微一蹙，显然也察觉到了异常，但并未出声催促，只是安静地等待她的决定。

陈千羽咬了咬下唇，贝齿在唇上留下浅浅的印子。

内心天人交战仅一瞬，那源于医者本心的悲悯终究占据了上风。

她不再犹豫，快步逆着稀疏下来的人流走过去，在少数几个尚未进殿、投来诧异不解目光的弟子注视下，小心翼翼地蹲下身，伸出双手，用最轻柔的力道，如同捧起一滴晨露般，将那只受伤的青色小鸟捧在了温热的掌心。

小鸟在她掌心微微颤抖着，羽毛拂过皮肤带来细微的痒意。它似乎连挣扎的力气都没有了，只是费力地抬起小小的头颅，那双琉璃般的眼眸望向陈千羽。

奇异地，那其中盈满的哀伤，竟在这一望之下稍微减退了些许，取而代之的是一丝微弱却清晰的依赖与祈求，它甚至试图将受伤的脑袋靠近她温热的指腹。

“你们先进去，我马上来。”陈千羽抬头，对白言冰和沈瑶说道，声音轻柔却坚定。

同时，她已运转起体内温润平和的元气，那气息带着她特有的、促进愈合的灵性。

一层淡淡的暖光，轻轻包裹住小鸟折翅的伤处，先稳住它不断流逝的生机与加剧的痛苦。

白言冰与沈瑶对视一眼，彼此眼中都掠过一丝了然。他们深知陈千羽外柔内刚的性情与那份对生命的执着，也相信她的直觉与判断。

“小心。”两人齐声道，不再多言，转身踏入了那深邃的殿门之中，身影迅速被黑暗吞没。

陈千羽怀抱这只奇异的青鸟，成为最后一位踏入大殿的修士。

当她迈过那道高高的青铜门槛时，身后沉重的殿门，仿佛有所感应般，发出一声低沉悠长的叹息，无声无息地缓缓合拢，将外界那片柔和的青光与广阔的广场彻底隔绝在外，只剩下门缝消失前最后一缕微光，映亮她沉静却坚定的侧脸。

\vspace{2em}

殿门之内是一个无比广阔、远超外部所见规模的空间：无数条通道、回廊、拱门纵横交错，向四面八方延伸，构成了一座极其复杂的巨大迷宫。

墙壁的材质不知是何物，触手温润，通体散发着均匀柔和的青光照亮前路，却又不显得刺眼。墙壁之上，镌刻满了巨大的、连贯的古老壁画。壁画的内容恢弘壮丽、栩栩如生，充满了灵动盎然的大道韵律，讲述着青鸾司掌生机、平衡万物的古老神职。

空气中弥漫着的，是浓郁到几乎化为实质的生机灵气，比之外界广场又不知强盛了多少倍。每一次呼吸，都仿佛饮下琼浆玉液，四肢百骸说不出的舒坦。

然而，与此同时，神识在此地如同陷入浓稠的泥沼，勉强能离体数丈便再难延伸；而那错综复杂的通道和看似有规律的壁画，更是在不断混淆、误导着人的方向感，使人难辨东西。

先进来的修士们早已被这巨大的迷宫和躁动的贪欲所分散，如同水滴汇入沙地，消失在各条通道深处。远处，不同方向不时传来惊喜的呼喊声、激烈的打斗声、以及灵力爆鸣或古老禁制被触动后发出的沉闷爆炸声，在这空旷的迷宫中回荡，更添几分混乱与危险的气息。

白言冰、沈瑶很快与吴刚等万仙盟的人，以及云隐别院其余人等在一处宽敞的十字回廊汇合。

众人苦迷路久矣，于是经过商议，选定一个方向，采用最稳妥的方式：每走过一段回廊，遇到门户便开门查探；每探索完一个房间，便在门上留下一个简明的记号，然后便转向下一个目标。

房间内的情形各不相同，有的空空如也，有的暗藏机括陷阱，猝然发动，需得小心应对；但也有的房间，里面存放着一些蒙尘的财富——或是堆积在墙角、灵气盎然的各色灵石；或是玉瓶中药力未曾完全散失的旧日丹药；或是样式古朴、灵光暗淡却依然可用的法器。

不久，周文彬带领的守旧联盟队伍也发现了这种方法的有效性，立刻有样学样，开始分散成数支小队，在迷宫中进行地毯式的搜索，同样在探索过的门户上留下标记。

然而，随着时间推移，无论是新生势力还是守旧联盟的修士，都陆续开始感觉到不对劲。

许多通道看似笔直通向更核心的幽深之处，满怀希望地走上一段后，却往往在拐过一个弯或穿过一道拱门后，愕然发现自己又回到了之前经过的、有着熟悉记号的岔路口，甚至干脆是冰冷的、刻画着壁画的死胡同。

墙上的壁画内容，初看时宏伟壮丽，看久了却觉得它们似乎在缓慢变化，又仿佛亘古如此，总在不经意间将人的思绪引向某种无始无终的循环。

不少弟子开始心浮气躁，如同被困在琉璃罩中的无头苍蝇。

明明宝藏可能就在附近，却不得其门而入，只能徒劳地在相似的廊道间打转，消耗着体力和耐心。

\vspace{2em}

唯独陈千羽，怀抱那只体温正在缓慢回升的青色小鸟，独自走在光影流转的迷宫之中，所经历的感受却与所有人都截然不同。

那干扰神识、混淆方向的无形力量，对她似乎影响甚微。

她心中有一种非常微弱、却真实不虚的模糊意识——似乎是一种若有若无的共鸣与牵引。

似乎……掌心那蜷缩的、依赖着她元气存活的小小生命，正用它微弱的气息，与她脚下这座古老殿宇的某处核心，产生着难以言喻的联系，并在冥冥中为她指明前路。

陈千羽没有抗拒这份感觉。

她自然而然地迈开脚步，不再刻意去记忆复杂的路径或辨识壁画的不同。

她走过刻画青鸾哺育幼雏的回廊，壁画上的神鸟眼眸似乎温和地望向她，轻轻颔首；她穿过一道拱门，门后本应是数条岔路，此刻却唯有其中一条，青光明亮而柔和，道路笔直地向前延伸，仿佛专为她而呈现。

怀中的青鸟，身体不再冰冷颤抖，折翅处在她持续不断输出的温润元气滋养下，那淤肿竟在缓缓消退，破损的皮肉下有微弱却顽强的生机在萌动、凝聚，试图接续断骨。

不知不觉间，她与白言冰、沈瑶等人走散了，也仿佛超脱了这座困扰所有修士的迷宫的固有规则。

她不再属于那些焦躁寻找、重复打转的人群。

\vspace{2em}

周围的景象开始发生微妙而持续的变化。

回廊愈发显得古老、静谧，墙壁散发的青光愈加纯粹，壁画上的青鸾形象愈发鲜活灵动，羽翼上的每一丝纹路都清晰可见，眼眸顾盼生辉，仿佛随时会发出一声清鸣，振翅从那墙壁之上飞腾而出，翱翔于宫殿穹顶。

最终，在迷宫的至深之处，她踏过最后一级泛着微光的台阶，眼前已无岔路，只有一扇孤零零的门户。

这扇门与其他任何房间的石门、木门都截然不同，它通体由凝实的青色光华构成，光晕流转，门扉上隐约有更复杂古老的纹路明灭，静静地矗立在一面刻画着青鸾归巢、敛翼安息的巨大壁画之前。

似是感应到她的到来，尤其是她怀中那只青鸟微弱却清越的一声低鸣，青色光门微微荡漾起来。

随后，它无声地、温柔地向内洞开。

门后流泻出的，是一片温暖、浩瀚、充满无尽生命喜悦的青色光辉。

\vspace{2em}

门后，并非她或任何人想象中的、堆满古老典籍或闪耀宝物的藏宝室，而是一座极为空旷、穹顶高远的圆形殿宇。

殿内空无一物，唯有地面、墙壁、穹顶都流淌着那种温润的青色光辉，将每一寸空间都映照得澄澈而圣洁。这里寂静得能听到自己血液流动的声音，却又仿佛充满了无声的宏大乐章。

而在殿宇的最中央，悬浮着一团无法用言语形容的青色光源。

它是“生机”这个概念最纯粹、最本源的显化，浓缩凝聚到了近乎实质的境地。它如同一颗宏伟的青色心脏，缓缓地、有力地搏动着：每一次舒张，便有无形无质却沛然莫御的生机道韵如水波般荡漾开来，充盈整个殿宇；每一次收缩，又仿佛将万物的喜悦与活力吸纳归一。

站在这团光辉前，陈千羽感到自己浑身上下的每一寸都在欢欣雀跃，多年行医积累的疲惫、先前跋涉时的紧张，乃至灵魂深处细微的尘垢，都被这纯粹的光辉洗涤、抚平。这不是力量的压迫，而是生命本源对同源者的呼唤与抚慰。

就在这时，她掌心传来轻微的悸动。

那只一直安静蜷缩、依赖她元气存活的青色小鸟，忽然抬起了头。它那双原本蒙着阴翳的琉璃色眼眸，此刻变得清澈无比，倒映着中央那团青色光辉，充满了孺慕、欢欣，以及一种归家的释然。

它发出一声清越悦耳、宛如冰玉相击的鸣叫，声音带着穿透灵魂的纯净喜悦。

挣扎着，用那双尚未完全恢复的翅膀，在陈千羽掌心轻轻一蹬，身体歪歪斜斜地飞了起来。仿佛初学飞翔的雏鸟，却又带着义无反顾的决绝，径直投向了殿宇中央那团最为炽盛纯净的青光。

小鸟的身形在接触光团的刹那，便如同水滴融入大海，悄无声息地消失了。

紧接着，那团青色光辉骤然膨胀、盛放！

无穷无尽的光淹没了陈千羽的视野，将她温柔而彻底地包裹。

一个古老、沧桑、却依旧温和慈爱的意念，跨越了漫长到难以想象的岁月长河，如同初春暖阳下解冻的涓涓溪流，自然而然地、涓滴不漏地流入陈千羽敞开心扉的灵台识海。

青鸾残念的声音，充满了时光沉淀的智慧与淡淡的忧伤。

“后来的生灵啊……你终于来了。”那意念的开端，像一声悠长的叹息，“我是此地主人最后残留的一点真灵印记，依托这秘境核心的生机本源而存，等候了太久太久。”

陈千羽心神震撼，在意识中回应：“您是……青鸾前辈？”

“称谓并不重要。重要的是，你看到了它，并选择了看到了它。”

那意念中流露出赞许的情绪，“当贪婪的目光掠过它，投向那扇被认为藏有至宝的门户时，唯有你，为一只看似平凡的受伤小鸟驻足。这种源于本心的、对生命本身的悲悯与尊重，正是我所执掌之道的基石。”

陈千羽眼前仿佛浮现出殿门外那混乱的一幕：匆忙践踏的脚步，冷漠无视的目光，还有那被随意踢开、瑟瑟发抖的渺小生命。她心中微微一疼。

“青鸾主生机，亦掌平衡。”那意念继续流淌，向她展开一幅浩瀚的图景，“生机并非一味滋生繁荣，平衡亦非僵死不变。春日勃发，秋日凋零，皆是循环；滋养万物，亦需抑制过盛，方为长久。”

无数玄奥的意象随之涌入陈千羽的脑海：种子破土时蕴含的爆破之力与柔韧之性，参天古木将阳光雨露转化为滋养的精密过程，伤病躯体中破损生机与修复力量的微妙拉锯，乃至一方天地间灵气的盈亏循环……这些不再是抽象的概念，而是一个个鲜活具体的道之显化。

在此道理之中，自然包含了无数早已失传于世的丹方配伍、药性医理、精粹炼药手法。

陈千羽感到自己熟知的那些药理知识，如同散落的珍珠被一根青色的丝线串联起来，焕发出全新的、更深层次的光芒。

那意念的声音渐渐变得空远，仿佛力量正在消散：“传承与你，并非要你成为另一个青鸾。你的道途与他们乃至我主都不同，没有那种强夺的本性。我只希望你能持此道，行你之路，在你所处的时代，护持那一点对生命的敬畏与平衡之心。”

这最后的印记即将消散。

“我散去后，这秘境也不能再维持多久了，希望你与你的朋友们能尽快离开这里。年轻的医者，善用这份馈赠……”

最后的意念如同温暖的潮水般退去，留下的是一片浩瀚无垠的理解。

不知过了多久，或许是一瞬，又或许是百年。

笼罩殿宇的磅礴青光渐渐收敛、黯淡，最终完全消失。殿宇中央空空如也，唯有空气里残留着令人心旷神怡的清新气息。

陈千羽静静站立在原地，双眸轻阖，长睫在白皙的脸颊上投下淡淡的阴影。

她周身的气息圆融流转，原本温润平和的元气质地未变，其中却多了一股浩瀚、灵动、仿佛能孕育万物的生机道韵，使她整个人看起来更加沉静深邃，宛如一株历经风雨却内含无尽生机的植株。

她眉心之间，一道极淡的、形似展翅青鸾的印记一闪而逝，随即隐没于光洁的皮肤之下。

在这段时间内，她自然从传承中找到了这只青鸾前辈的生前修为：圣魄境不朽期。

这个修为具体有多高，她对此没有什么概念。

不过这不重要。

她缓缓睁开双眼，眸底似有清澈的青光流转了一瞬，旋即恢复为原本的明净；但那明净之中，却多了几分洞彻生命规律的深邃。

传承已毕，几乎无需消化，已然成为她认知与道基的一部分。不过，她并没有试图当场突破境界。

眼下，离开这个迷宫，才是当务之急。

没有丝毫犹豫，心念一动，触动了怀中一枚简易的传讯符。

\vspace{2em}

一道清晰而简短的信息，传向迷宫之中与她气息相连的同伴：“传承已得，速退。”

迷宫深处，白言冰刚以精妙剑势引偏了一具古老铠甲傀儡的重击，沈瑶的刀光趁机切入其关节缝隙，吴刚则咆哮着用一柄刚捡来的重锤砸碎了它的核心。几人配合无间，正准备查看这房间内是否有遗漏之物，怀中传讯符同时传来了微热与波动。

信息入心。

“得了？”沈瑶挥刀震开傀儡碎片，清丽的脸上第一次露出近乎懵然的诧异，她看向白言冰，重复道，“千羽她……找到了核心？这么快？”

这效率远超所有人最乐观的预估。

白言冰的动作也是一顿，眼中瞬间爆发出惊喜的光彩，但那光彩立刻被更强大的决断力所取代。

他对陈千羽毫无保留的信任让他立刻做出反应：“信千羽！所有人，向入口方向撤离，沿途汇合，准备离开秘境！”他的声音沉稳有力，瞬间压过了刚刚战斗的余响。

没有丝毫留恋，他们甚至不去看那傀儡残骸中是否还有灵材，立刻转身，按照记忆中留下的标记和大致方向，朝着青鸾殿大门处疾退。

途中遇到其他仍在迷茫探索或与零星禁制纠缠的新生势力弟子，无论是云隐别院还是万仙盟的，他们立刻言简意赅地告知：“陈师姐已得传承，速退！”

惊愕浮现在每一张脸上，但无论是出于对白言冰权威的信服，还是对陈千羽机缘的震惊与信任，没有人质疑。求生的本能和对大局的判断让这些弟子迅速压下贪念，加入撤退的洪流。

队伍像滚雪球般壮大，行动却愈发迅捷有序。

另一边，守旧联盟的弟子很快察觉到了异常。

原本如同无头苍蝇般在迷宫中偶尔碰面、彼此戒备的云隐别院和万仙盟的人，忽然都朝着同一个方向快速移动，行为坚决，毫不拖泥带水。

几个小队将相同的消息迅速汇总到周文彬处。

周文彬正为迷宫所困而烦躁，闻讯先是一愣，随即脑中闪电般划过进入殿门时那最后一幕——那个云隐别院的女医修，逆着人流，弯腰捧起了一只该死的、碍事的伤鸟！

他的脸色“唰”一下变得铁青，额角青筋跳动，一股混合着极度嫉妒、被戏弄的羞辱以及计划落空的暴怒直冲顶门，烧得他双眼泛红。

他几乎是从紧咬的牙关里，一字一字地挤出充斥着毒恨的话语：“定是那陈千羽！最后抱着鸟进去的那个女人！她得了传承！”

他猛地转向周围同门，声音因激动而尖利，“不能让他们带走！拦住他们！不惜一切代价！”

然而，迷宫的复杂在此刻成了新生势力最好的掩护。众人撤退路线明确，彼此间依靠简单的信号和“三三制”的小组配合互相掩护、交替断后，行动高效得像演练过无数遍。

守旧联盟队伍本就分散，而在他们自己都不熟悉的迷宫之中捉人更是困难。

仓促间的拦截不是扑空，就是被对方默契的配合轻易化解或摆脱。他们只能眼睁睁看着白言冰、沈瑶等人与从核心区域方向匆匆赶来的陈千羽顺利汇合，然后把她簇拥在中心，汇聚成一股洪流，头也不回地冲向迷宫出口。

“废物！追！出去再说！立刻联系长老，告知情况！”周文彬气得浑身发抖，一脚踹在旁边的发光墙壁上，怒吼着带领手下，也疯狂地追了上去。

那扇厚重的青铜殿门，不知何时已无声洞开，仿佛专为离去者让路。

两派人马前一后，如同两道逆向的激流，冲出青鸾殿，掠过那片依旧笼罩在柔和青光下却已感觉不同的广场，向着秘境入口的方向全力疾驰。

沿途，秘境中那些原本稳定的禁制和生机脉络开始出现紊乱，时而凭空爆发出一股混乱的、充满撕裂感的生机冲击，或是地面骤然长出尖锐的荆棘，给这场生死时速般的逃亡与追击增添了无数变数与凶险。

当那熟悉的七彩漩涡入口终于在望，并且开始出现不稳定的、时缩时胀的波动时，双方人马皆是狼狈不堪，不少人身上添了新伤，但速度却丝毫不减。

新生势力众人护着陈千羽，率先一头扎进了漩涡；守旧联盟弟子红着眼，紧随其后冲入。

\vspace{2em}

身体穿过空间通道的短暂晕眩过后，脚下传来了内域熟悉的、带着碎石的土地触感。

然而，等待新生势力众人的，远远不是约定的“界外不得干涉”的平静，更不是师长接应的欣慰，而是一场早有预谋、卑鄙到极致的血腥伏击！

守旧联盟的弟子在冲出秘境七彩漩涡的第一时间，就毫不犹豫地捏碎了早已扣在掌心的紧急传讯符箓。

忽而，峡谷两侧原本看似平静的山岩、树林后方，骤然爆发出数十道强烈的灵力波动！

“陈千羽！交出青鸾传承，饶你同道不死！”

陈炎那熟悉的、充满灼热霸道的怒吼声如同惊雷炸响，伴随着其声浪，漫天炽热狂暴的火球、烈焰长枪如同陨石雨般轰然砸向刚刚立足未稳的新生势力人群！

几乎同时，刘长靖冰冷凌厉的剑气化为一张巨大的光网，封锁了他们侧翼与后方的退路；丹鼎阁门下高手释放出的墨绿色毒瘴如同活物般从地面蔓延开来，更有隐藏的困阵光华亮起；更有四五名同样散发着真符境强大气息、服饰各异的修士，显然是守旧联盟邀约或麾下的客卿、附属势力长老，带着大量眼神凶狠的低阶修士，如同捕猎的群狼，从峡谷两侧嘶吼着冲杀下来！

瞬息之间，新生势力众人便陷入了十面埋伏的绝境。

“卑鄙！你们违背约定！”峡谷上方，传来了吴东浩惊怒交加的吼声。他与叶牧谦本是按约定前来接应，万万没想到对方竟敢如此公然践踏规则。

两道强横的身影立刻从高处扑下，叶牧谦心念领域如一张无形巨网立即张开，直取最为狂暴的陈炎；吴东浩则啪地一声打开折扇，法术波动悍然拦向刘长靖的剑气网络，力图为下方陷入绝境的弟子们分担最大的压力。

随即，最强的陈炎和刘长靖被叶牧谦和吴东浩两人强行缠住，但更多的真符境摹形期、乃至玄脉境界巅峰的修士们开始向白言冰等人攻去！

“结阵！突围！”白言冰目眦欲裂，看着呼啸而来的致命攻击和密集的敌人，一股热血直冲头顶，狂吼出声。

他毫无保留，见独领域全力展开，那圆融澄明、映照己心的无形力场瞬间扩张，将他与身边众多同伴笼罩。他手中长剑发出清越长吟，剑光化作一道坚韧无匹的游龙，主动迎向那最炽热、最狂暴的几道火焰陨石，剑势浑圆连绵，竟能以一己之力卸开真符境长老的含怒一击。

沈瑶的反应同样快到极致，她清叱一声，周身气息陡然变得飘忽诡谲，人随刀走，化为一道捉摸不定的黑色流光游走在侧翼。刀光配合元气激发的幻术，干扰、误导、甚至将数道袭向修为较低同伴的攻击引偏。

吴刚、陈默等散修在最初的震惊后，爆发出亡命之徒般的狠劲与怒吼。他们自然不能像白言冰与沈瑶两人搞出什么精妙的幻术或领域，全靠血勇与平日里刀头舔血的默契。

散修们也三五成群，迎着扑上来的守旧联盟伏兵厮杀在一起，用身体和简陋的武器构筑起一道混乱却顽强的防线。

陈千羽站在众人保护的中心，脸色微微发白。

她刚刚接受完那浩瀚如海的传承，修为也未曾突破。但看着周围瞬间陷入血火的同伴，她毫不犹豫地抬起双手，十指如穿花蝴蝶般结出一个古朴玄奥的手印——这是传承中最基础、却也最核心的生机引导与护持法门。

春——和畅！

柔和的青色光晕以她为中心，如同水波般迅速扩散开来，笼罩住己方众人。

这光晕是她以自身精纯元气为引，强行模拟、构造出的一个临时性的生机领域，比真正的见独领域更耗费心神与灵力；它无法阻挡飞剑火矢，却蕴含着青鸾传承独有的、滋养万物的奇异生机。

光晕笼罩之下，受伤同伴伤口流血的速度肉眼可见地减缓，剧烈消耗带来的疲惫与灵台昏沉为之一清，仿佛干涸的土地得到了春雨的滋润，一股顽强不屈的生命力从心底被点燃。

然而，敌我力量悬殊得令人绝望！

伏击者不仅人数占优，更有多位真符境长老级人物坐镇指挥、出手。对于绝大多数修为仅在玄脉境或养气境的新生势力弟子而言，真符境修士随手一击的余波，都足以让他们骨断筋折，甚至当场殒命。

整个探索队伍中，唯有白言冰凭借见独境的修为能硬抗真符，沈瑶凭借幻术能勉强在真符境的攻击下周旋、抵挡片刻；其他人则如同狂风中的落叶。

顷刻之间，残酷的现实便撕开了希望的口子。

一名万仙盟的散修躲闪不及，被一道逸散的烈焰擦中，半个身子顿时焦黑，惨叫着倒地；一名云隐别院的弟子奋力格开当面之敌的长剑，却被侧面袭来的毒镖射中大腿，麻痹感瞬间蔓延，随即被乱刀砍中；更有甚者，直接被强横的灵力冲击震得口喷鲜血，倒地不起……

猩红的血液，迅速在峡谷灰黑色的地面上蔓延开来，刺鼻的血腥味混杂着焦糊味、毒瘴的甜腥气，构成一幅地狱般的图景。

“走！不要恋战！快走！”叶牧谦充满压抑怒火的声音，如同惊雷般穿透混乱嘈杂的战场，清晰传入每个苦苦支撑的弟子耳中。

他本人虽然见独圆满、修为高绝，但被陈炎与刘长靖两人死死缠住，根本无法分身救援下方。

“冲出去！”白言冰嘶吼道，他浑身浴血，青衣早已被染成深褐色，分不清是自己的还是敌人的。

他的领域在数名真符境修士的冲击下剧烈波动，却依旧坚韧地支撑着。

他一剑荡开面前之敌，剑势陡然变得惨烈决绝，化为一道锐不可当的锋芒，朝着包围圈相对薄弱的一处，悍然突进！

与其说是斩，不如说是刺！

带着全身重量、见独领域，以及不屈不挠的守护意志，全身撞击过去的刺！

“拦住他！”

两侧的两名守旧联盟修士——玄脉境融贯期——急忙一前一后，试图阻挡白言冰的突围。

没想到，这临时变出的一招，竟然蕴含着不可阻挡的威势。

站在前面的那位修士，白言冰的剑尚未及身，便被那无形的领域一撞！

踉跄着后退几步，几乎要撞到后面那人身上。

两人匆忙防御，却发现白言冰竟然不在他身前！

那他哪去了？

下一刻，那位修士忽然觉得后背被另一个人狠狠地撞了一下。

随即，便是剑尖刺穿胸口的感受。

什么！

白言冰自己见到剑上穿成一串的两个死不瞑目的守旧联盟修士，自己也有点愕然。

他刚才刺前人是伪，闪到后人身后背刺之是真。只是没想到，自己见独领域竟然也能用来撞人？

不过眼下并没有什么思考的余裕。

“包围圈已经撕破，快撤！”

沈瑶默契地紧随其后，刀光幻化成一片令人眼花缭乱的死亡光幕，掩护他的侧翼。

吴刚咆哮着，用以伤换伤的打法，将另外两名试图合围的敌人硬生生逼退。陈千羽咬牙维持着生机光晕，同时搀扶起身边腹部受创、血流不止的陈默。

一行人护着几名尚能行动的伤员，靠着“三三制”瞬间爆发的默契配合、陈千羽不断提供的生机支撑，残存的人马如同离弦之箭，从那缺口处冲出，朝着云霞山、云隐别院的方向亡命奔逃。

叶牧谦与吴东浩亦爆发出全力，心念领域中吴东浩的灵气锋刃纵横交错，且战且退，牢牢挡住身后追兵最凶猛的那些攻击，为弟子们的逃亡争取那宝贵的一线生机。

箭矢尖啸，法术轰鸣，剑气纵横，从身后不断袭来。不断有人中招倒地，发出闷哼或惨叫，又立刻被身旁的同伴拼死拽起，拖拽着继续前进。

鲜血，点点滴滴，连成触目惊心的红线，染红了这条用生命铺就的逃亡之路。

每一次回头，都可能看到熟悉的面孔永远倒下；每一次喘息，都带着铁锈般的血腥味。绝望与悲痛如同冰冷的潮水，一次次试图淹没他们，却又被更强烈的求生欲与愤怒点燃的斗志狠狠压回。

当云隐别院那熟悉的亭台轮廓、以及笼罩其上的淡淡防护阵法光晕终于映入眼帘时，出发时三十多名精锐，只剩下十一二人还能勉强站立，还有几人倒地不起、重伤昏迷。

每个人身上都带着或轻或重的伤，衣衫褴褛，血迹斑斑，气息萎靡到了极点，模样凄惨如同从地狱血池中爬出。

身后，喊杀声与追击的灵力波动在靠近别院范围后，终于渐渐减弱、停止。

守旧联盟的伏兵似乎终究对直接攻击云隐别院本体心存顾忌，依然未敢逾越最后的界线。

最后几步路，几乎是用尽最后一丝力气完成的踉跄。

\vspace{2em}

当终于冲入九宫迷踪的范围，确认那致命的攻击不再袭来时，紧绷到极致的神经骤然松弛。

扑通、扑通……幸存者们如同被抽掉了全身骨头，顿时瘫倒一片，剧烈的喘息声、压抑的痛哼声、劫后余生的哽咽声响成一片。

陈千羽自己也到了强弩之末，识海因过度强行模拟见独领域而阵阵抽痛，元气几乎枯竭。

但她只是擦了擦额角的冷汗与血污，便毫不迟疑地跪坐在重伤员之间，苍白的手指迅速而稳定地检查伤口，掏出随身携带的、所剩无几的药瓶与银针，开始紧急施救。

柔和微弱的青光再次从她指尖亮起，虽然远不及之前，却依旧带着生的希望。

白言冰以剑拄地，单膝跪倒，胸膛剧烈起伏，每一次呼吸都牵扯着身上的伤口，带来火辣辣的疼痛。

他脸上溅满血污，俊朗的面容此刻只剩下鏖战后的疲惫与狼藉，但那双向来沉静的眼眸里，劫后余生的庆幸只是一闪而逝，随即被一种深沉到极致、几乎要喷薄而出的熊熊怒焰所取代。

沈瑶靠坐在不远处一段残破的矮墙边，她左边的肩膀，一道深可见骨的剑伤皮肉翻卷，鲜血仍在汩汩外渗，染红了半边黑衣。她清丽的脸上没有太多表情，只是紧抿着失色的唇，额角渗出细密的冷汗。

她咬着牙，用还能活动的右手，颤抖而坚定地给自己洒上金疮药，然后用牙齿配合，撕下一条相对干净的衣襟，试图包扎。

叶牧谦和吴东浩两人的身影无声地落在庭院中央。这两人倒是没有受到看起来过重的伤，却更显凶险。

真符境只是个玻璃大炮。这个阶段的修士之间打斗，要么无伤，要么陨落。

叶牧谦的一袭青衫染满尘埃与暗红色的血点，头发有些散乱。他的面色冰冷如万载玄铁，没有丝毫表情，目光缓缓扫过庭院中横七竖八、伤痕累累的弟子们，扫过地上那来不及擦拭的斑斑血迹，扫过空气中弥漫的浓重血腥与悲愤气息。

他就这样静静地站着，许久，许久，一言未发。然而，任谁都能感受到，在那死寂般的沉默之下，是压抑到极致、即将冲破一切束缚的滔天怒意与冰冷风暴。

陈千羽用最快的速度，将几名最危重伤员的伤势暂时稳住，吊住了他们最后一口气。做完这一切，她已是摇摇欲坠。她强撑着站起身，对叶牧谦、白言冰等人微微颔首，没有多说什么，眼眸中交织着疲惫、悲恸，以及一丝刚刚领悟的、关于生命重量的明悟。

随即，她便转身，径直走向静室，关上了门。她需要时间，需要一个绝对安静的环境，去整理那浩瀚的传承，去冲击那早已松动的瓶颈，去获得足以应对接下来一切的力量。

沈瑶亦寻了一处静室闭关，她也觉得自己的领域快要成型了。

\vspace{2em}

庭院中的沉寂，持续了整整一日。

次日，当晨光再次洒满云隐别院，两扇静室的门，一前一后地被从内轻轻推开。

沈瑶先从她的静室踏出，身上伤口已经用绷带强行止血，但绷带上渗出、干结的血迹依然触目惊心。

她创新性地将自己的领域与幻术相结合，用元气刻画幻阵。从而，在节省布阵珍贵材料的同时，还能让幻阵随着她的移动而移动。

这幻术若与她的灵动步法相结合，对敌人而言将会是一场灭顶之灾。也正因此，沈瑶的脸上罕见地露出了一丝自得。

她的幻术领域，自名为“千幻流光”。

随后就是陈千羽缓步走了出来。

她整个人的气质，已然发生了天翻地覆的蜕变：之前的温婉宁静依然存在，然而在这底色之上，却多了一份源自生命本源的深邃与厚重，一种洞察生灭、触摸平衡的从容。

她自己也成功地将青鸾传承那浩瀚如星海的感悟，与自身坚定不移的求索根基、与那颗医者仁心相互印证、水乳交融；而她的生机领域，更是令人惊诧。

领域悄然展开的范围内，陈千羽能异常清晰地感知到万物内部生机的流转状态、盈虚变化，乃至那些细微的、可能导致崩溃的失衡之点。

此领域竟然包含两重变化：

其一，自然是先前在追逐战中强行模拟出来的“和畅”形态：陈千羽能以自身精纯元气与领悟的生机之道为引，强行调和、重新平衡那些破损的生机脉络，甚至只要对方尚存一息，她便能以损耗自身元气为代价，为其强行续命；

其二，便是尚未出现过的“萧瑟”形态：与“和畅”恰恰相反，这一形态之下，她能主动去破坏敌人体内的生机平衡，从而大大压制敌人的行动能力，甚至敌人死于这萧瑟秋风之中也不无可能。

这已超越了寻常医术的范畴，是在生死边缘行使恩赐与判决的惊世之能，是连丹鼎阁主王丹那种层次的人物都未必触及到的、近乎于“道”的层次，将陈千羽的医道，推上了一个全新的、令人仰望的峰巅。

白言冰倒是没什么好说的，他忙着帮叶牧谦、吴东浩整理秘境中其余所得，并安葬那些因各种原因把生命留在秘境中的弟子的骨灰。

只是把强行突围的那一招起了个名，叫“孤弓”。

\vspace{2em}

然而，境界的突破、领域的领悟、招式的变幻，并未给云隐别院的任何人带来丝毫喜悦，只有更沉默的压抑。

庭院中尚未散尽的药味与血腥味，同伴们苍白憔悴的脸庞和身上缠裹的绷带，都在无声地诉说着昨日的惨烈与卑鄙。

云隐别院方面一直秉持着克制与发展并重的原则，在诸多摩擦中甚至做出了不少让步，只为争取那来之不易的和平发展时间。

但是，守旧联盟一而再、再而三的倒行逆施，从秘境入口的暗算，到秘境中的偷袭，再到最后这公然违背约定、不惜动用长老级力量进行无耻伏击、造成惨重伤亡的行径，已经彻底、完全地践踏了所有的底线与规则。

昨日那场伏击的每一个细节，同伴倒下的每一个身影，鲜血染红大地的每一寸画面，都在幸存者眼中反复灼烧。

是可忍也，孰不可忍也？

那压抑到死寂的庭院之中，风暴，已然在平静的表象下酝酿到了极致。

沉默啊，沉默啊！

不在沉默中爆发，就在沉默中灭亡！

\vspace{2em}

次日，一纸堂堂正正、却字字如刀的檄文，猝然划破了内域维持许久的冷战僵局。

檄文由叶牧谦一人所撰，以云隐别院之名发出；以刀削斧凿的文字，以其理念为核心，向整个内域发出雷霆般的号召：所有认同“问道途”之念、苦于守旧联盟压迫与盘剥已久的修士与凡民，当共举义旗，涤荡寰宇，争一个公理，求一条活路！

檄文在内域各处被飞速地传抄，其文辞激越昂扬，理直气壮，更因蕴含着一股受尽屈辱后迸发出的决绝意志，而具有了灼人的力量。

其中两句，尤为耀眼夺目，如同烧红的烙铁，狠狠地烫在了每一个读到它的人心头上：

“我未失礼于彼，彼乃无礼横行，专恃兵坚器利，自取决裂如此！”

“与其苟且图存，贻羞万古，孰若大张挞伐，一决雌雄？”

檄文精准地道出了所有被压迫、被孤立、被践踏者心中积郁已久、几乎要喷薄而出的愤懑与那份置之死地而后生的决绝。

\vspace{2em}

在那座象征着内域传统权力巅峰的紫阳宗凌霄殿内，副宗主、守旧联盟此战的实际统帅陈炎，同样手持一份抄本。

他端坐于主位之上，指尖漫不经心地轻轻敲打着玉简，脸上竟也浮现出一丝复杂的、近乎玩味的笑意，低声自语道：“好骂。”

但这声赞叹里，浸透了居高临下的冰冷轻蔑与毫不掩饰的嘲讽。

大张挞伐？一决雌雄？

陈炎几乎要嗤笑出声。

在他看来，这种宣战甚至完全是有利于他们一方的：以前或许还有所顾忌、不敢轻易越雷池一步；但战书一下，双方自然可以堂堂正正、短兵相接了。

檄文越激愤，在陈炎的眼中，便越显得可笑而悲凉，如同困兽最后的哀鸣。

他的自信源于冷酷到极致的现实计算：

紫阳宗有着久经操练、传承有序的精锐弟子，他们演练纯熟的战阵攻防一体；丹鼎阁有着庞大的资源储备仓库，其门下擅长辅助与诡异毒术的药修、毒修被大量编入作战序列；而已经被刘长靖彻底掌控的无极剑派及其众多附属势力，则能够提供最为锋利、擅长攻坚与突袭的剑锋。

而一个崛起不过数载、根基浅薄、全靠一群被排挤的散修和少数几个叛逆宗门弟子支撑起来的新生势力，其底蕴拿什么来抗衡紫阳宗、丹鼎阁、无极剑派这三大巨头数百上千年积累的庞大资源、深厚功法和无数强横修士？

就算己方不主动进攻，仅凭旷日持久的资源消耗和硬实力的绝对碾压，也足以将云霞山上那点可怜的反抗力量慢慢磨死、困死。

\vspace{2em}

云霞山上。

虽然说叶牧谦和吴东浩形成了共识，不能再继续绥靖以紫阳宗为首的老大势力的联合绞杀，必须要重拳回击，他们在诸弟子前往青鸾秘境的过程中，在外界也确实拉起了一支由散修构成的义军；但在如何回击上，双方确实产生了分歧。

此刻，云隐别院之内，关于具体战术的实施，正在进行一场争论。

与其说是争论，不如说是一场理念与情绪的直接对撞。

青石板周围聚集了三四十人，除了叶牧谦、白言冰、陈千羽、沈瑶，以及很早就跟随叶牧谦随修的陈默等人，其余皆是吴东浩凭借自身威望与吴刚的情报网网罗来或主动汇聚来的、自愿组成义军的散修头领或代表人物。

洞府内的空气因灵灯与众人身上散发的驳杂气息，显得有些浑浊。

叶牧谦本人虽然没有完整地读过《论持久战》，但是其精髓倒是早早地刻进了他的心中。

而他的观点自然也是非常朴实无华的：“‘问道途’有民心汇聚之功，此战长远来看，我方必胜；但按如今实力来看，敌强我弱，我义军不能正面撄其锋芒，故不能速胜。叶某认为，短期来看，义军主力必须固守少数据点；其余人等化整为零，潜入广大乡村地区，对敌军游而击之，并广泛传播我问道途理念，团结更多群众于我旗帜之下，陷敌人于汪洋大海，然后胜之。”

吴东浩还没发话，草图另一侧，一位一身劲装、身形挺拔如枪的另一位散修头目立即插言。

头目名为尹壮图，修为正在玄脉境通络期。

他声音洪亮，斩钉截铁：“叶先生的方略，深远周全，尹某佩服。可我等举义，聚拢人心，靠的不是画地为牢、被动挨打！这么多年了，兄弟们从四面八方赶来，心里头都憋着一股火呢！”

他猛地一挥手，目光灼灼地扫过在场许多面露赞同、甚至隐含焦躁的散修头领：“我们必须要一场胜仗——一场堂堂正正、能让所有人都看见的胜仗！一场能把陈炎那老匹夫的嚣张气焰打下去的胜仗！提气，壮胆，更要告诉全内域的兄弟姐妹，我们不仅敢立旗，更敢拔剑，能砍下那群老爷们的脑袋！”

他的话语如同火星溅入干草堆，瞬间点燃了许多人的情绪。

这些散修大多是在旧有秩序下备受排挤、挣扎求存的个体，能修炼至今，靠的便是一股子狠劲和不择手段的拼搏；话语更是直接、血性，充满了草莽豪气，远比叶牧谦那些迂回而忍耐的策略更合胃口。

“尹兄说得在理！”

“对！咱爷们儿聚在这儿，不是来当缩头乌龟的！”

“总得干他娘的一票大的，让紫阳宗知道疼！”

“是啊，叶先生，咱们的实力未必就不如他们一股偏师，找个机会，狠狠咬下一块肉来！”

附和声此起彼伏，不少人的眼睛亮了起来，仿佛已经看到了伏击成功、凯旋而归的场景。

那种渴望用胜利证明自己、用敌人鲜血洗刷憋屈的心情，几乎要凝成实质。

叶牧谦静静地听着，脸上依旧平静，只是眉宇间掠过一丝几不可察的凝重。

待声浪稍平，他才缓缓开口，声音一如既往的平稳，却带着一种穿透喧嚣的清晰：“尹兄所言，是为提振士气，凝聚人心，此心可鉴。然，战争非意气之争，更非江湖械斗之放大。”

他向前半步，手指再次落回青石板，沿着代表敌军推进方向的几道箭头移动。“敌军大致分三路而来，看似有隙可乘。陈炎用兵，看似笨重迟缓，实则步步为营，稳扎稳打，其战阵互为犄角，灵力探测法阵与巡逻警戒网交错密布，几乎没有留给大规模伏击部队隐蔽接近、从容设伏的空隙。”

他顿了顿，看向诸多散修，“即便我们能集结一支足以撼动其一路偏师的精锐力量，成功突袭，最快能多久结束战斗？半个时辰？一个时辰？”

尹壮图立即被噎住了。

吴东浩则思考了一下：虽然说支持义军的声势确实浩大，但是这么长时间以来，真正上山、愿意拼命作战的不过两万左右，而修为稍高、到达玄脉境界者区区数千。

但数千玄脉境，再怎么说也已经不是一个小数目了。

吴东浩一估计，得出了结论：“若谋划得当，突袭迅猛，速战速决，一个时辰内当可击溃其一部前锋！”

“一个时辰。”叶牧谦重复道，指尖在敌阵后方一点，“足够另外两路，甚至中军主力做出反应。届时，我方伏击部队大概率陷入夹击或苦战；即便伏击取胜，缴获些许物资，杀伤数百敌玄脉境界及以上的修士，于敌军整体而言不过九牛一毛。然我军可能付出的代价，将是数百甚至上千最精锐、最敢战之士的折损。”

他抬起头，目光扫过那些情绪激昂的头领：“以我有限之血勇精锐，赌敌军无穷之资源消耗；况且，败则伤筋动骨，胜亦无济于事，此实是下下之选！”

另一位散修头目，叫曹润生的，脸色有些涨红，争辩道：“叶先生！打仗哪有不冒险的？若是瞻前顾后，何时能打开局面？我们一味退守，只会让他们越发骄狂！兄弟们心里这口气不顺，修为都难有寸进！我并非要浪战，而是寻机歼敌一部，哪怕两败俱伤，也要打出我‘问道途’的血性！让天下人看看，我们不是任人拿捏的软柿子！”

叶牧谦自己也有些忍不住了：“你们这是胡闹，用好不容易聚集起来的主力去胡闹！”

他再次指向草图上那片代表敌人腹地、统治薄弱的地域，“决胜之力，在彼不在我。过早暴露我们有限的主力精锐，将其消耗在正面碰撞中，乃是舍本逐末。”

不同的主张清晰地摆在了所有人面前：诸多散修要的，是一场立竿见影、提振人心的正面战斗，哪怕代价惨重；叶牧谦坚持的，则是隐忍固守，积蓄力量，将胜负手放在更深远、也更需要时间的敌后发动与群众争取上。

会场的氛围变得微妙而凝重。散修们大多面露不甘，他们觉得叶牧谦太过书生气，太过保守，甚至有些怯战。

在他们混迹底层的经验里，“百无一用是书生”，那种书生气是要不得的，必须靠着一股不要命的狠劲杀出条血路才行！

退让和隐忍往往意味着失去机会，随即人为刀俎、我为鱼肉！

叶牧谦描绘的“星火燎原”固然美好，但太过遥远：远水难解近渴。

眼下，他们需要的是胜利，是证明，是宣泄！

而叶牧谦这边，除了白、陈、沈三位无条件相信他的判断之外，吴刚也拧着眉头，看看正气势汹汹的尹壮图、曹润生等，又看看叶牧谦，似乎内心也有些挣扎。

他理解那些散修的躁动，因为那也曾是他的心态——他变得长袖善舞，那已经是后话了。

眼见得会场内部隐隐有了分裂的迹象，吴东浩深吸一口气，知道言语上的辩论已难见效。这时候，便是他发挥能力的时候了。

和稀泥！

他转向众人，抱拳道：“诸位兄弟，是非曲直，难以论说。我等既举义旗，便是军中同袍。军中最重号令统一，而号令出自众议。我提议，就此战术方略进行表决。在场诸位，皆有表决之权。无论结果如何，吴某必当遵从！”

此言一出，附和声再次响起，且更加响亮。对于这些散修而言，这种“公投”式的方式，远比听从某一人——尤其是叶牧谦这样“不通实务”的学术领袖——的决断，更让他们感到被尊重。

叶牧谦眼帘微垂，心中无声地叹息。

这就是人心！

在巨大而切近的压力下，在渴望即时反馈的情绪驱使下，更具煽动性、更符合直接经验的主张，往往能压倒更理性、更深远但见效慢的计划。

他也没有反对表决，因为强行压制只会离心离德。

“行。便依吴兄所言。”

表决的过程很快，几乎毫无悬念。

白言冰、陈千羽、沈瑶，自然站在叶牧谦一边；

吴刚、吴东浩弃权。

而其余三十余名散修头领或代表，超过九成，齐刷刷地举手，或者出声，明确表示支持“主动寻机出击”的方略。他们的眼神热烈，带着对胜利的渴望。

相比之下，叶牧谦那套守正待机、“决胜于千里之外”的理论，在他们听来有些缥缈，远不如“带兄弟们砍翻他们”来得实在。更何况，叶牧谦温文尔雅，更像一位令人尊敬的师长、学者；而尹壮图、曹润生则浑身散发着刀头舔血、敢打敢拼的“自己人”气息。

结果显而易见。

尹壮图再次抱拳，声音铿锵：“承蒙诸位兄弟信任！尹某在此立誓，必当竭尽所能，寻隙而击，为咱们‘问道途’义军，打出第一场漂亮的胜仗！”

欢呼声在洞府内响起，充满了对即将到来战斗的期待。

叶牧谦安静地站在原地，青衫素雅，与周围热血沸腾的氛围显得有些疏离。白言冰担忧地看向师尊，陈千羽抿了抿唇，沈瑶则目光微冷地扫过那些欢呼的散修。

待声浪稍歇，叶牧谦才缓缓开口，声音依旧平静无波：“既然众议已决，那……吴兄便为义军总指挥，节制诸人；尹壮图、曹润生则为副手，总揽寻机出战事宜。”他看向吴东浩，又扫了一眼尹、曹二人，目光深邃，“望你们审时度势，慎选战机，保全我义军元气为要。”

尹、曹二人郑重抱拳：“叶先生放心，我等省得！”

叶牧谦点了点头，不再多言。

打一打，也好。

实践是检验真理的唯一标准，即使实践存在代价。而搞一场会战，也确实能真正摸一摸敌人的底细。

当然，在叶牧谦本人看来，这场会战必然失败。他相信吴东浩的镇场能力，只希望尹、曹二人能别把义军主力糟蹋得元气大伤了。

\vspace{2em}

数日之后，义军开拔，一万余人浩浩荡荡地下了山，本来因为“山林啸聚”一番而鼎沸的云霞山顿时冷清了不少，只剩下少数后勤部。

叶牧谦、白言冰、陈千羽、沈瑶均未随军。

望着那七拼八凑但气冲斗牛的、渐行渐远的义军部队，叶牧谦深深地叹了一口气。

“瑶儿，走吧，再去加固一下九宫迷踪阵。言冰，你去协助那些后勤部队，在敌人可能攻上来的地方埋设陷阱、建设工事。千羽，你去看看药田，保证能够按时收获。”

四人各个心乱如麻。

“未来大概率会有一场守山的苦战了。”叶牧谦叹气道。

\vspace{2em}

这是一片位于云霞山以北约三百里、两片低矮丘陵夹峙下的宽阔浅谷。

按照尹壮图与几位心腹部下反复侦察推演得出的结论，这里是守旧联盟西路一支偏师——由紫阳宗某位外门长老孙翊率领、约八千人的战营——前往预定集结地的必经之路，且因其地形相对闭塞（虽然也没有那么闭塞），大军难以迅速展开，是打一场迅猛伏击的理想场所。

吴东浩本人倒也谨慎，一万多人的义军，留下了两千士兵在更远的地方试图接应、作为预备队。

前来埋伏的，大约万人。

尹壮图站在沟壑一侧的隐蔽高点上，劲装被山风吹得紧贴身躯，猎猎作响。

他目光如鹰隼，扫视着下方逐渐进入伏击圈的敌军队伍。那支队伍军容严整，士卒身着赤红的制式灵甲，行动间隐隐有灵力共鸣的微光流转，步伐虽然因山路而稍显缓滞，但前后呼应严密，斥候游骑交替前出，显示出良好的训练素养。

他身边，曹润生半蹲着，手中一枚符箓微微闪光，低声道：“尹兄，看旗号，是紫阳宗的一部，主将是那个叫孙翊的长老，听说修为在玄脉后期。他们队列中央那些覆盖着毡布的大车，灵力波动不弱，怕是运送的补给或某种法器。”

尹壮图舔了舔有些干裂的嘴唇，眼中非但没有惧色，反而燃起更炽烈的战意。

“紫阳宗的赤皮狗，装备倒是光鲜。正好，宰了这条大鱼，扒下他们的皮，给兄弟们换换装！告诉各队，按原计划，等他们前军过了一半，中军那些大车进入沟底最窄处，听我号令，一起动手！先集中火力轰击车队，打乱他们阵脚，分割前后，然后近身绞杀！”

在他们的规划下，中军将是主要的突破口，在这里安排了六千人之多。而在突袭发起之后，袭扰前军和后军的，分别不过千人。

命令被悄无声息地传递下去。埋伏在沟壑两侧乱石、灌木后的上万义军修士，此刻屏息凝神。他们大多眼神狂热，紧握着自己五花八门的兵器——有的只是凡铁打造、附加了基础锋锐符文的刀剑，有的是品相不一的低阶飞剑，更有拿着奇门钩锁、淬毒短弩的。

虽然装备驳杂，但那股子被压抑了许久的、渴望用敌人鲜血证明自己的杀意，却几乎凝成实质，让空气都变得粘稠而躁动。

吴东浩隐在另一侧的山崖后，眉头紧锁。

他同意了这次出击，但内心深处那丝不安始终挥之不去。

他看得比尹壮图更清楚：敌人数量虽略比我方少，但其阵型严谨，反应速度绝不会慢。

义军人数虽多，可一旦第一波突袭不能达到预期效果，陷入缠斗，后果不堪设想。

“盯紧侧翼和后方，一旦情况有变，准备接应尹壮图他们撤出。”

此刻，那支队伍，如同一条赤红色的巨蟒，缓缓游入黑云沟的腹腔。阳光被两侧山崖遮挡，沟内光线昏暗，只有甲胄和兵刃偶尔反射出冰冷的幽光。队伍中隐约传来军官的呵斥和车轮碾过碎石的声响。

前军过了谷口，辎重正好进入沟底最深处。

时机到了！

尹壮图猛地站起，灵力毫无保留地爆发开来。他手中那柄门板似的厚重巨剑高高举起，剑身嗡鸣，亮起土黄色的沉重光芒。

“兄弟们！随我杀——！”

“杀啊——！！！”

“为了问道途！为了活路！”

积蓄已久的怒吼声，如同山洪暴发，从沟壑两侧轰然炸响！刹那间，成千上万道灵力光芒亮起，虽然驳杂不纯，却汇聚成一片令人目眩的死亡潮汐，朝着沟底的敌军倾泻而下！

最先落下的是各种远程攻击。火球、冰锥、风刃、雷符、毒镖、飞针……密密麻麻，如同夏日狂暴的冰雹，劈头盖脸砸向尚未完全反应过来的紫阳宗队伍。

尤其是针对那十几辆覆盖毡布的大车，更是受到了重点关照，超过百道攻击汇聚成数股洪流，狠狠撞了上去！

轰！轰轰轰——！

剧烈的爆炸声接连响起，火光冲天，浓烟翻滚。几辆大车瞬间被炸得粉碎，里面的灵石、矿石、甚至尚未组装完成的弩炮零件四处飞溅。

拉车的驯化妖兽发出凄厉的哀嚎，受惊之下横冲直撞，进一步搅乱了敌军队列。

“敌袭！御！”

紫阳宗队伍中响起了尖锐的警哨和军官声嘶力竭的吼叫。

训练有素的精锐在这一刻展现了价值，尽管遭遇突袭，出现了一些伤亡和混乱，但士卒仍然在极短时间内做出了反应。赤红色的灵光从他们甲胄上升腾而起，迅速连成一片，形成一面巨大的、半弧形的火焰护盾，将后续的大部分攻击抵挡在外。

盾牌铿锵组合，长矛如林竖起，瞬间构筑起一道坚固的防线。

但义军的突袭势头太猛，第一波打击也确实造成了可观杀伤。近百名紫阳宗修士在最初的混乱中非死即伤，队伍被拦腰截断，前后失去了联系。

“就是现在！冲下去！近战！杀光他们！”尹壮图身先士卒，如同一头发狂的猛虎，从数丈高的崖壁上一跃而下，巨剑带着开山裂石之势，狠狠斩向一名刚刚组织起小队、试图反击的紫阳宗什长。

那什长举盾格挡，但尹壮图含怒一击何等狂暴，只听“咔嚓”一声脆响，精铁盾牌连同其后的一条手臂应声而断，鲜血喷溅！

尹壮图去势不减，巨剑横扫，又将旁边两名惊愕的紫阳宗修士拦腰斩飞！

“尹首领威武！”

“杀啊！”

主将的悍勇极大鼓舞了士气，无数义军修士红着眼睛，从埋伏点冲下，如同决堤的洪水，涌向被分割开的敌军段落；灵枢境的找灵枢境的大头兵，玄脉境的找玄脉境的长官，捉对厮杀。

短兵相接，血肉横飞的惨烈战斗，瞬间在黑云沟狭窄的底部上演。

个人武艺与悍不畏死的精神，在这一刻似乎取得了优势。

许多散修出身、在底层摸爬滚打练就了一身厮杀本领的义军，三五成群，凭借对地形的熟悉和一股不要命的狠劲，围着那些试图结阵自保的小股紫阳宗修士猛攻。

刀光剑影，符篆乱飞，惨叫与怒吼混杂在一起，浓烈的血腥味迅速弥漫开来。

曹润生也杀入了战团，他使得是一对分水峨眉刺，身形滑溜如鱼，专挑敌人防护薄弱处下手，出手狠辣刁钻，几个照面便放倒了两名紫阳宗士卒。他看到战局似乎正向有利于己方的方向发展，脸上忍不住露出一丝兴奋的潮红，高喊道：“兄弟们加把劲！吃掉他们！”

分配给敌人前军和后军的小股部队也开始袭扰起来，迟滞着敌人的回援，短期内成果斐然。

就连远处观战的吴东浩，看到初期战果，紧绷的心弦也略微松了一分。

难道……叶先生的判断还真是过于保守了？

我们真能啃下这块硬骨头？

然而，他们高兴得太早了。

孙翊此刻已从最初阵型被搅乱的惊恐中镇定下来，冷哼一声，命令通过灵力传遍己方尚能联系的部队。

“传令！中军抛弃辎重，自行结阵固守！前后两部，起焚城大阵，左右卫护，缓步压上，与中军靠拢！”

想了想，又对近卫补充道：

“发求援焰火，告知西路周圣子：我部遇伏，敌军近万，请求合围！”

命令迅速被执行。

只见那些尚未被义军彻底缠住的中军紫阳宗修士，迅速放弃与冲下来的散修纠缠，开始向后收缩。他们彼此靠拢，大约每百名修士便迅速聚集成一个紧密的圆阵，外围盾牌层层叠叠，长矛从缝隙中刺出，内里的修士则开始整齐划一地掐动法诀，赤红色的火行灵力在他们之间疯狂汇聚、共鸣。

“不好！他们要结战阵！”一名有见识的义军小头目骇然变色，大声示警。

但已经晚了。

中军的一个个防御圆阵已经结成，变成铁桶，打不动了！

更要命的是，前军和后军已经结成了更具有攻击性的战阵！

“烈焰焚城，焚尽八荒！疾！”

孙翊暴喝一声，手中令旗猛地挥下！

嗡——！

众多紫阳宗修士组成的战阵中心，空气剧烈扭曲，一团团直径超过十丈、炽白到让人无法直视的巨大火球骤然成型，散发出恐怖的高温，连周围的岩石都开始融化、流淌！

下一刻，这团毁灭性的火球轰然炸开，化作无数道手臂粗细、灵动如蛇的炽白火线，以战阵为中心，呈放射状向着四面八方疯狂蔓延、抽打、绞杀！

这正是紫阳宗核心战阵之一的威力！

嗤嗤嗤——！

火线所过之处，无论是义军修士仓促撑起的护体灵光，还是他们手中的低劣法器，甚至是坚韧的肉身，都如同热刀切猪油般被轻易撕裂、洞穿、点燃！

“啊——！我的腿！”

“火！扑不灭！”

“退！快退开！”

惨叫声此起彼伏，原本气势如虹、扑到近前的义军，顿时遭受重创。

数十名冲在最前面的悍勇之士，瞬间被火线穿透，变成一个个燃烧的火人，凄厉地翻滚着，很快化为一堆焦炭。更多的人被火线擦中，护体灵光剧烈闪烁后破碎，身上燃起难以扑灭的灵火，惨叫着向后逃窜。

原本还算有序的进攻浪潮，被这突如其来的、覆盖性的恐怖打击彻底打懵、打散！

而这，仅仅是个开始。

孙翊显然意图将这股胆大包天的伏兵全部留下。

战阵在第一波齐射后，开始如同一个个燃烧的刺猬，缓缓向前、向沟内压迫移动。火线不断从战阵中衍生、抽击，清扫着前方一切障碍。

同时，战阵内部的修士轮换施法，一道道范围稍小但同样致命的大片法术被抛射出来，落在义军人群相对密集的区域。

黑云沟内，顿时化为了火焰地狱。原本昏暗的沟底被映照得一片通红，炽热的气浪扭曲升腾，将空气都烧得滚烫。焦臭味、皮肉烧灼的可怕气味混合着血腥，形成了令人作呕的死亡气息。

义军最大的劣势——缺乏统一指挥和应对集团战阵的经验——在此刻暴露无遗。

他们习惯了小规模的混战和游击，也对此颇有自得；何曾见过如此恐怖、如此有组织性的灵力覆盖打击？

慌乱之下，众多小队队长的命令也混乱起来：

“顶上去！”

“跑啊，等死呢！”

“我听哪个！”

命令互相冲突，导致局部更加混乱。

有小队想从侧翼绕过战阵，但那战阵移动间毫无破绽，很快全队死的死、逃的逃。

有人想远程攻击打断施法，可零星的法术打在厚重的联合灵力护盾上，连涟漪都难以激起。

尹壮图目眦欲裂！

他亲眼看着自己麾下几名得力干将，在试图组织一波反冲锋时，被交叉扫过的七八道炽白火线直接气化，连渣都没剩下。

他狂吼一声，将巨剑舞得密不透风，磕飞了几道袭向自己的火线。虽然尹壮图本人修为较高、扛下了众多低阶修士组成的超大型战阵的攻击；但那火线蕴含的灼热灵力和冲击力，震得他气血翻腾，手臂发麻。

“不能退！退了就全完了！跟我上，杀了孙翊那家伙！！”

尹壮图知道，不击破中军、不停止孙翊的动作，今日所有人都要葬身火海。

尹壮图与曹润生集结了身边还能指挥的约两百名最精锐、最不怕死的修士，鼓起残存的灵力，化作一股决死的尖刀，试图强行凿穿战阵外围的防御，强行换死孙翊，停止他对毁灭性超大战阵的指挥。

这是一场鸡蛋碰石头的自杀式冲锋。

“螳臂当车！”孙翊在阵中看得分明，嘴角泛起残忍的冷笑。

“剑修！给我撕了他们！”

早已在战阵侧后方蓄势待发的五百名无极剑派外剑脉剑修，如同一群闻到血腥味的鲨鱼骤然杀出！他们的任务，就是应对这种敌方精锐的突围或反击。

这些剑修身法极快，行动间带着森然的统一性。他们没有结什么复杂战阵，只是简单分为数支小队，如同几柄淬毒的匕首，倏忽之间便切入尹壮图组织的冲锋队列侧翼和薄弱衔接处。

剑光！

冰冷、迅疾、精准、高效的剑光！

他们的剑法远远不如已被清理的心剑脉那般灵动缥缈、直指本心，但在杀伐效率上，却达到了某种极致。

每一剑都直奔要害，配合默契，往往两三人一组，瞬间便能绞杀一名冲在前面的义军好手。

“呃啊！”曹润生一个不慎，左腿被一道刁钻的剑气划过，深可见骨，鲜血狂喷。

他闷哼一声，险些栽倒，被旁边一名修士拼死拖住。

尹壮图的冲锋势头被硬生生遏止，陷入了与这些精锐剑修的死斗之中。他巨剑势大力沉，每一击都足以开碑裂石，但那些外剑修士滑溜异常，根本不与他硬拼，只是游斗、袭扰、分割，将他身边稍不留神的战友们一个个快速点杀。

他空有一身勇力，却仿佛陷入了一张无形而坚韧的剑网，憋屈得几欲吐血。

就在尹壮图部陷入苦战、整个义军战线因为核心战阵的碾压和外剑修士的切割而开始动摇、溃散的危急关头——

呜——！

低沉雄浑的号角声，从黑云沟的两端入口处，同时响起！声音苍凉悠远，却带着令人心悸的肃杀之意。

紧接着，地平线上，烟尘大作！旌旗招展，灵力冲霄！

西路讨伐军的真正主力，到了！

而且不止一路！从沟壑的东、西两个出口方向，同时出现了滚滚而来的钢铁洪流。东面是更多的紫阳宗战阵，赤旗如血；西面则是丹鼎阁与部分附属宗门的联军，旗帜花色繁杂，但阵型同样严整。他们显然接到了孙翊的求援，并做出了最迅速、也最致命的反应——反包围。

周文彬本人没出现，似乎是算准了用不着他出手。

原本计划打一场速战速决伏击战的义军，此刻反而成了瓮中之鳖。

“完了……”吴东浩远远看到这一幕，一颗心沉到了谷底。

最坏的情况发生了。

敌人的反应速度和兵力调动效率，远超他们的预估。陈炎“稳扎稳打，互为犄角”的策略确实是将强大的实力化为了难以撼动的整体优势。

他们这支孤军深入的义军主力，已经成了送上门的肥肉！

“吴首领！东、西两面都有大队敌军出现！我们被包围了！”浑身浴血的斥候连滚爬爬地冲过来报信，脸上满是绝望。

吴东浩手青筋暴起，他猛地看向沟底仍在拼死搏杀的尹壮图部，又看看周围已经开始出现恐慌、骚动的义军部队。他知道，必须立刻做出决断，能撤出去多少是多少！

“传令！各部不要恋战！以小队为单位，自行突围！能走多少是多少！尹首领那边……我去接应！”吴东浩咬牙。

他知道凭自己真符境摹形期的修为，在许多个成熟战阵中救出尹壮图希望渺茫；但作为义军总指挥，他不能眼睁睁看着主将战死而不救，那会让士气彻底崩溃。

然而，战场局势的崩坏速度，比吴东浩预想的还要快。

当两路生力军加入战场，并且迅速在外围构建起更严密的封锁线，开始用远程法器和战阵技能向沟内进行覆盖式轰击时，义军的崩溃便不可避免地发生了。

丹鼎阁的修士隐匿在侧翼，将大片大片的甜腻毒瘴借助风势吹入沟内，许多修为较低的义军吸入后立刻手脚酸软，灵力运转停滞，成了待宰羔羊；地面在土行法术的干扰下变得泥泞不堪，甚至突然塌陷，严重阻碍了突围的脚步。

绝望像瘟疫一样蔓延。原本凭借血勇支撑的斗志，在绝对的实力差距和绝境面前，迅速消融。

有人扔下武器，跪地投降，但换来的往往是毫不留情的一剑。

更多人像没头苍蝇一样乱撞，然后被各色法术光芒吞噬。

尹壮图身边的战友越来越少，最后只剩寥寥十余人，被超过百名的紫阳宗精锐和外剑修士团团围住在一个小山坡上。

他浑身是伤，巨剑拄地，大口喘着粗气，眼神却依旧桀骜凶狠，死死盯着不远处被重重保护的周焱。

“贼子！可敢与尹某单挑！”他嘶声吼道，声音沙哑。

孙翊嗤笑一声，甚至懒得回答，只是摆了摆手。

顿时，更加密集的法术和剑气向那小山坡倾泻而去。

“保护首领！”最后几名亲卫怒吼着扑上，用身体挡住了大部分攻击，瞬间被淹没在光芒中。

尹壮图狂吼，压榨出最后一丝灵力，巨剑爆发出最后的黄光，一式力劈华山，斩碎了袭向自己的几道剑气和火球。但他已是强弩之末，护体灵光黯淡如风中残烛。

噗嗤！噗嗤！

两柄飞剑，一左一右，如同毒蛇吐信，趁机突破了他巨剑的防御范围，狠狠贯入了他的胸膛和腹部！

尹壮图高大的身躯猛地一震，他低头看了看透体而出的剑尖，脸上闪过一丝茫然，随即化为无尽的愤怒与不甘。

“悔……不该……”他喃喃吐出几个字，眼神的光芒迅速黯淡下去，最终轰然倒地，激起一片尘土。那双曾经燃烧着炽热战意的眼睛，至死未曾闭合，直直望着云霞山的方向。

几乎在尹壮图倒下的同时，吴东浩才杀入重围。

见到尹壮图已经陨落，也只能长啸一声，扛起尸体再次强行杀出。

真符境毕竟不是法相境，虽然说这么多战阵留不下他，但要说硬碰硬，那还是远远不够的。

另一边，曹润生早在腿受重创时就被亲信架着后撤，但突围路上同样遭遇了数次截杀，身上又添了几处新伤，最重的一处在后背，是被一道散射的剑气所伤，差点伤及灵枢，人也因失血过多而昏迷过去。

主将战死，副将重伤。吴东浩此刻留下的后手——在远处接应的数千修士——确实起到了自己的战术作用，成功发动了突袭、搅乱了些战场局势，但已经不能改变溃败的结果了。

而原来埋伏敌人的、剩下的修士彻底失去了组织，只能凭借本能亡命奔逃。

那已经不能称之为撤退，而是一场彻头彻尾的溃败。

守旧联盟的军队则如同经验丰富的猎手，开始有条不紊地追杀、分割、歼灭溃逃的残敌。

一开始参与埋伏的近万人，折了五千有余，剩余的拼死拼活才逃了出来；预备队伤亡倒是不大，但也心有余悸。

不少人本因为自己被编入预备队而感到不忿，但此刻只剩下更大的庆幸。

夕阳如血，缓缓沉入西边的山峦，将这片修罗场映照得一片凄艳。

喊杀声渐渐平息，取而代之的是伤者压抑的哀嚎、胜利者打扫战场的喧嚣，以及食腐妖兽被血腥味吸引而来的低沉嘶吼。

曾经气势汹汹而来的一万两千余义军修士，最终能够狼狈逃出的，不足八千。而且绝大多数带伤，建制彻底打散，士气低落到了极点。

更惨重的损失是高层：义军骁将尹壮图确认战死；曹润生重伤昏迷，脊背的伤势可能留下永久性的隐患，即便醒来，战力也大不如前。吴东浩本人修为高绝，在背负尹壮图尸体回来的途中受了些轻伤，所幸无大碍。

其余在战斗中表现出色、有一定威望的中下层头目，伤亡更是超过三分之一。

经此一役，云霞山义军可谓元气大伤，不仅损失了大量最敢战、最有经验的骨干，更重要的是，那口下山时“欲与天公试比高”的心气，被彻底打没了。

恐惧、沮丧、质疑、乃至对叶牧谦当初劝阻的后悔情绪，开始在残存者中悄然蔓延。

而对守旧联盟而言，此战役则是一场预料之中的胜利。虽然己方阵亡修士也大约二千有余，但比例上看并不算多；更重要的是，它以极小的代价重创了对手的有生力量，还验证了己方战略战术的正确性，极大地鼓舞了士气。

陈炎在接到战报后，只是淡淡地评价了四个字：“乌合之众。”

通往云霞山的最后屏障，正在这场血腥的大败后，缓缓洞开。真正的铁壁合围与攻坚，即将开始。而云霞山上，叶牧谦站在反复加固过的阵眼枢纽处，接到染血的败报时，脸上并无多少意外之色，只有一片深沉的凝重与决绝。

实践确实检验了真理。只是这代价，未免太过沉重。

随后，在战役失败的总结与反思会议上，叶牧谦自然而然地接过了义军的总指挥权。

叶牧谦不禁喟叹：每次他接手的总是一片烂摊子！

但斯大林格勒从来不相信眼泪。

\vspace{2em}

大规模会战宣告失败后的一年内时间，便在无数场规模不等、却同样浴血的小规模接战，在一次次带着不甘与悲愤的撤退，在一声声无奈放弃经营许久据点的叹息中，迅速流逝。

吃下一地，稳固一地，再吃下下一地！

一年，原问道途所有地下情报与走私网络，近乎全毁！

没毁的也不敢再出来活动了！

守旧联盟的大军，如同一个巨大而沉重的石磨，缓慢、笨重，却带着无可抗拒的力量，一寸一寸地碾过了问道途势力范围的外围城镇、资源点和战略据点。

最终，这碾压一切的兵锋，携带着胜利者的傲慢与肃杀，直抵那最后的、也是最高的屏障——云霞山。

这一日，陈炎亲率中军主力，旌旗蔽空，终于兵临云霞山下。

站在山巅俯瞰，景象足以让最坚韧的战士也感到一阵窒息般的绝望。

山下，守旧联盟的营寨依着地势，连绵铺展开去，一眼望不到边际，仿佛给大地铺上了一层铁灰色的毡毯。

各色代表不同宗门、不同战营的旌旗如林般矗立，在风中猎猎作响，几乎遮天蔽日。

从那些规划严整的营寨中，升腾起一道道、一片片色泽各异的灵力光晕。

赤红的火行灵力灼热逼人，青木的生机中暗藏毒瘴，白金剑气锋锐刺目，土黄厚重如山，水蓝绵密森寒……

这些光晕交织、融合在半空，形成了一片瑰丽奇诡却充满纯粹杀伐之气的灵力光云，冲霄而上，将云霞山半边天空都渲染得变了颜色，连日光都显得黯淡昏沉。

出入山林的每一条大小要道，都被鹿砦、壕沟以及时刻运转的警戒与探测法阵几乎彻底锁死，巡逻修士的队伍甲胄鲜明，气息精悍，交错往复，目光如鹰隼般扫视着山上的动静。

山下，是阵列森严的钢铁森林，是汹涌澎湃、近乎无限的灵力海洋，是志在必得、以泰山压顶之势而来的碾压洪流。山上，是历经一年苦战、人人面带疲惫却眼神依旧不屈的战士，是日益见底的丹药瓶与灵石箱，是寄托了内域无数受压迫者最后希望与信仰的险要壁垒。

一副瓮中捉鳖的恶劣态势，已然在战场层面彻底形成。

战争的初级阶段，以守旧联盟凭借绝对硬实力的无情碾压，和问道途一方被迫陷入艰苦卓绝的战略防御、乃至被重重围困的绝境而告终。

那赤裸裸的、令人绝望的实力差距，在这场从一开始就极不对等的生死较量中，展现得淋漓尽致，冰冷而残酷。

然而，正是在这山下连营的肃杀之气几乎要凝结成实体、如同铅云般沉沉压上云霞山巅、让人喘不过气来的时刻，坐镇于云隐别院深处、那掌控整座山防御体系核心阵眼处的叶牧谦，其神色却异乎寻常地平静，甚至可以说得上是淡漠。

他独自立于阵眼枢纽的灵光之中，青衫素雅，身形挺拔如孤峰。

那双仿佛能洞悉世间一切虚妄与表象的眼睛，平静地望向山下那密密麻麻的营寨，望向那冲霄而起、炫耀着武力的斑斓灵力光云。

眸底深处，并无半分寻常人应有的慌乱、恐惧或是愤怒，只有一片如同古井深潭般的、深邃的了然。

过去两年多看似平静实则暗流汹涌的冷战岁月里，过去一年正面战场被不断推进、绞杀的日子里，除了指挥义军撤退，他倾注最多心血、智慧与珍贵资源的，便是此云霞山本身了。

他有信心，九宫迷踪能够将脚下这座云霞山，从一座普通的险峻山峰，经营成一座纵使面对千军万马、拥有雄厚资源的敌人，也无法被轻易碾碎、啃下的战争堡垒。

而作为整个堡垒核心的，自然就是他的得意之作：九宫迷踪阵！

“传令各部，”叶牧谦的声音通过嵌入山体灵脉的特定传讯法符，清晰、稳定、不带丝毫波澜地回荡在山上各处要害据点守备负责人的耳中。

“放弃外围所有隘口、前出哨所，全部撤回云霞山主峰预设防御圈内。”

他略一停顿，那平静的声音却下达了一个将决定接下来无数人生死的命令：

“即刻起，开启九宫迷踪阵第一重变。”

命令如山，早已在无数次演练中刻入骨髓的守军，此刻如同上好了发条的精密器械，开始沉默而高效地运转。

散布在山脚下各处险要地点，原本准备誓死阻击、拖延时间的修士小队，毫不犹豫地放弃了经营许久的阵地。

埋设预设的阻挠陷阱后，小队们沿着早已勘测熟稔的隐秘小径，迅速而有序地向主峰核心区域收缩。

在敌我实力差距悬殊到令人绝望之时，主动将伸出的五指攥紧成拳，收回胸膛，看似退缩，实则是为了将力量凝聚于一点。

将那看似柔软的棉絮层层包裹内部最坚硬的铁芯，准备在敌人最意想不到、也最能发挥己方优势的位置，给予其最凶狠、最致命的反击。

这，是叶牧谦前世阅读那些波澜壮阔的历史时，从一次次以弱胜强、于绝境中开辟新生的经典战例中，领悟到的精髓。

叶牧谦自认无法像前世的伟人一样有着极高的军事天赋与才能，而现在的条件也不允许：正面兵力不足以打出如雷霆之势摧枯拉朽的三大战役，他自己的指挥能力也不足以打出四渡赤水这样精彩且精确的战役。

但如果仅仅是打一个反围剿，那似乎容易不少，至少他有很多现成的答案可以抄。

\vspace{2em}

云霞山，本就以山势险峻、路径错综复杂著称。

强行仰攻本就是兵家大忌，而加上九宫迷踪阵，对攻山的守旧联盟而言，情况只会更差。

就在守旧联盟的先锋部队，以严整战阵为依托，沿着几条主要山道气势汹汹、志在必得地向上推进之时，原本清晰可辨的山路、熟悉的溪涧与巨石，仿佛被一只无形而庞大的手轻轻搅动、抹去，然后重新排列。

浓密得化不开、伸手不见五指的白雾，不知从岩缝中、树根下还是地脉里汹涌而出，顷刻间如同涨潮的乳白色海水，弥漫、吞没了整个山林外围区域。雾气中蕴含着被阵法强行扭转的紊乱灵气，如同无数细密的针，疯狂刺扰着闯入者的神识探查：

斥候刚刚回头，信誓旦旦地报告前方路径正常，转瞬间，身后的主力部队就骇然发现，方才走过的坚实山路已然消失，脚下竟是云雾缭绕的陌生断崖；

明明按照记忆和地图该向左急转的路口，踏出去的队伍却直通一个布满苔藓湿滑、且犬牙交错的天然石刺绝地。整片外围山区，在短短半盏茶的时间内，便化作了一座庞大、诡异、时刻变化的活迷宫。

这，还仅仅是令人头痛的序曲。

九宫迷踪阵的扭曲空间效应，与预先巧妙埋设在关键节点的无数陷阱相结合，产生了叠加的恐怖效果：

当敌军队列在浓雾中如盲人般艰难摸索时，脚下看似坚实的落叶覆盖的地面可能突然无声塌陷，露出下方深达数丈、底部插满淬毒精钢刺的陷坑，瞬间吞噬数名躲闪不及的修士；

头顶上方那仿佛亘古存在的岩壁会毫无征兆地整体松动，数十根被符文加持过、裹挟着万钧之力的合抱滚木与棱角狰狞的礌石轰然砸落，破风之声如同死神的咆哮；

更隐蔽且防不胜防的，是那些触发式的小型连环灵力爆符，或是被特殊法术催生的、见血便疯狂滋长的铁线荆棘藤蔓种子。

这些陷阱虽难以对结成战阵、有层层护身灵光笼罩的高阶修士造成瞬间致命伤，却足以在关键时刻打乱他们严谨的推进节奏，迫使整个战阵停滞或变形，更在不断消耗他们宝贵的灵力以维持防御，并在每一名敌军心中持续施加无形的压力，制造挥之不去的混乱与随时丧命的恐慌。

守旧联盟想象中摧枯拉朽的推进速度，仅仅在山脚之下，就被这层黏稠、危险、无所不在的迷雾沼泽极大地迟滞、拖住了。

\vspace{2em}

天师府。

作为内域名义上地位超然的中立势力，其总坛不断向冲突双方发出正式文函，呼吁“秉持天道好生之德，即刻停战，以和谈化解干戈”。

紫阳宗为首的守旧联盟自然对此嗤之以鼻，势在必得；云霞山一方更无妥协余地。

但天师府的举动，远不止于这些流于表面的公文往来。

总坛深处，幽静的观星阁内，府主沈天穹一袭朴素道袍，正凭栏远眺，目光似乎穿过了千山万水，落在了烽火连天的云霞山方向。他的指尖无意识地在冰冷的玉石栏杆上轻轻敲击，那节奏缓慢而深沉，如同他此刻反复权衡的内心。

对叶牧谦此人，沈天穹的情感颇为复杂。自女儿沈瑶拜师，以及后来那次开诚布公的长谈后，他便一直在关注这个看似年轻之人的动向。

叶牧谦所倡的“问道途”，那种打破陈规、务实求真、愿为更多修士乃至凡人开辟前路的理念，在沈天穹看来，固然激进，甚至有些理想主义的天真；但其内核的精神，却与他年轻时所怀的某种抱负隐隐共鸣。

那是一种不愿被旧有秩序完全束缚、渴望看见新变化的微妙欣赏。

然而，身处他这个位置，个人的好恶必须让位于全局的权衡。

沈天穹收回远眺的目光，转身望向阁内悬挂的内域堪舆巨图：图上写写画画，已经很旧了。

云霞山的位置被一枚红色的符文标记着，而代表守旧联盟兵锋的黑色箭头，正从四面八方向其合围。

他看得十分透彻：陈炎所率联军，实力占据绝对上风，若任由其速战速决，以雷霆万钧之势碾碎云霞山，那么此后内域的格局将彻底失衡。

一个迅速膨胀、吞并了革新派所有遗产——无论是有形的资源还是无形的声望——且毫无制衡的紫阳宗、丹鼎阁、无极剑派联盟，对天师府而言绝非福音。

那将意味着，天师府这超然中立的地位将受到前所未有的挑战与挤压，甚至可能被卷入更危险的漩涡。

一个僵持的、相互消耗的云霞山战场，才更符合天师府的长远利益——既能持续牵制、削弱陈炎一方的锐气和力量，又能为内域其他观望势力留下反应与抉择的时间。

更何况……

他的女儿，沈瑶，就在那座被围得铁桶一般的山上。尽管父女因理念不合早已疏离多年，当年激烈的争吵与沈瑶的毅然离家仍历历在目，那份血缘牵连与深藏的关切却从未真正断绝。

他得到的情报显示，山上战况惨烈，攻防每日都在见血。

作为一个父亲，他无法想象，若云霞山防线崩溃，沈瑶会遭遇什么。

\vspace{2em}

数日后，天师府总坛传出的命令，便带上了这多重考量混合后的温度。沈天穹以“苍生罹难，修士亦为生灵，不可坐视伤亡惨绝”为由，正式对外宣布，将对云霞山交战区域进行“人道援助”。

这个理由确实冠冕堂皇，任谁也挑不出大的错处，完美地遮掩了其下涌动的战略算计与私心牵挂。

命令下达，天师府这架庞大而古老的机器开始以特有的效率运转。一队队打着天师府杏黄旗号的车队，从总坛及各处分观出发。这些车队主要由未经历练的低阶弟子和谨慎挑选的世俗忠仆组成，车辆普通，护卫力量也仅足以应付寻常匪患，摆明了姿态——我们只是运送救济物资。

车上装载的，是“用以救治伤患、保全元气”的基础物资：品阶不高但数量可观的止血散、养气丹，洁净的灵纱布，以及大量不易腐坏、能勉强补充灵气消耗的低灵谷物。

它们无法决定战局，却足以维系最基础的生命线。

对于兵精粮足、后方补给源源不断的守旧联盟而言，这些“低品”物资不过是隔靴搔痒的锦上添花。他们接收得漫不经心，负责查验的军需官甚至略带嫌弃地检查着丹药的成色，挥挥手便让人将东西丢进庞大的后勤仓库角落，仿佛接收了一样碍眼的累赘。

然而，这些打着杏黄旗的车队试图靠近云霞山时，情形便截然不同了。

守旧联盟的外围哨卡如临大敌，通往山区的每一条小路都已被严密封锁。车队在距离防线尚有数里之地便被凶神恶煞的巡逻队拦下。

“站住！天师府的车队？此地乃军事禁区，闲杂人等不得靠近！”一名身着紫阳宗服饰的小头目按着剑柄，眼神不善地扫视着车队。他身后的士兵们早已结成简易战阵，灵力隐隐流转，封锁了前路。

带队的天师府低级执事是个面相憨厚的中年人，他陪着笑脸，取出盖有沈天穹印鉴的正式文书：“这位道兄，我等奉府主之命，运送人道援助物资，救治伤患，此乃文书凭据，符合天道常伦，还请行个方便……”

“方便？”小头目嗤笑一声，粗暴地打断他，“谁知道你们车里装的是丹药还是爆符？是粮食还是兵刃？前线正在鏖战，谁知道你们是不是打着援助的幌子，给山上的逆贼输送给养？退回去！”

类似的刁难，在各个方向的哨卡反复上演。车队被无故扣留盘查数个时辰，货物被要求全部卸车，一粒粒谷物、一颗颗丹药都被反复用神识探查甚至刮开检验；道路“恰巧”被落石或临时布置的简易法阵阻断；通关文书被挑出各种格式或措辞上的“问题”……

守旧联盟的基层军官们，显然得到了上峰的某种默许，竭尽所能地给这些杏黄旗车队制造麻烦，拖延、阻碍物资的送达。

但他们终究不敢真的撕破脸皮。

当带队执事据理力争，甚至抬出府主沈天穹的名号，并暗示若物资无法送达、伤患死者众，天师府必将质询其违背“天道好生”之责时，那些气焰嚣张的小头目们往往脸色变幻。

他们接到的指令，也确实是“刁难”与“拖延”，“尽量阻止物资上山”。

真正阻挡所有物资，等同于公然打沈天穹和整个天师府的脸。沈天穹本人作为法相境修士的强悍实力，像一座无形的大山，压在这些执行具体任务的修士心头。

他们可以刁难，可以扣留一部分，却绝不敢真的下令攻击车队或彻底没收所有物资。

于是，在反复的拉锯、交涉、妥协之后，总有一部分车队，在经历了重重盘剥与“损耗”后，被允许通过最后一道关卡。

对于被重重围困、资源如沙漏般飞速流逝的云霞山而言，这些历经千难万险才得以流入的援助，其意义远不止于物资本身。

当一袋袋染着尘土甚至血污的低灵谷物，一箱箱品相普通的止血散和养气丹被运回山上时，引发的确是绝境中难以言喻的激动。

久旱逢甘霖，他乡遇故知！

他们并未被外界彻底遗忘和抛弃，有一条极其微弱却真实的生命线，正穿透铁桶般的封锁，连接着山外！

叶牧谦与沈瑶对此心照不宣。沈瑶抚摸着短刀冰凉的刀柄，望向山外遥远的天师府方向，清冷的眸子里情绪复杂难明。

她知道父亲此举，公私掺杂。既有对其理念某种程度上的认可与对平衡局面的算计，也必然少不了对她这个叛逆女儿的担忧。

这份来自冰冷规则与家族隔阂之外的、别扭而隐晦的关切，让她心中某处坚冰，悄然融化了一丝裂缝。

尽管杯水车薪，但这股细微却持续不断的“输血”，实实在在地帮助山上撑过了最初、也是最艰难的消耗阶段。

它让更多重伤员在简陋的救护所有了生还的机会，让筋疲力竭的战斗修士在轮换的间隙，能多咽下一口蕴含灵气的食物，多炼化一颗丹药，从而多恢复一丝宝贵的、可能决定下一次生死交锋的元气。

这是在绝望的钢铁壁垒之中，来自冰冷战争规则外的一丝微弱却真实的暖意，它维系的不仅是生命，更是摇摇欲坠却始终不肯熄灭的希望之火。

\vspace{2em}

过了近一个月。

接连强行冲山的守旧联盟精锐的伤亡、不眠不休推算阵法演进的阵法师的努力，终于使得他们在九宫迷踪的第一重迷雾中，勉强撕开了几条断断续续的通道，总算爬上了半山腰。

然而，这里的迷雾固然变得稀薄，视野稍清，迎接他们的却也不是预想中的开阔地，而是更加赤裸、更加森严的死亡壁垒。

九宫迷踪阵第二重变：借阵法和预先建构的有利地形，迟滞敌人，增强自己。

地形已被改造得面目全非，山路崎岖破碎，迫使他们的百人级大型战阵彻底瓦解，化整为零。

空气中，弥漫着泥石、钢铁与隐隐血腥混合的气息。

就在这里，在这片被改造过的山林里，游击战的铁律开始用最残酷的方式，书写自己的注脚。

叶牧谦的命令非常简洁高效：

——一发子弹消灭一个敌人；

——打一枪换一个地方。

其余的，自由出击。

最初的战斗，并非所有人都能立刻适应。前一句话非常好理解，是告诫他们要节约灵力，应对持久化的战役；后一句话，是什么意思？

对“打一枪换一个地方”这句话不乏反感之人，特别是那些后来加入万仙盟、以勇猛凶悍著称的草莽散修。

一个叫石虎的灵枢境巅峰散修，带着他两个同样脾气火爆的结拜兄弟，被分配在吴刚负责的一片乱石区。

第一次遭遇敌情，是三人的巡哨小组发现了一小队五名紫阳宗斥候，正沿着一条干涸的溪床小心翼翼地上探。

石虎当时就红了眼。“才五个！老子一人就能劈了他们三个！弟兄们，上，拿了这头功！”

他不等队友完全响应，便如猛虎出闸般从掩体后跃出，手中鬼头刀抡圆了，带起一片腥风，直扑最近的那名斥候。

他的突袭确实迅猛，一刀便劈飞了那名斥候匆忙格挡的长剑，随即把他砍倒在地。他的两个兄弟也嗷嗷叫着冲了出去，凭个人勇武各自与剩下几个修士缠斗起来。

战斗起初很顺利，他们凭借个人武力，瞬间压制了对方。

然而，那名被重创的斥候在失去生命前，拼死捏碎了一枚示警符箓。尖锐的灵力啸音划破山林。

“大哥！他们叫了援兵！快走！”一名兄弟格开对手的剑，急声喊道。

“走什么走！还剩两个，宰了他们再走！”

杀红了眼的石虎看着眼前剩下的两个惊恐后退的斥候，觉得煮熟的鸭子不能飞了。

就是这短短几个呼吸的贪功恋战，要了他们的命。

溪床上方，一处看似天然、布满苔藓的岩壁后，突然转出整整一个十二人的紫阳宗快速反应小队，为首的是一名玄脉境开脉期的修士。冰冷的箭矢与灼热的火球瞬间覆盖了溪床区域。

石虎的怒吼戛然而止，他被三支破灵弩箭贯穿了胸膛。他那两个兄弟，一人被火球正面击中，化作焦炭；另一人想逃，却被那灵枢境修士御使的飞剑追上，从后心透入，钉死在溪床的石头上。

他们消灭了三名敌人，却赔上了自己三条命，还暴露了这片区域的防御存在。

消息传回，吴刚脸色铁青。

“我方四十七号哨位已经暴露，短期内不得前往该哨位！”

类似的情况，在最初的几天里并非孤例。

有的小组伏击得手后，觉得敌人慌乱溃逃，有机可乘，多追了几步，结果落入了敌方后续梯队的反伏击圈。

有的小组则固守着某个视野不错的狙击点，连续射杀了几名敌人后，不肯及时转移，很快便被锁定，招来覆盖性的集团法术轰击，连人带工事被炸上了天。

鲜血与瞬间消逝的生命，成为了最冷酷的教官。

叶牧谦说的确实不算数，但一条条活生生的性命，用他们的牺牲反复验证着叶牧谦那两条看似简单的铁律：

——轻易不出动，出动则至少消灭一个敌人；

——出手之后，无论有没有被人发现、有没有叫来援兵，立即换一个地方潜伏。

散修们开始变了。

适应环境，自然也是散修最常见的美德。

他们可能依旧不习惯这种偷偷摸摸、打了就跑的打法，觉得不够痛快，但他们更不想死。

现实的伤亡数字，逼着他们将质疑、将个人勇武的骄傲，死死压在心底。

\vspace{2em}

于是，真正的、高效的游击开始上演。

一个由一名原云隐别院剑修、一名善使符箓的散修和一名猎户出身的神射手组成的三人小组，已经在一处视野极佳、伪装巧妙的树冠狙击点潜伏了整整六个时辰。

他们下方，是一条被落石半掩的废弃小径。

黄昏时分，一队八人的守旧联盟巡逻队沿着小径走来，队形并不算紧密，神色疲惫。树冠上的神射手，将一支刻着破甲符文的铁箭稳稳搭在长弓上，向其中灌注了灵力，呼吸瞬间变得缓慢而深长。

他瞄准的不是为首那个看似指挥者，而是队伍中间那名唯一穿着丹鼎阁服饰、腰间挂着数个药囊的修士。

“咻——”

箭矢离弦的声音微不可闻，下一刻，那名丹鼎阁修士喉咙便被贯穿，哼都未哼一声便扑倒在地。

队伍瞬间大乱。

“敌袭！在树上！”有人指着箭矢来向大喊。

几乎同时，那名忙着画符的修士从另一侧早已看好的位置，扬手打出了三张符箓——两张爆炸符，一张烟雾符，炸伤了几个人，造成了更大的混乱。

“走！”剑修低喝。

三人没有丝毫犹豫，甚至不看战果。

神射手将长弓背好，抓住一根垂下的藤蔓，第一个滑向数十丈下方早就探好的一处岩缝。符修紧随其后。剑修最后一个离开，离开前，他用脚尖触发了一个旁边的机关。

当巡逻队从混乱中恢复，无大碍的五名修士愤怒地朝着原本的树冠位置施展术法、投掷符箓时，那里早已空无一人，只有燃烧的枝叶。

而他们试图追击时，却又触发了疯狂滋长的铁线荆棘，再次被拖延、刺伤。

另一侧。

在泥泞障碍区边缘，另一个小组演绎了另一种配合。

他们故意露出一个破绽，让一支试图绕行的敌方精锐小队发现了他们“仓促”撤退时留下的痕迹。

那支小队果然中计，兴奋地追入一片表面覆盖着落叶、看似坚固的低洼地。

结果，落入了一个连环陷阱。第一个踩上去的修士脚下的“地面”瞬间塌陷，露出一个不算深但布满黏稠淤泥和隐藏水刺的陷坑。

紧随其后的两人急忙想救，却被埋伏在侧翼的陈千羽甩了四根钢针，钢针深深刺入了他们的脚踝，挑断了其脚筋。两人脱力，立即向前栽倒，于是也掉进了泥坑里面。

与此同时，小组中那名远程支援者，冷静地用一枚低阶爆炸符在陷坑旁引爆，巨大的噪音和灵力冲击让陷坑里的人更加慌乱挣扎。

在泥潭中，挣扎反而会使得自己陷进去；只有冷静下来，才能慢慢上浮。

当然他们在泥潭中还受了皮外伤，那救回去八成也会感染——哦，感染在修仙界不算太大的事情，但也能消耗一些后勤。

剩下三人甚至没看清发生了什么事，只能站在坑边，试图拉起三个掉进坑里的家伙。

陈千羽的组员以迅雷不及掩耳之势把这三个家伙踹进泥坑之后，立即远遁，看都未多看那些在泥泞中惨叫的敌人一眼；借助对地形的熟悉，迅速消失在错综复杂的巨石和湿滑的坡地后。

等敌方的支援赶到，只看到六个在泥坑里狼狈不堪、失去战斗力的同袍，其中三个人已经半死不活，敌人早已无影无踪。

——一发子弹消灭一个敌人；

——打一枪换一个地方。

在一次次生死边缘的实践后，这两条律令被山上每一个幸存者用肌肉记忆刻入骨髓。

那些曾经对此不屑一顾的散修，在经历了最初的惨痛损失，又亲眼看到、亲身参与了这种高效而冷酷的杀伤后，心态发生了微妙而彻底的变化。

一次战斗间隙，几个浑身泥污血渍的散修挤在狭窄的掩体里分食着一块硬饼。

其中一个脸上带着新鲜刀疤的汉子，狠狠咬了一口饼，含糊不清地叹道：“娘的，以前觉得叶先生这打法忒不痛快，现在想想……真他娘的高！就这么一点点磨，看着那些紫阳宗的龟孙子明明人多，却像没头苍蝇似的挨揍，真解气！老子这条命，就是听了令，没贪刀，才捡回来的。”

旁边一个年纪稍大的散修，默默磨着自己因为砍人而略有卷刃的剑，幽幽接了一句：“何止是高明……老子活了几十年，打过不少架，从没见过这种打法。让你有力没处使，有火没处发。”

“叶先生……真神人也！”

最后那一声叹息般的“真神人也”，在狭小的空间里回荡，没有豪情万丈，却充满了一种劫后余生、心服口服的感慨。

鲜血与生存，永远是最有说服力的老师。

当那些不信任的、质疑的人都变成了尸体，而严格遵守这些“简单”规则的人却活了下来，并不断给强大敌人制造着伤亡和恐惧时，所有的怀疑都烟消云散，只剩下对制定规则者近乎敬畏的信服。

白言冰、陈千羽、沈瑶、吴东浩、吴刚等核心人物，便如同定海神针，钉在各处要冲，他们自身就是这种战术最坚定的执行者和示范者。

而整条战线，就在这无数个精干如匕首的小组神出鬼没的袭扰下，变成了一架高效而沉默的杀戮机器。

守旧联盟空有数倍兵力，却感觉自己像是在用沉重的铁锤，愤怒而徒劳地击打着无形无影的流沙和铁蒺藜。

他们磅礴的力量被主动分散、被巧妙引导，每一次蓄力的重击，要么打在空处，要么只能换来零星却绝无间断的精准回刺。

看着同袍昨日还鲜活的面孔，今日便成为一具被拖回的冰冷尸体，或者干脆失踪于山林再也寻不回，一股深入骨髓的无力与挫败感，如同这山间弥漫的、无法驱散的湿冷寒气，悄然侵蚀着攻城大军原本高昂的士气。

\vspace{2em}

山下中军大帐内，陈炎接到前线接连传来的受挫战报与日益增长的伤亡数字，脸色阴沉得能滴出水来。 他意识到，单靠中低阶修士组成的常规军队推进，恐怕真要在这座被对手经营得如同铁刺猬般的山上，撞得头破血流，甚至损兵折将、威望扫地。

速胜的耐心被这座顽山消耗殆尽，他眼中寒光一闪，决定不再保留，动用高端力量强行破局，以绝对的实力碾压一切诡计。

“命令周文彬，率本部真符境精锐，自东侧山脊强行御空飞渡，直扑别院核心区域！”陈炎的声音如同金铁交击，冰冷而不容置疑，“告诉他们，不惜代价，给我找出并摧毁其阵法核心节点！”

同时，他转向身旁侍立的随军高层，厉声指示：“让丹鼎阁的高阶药师与紫阳宗的阵法师即刻配合，启动‘焚天’与‘瘴海’预案！我要这半山腰以上，再无一片完整的树叶，再无一处可资利用的地形！”

即便不能直接杀伤那些隐藏极深的老鼠，他也要用覆盖性的烈焰与毒瘴，将山体烧成焦土，将阵基腐蚀瘫痪！

命令下达，山下营地骤然爆发出数道磅礴凶悍的气息，如同沉睡的凶兽苏醒。

以周文彬为首，总计七八名真符境高手周身灵光暴涨，化作数道刺目流光，如逆射的陨星般拔地而起，悍然刺破云雾，直扑山腰以上！

真符境修士已能较长时间御空飞行，足以规避绝大部分恼人的地面陷阱。

在他们身后，数十名玄脉境融贯期的好手也各展手段，或借助法器，或彼此借力，紧紧跟随，意图凭借这石破天惊的一击，从空中强行撕开一条通往胜利的血路。

与此同时，山下数个早已筑好的高大法坛骤然亮起刺目欲盲的光芒，狂暴暴虐的火行灵力与诡异粘稠的墨绿色毒雾开始疯狂汇聚，规模之浩大，灵压之骇人，让远处山巅观看的守卫者们都不由得面色发白，仿佛看到了天灾降临。

然而，中枢阵眼之内，始终如古松般盘坐的叶牧谦，感受到那自天空与山下同时传来的、足以令寻常见独境修士心神摇曳的恐怖灵力波动，缓缓睁开的双眼之中，却掠过一丝“终于来了”的锐利精光。 他似乎，等待的正是对手按捺不住、祭出高端力量的这一刻。

他双手平稳而坚定地虚按在身前——那里是一个由无数道细微光线精准勾勒出的、惟妙惟肖的云霞山脉灵气流动虚影。

磅礴的精神力与他自己独有的、能与周遭天地产生深层共鸣（至于为什么只有他能引起共鸣，叶牧谦并未能深究明白，于是也不去想它了）的精纯元气，如同决堤洪流，奔涌而出，透过身下与整座大山相连的阵法核心脉络，瞬间沟通、唤醒、引导了整座云霞山沉寂却浩瀚的庞大地气。

“地脉如龙，听我号令……起！”一声清喝，虽不高昂，但仿佛蕴含着律令山川的威严。

轰——隆——

九宫迷踪第三重变！

叶牧谦毕竟出身于现代社会，哪里会不知道防空的重要性！

刹那间，整座云霞山主峰，从山基到山巅，似乎被注入了生命，发出一阵低沉雄浑、仿佛源自大地深处的嗡鸣与震颤。

以叶牧谦所在中枢为原点，一层肉眼难以清晰捕捉、却给人以无比坚韧厚重之感的淡金色灵力屏障，如同一个巨大无比、倒扣而下的琉璃巨碗，伴随着嗡鸣声迅速向上蔓延、合拢，最终将山腰以上的核心防御区牢牢笼罩在内！

这屏障中蕴含着叶牧谦对阵道导引、分流、化解的玄妙领悟。

数位真符境长老联手轰出的、足以融化金铁的赤红咆哮火柱，与撕裂空气的凌厉无匹剑气长龙，几乎同时狠狠撞击在淡金色的屏障光幕之上，但预想中的惊天爆炸与屏障破碎并未发生。

那声势骇人的攻击，如同泥牛入海，其大半毁灭性能量被那层流转着奇异韵律的屏障表面巧妙地偏转、卸开、引导。

炽热的火流被引偏斜射向高空无害消散，锋锐的剑气被分化成数百道细流，汇入山体周边呼啸的狂风。余下的威力，则被屏障均匀分散，传导至下方绵延的山体基石，被厚重的大地无声吸纳。

山下法坛蓄势完毕，释放出的、如同赤红浪潮般的烈焰风暴与遮天蔽日的墨绿色毒瘴之云，在触及这层淡金色屏障时，也同样被大幅削弱、稀释、阻隔，未能对屏障后的山体植被与工事造成预期中焦土千里的毁灭效果。

几乎就在淡金色屏障光华大盛、成功抵御住第一波毁灭性打击的同一刹那！

东侧一处陡峭制高点的鹰嘴岩上，沈瑶的身影如鬼魅般悄然浮现。山风猎猎，吹拂着她束起的长发与紧身的玄色劲装，勾勒出凛然如刀的身形。她手中那柄弧度优美的短刀正在微微轻颤，发出几乎听不见却直抵灵魂的细微嗡鸣。

面对数名试图趁乱从侧翼低空掠过、寻找阵法屏障薄弱点或能量节点的敌人，她不退不避，架势倏然展开——

“千幻流光”！

铮——！

刹那间，以沈瑶所处岩石为中心，方圆数十丈内的雾霭、林间疏漏的光斑、甚至屏障流转的淡金微光，仿佛都被强行抽取、熔炼、再造！

数十上百道虚实交织、真假难辨、绚丽夺目却又杀机凛然的璀璨刀光，凭空浮现！

这些刀光如同被赋予了独立的生命与狩猎本能，化身为一条条灵动无比、轨迹完全无法预测的光之游鱼、疾电或幽魂，时而如仲夏夜的流星暴雨，毫无规律地倾泻覆盖一片区域，逼得敌人手忙脚乱地防御；时而又骤然汇聚，凝成一道凝实到极点的深寒刀芒，如同毒蛇吐信，精准刺向因攻击受挫而出现刹那松懈或灵力衔接空白的敌人要害。

真幻交织，虚实互变，令人心神俱疲。

一名紫阳宗玄脉境巅峰的执事，怒喝着挥动火焰盾牌格开了迎面三道看似凶悍的刀光，却冷不防被一道从脚下阴影中悄然钻出、看似虚幻的刀光掠过小腿，护体灵光竟如纸糊般被割开，血光迸现！

沈瑶一人一刀，凭借对地形的极致利用、阵法屏障的掩护以及自身诡谲莫测的领域，竟如一道无形而坚韧的壁垒，生生牵制、迟滞住了数倍于己的敌方高手，令他们无法专心寻找破阵之法，反而不得不分出大量精力应对这无所不在、防不胜防的袭杀。

云霞山，这块被守旧联盟上下视为大军一到便可唾手可得的“骨头”，出乎所有人意料地，在鲜血与火焰的洗礼中，彻底显露出其内里——这根本就是一块坚硬无比、布满倒刺、并且会主动反击的铁桶！

陈炎后续又陆续调集生力军，不惜代价发动数次规模更大的猛攻，前后历时月余，大小血腥接触战数十场，自身伤亡累积已达到一个令中军帐内气氛凝滞的惊人数字，却始终无法突破那最后一道、仿佛拥有生命的核心防线。

第一重变目前已经被全部瓦解，但那毕竟是无人组成、无人参与的；第二重变和第三重变，山上万众一心，此刻根本破不开！

更要命的是，就算他们试图放火烧山，陈千羽的和畅领域一来，草木一夜之间便又长出来了！

这可真真正正是“野火烧不尽，春风吹又生”！

浩浩荡荡的联盟大军，被牢牢拖在了云霞山陡峭的山麓与山腰，进退维谷，骑虎难下。

山下，速战速决、一鼓定乾坤的意图彻底落空，战争无可避免地滑向了他们起初并不愿意但也并不十分介意的消耗泥潭。

陈炎此刻也有些泄劲了。

他的战略也变了：围而不打。

我给你围起来，让你消耗。

你们总不能不吃饭吧？

沈天穹的输血虽然存在，但陈炎冷酷的计算之下，显然不够山上消耗的！

但他显然低估了叶牧谦等人的韧劲：叶牧谦主持下，冷战期间就在山上开凿了水井、开辟了药田、灵田，辟谷丹（粮食）、基础草药均基本自足；天师府又有输血，虽然不多，但蚊子腿也是肉。

要支撑这么多人，勒紧裤腰带一段时间也并非不可。

但能勒紧一时，还能勒紧一世吗？

山上，尽管防线在不少将士以生命为代价的捍卫下暂时稳如磐石，但每个幸存者都能清晰地感受到，那来自山下无穷无尽压力的迫在眉睫。

改变，无论是对于渴望破局的守旧联盟，还是对于必须寻找生路的“问道途”而言，都已迫在眉睫。但至少，死守云霞山的战略目的，已经完成了一部分。

单纯的防御，无论阵法多么精妙，战术多么灵活，终究是在消耗自身本就有限的“血液”。每一块灵石的耗尽，每一瓶丹药的见底，每一处工事的破损，都在削弱这座孤山抵抗的资本。

\vspace{2em}

叶牧谦作为主帅，比任何人都更清醒地认识这一点。

就在山下攻势最为猛烈、山上承受压力仿佛达到顶点的一个深夜，他独自立于别院后方残破的观星台遗址。头顶的夜空，被己方阵法流转的淡金灵光与远方敌军营地连绵的篝火、法坛光芒映照得一片混沌，不见星辰。

然而，前世那些波澜壮阔的记忆碎片，却如同穿透了无尽时空的璀璨星辰，清晰无比地倒映、串联起来——关于在强大敌人围困下依然屹立的红色根据地，关于以弱势兵力周旋并打破一次次“围剿”的惊人战术，关于深入敌人统治腹地、以点点星火最终燎原的传奇。

“坚守要点，以空间换取时间……发动群众……陷敌于人民战争的汪洋大海……”

他低沉的自语声消散在夜风里，轻不可闻，但那双望向山下广袤黑暗地域的眼眸深处，却骤然闪过一丝足以劈开迷雾的决然锐光。

僵局必须打破，而破局的关键力量，不在山上这有限的不到一万人中，恰恰在那山下——那片被守旧联盟视为统治基石，却又在实际上被他们视为可以随意榨取之草芥的、广袤而沉默的底层疆域之中。

\vspace{2em}

翌日，一次参与者极少、保密级别最高的核心会议，在云隐别院深处一座被层层阵法灵光彻底隔绝的静室内召开。空气中弥漫着聚灵阵运转的低微嗡鸣与紧张的气氛。

叶牧谦在一块平整的青石板上，以元气光线勾勒出一幅简陋却关键的内域势力分布草图。他的手指，稳稳地点在了代表着云霞山外围，那一片片连绵的、标注着中小型城镇、散修聚集点和凡人村落符号的区域。

“固守云霞，是为中流砥柱，不可有丝毫动摇。”叶牧谦的声音平稳而有力，仿佛在陈述一个天地至理，“这里将会，并且在未来一段时间内，一直吸引敌军主力。它更是整个‘问道途’的主心骨、精神旗帜，是我们的基本盘，不能放弃。”

他的目光转向侍立一旁的两位大弟子，“白言冰，陈千羽，你二人随我继续坐镇中枢，白言冰统筹防务，临机决断；陈千羽保障救治，维系生机。山在，人在，旗就在。”

白言冰与陈千羽没有丝毫犹豫，同时抱拳，沉声应诺。白言冰眼中是剑锋般的坚定，陈千羽眸子里则是柔韧如青藤的沉静。

然而，问题如同阴云悬在每个人心头。如此被动固守，与敌人拼消耗，哪里能看到胜利的曙光？革新派一方的后勤储备与资源底蕴，远远无法与树大根深的守旧联盟相比，拖延下去，最终的结局几乎是注定的慢性失血而亡。这不仅是所有人的疑问，也正是叶牧谦近期全部思考所聚焦的核心。

“真正的决胜之机，不在山上，”叶牧谦的手指从云霞山图标移开，稳稳地划过草图，落在了那些标志着守旧联盟腹地、但其直属统治力量相对鞭长莫及或较为薄弱的区域，“而在于此处。”

他的目光抬起，落在另外两人身上：“沈瑶，吴刚。”

被点名的两人身躯瞬间绷直。沈瑶清丽的面容上，眸光冰澈凛冽，如同雪原上即将出鞘的剑；吴刚则收敛了平日里那副万事皆可调侃的玩世不恭，粗犷的脸上流露出罕见的郑重与肃穆。

“你们二人，各自挑选一支绝对可靠、精干灵活的小队，人数贵精不贵多。” 叶牧谦的手指关节轻轻敲击着青石板上的那些区域，“目标，不要强行攻打任何重镇要塞，而是潜入这些地方。”

他的指尖重重一顿，仿佛要将决心钉入石板：“那里，有被宗门与家族层层加码的苛捐杂税压得脊梁弯曲、喘不过气的散修；有面朝黄土背朝天，辛勤劳作一年所得大半被盘剥、仅能勉强糊口乃至挣扎在饿死边缘的凡人百姓；还有无数因出身、机缘或无背景，而对现状满怀愤懑却又敢怒不敢言的失意者。”

叶牧谦的语调，带着一种洞悉本质的寒意与热度，“他们是守旧联盟维持统治、汲取养分的根基土壤，但也恰恰可能是这块看似坚固巨石上，最易松动、最先崩裂的一块砖。”

他详细阐述了行动的方略核心：依托对云霞山周边地形地貌的绝对熟悉，借助“九宫迷踪阵”尚未被敌军完全勘破的变幻特性作为掩护，精确计算敌军巡逻队伍的间隙，选择夜色最深浓、戒备最易松懈的时刻，如同水银泻地般悄然渗透出去。

不求一次歼灭多少敌人，不求破坏多少后勤运输，更不贪图占领一城一地。

唯一的核心目的，是将“问道途”追求公道、反抗不公的理念，将那“人人皆可凭自身努力追寻大道、改善生计”的希望火种，尽可能广泛地播撒到那些沉闷而压抑的土地中去。

“具体如何行事，你们可根据实际情况临机决断。”叶牧谦的目光如同实质，扫过沈瑶和吴刚，“但必须牢记我们的根本行动纲领：发动群众，建设组织，积蓄力量。”

他格外强调，“其中，‘发动群众’是重中之重，是源头活水！这场战争注定旷日持久，从长远来看，无数凡人百姓与低阶修士，既是我们的根基所在，也是他们自身改变命运的希望所系！”

“我们能做且必须去做的，”叶牧谦越说越兴奋，语调也越来越激昂，“就是让更多的平民百姓、低阶修士，从心底里理解我们、认可我们、最终愿意支持并追随我们！”

他的视线最终牢牢锁定沈瑶与吴刚，字句如锤：“我要你们，成为插入敌人看似稳固之后方的两把利剑！一把，斩断其汲取民脂民膏、输往前线的血管；另一把，点燃其统治根基之下，那早已遍布干柴的星星之火！”

这绝非一时冲动的冒险，而是叶牧谦结合前世智慧与此世现实，经过深思熟虑后得出的破局之策，得益于他初入内域时，在青石城那长达一年的生活与细致观察。

那一年，让他明确了一个看似简单却至关重要的道理：这个世界虽存在移山倒海的超凡个体，但在大规模战争状态下，高阶修士很大程度上被“兵器化”、脱产（甚至平常情况下也大多是脱产的），而维持这支“兵器”运转的庞杂后勤——从粮秣布匹到基础矿材，从简单工事修筑到短途物资转运——其真正的主力，恰恰是数量庞大的凡人与低阶修士。

也正因洞悉了这一战争维度的本质，他才敢于借鉴并化用那经过历史检验的智慧，坚定地走向这条发动最广大底层力量的群众路线，开辟决定胜负的第二战场。

行动计划迅速在高度保密中敲定。

大多数人可能并不理解为什么修士打架要去管平民百姓死活。但沈瑶有着长达十年的草根经历，吴刚也在万仙盟做过很多低阶散修的工作，所以关于叶牧谦所言“发动群众”，他们是有经验的，也是发自内心认同的。

沈瑶精心挑选了十人，其中多为原无极剑派心剑脉中剑法以轻灵迅捷见长、尤其擅长隐匿潜伏的弟子，再搭配数名在加入万仙盟前便是独行游侠、极擅野外生存、陷阱布置与情报侦查的散修老手。

吴刚的队伍则更侧重综合能力与适应性，既有万仙盟中熟悉各地黑市门路、人情往来与社会三教九流的“路路通”，也有像陈默这样早期转修问道途、能作为鲜活榜样吸引同类的修士，还包括两名思维敏捷、口齿伶俐、善于在不同场合进行说服与鼓动的特殊人才。

月黑风高，正是潜龙腾渊时。

在云霞山防御部队又一次浴血击退敌军组织的夜袭、战场迎来短暂死寂的间隙，两条比夜色更加深邃、更加沉默的“细流”，借助山间尚未散尽的雾气与复杂崎岖地形的完美掩护，从守旧联盟包围圈两个因地形限制而形成的、极为细微且不易察觉的结合部缝隙中，悄然滑了出来。

沈瑶一袭不起眼的深蓝近黑劲装，将姣好的身形完全融入阴影。她甚至未曾全力催动见独领域，仅仅让其维持在周身尺许，完美地扭曲光线、吸纳气息。

她如同夜幕中飘忽的领路幽魂，身后两支小队成员个个屏息凝神，将一切行动声响降至最低，如同真正的鬼魅般紧跟着她，穿行在林木与陡崖的阴影之间。很快，那旌旗如林、灵光冲天的连绵营寨被远远甩在身后，模糊成一片背景光晕。

他们的面前，是沉沉睡去的广袤山林，以及山林之外，守旧联盟统治下那些光影交织、苦难与麻木并存的灰色地带。

两队人马在一个早已标记好的三岔口悄无声息地分开，互递一个坚定的眼神后，便向着不同方向，隐没于更深沉的黑暗之中。

\vspace{2em}

沈瑶带领的小队，彻底化为了一柄经过千锤百炼、敛尽所有华光的漆黑匕首，沉默、精准、致命地刺向敌人后方。他们昼伏夜出，完全避开官道与任何稍具规模的城池，专挑那些小路行进。

目标极其明确：那些由守旧联盟附属中小宗门或地方家族控制、驻防力量有限、看似不起眼，却实际承担着为前线大军输送物资、同时层层盘剥底层以供自身享乐的城镇节点。

七日后，经过周密迂回，他们抵达了第一个目标——位于紫阳宗势力范围边缘地带的灰谷镇。此地因出产一种用于低阶法器淬火增韧的辅料“灰纹铁”而得名，镇民结构简单，多为在矿洞中讨生活的矿工、往来运输的贩夫走卒，以及少数依附这条产业链生存、修为低微的散修。

镇守此地的，是紫阳宗某外门长老的一个姻亲家族，修为最高者不过玄脉境中期，但凭借后台，在镇内作威作福，苛捐杂税名目繁多，征敛无度，是标准过头的土皇帝。

沈瑶没有丝毫急躁。她亲自带领两名最擅长潜行侦查的队员，如同一缕青烟般在镇内外潜伏、观察了整整三天，摸清了守卫懒散的巡逻规律、镇内几处关键仓库与货栈的具体位置与守备情况，以及镇上矿工、散修们主要聚集闲聊的几个茶摊、酒肆。

她清晰地看到，因为战争，紫阳宗对灰纹铁的收购价格被压得更低，而经由联盟商会控制的粮食、盐、粗布等生活必需品的价格却飞涨。镇民们脸上菜色更重，眼神麻木中压抑着怒火，私下怨声载道，但面对镇守家族那几个嚣张跋扈的护卫和背后的紫阳宗，无人敢带头反抗。

行动日，选定在镇上一次旬日小集的日子。正午时分，狭长的集市街道略显拥挤，充斥着讨价还价声、牲畜叫唤和粗重的咳嗽声。沈瑶命令所有队员——也包括她自己在内——换上与当地矿工、农夫无异的粗布短打，甚至故意蹭上些尘土泥污。

大部分队员化整为零，悄然散布在集市外围关键路口，如鹰隼般警戒。她自己则仅带两名最得力的队员，步履自然地穿过人群，径直走向位于集市中段、那间由镇守家族开设、兼做着收税与强买强卖勾当的最大粮铺。

粮铺掌柜是个脑满肠肥的中年人，正尖声呵斥着一个蜷缩在柜台前、颤抖着掏不出几块下品灵石购买全家半月口粮的老矿工，唾沫星子几乎喷到老人脸上。忽然，掌柜浑身肥肉一僵，感到脖颈侧面传来一阵深入骨髓的冰凉，一柄弧度优美、却散发着森然死意的幽暗短刀，已无声无息地贴在了他的皮肤上。

“所有人，不许动，原地蹲下！”

沈瑶清冽的声音并不高昂，却奇异地压过了集市所有的嘈杂，清晰地钻入每个人耳中。与此同时，她身上那股修士特有的、凝练如精钢的气息微微释放了一瞬。

如同寒风掠过，顿时让粮铺内几名仅有灵枢境修为、平日里狗仗人势的护卫双腿发软，脸色煞白，僵在原地动弹不得。

集市上的人群先是一静，随即爆发出本能的骚动与惊呼，但在外围队员冰冷而有序的维持手势，以及沈瑶身上那似有若无、却足以让普通人灵魂战栗的凛冽气场笼罩下，骚动迅速被压制下去，化为一种死寂般的、充满恐惧与惊疑的安静。

沈瑶看也不看那抖若筛糠、面如土色的掌柜，她的目光平静地扫过集市上那一张张或麻木、或惶恐、或畏缩、又或在不经意间流露出一丝好奇与期盼的脸庞——那是被生活重压碾磨过的底层百姓最真实的面容。

她深吸一口气，朗声开口，声音在元气的包裹下传遍半个集市：

“我乃云隐别院，叶牧谦叶先生座下弟子，沈瑶。今日来此，不为杀戮，不为劫掠，只为替各位父老乡亲，取回本就属于你们的活命之物！”

话音未落，她手腕看似随意地一抖，一道凝练的元气刀光掠过，粮铺柜台上那本记录着巧取豪夺与高利盘剥的厚重账册，应声被挑飞至半空，纸张哗啦散开，如同片片控诉的雪花。

几乎在同一时刻，她身旁一名队员如同蓄势已久的猎豹般窜出，身法快得只在众人眼中留下一道残影，几道精准如尺量般的剑气“叮叮”数声轻响，粮铺通往后院仓库那扇厚重木门上的精铁大锁应声断裂。

第三名队员则早已如同鬼魅般贴近了闻讯从隔壁酒楼仓皇冲出的几名镇守家族护卫，简单利落的几记手法，便让他们瘫软倒地，暂时失去了行动能力。

“去，打开仓库！将里面囤积克扣的粮食、盐、布匹，全部搬出来，就在这集市空地上，按户按人头，公平分发！”沈瑶的命令清晰果断。

她的队员们毫不犹豫地执行。更令人动容的是，那几个原本被掌柜呵斥、蹲在墙角的老矿工，此刻眼中骤然爆发出一种近乎凶狠的光芒，那是对生存最本能的渴望被点燃的火焰！他们低吼一声，仿佛瞬间年轻了十岁，跟着问道途的队员们一起，冲进了那座平日里他们连靠近都不敢的仓库。

当第一袋沉甸甸的、散发着谷物清香的麻袋被扛出来，随后是第二袋、第三袋……当一匹匹虽然粗糙却厚实的土布、一罐罐洁白如雪的盐巴被陆续搬出，在集市空地上逐渐堆积成一座小山时，死寂的人群中终于爆发出了难以置信的、混杂着狂喜与哽咽的低呼与骚动。

许多人的眼睛死死盯住那些物资，眼眶迅速泛红。

沈瑶知道，火候已到。她的声音再次响起，依旧清冽，却注入了一种铿锵的、直指人心的力量：

“乡亲们看好了！这些粮食、这些布匹、这些盐，哪一样不是你们的汗水换来？哪一样不是守旧联盟那些高高在上的老爷们，从你们嘴里夺去、从你们身上剥下，用以填满他们自己的仓库，供养前线那些用来继续压迫更多像你们一样百姓的军队？！”

她的话语如同匕首，剖开了血淋淋的现实，“他们何曾将你们当人看待？在他們眼中，尔等不过是可以随意驱使、榨取、丢弃的牛马与蝼蚁！”

她停顿下来，目光如电，扫过人群。她看到，许多双原本麻木的眼睛里，开始燃起压抑已久的火焰；许多紧抿的嘴唇在颤抖；一些年轻人捏紧了拳头，指节发白。那是积郁多年的悲苦与愤懑，被这番话说穿、点燃后的共鸣。

“我师叶牧谦，观天地不公，悯众生之苦，另辟‘问道途’！”沈瑶的声音陡然拔高，带着一种宣告般的穿透力，“此道，不假外求，不拜神佛，只修己身！无需依赖大宗门垄断的昂贵灵石、稀缺丹药、天价法器，凡人亦可凭自身毅力，引气入体，踏上修行之路！即便是资质寻常的低阶修士，修习问道途后，其劳作之能、体魄之健、乃至心志之坚，都远非寻常灵枢径修士可比！这意味着什么？”

她自问自答，话语斩钉截铁：“这意味着，跟着我们，你们或许无法立刻飞天遁地，但你们能凭自己的双手，开垦更多荒地，收获更多粮食；能凭自己的力气，从事更繁重的劳作，换取更安稳的生活；能让你们的子女，有机会不花一个冤枉钱，就接触到真正的修行法门，拥有改变命运的可能！而不是世世代代，被锁死在矿洞、田垄和最低贱的活计上，永无出头之日！”

这番话语，没有任何虚无缥缈的大道理，而是用最朴素、最直接的方式，描绘了一个触手可及的、关于“更好生活”的图景。

对于这些挣扎在生存线上、对高阶修士世界如同仰望星空的底层百姓而言，什么长生久视、什么大道法则都太过遥远；唯有实实在在的粮食、衣物、盐巴，以及让下一代人过得比自己稍好一点的希望，才能真正撼动他们坚固如铁石的心防。

“守旧联盟为何要不惜代价，围困云霞山，打压我等？”沈瑶的语调转为凌厉的质问，“正是因为他们恐惧！恐惧一旦人人皆可凭自身努力修行、改善生活，无需再仰仗他们的鼻息、购买他们垄断的资粮，他们那作威作福、骑在所有人头上吸血的特权，就将如同烈日下的冰雪，彻底崩塌消散！”

“我们一退再退，他们却步步紧逼，甚至不惜设下埋伏，屠戮我同道！试问，我‘问道途’求一个公道，何罪之有？你们这些终日劳作不得温饱的百姓，又何罪之有，要承受这般无休止的盘剥与轻贱？！”

此时，物资分发已经在队员和一些胆大起来的镇民主持下，开始有序进行。

当第一个头发花白的老妇人用颤抖的双手接过一小袋粮食和一小块盐巴时，浑浊的泪水瞬间滚落；当一个面黄肌瘦的孩子紧紧抱着一块粗布，脸上露出罕见的、属于孩童的笑容时……整个集市的气氛变了。

先前对沈瑶这位“修士”的恐惧与隔阂，在她身穿粗布衣、口吐百姓言、行事为百姓夺回生存之资的对比下，迅速转化为一种难以言喻的信任与激动。

那是一种最质朴的认知：这个“仙子”和以前见过的、那些高高在上、冷漠索取的人们，不一样。

沈瑶敏锐地捕捉到了这种变化。她向人群中几个一直紧紧盯着她、眼中除了激动还有某种更深层渴望的年轻散修招了招手。

这几人修为低微，多在灵枢初、中期徘徊，衣衫陈旧，显然是镇中最底层的那批散修。她接过队员递来的几本纸张粗糙、装订简易的薄册子。

“此乃我师亲撰、简化优化的问道途《养气篇》最基础入门法诀。”沈瑶的声音缓和了些，带着一种郑重的意味，“文字浅白，路径清晰，于资源要求降至极低，唯重心性与毅力。此乃正道之基，虽简，却直指修身自强之本意。有意者，可自取观之、习之。能否入门，踏上此途，全在个人抉择与坚持。”

她将册子亲手递给那几名眼神灼热的年轻散修，又额外留下几份，放在分发物资的木台上。

“我辈修士，生而在世，当有自强不息、不仰人鼻息之志！而非浑浑噩噩，任人鱼肉宰割！”

一位衣着最破旧、手上满是矿洞污渍却洗得发白的年轻散修，几乎是颤抖着接过册子，急不可耐地翻开第一页。扉页之上，并非什么高深功法总纲，只有十个大字：

“仰不愧于天，俯不怍于人。”

年轻散修浑身一震，呆呆地看着这十个字，眼眶瞬间通红。他身边另一个同伴凑过来看，也是身躯微颤，低声喃喃：“不愧本心……当今之世，能做到的，有几人？”

修士与凡人百姓的关注点，在此刻微妙地分化开来。

大多数百姓接过救命的粮食布匹，感激涕零，心中记住的是“叶先生”、“沈仙子”给了他们活路，对那十个字或许感触不深，只觉得是“好人该有的样子”。

而这些挣扎在修行底层、看多了蝇营狗苟、自身也常因资源匮乏而不得不做出些违心之举的低阶散修们，对这十个字的共鸣却如同惊雷，直击灵魂深处。

它触碰的，是他们身为修士，却活得不如意、甚至时常自我质疑的那份尊严与骄傲。

整个分发过程高效而迅速，持续了不到一个时辰。

在镇守家族可能的主力修士闻讯从烟花之地或修炼静室匆忙赶来之前，沈瑶小队早已带着缴获的部分关键账册——作为守旧联盟盘剥的罪证——和少量便于携带的灵石，如同他们出现时一样，悄然撤离，瞬息间便融入镇外山林，消失得无影无踪。

留给灰谷镇的，不仅是家家户户或多或少分到的、足以支撑一段时日的活命物资；不仅是那几本在暗地里被无数人渴望、偷偷传抄临摹的《养气篇》入门册子；更重要的，是“沈瑶”“叶牧谦”“问道途”“仰不愧于天，俯不怍于人”这些名字、话语和理念，如同巨石终于投入了这潭被压抑已久的死水，激荡起层层叠叠、再难平息的思想涟漪与渴望的波澜。

关于云霞山的故事，关于另一种活法的可能，开始在这个小镇的街头巷尾、矿洞内外，以各种被添油加醋却核心不变的版本，秘密而顽强地流传开来。

此后月余时间里，沈瑶带领这支精干小队，如同技艺高超的弈者，在守旧联盟广袤的统治棋盘上，落下了一枚枚看似微小、却位置关键的棋子。

他们如同飘忽不定、却总能精准命中要害的幽灵，连续“光顾”了四五个与灰谷镇情况类似的城镇据点，有时会选择在当地守旧势力头目正在大摆筵席、醉生梦死时突然现身，以雷霆手段控制全场，在众目睽睽之下揭露其奢靡与底层困苦的对比；有时则会暗中联系当地一些早已对压迫心怀不满的小型散修团体或受排挤的家族旁支，里应外合，事半功倍；在一次特别行动中，他们甚至成功伏击了一支押送队伍，解救了一群被守旧联盟强行征发、准备送往最危险的前线地段充当肉盾“炮灰”的散修苦力，并将他们妥善安置、吸纳……

随着这些行动如水波般扩散，“问道途”以及其核心人物叶牧谦的形象，在广大的底层区域，开始以一种复杂而多元的方式被认知、被谈论。

对于绝大多数挣扎在温饱线上的凡人百姓而言，“叶先生”“云隐别院”这些词汇，最初或许模糊而遥远。但经由沈瑶等人一次次实实在在分发粮食、布匹、盐巴的行动，“问道途”逐渐与“能让我们吃上饭”、“不欺负我们”、“替我们出头”这些最朴素、最根本的利益诉求画上了等号。

他们或许不懂“不假外求”的深奥，也难以理解“仰不愧天”的境界，但他们心中有一杆最朴素的秤：谁对他们好，他们就记着谁的好。叶牧谦的形象，在他们口中传颂时，常常被赋予一些朴素的想象，有时是“下凡救苦的神仙”，有时是“法力无边、专打坏蛋的大侠”，有时甚至就是“一个特别厉害、心肠好的读书人”。

这种认知虽然不够准确，甚至有些神化或简单化，却无比真实地反映了底层民众最直接的感激与期望。

所有平民百姓，实际上都渴望一个能带来公平、让他们活下去的救星或青天。

而对于数量稍小、处境稍好但也极其有限的低阶散修，以及部分出身寒微、在守旧联盟体系内不得志的低阶修士而言，他们对叶牧谦和“问道途”的看法则要复杂、深刻得多。沈瑶留下的那本《养气篇》册子扉页上的十个字，如同黑暗中的一盏孤灯，照亮了许多人内心不曾明言的委屈与坚守。

“问道途”宣称的“无需依赖垄断资源”“仅凭自身努力亦可修行进阶”，更是切中了他们修行路上最大的痛点——资源匮乏与上升无门。

叶牧谦的形象，在他们心中更接近于一个“开宗立派、为吾辈寒微者开辟新路的大贤先师”。

他不仅提供了一种可能更公平、门槛更低的修行路径，更重要的是，他提供了一种精神上的认同与尊严。他们讨论叶牧谦时，会探讨他那奇特的“养气”“见独”境界究竟有何奥妙，会争论“问道途”与“灵枢径”在根本理念上的差异，会羡慕白言冰、沈瑶等人能追随这样的师长。

叶牧谦之于他们，既是现实摆脱困境的希望，也是一种理想人格的投射。

\vspace{2em}

吴刚那面采取的行动，和沈瑶则大相径庭。

如果说沈瑶行动在明，如同撕裂黑夜的雷霆与烈火，以直接而凌厉的方式发动群众，凭借夺回物资、公开宣讲，在底层百姓心中迅速点燃对守旧联盟的憎恶与对“问道途”最朴素的认同；那么，在另一条更为隐蔽的战线上，吴刚所执行的，则是战略的另一个精妙侧面：团结一切可以团结的力量。

这显然是基于两人截然不同的经历与资源。

沈瑶十年游历，深谙草根疾苦与民间生态；而吴刚的优势，则根植于万仙盟那张早在冷战时期就已存在、并在之前经济封锁中历经考验而变得越发坚韧细密的地下网络，以及极其善于揣摩人心、权衡利弊、在复杂利益格局中找到平衡点与共鸣处的特殊才能。

虽然说先前那个网络在守旧联盟巨象踏步一般的前期推进之中，不少敢于露头的地下网络联络员都被抓到，面临监禁，甚至当场牺牲；但依然有着少部分原地下网络联络员依然保存了自己的上下游联络方式、情报和性命。

不过已经建立过一次地下网络，在原网络的废墟之上重建一个网络，那难度相形之下会稍低些。

吴刚选择的切入点，自然也是散修比例更高、商业活动更活跃、各方势力交错、情况也更为微妙的区域——情况越混乱，乡土气息越浓厚，那么守旧联盟的触角就更难管辖到哪些地方，重建网络、团结地方的操作自然也相对容易一些。

他率领小队潜出重围后，抵达的第一个重要节点，便是这样一个地方——名为“黑石集”的大型散修贸易据点。

\vspace{2em}

黑石集坐落在一片盛产低阶炼器矿石与药材的山脉出口。建筑杂乱无章，街道上弥漫着矿石粉尘、药材腥气与汗臭混合的味道。

这里名义上受一个依附于丹鼎阁的中等势力管辖，但天高皇帝远；真正的秩序维持与利益分配，很大程度上依赖于散修们自发形成的潜规则，以及几个在此地盘踞多年、彼此制衡的本土小帮派。这里是信息的集散地，也是资源的流转池，更是无数散修讨生活、碰运气的江湖。

吴刚如同一条回到熟悉水域的游鱼，通过队伍里的本地散修，只用了不到两天，便在一家充斥着劣质酒气与粗豪喧哗的嘈杂酒馆后间，与黑石集本地一位颇有威望、人称“老黑石”的散修头目接上了头。

老黑石是个年长的矿工头目，人如其名，皮肤黝黑粗糙如历经风霜的岩块，一双眼睛却锐利依旧。他打量着一身寻常散修打扮、笑容里带着三分市侩七分诚恳的吴刚，以及吴刚身后那位气质沉静、与周遭环境隐隐格格不入的陈默，并未立刻表态。

吴刚也不废话，寒暄两句后，直接伸手入怀，取出的是一卷看似普通的账册抄录。

他将抄录在油腻的木桌上摊开，手指点向上面几行迥异的数字。

“老黑石，您是明白人。咱们不讲虚的，看这个。”

吴刚的声音压得很低，却字字清晰，“您看仔细了。这一行，是丹鼎阁今年在咱黑石集收购黑须草的‘官定’收购价，明码标价，童叟无欺，对吧？”

黑须草算是一种不算多难培植的灵草，能够代替更常见的凝气草，在黑石集附近地区繁盛。这是黑石集附近部分从事农业散修的主要生产对象。

老黑石虽然是矿工，但是家里不乏种植黑须草的，所以行情也大概了解。

这个价格，和他的听闻差不多。

吴刚指尖一滑，落到下一行：“您再瞧这一行。这是同一批黑须草，被他们运出去，简单炮制一番后，转手卖给前线各宗，或者是炼制成最低阶的回气散后，在联盟控制的大城药铺里的售价。”

老黑石鼻翼微微翕动，盯着那串数字的目光仿佛要将其烧穿。

他仔细地看了一遍账本，却并未发现什么作假之处。

吴刚啧啧两声，摇着头，脸上写满了替人不值的唏嘘，“这还仅仅是一项，要是加上矿业、养殖业，中间的差价……嘿，我粗算了一下，刨去成本，剩下的利润，够把咱们黑石集老老少少、所有散修兄弟的吃喝用度，舒舒服服提升好几个档次，甚至能供养几个有天赋的后生安稳修炼到灵枢后期了。”

老黑石脸色本就黝黑，此刻更显得阴沉了几分。

这些事，他何尝不知？

黑石集的散修们起早贪黑种植的黑须草，冒着被妖兽袭击、矿洞坍塌的风险采集来的药材矿石，大半利润都被上游的丹鼎阁及其附庸家族轻易攫取。他们就像守在泉眼边的挑水工，拼死拼活挑上来二十桶水，自己能留下的，不过区区数桶。

吴刚观察着他的神色，又适时递上一张新的纸，上面是紫阳宗最新下达的征税令。“还有这个，紫阳宗的‘特别军税’，这个月额度比上个月又硬生生抬高三成！理由是‘云霞山战事吃紧，需加大投入’。”

他语带嘲讽，“可咱们得了什么好处？前线死的散修兄弟，抚恤金可有一枚灵石足额发到他们家人手里？咱们黑石集被征调去的后生，是编入了紫阳宗、丹鼎阁的核心战阵享福，还是被充作冲锋在前、试探陷阱、消耗对方箭矢的仆从军、敢死队？”

老黑石腮帮子咬得咯吱作响，眼中闪过一丝痛楚与无力。

这些都是血淋淋的现实，他手下就有几个兄弟一去不回，家人只等到一张轻飘飘的“阵亡通知”和几块打发叫花子的灵石。

更让他内心备受煎熬的是，面对丹鼎阁和紫阳宗的层层压力，他这个本地头目有时也不得不协助征调、催缴税款，明知是火坑，也得推着乡亲们往里跳，背地里不知挨了多少骂，听了多少戳脊梁骨的话。

虽然市侩，但他至少还有良心。

吴刚身体前倾，声音压得更低，却如淬毒的细针，直刺老黑石的心防：“老哥，咱们敞开了说。在紫阳宗、丹鼎阁那些高高在上的大人物眼里，咱们散修是什么？”

随即自问自答，每个字都透着冰冷的寒意，“是耗材！是用完即弃的灵符箭矢！是维持他们战阵运转、可以随意填补进去的数字！是需要时肆意强征暴敛、无用时便弃如敝屣的蝼蚁！”

“可我万仙盟，为何最终铁了心，跟云隐别院、跟叶先生站在一起？”吴刚话锋一转，语气带上了一种同道相认的诚挚，“就是因为看透了这一点！咱们散修，不想永远当耗材，当蝼蚁！”

他指向身旁一直沉默不语的陈默：“陈默道友，您或许听说过。原先也就是个跟咱们差不多的普通散修，机缘巧合得了叶先生指点，转修‘问道途’。您看看他现在，修为精进多少暂且不说，这精气神，这眼神里的笃定，您感受感受！在云隐别院，凭本事说话，有功则赏，叶先生传道解惑，何曾收过弟子一枚灵石的孝敬？”

陈默适时地微微颔首，周身那股沉静而内蕴锋芒的气度，与他朴素的衣着形成奇特的和谐，本身就是最具说服力的无声证言。

老黑石的目光在陈默脸上停留片刻，又回到吴刚身上，眼中的警惕与犹豫终于开始松动，取而代之的是一种强烈的、对现状不甘的共鸣，以及对另一种可能的好奇。

吴刚知道火候已到，不再停留于控诉，而是抛出了实实在在的合作方案：“老哥，不瞒您说，云霞山是被重兵围着，但您也是老江湖，该看得出，那不是一时半会儿能啃下来的硬骨头。 长久拖下去，守旧联盟这般涸泽而渔，四处树敌，内里能支撑多久？反观问道途这边，得道多助，路子越走越宽。 我们不求黑石集的兄弟们立刻明火执仗地造反，那是把大家往绝路上逼。但有些事，咱们可以暗中携手，互惠互利。”

他条理清晰地列出几条建议：黑石集的散修可以暗中将一部分采集到的药材、矿石，以远比丹鼎阁“官价”公平的价格，出售给吴刚联系的秘密渠道；对于联盟下达的征调令，可以采取阳奉阴违、拖延敷衍的态度；在不暴露自身核心的情况下，可以传递一些关于本地守军动向、物资运输路线等非绝密但对问道途了解后方态势有益的情报；甚至，可以暗中关照、保护那些在黑石集内，因听了传闻而对“问道途”产生兴趣、偷偷尝试修炼那本《养气篇》基础法诀的年轻人。

“作为回报，”吴刚的承诺同样具体而务实，“我们万仙盟的网络，可以帮黑石集的兄弟们，开辟一些……嗯，守旧联盟控制区之外的‘贸易选择’。你们需要但被联盟卡脖子的特定药材种子、某些特殊的工具法器图谱、甚至是一些外界的信息，我们可以想办法。 我们还可以分享一些，如何‘合理’规避某些过于严苛征役的小窍门；万不得已时，甚至能为一些兄弟的重要家眷，提供远离这战区边缘的临时安置渠道。”

当然，吴刚深谙驾驭之道，明白纯粹的利诱不足以长久，必须辅之以必要的威慑。

“当然，丑话说在前头。合作归合作，底线不能破。” 他脸上的市侩笑容收敛，目光变得锐利，“若是黑石集的哪位，借着咱们合作的机会，转头去欺压更底层的矿工、盘剥普通百姓，那我的人也会一笔一笔记得清清楚楚。将来若是时移世易，大局有变，我军真的得了势，到了清算的时候，这些人会是什么下场，我可不敢保证。”

最后，他总结道，语气恢复了那种令人容易接受的、权衡利弊的商人式诚恳：“风险可控，利益可见。 他们视咱们为蝼蚁，随时可以牺牲；而我们，却愿视黑石集的诸位为潜在的同道，为将来大业的基石。今日在此埋下合作的种子，浇水施肥，来日若得天时，未必不能生长为庇护一方、让兄弟们活得更有尊严的参天大树。即便最坏的情况，咱们没能成事，对黑石集而言，多一条隐秘的退路和财路，总比只在守旧联盟这一棵树上吊死、被一点点榨干骨髓来得强，对吧？”

这番立足于冰冷利益计算、却又巧妙描绘出远景与尊严的剖析，没有空泛的口号，每一句都敲在老黑石和类似头目们最关切、最疼痛的点上。它揭示了被剥削的现状，提供了切实的替代方案，指明了潜在的风险与远期的希望。

老黑石沉默了许久，最终，重重地将杯中残酒一饮而尽，酒杯落在桌上发出一声闷响。

这，便是应允的信号。

合作，就此在不见光的暗处生根。

虽然这种操作在沈瑶这种磊落的人看来并不光彩，但确实是非常务实而有效的了。

吴刚行事极为周密。

会谈之后，队伍中那名原本就出身黑石集附近、对当地人情脉络了如指掌的散修，便以“游子厌倦漂泊，归乡定居”为由留了下来。

明面上是落叶归根，暗地里，则承担起联络、协调，以及某种程度上的“观察”职责，确保合作沿着既定轨道进行，防止有人首鼠两端或趁机作恶。

另一方面，那个散修钉进黑石集之后，新的情报网络的最初节点自然也形成了。先前因躲避危机而停用、休眠的旧网络，也自然会随着新网络的启用而慢慢复苏。

这张无形的网，开始成为散修们共享信息、协调行动、彼此声援的地下生命线。联盟哪里又加了新税？哪支巡逻队换了凶狠的头目？哪个家族的仓库最近看守特别松懈？

这些零碎却至关重要的信息，开始在网络中悄然流动。

\vspace{2em}

吴刚更深的谋划，在于人的塑造。

他的眼光毒辣如鹰隼，目标不是那些热血上头、叫嚷最响的年轻人，也不是那些早已麻木、逆来顺受的老油子，而是那些本身在一定范围内具有威望、头脑相对清醒、对现状积累了足够不满、内心深处依然渴望改变的“中间层”。

“老刀子”秦厉，便是这样一个被吴刚“相中”的人。

他是黑石集南区矿洞的一个矿工，人如其名，脾气又硬又直。

但老刀子近来越发沉默，眼神里时常流露出一种深沉的疲惫与郁结。他侄子三个月前被紫阳宗的征调令强行带走，说是补充前线辅兵。一个月后，只传回一块破碎的身份玉牌和五块下品灵石的抚恤。他偷偷托人去打听，隐约得知侄子是被派去执行一次危险的试探性冲锋，死得毫无价值，尸骨都没找回。

本来骂了两句问道途，这事情就应该结束了；但他之后又听说，同期被征调的、几个依附于丹鼎阁的宗门子弟，经上下打点，却大多被安排在相对安全的后勤位置。

这件事像一根毒刺，深深扎进老刀子心里。

自此，他开始平等地怨恨问道途和守旧联盟，但更多的是无力。

他知道自己这点力量，对抗不了庞然大物，曾经的热血和硬气，在冰冷的现实面前，似乎只剩下每日收工后，在烈酒中寻求片刻麻木。

吴刚注意到老刀子，是通过那名留在黑石集的联络员。几次“偶遇”和酒馆里的观察后，吴刚判断，这是一个有底线、有担当、且内心痛苦正在寻求出口的人。他需要的不是煽动，而是一个支点，一点火星，以及一条若有若无、却切实存在的路。

第一次接触，不是在酒馆后间，而是在老刀子收工后，独自坐在矿洞外一块大石上，对着夕阳发呆的时候。吴刚扮作一个路过歇脚、收购零散矿石的行商，很自然地坐到了不远处，递过去一袋水，随口聊起了今年矿石的品相难寻，物价飞涨。

老刀子起初戒备，但吴刚谈论的都是最实际的行情，语气里没有高高在上的修士架子，反而像是个精明的生意人，偶尔还骂两句丹鼎阁收购官压价太狠，盘剥得太凶。这些抱怨说到了老刀子心坎里，两人的话匣子慢慢打开。

吴刚没有一上来就提“问道途”——他也知道老刀子对问道途的印象好不到哪去——他只是像一个见识广博的旅人，讲述着沿途见闻：“往南走，过了青岚江，那边的散修日子也不好过，但听说有些地方，散修自己抱团，搞了些小型的‘互助会’，一起接活儿，一起跟收购商讨价还价，虽然也艰难，但听说被盘剥得没咱们这儿这么狠……唉，也是听说，不知真假。”

老刀子听着，浑浊的眼睛里闪过一丝微光，闷声道：“抱团？谈何容易。没个领头扛事的，没点底气，一盘散沙，一吓就散。”

吴刚深以为然地点点头，状似无意地提起：“也是。不过我还听说，云霞山那边的散修，好像有点不一样。那边被围得铁桶似的，但里头的散修听说反而挺硬气。”

“别他妈的跟我提云霞山那边，闹心。”老刀子道。

“哈哈，不是夸他们。我又听说，有个叫什么陈默的，以前也就是个普通散修，现在可是云隐别院叶先生眼前的红人，又没打点又没干啥的，但修为精进不说，管着不少事儿呢。啧，这人跟人，命不同啊。”

“陈默？”老刀子似乎对这个名字有点印象。

黑石集消息杂乱，关于云霞山的只言片语里，好像确实提过这么个人物。他以前只当是谣传，此刻从一个看起来实在的“行商”嘴里说出来，分量似乎重了一点点。

“他不是吹出来的？”老刀子还是不相信。

“谁知道呢，道听途说。”吴刚摆摆手，结束了这个话题。

转而，说起一种能缓解矿工长期劳作筋骨酸痛的廉价药油配方，说是南方某个小地方流传的土方子，材料便宜，就是配伍有点讲究。

他很简单地说了说主要材料和大概比例，然后就像闲扯完一样，拍拍屁股起身，说还要赶路去下一个集子。

老刀子将信将疑，但矿工们的劳损病痛是实实在在的，他留了心。

几天后，他试着按那模糊的方子配了点药油，给几个老伤痛的兄弟用了用，反馈居然不错，虽然比不上真正的丹药，却也能缓解不少痛苦。

这让他对那个神秘“行商”的话，又信了一分。

当然，纵然如此，想要改变一个人对问道途的看法，道阻且长。

这只是开始。

这个行商倒是见多识广，有时候三天来一次，有时候十天半月不来一次；但每次和他闲聊，总能套到些不错的消息。

有时是一批市面紧俏、能有效防护矿尘的低阶面罩材料来源信息；

有时是关于紫阳宗某个税吏喜好和弱点的几句提醒，让老刀子能在交涉时少吃亏；

有时则是一些外界流传的、关于守旧联盟内部倾轧、后勤吃紧的模糊消息；

吴刚有时候也会装作不经意间提到问道途，讲的倒是半真半假，也算客观；又说有个女侠据说出自问道途，这几天四处乱窜，到处破坏敌人统治，乃至开仓放粮赈济灾民……

这些事情倒是真的，老刀子倒也听说过不少。

在十分话语中只有一分提及问道途，是吴刚的策略——而这一策略自然也是卓有成效的。

潜移默化间，老刀子看问题的角度和深度，悄然发生了变化。

他不再仅仅盯着眼前的矿石和工钱，开始下意识地思考丹鼎阁定价的不公根源，思考紫阳宗征调政策背后的冷酷算计。

更重要的是，他作为南区矿工中有威望头目的身份，使他的变化和偶尔透露出的只言片语，开始产生影响。

他依然不可能公开说什么——实际上对问道途的印象虽然好了些，但也确实没太好。

但在矿洞休息的间隙，在劣质酒馆的一角，当手下兄弟或相熟的散修抱怨“这日子没法过了”、“上面心太黑”、“后生们没出路”时，老刀子闷头喝一口酒，可能会沉沉地叹口气，无意间接上一句：

“听说云霞山那边，散修死了，抚恤是足额发的，家里人也有人安置。”

或者，在有人羡慕某个依附家族子弟又得了好处时，他擦擦脸上的灰，冷不丁冒出一句：“依靠别人赏饭，终究不稳当。我好像听说，云霞山那个叶先生，传道不怎么看重出身和灵石，有个叫陈默的，以前也就是个跑单帮的，现在可不一样了。”

他甚至连表演的痕迹都没有——因为根本没有表演，全都是真正拉闲散闷的口气。

配合着他本身硬气、实在的形象，反而比任何激昂的演说更具说服力。

听者会愣一下，然后追问：“真的假的？老刀子你听谁说的？”

老刀子通常只会摇摇头：“道听途说，谁知道呢。”

但那种子，已经借着他自己都没意识到有多不经意的话语，种进了听者的心里。

他们会私下讨论，会去打听，会结合自身遭遇去想象——“如果那是真的……”。

\vspace{2em}

吴刚塑造的，不止一个“老刀子”。

“老刀子”们或许不清楚背后运作的万仙盟网络全貌，甚至不清楚吴刚具体是个什么东西，也不知道哪个“到处乱窜的女侠”是何许人也。

但他们能模糊地感知到，有一群不同于守旧联盟的、带着些许温度的家伙——甚至有时候自己都不理解这些。

于是，这些人便成了吴刚网络中最基层、也最稳固的无形节点。

他们甚至不会打出旗号，不会聚众闹事；但在日常生活的每一个缝隙里，他们偶尔泄露出的一两句“牢骚”或“感慨”，持续地瓦解着对守旧联盟的敬畏与顺从，播撒着对另一种可能的好奇。

而陈默，这个活生生的“榜样”，则像巡回展示的珍宝，在绝对安全的掩护下，被吴刚带到几个类似黑石集的关键节点，与这些“节点”人物进行更深入的接触。

在伪装过的密室中，陈默没有任何夸夸其谈。他只是平静地讲述自己作为一个普通散修曾经的窘迫与迷茫，讲述在云隐别院第一次听到叶牧谦讲解“仰不愧于天，俯不怍于人”时的震撼，讲述转修问道途后，那种虽然缓慢却扎实的进步，以及内心逐渐明晰的笃定感。

他展示的是一种迥异于寻常散修的精气神——沉稳、内敛，眼神清澈而坚定，没有长期被盘剥压榨者的畏缩与怨毒，也没有骤然得势者的骄狂。

老刀子这样的人，看人极准。他们能看出陈默不是装的，那种从骨子里透出的踏实与尊严，做不得假。

陈默的存在本身，就是最有力的证据，证明了“问道途”和云隐别院，至少为他们这样的人，提供了另一种活法——或许依然艰难，但更有尊严，更有希望。

“陈默道友，”老刀子在一次秘密会面后，忍不住低声问，“叶先生……真如传闻所说，传道不论出身，不索贿赂？”

陈默认真点头：“老师有教无类。在别院，一切用度、功法，皆按需分配，按劳、按功获取。欺压同门、盘剥弱者，是首要大忌。”

老刀子沉默良久，重重吐出一口浊气。

“妈的，这问道途还真有点东西。”

\vspace{2em}

于是，发动群众的效果，开始以另一种更绵长、更深刻的方式，在黑石集以及类似的土壤上显现出来。

效果首先体现在最直接的经济层面——这自然是最快的，拜此地散修的实际领袖、与吴刚秘密签订条约的老黑石所赐。

丹鼎阁又派人来收购药材矿石了。这个精明的管事来这里也不是一回两回。却发现这个月黑石集上交的配额，虽然面子上勉强够数，但品质上佳的“尖货”比例明显下降，且总数量竟然比预估的“正常产出”少了那么一成半成。

“怎么回事，老黑石，你们这些家伙是不是偷懒了，没好好替我们收货？”

老黑石一脸愁苦地摊手：“仙师明鉴啊！今年矿脉贫瘠，妖兽闹得也凶，兄弟们实在是拼了老命了……”

“抗命是要杀头的。”管事无奈地道，“我也是有命令在身的，收货物收不上来我也不好做啊。”

老黑石也愁眉苦脸地递上来一张纸。

“您看，伤亡名单都在这呢。”

管事也没什么可说的，毕竟都到这种程度了，死者为大；况且收货的面子上也够数了，他对上级也有个交代。

而与此同时，吴刚的秘密渠道，却总能以比官价高出三到五成的“公道价”，稳定收到一批成色不错的货物。散修们拿到实实在在的更多报酬，对老黑石和那神秘的“外路朋友”自然心存感激，心照不宣。

\vspace{2em}

其次体现在对联盟统治的消极反抗上——这个自然要慢得多了，毕竟改变人心是一个更加缓慢的过程。

当紫阳宗的传令修士再次来到黑石集，要求紧急征调三十名灵枢境散修去填补前线某个危险区域的防务缺口时，响应者寥寥。

以往靠利诱和威逼总能凑齐的人数，这次磨蹭了三天，也只来了十几个老弱或实在走投无路之人。

传令员自然是不信老黑石等人“动员不上来”的说辞的，他倒也务实，选择亲自下去动员散修。

老黑石同样跑前跑后，跟着传令员“苦口婆心”动员，却挨了老刀子一通臭批：

“驴*的紫阳宗的狗东西，你先还我侄儿再说！老黑石你他妈的也是，怎么还敢来这，再敢来我把你揍出去！”

一连几处，皆是如此。老黑石和传令员也是挨了不少唾沫星子。

最终，老黑石也只能无奈地回禀：“仙师，您也看到了，人心涣散，惧怕前线死地，小老儿实在是尽力了……”

传令员灰头土脸，挨了一天骂饶是谁也都不好受，也只能带着这十几个人回去复命了。

这自然也仅仅是个缩影，但随着时间流逝，征调令的执行效率确实在大打折扣。

联盟在黑石集及周边区域的驻军换防时间、巡逻路线、仓库守备的薄弱时段……这些信息，开始通过那张地下网络，悄然流向吴刚，经过汇总分析后，又变为更有效的行动指导，流向沈瑶等人。

而守旧联盟对黑石集等地内部的真实人心动向、暗流涌动，却如聋似瞎，只知道收税、招人不太好办了。

\vspace{2em}

最后，也是最重要的，是一种人心氛围的微妙转变。

集市上，茶摊酒肆里，公开咒骂联盟的声音或许还不敢有；但那种麻木的顺从在减少，取而代之的是一种闪烁的、心知肚明的沉默，以及私下交流时眼神中蕴含的别样意味。

当有联盟修士趾高气扬地走过时，背后投去的目光里，敬畏在减少，隐忍的厌恶与一种近乎怜悯的疏离感在增加。

那句“仰不愧于天，俯不怍于人”，开始像一道微弱却执拗的光，照进一些不甘沉沦的心灵深处。

吴刚领导的这张网，如同生长在守旧联盟庞大躯体深处的无形菌丝。它不发动正面的攻击，不攫取显赫的城池，却以惊人的耐心与韧性，无声地缠绕、渗透进联盟基层统治最细微的根系之间。

它缓慢而持续地吸走本该属于联盟的养分，悄然松动其统治的土壤，并将另一种关于生存方式、关于尊严与未来的可能性，如同孢子般，悄无声息地播种在无数散修的心田。

\vspace{2em}

当然，暗流涌动之下，必有波澜。

守旧联盟能做到现在这种大小，如果其说其中枢是蠢物那有些太不负责任了；只是，其臃肿的躯体对发生在末梢神经的细微刺痛，反应往往迟钝而傲慢。

这件事一开始甚至没有捅到陈炎那里，

紫阳宗前线那位负责后勤协调的中级执事，翻看着下面陆续呈上来的卷宗，眉头拧成了疙瘩——这个月从黑石集等几个固定供应点收上来的低阶矿石和药材，平均品质比上月低了半成，总量也差了那么一丝。

呈报上来的理由千篇一律：矿脉贫瘠，妖兽滋扰，人手折损。

“废物！”执事低声骂了一句，随手将卷宗丢到一边。

在他眼中，这只是下面那些贪婪又无能的附庸小家族或散修头目惯用的伎俩，想借此讨要些补贴或降低配额。

相较于千头万绪的战报和日常管理，这等小事尚不足以惊动更高层。

要是这种鸡毛蒜皮都上报，那免不了挨一顿批评：“此等小事还拿来犯我，你能力是不是有点过于有限了？”

是啊，这些下级，能力是不是过于有限了？

这个执事有些恼火，批了一句：“严加督促，务必足额，否则以延误军机论处！”便将这份报告归档了。

几乎同时，另一份来自某个地方驻军的报告也送到了他的案头，报告称辖区内“流言频起”，有低贱散修私下议论前线战事不利，甚至妄言紫阳宗苛待仆从军。

执事的火气更大了，提笔批复：“妖言惑众者，一经查实，立斩！驻军加强巡防，以儆效尤。”

他依然没有将这几件“小事”联系起来。

在他，乃至守旧联盟大多数中上层修士的思维定式里，后方的些许杂音，无非是管理懈怠或愚民无知，派几队更凶悍的巡逻队下去敲打一番，杀几只鸡，猴子们自然就老实了。

他们尚未意识到，这早就不是简单的懈怠或愚昧了；而是一种无声的、织网般蔓延的抵抗。

真正让“小事”开始引起区域层面注意的，是沈瑶那边愈演愈烈的打砸抢。

距离黑石集数百里外，一处由丹鼎阁某个附属家族控制的灵谷转运站，在某个深夜燃起了冲天大火。不仅转运站本身被焚毁、数千担灵谷不翼而飞，更有数名家族修士被当场格杀。

现场留下用血迹涂抹的大字：“取不义之粮，还于饥民！”

只知道那名——或者那伙——“贼人”擅长短刀，轨迹极为精确，基本上全奔着一刀毙命去的。

几乎每隔十天半月，类似的事件就会在守旧联盟控制区域的薄弱地带上演。

有时是押送税款的队伍在半道被劫，护卫被杀，灵石不翼而飞；有时是某个盘剥甚重的小家族仓库被撬开，里面的物资一夜之间消失大半，只留下些散给附近穷苦人家的痕迹。

行事者来去如风，手段狠辣精准，绝不纠缠，一击即走。渐渐地，一个模糊的代号在底层流传开——“那伙煞星”，或者更具体点，“那个穿蓝衣服、刀快得像鬼一样的女人”。

\vspace{2em}

于是这些事件终于不再是“小事”。

损失实实在在，影响恶劣，更重要的是，它们指向了一个明确的、有组织的势力活动。

不知道这贼人是哪来的，但是有贼人到处活动这事情是真的。

几位主事者们坐不住了。

一次紧急议事中，负责本区域防务的紫阳宗执事面色不善地敲着桌子：“不能再姑息了！先是物资短少，征调不力；现在匪患公然劫掠，袭杀我修士！这已不是简单的散修和老百姓了，必须要出重拳！”

另一位丹鼎阁的代表则更关心实际利益，阴恻恻地说：“尤其是黑石集一带，近来颇有古怪。收缴上来的货物品相持续走低，必有猫腻。说不定，那里就是匪徒销赃、补充给养的窝点之一！”

他们迅速达成了共识，双管齐下。

一方面，调集精锐，加大对那支“流窜匪帮”的清剿力度，重金悬赏其人头；另一方面，对黑石集这类“可疑区域”施加高压，派加强的巡逻队进驻，明查暗访，务必揪出内鬼，掐断任何可能与“匪帮”联系的渠道，重新树立联盟的绝对权威。

\vspace{2em}

数日后，第一支背负着震慑与清查双重命令的联盟巡逻队，开进了黑石集。

这支队伍由二十名修士组成，远超平日巡防的规模，其中领队是一名玄脉境通络期的紫阳宗资深弟子，面色冷峻，眼神里带着不加掩饰的审视与倨傲。队伍中还有两名来自丹鼎阁、专管财务的执事，专门负责查验货物账目。

集市主干道上，人群下意识地分开，原本嘈杂的讨价还价声低了下去，无数道目光或畏惧、或麻木、或带着隐晦的疏离，从街道两旁、摊位后面投来。

巡逻队队长很满意这种效果，他要的就是这种令人窒息的威压感。

他们首先直奔老黑石的“公事房”——一间简陋的石屋。老黑石早已得到风声，恭敬地迎了出来，腰弯得很低，脸上堆满了讨好的、属于小人物的愁苦笑容。

“仙师们大驾光临，有失远迎，恕罪恕罪！”老黑石搓着手，语气卑微，“可是上峰有什么新指令？小老儿一定尽心竭力去办！”

队长没理会他的客套，径直走入石屋，目光如鹰隼般扫过每一个角落。两名丹鼎阁执事则毫不客气地开始索取近三个月的出货账册、矿工伤亡记录、以及所有与丹鼎阁交易的明细凭证。

老黑石早有准备。账册记得清清楚楚，明面上的出货量、品质分类完全符合“官方要求”，只是旁边用朱笔小字备注的“矿脉贫瘠预估减产”、“妖兽袭击损失”、“新矿工伤亡抚恤支出”等条目，密密麻麻，触目惊心。

伤亡名单更是厚厚一摞，每一个名字后面都按着手印或简单的标记，透着一股沉甸甸的、令人不悦的真实感。

一名执事皱着眉头翻看，试图找出涂改或作假的痕迹，却一无所获。他抬起头，狐疑地盯着老黑石：“老黑石，这数目，可经得起查验？若是虚报伤亡，侵吞物资，你知道后果。”

老黑石脸上的愁苦简直要滴出水来，他指着账册，声音都带上了哽咽：“仙师明鉴啊！借小老儿一百个胆子也不敢啊！实在是……唉，仙师们若是不信，大可去矿洞看看，去那些死了男人的家里问问！这日子，难啊！”

他表演得如此自然，将一个小人物在夹缝中求生存、既要完成任务又无力回天的形象刻画得入木三分。至于这些账目，在吴刚等人的参与和指导下，七分真、三分假，假的地方大致上也算天衣无缝。

巡逻队队长哼了一声，对这种诉苦把戏不感兴趣。

他手一挥，两名丹鼎阁财务执事立即开始上前查账，算盘打得噼里啪啦。

当然，巡逻队队长更关心的是“匪患”线索。

“最近集市上，可有什么生面孔？有没有人私下打听云霞山，或者交易来路不明的物资？”队长冷声问道，目光锐利如刀，仿佛要剖开老黑石的脑子看看。

老黑石茫然地眨眨眼，做努力思索状：“生面孔……这黑石集哪天不来几十号生面孔？都是讨生活的苦哈哈。云霞山？哎哟，那可不敢提，那是叛逆之地啊！至于来路不明的货……”

他压低了声音，显得小心翼翼，“仙师，咱们这儿规矩，来历不明的东西不收，怕惹麻烦。倒是……倒是偶尔听说，有些胆大包天的家伙，会偷偷把一点私货卖给那些行踪不定的游商，价格能高些，但风险也大，小老儿是严令禁止的！”

他把“偷偷”和“游商”咬得很重，既暗示了可能存在的地下交易——这符合联盟对平民百姓、底层散修的刻板印象，又把自己撇得干干净净。

与此同时，巡逻队的其他成员已经分散开，在集市上进行抽查。他们闯入几家较大的药材铺和矿石行，翻箱倒柜，盘问掌柜和伙计，语气凶狠。

集市上一时间鸡飞狗跳，人心惶惶。

然而，他们很快发现，面对的是一个无比配合却又无比空洞的群落。

所有的家伙此刻都和常规没有两样。老刀子秦厉正光着膀子，和一群矿工蹲在街角啃着干粮，满脸灰尘，眼神和其他人一样警惕而麻木。

巡逻队员过来盘问，他夹枪带棒地骂了两句，表示什么都不知道，只念叨着工钱难挣，饭不够吃，“最好少收点税”。

所有的关键交易，早已转入更深、更隐秘的渠道。

吴刚的秘密运输线采用了更灵活的单线联系和分段接力，物资在联盟巡逻队到来前，就已通过不同的小路，化整为零地转移或藏匿起来；然后再分批混在更多的货物中，或采取其余手段，把这些“灰色”货物逐渐地变成经得起查验的好货。

通俗来说，那叫“洗钱”。

洗钱归洗钱，至少他这个洗钱确实起到了积极效果。

现在集市上流通的，都是经得起查验的清白货物。

巡逻队折腾了大半天，近乎一无所获；查账的两个财务人员也没看出什么蹊跷。

但他们想找的内鬼、联络点、大批违禁物资，一样也没发现。

领队的队长脸色越来越难看。他凭直觉感到不对劲，这里的气氛太过正常，正常得过了头。

可他抓不到把柄。

最终，他只能抓了几个平日里就小偷小摸、或是酒后胡言乱语说了几句怪话的倒霉蛋交差。

临行前，巡逻队队长严厉警告老黑石和集市上几个有头脸的人，必须严密监视可疑人等，及时上报，否则以通敌论处。

然后，带着一丝悻悻然和犹疑，巡逻队离开了黑石集，前往下一个“可疑”地点。

\vspace{2em}

一连几处，皆是如此。

\vspace{2em}

然而，当后方“物资损耗异常”、“征调不力”、“流言频起”乃至“匪患劫掠”的报告太多、无法处理，最终如同雪片般从各个方向汇集，甚至有几份直接呈送到了坐镇前线的陈炎案头时，事情的性质就变了。

尽管陈炎的主要精力仍被云霞山这个“硬骨头”牵制，但后方此起彼伏的烽烟，尤其是那些精准打击补给线、动摇统治基础的行动，让他无法再将其视为简单的“流窜匪患”或底层散修的抗命。

一道带着陈炎冰冷意志的命令被迅速下达：后方各区域，必须严加整肃，加大巡查与清剿力度，务必在最短时间内扑灭这些“星星之火”，保障后方安稳与前线的物资供给。

压力如同不断收紧的绞索，开始套向黑石集这样的地方，也迫使沈瑶和吴刚的行动更加艰难和险象环生。

双方的探子、小股部队在城镇外的荒野、山林间的遭遇变得频繁。

在一次对某处小型联盟物资中转站的侦察任务后，沈瑶带着游击队的三名精锐队员，趁着夜色撤离，准备返回位于深山中的临时营地。

他们行动如狸猫，借着山林阴影和沈瑶以千幻流光诀制造的细微视觉扭曲，几乎与夜幕融为一体。

就在穿过一片稀疏的林地，即将进入更茂密的山岭时，前方突然传来杂乱的脚步声和低沉的呼喝声，还有几道神识之光如同探照灯般扫过林间空地——是一支约十五人的联盟巡逻队，似乎是在执行例行的夜间巡防，领头的是一名神色警惕的玄脉境修士。

双方距离不过百丈，且处于相对开阔地带，遁走极易引发灵力波动而被察觉。

硬闯？对方人数占优，一旦纠缠，附近很可能还有其他巡逻队。

现在还没被发现，还有时间伪装！

电光石火间，沈瑶眼神一凛，手势微动。

队员们心领神会。

只见四人几乎同时，发生了细微却显著的变化：沈瑶身上那属于精锐修士的凌厉气息瞬间收敛殆尽，换上的是一种长途跋涉后、饱经风霜的疲惫与谨慎，她甚至下意识地弯了弯腰，仿佛背负重物；其他三名队员则迅速将身上的衣装弄得凌乱，脸上也抹上些尘土，眼神变得茫然又带着点小民的畏缩。

一切准备完成，几人调整方向，略显慌乱地朝着巡逻队即将经过的方向“撞”了过去。

常言道：最危险的地方，就是最安全的地方。

“什么人？！站住！”巡逻队立刻发现他们，数道神识和兵刃瞬间锁定。

沈瑶立刻举起双手，用带着浓重口音、有些颤抖的方言喊道：

“军、军爷饶命！俺们是北边野猪岭的采药人，进山寻点山货，迷了路……这、这才转出来……”

领头的玄脉境修士眉头紧皱，神识仔细扫过四人。

修为最高者不过灵枢中期，气息虚浮杂乱，衣着普通甚至寒酸，身上带着淡淡的草药泥土味和汗味，确实像极了常年混迹山野的底层采药散修。

他目光锐利地扫过沈瑶的脸——那是一张粗糙、黝黑、带着惶恐的农妇面孔，毫无特点。

但野猪岭离这边有五十里地！

“野猪岭？跑到这五十里外的林子来采药？当老子好糊弄？！”

巡逻队长厉声喝道，同时示意手下散开，形成半包围。

沈瑶身体抖得更厉害，带着哭腔道：“军爷明鉴啊！俺们……俺们也是没办法！野猪岭那边被征了‘入山税’，俺们交不起，听说这边林子深，没人管……就、就冒险过来了……谁知道这鬼林子这么大……”

她一边说，一边从怀里颤巍巍掏出一个破旧的布袋，倒出几株品相普通、还沾着泥土的“灰线草”和“地根藤”，都是黑风林附近常见的低阶药草。

“军爷，这是俺们今天采的……都、都孝敬军爷，只求军爷指条明路，放俺们回去吧……”

巡逻队长看着那几株再普通不过的药草，又看了看眼前这四个采药人瑟缩恐惧的模样，心中的疑窦去了大半。

这种为了逃避苛捐杂税、冒险进入陌生区域讨生活的散修他见多了，都是些穷得叮当响、修为低微的可怜虫。

抓起来压榨，也榨不出什么油水。

更何况那农妇容貌根本入不了他的脸。

他嫌恶地挥挥手，像驱赶苍蝇一样：“滚！以后不许再靠近这片区域！再让老子看见，按奸细论处！”

沈瑶四人如蒙大赦，连连鞠躬道谢，然后互相搀扶着，脚步踉跄却速度不慢地“逃”向了与真实营地相反的方向，迅速消失在黑暗的林地中。

直到彻底脱离对方神识范围，四人才停下，互望一眼，都从对方眼中看到一丝紧绷后的放松。

沈瑶面容一阵水波般的模糊，恢复了原本的模样。

这次蒙混过关，靠的是对底层生态的精准模仿、毫无破绽的伪装，以及沈瑶那足以完美控制自身气息甚至在一定程度上干扰低阶修士感知的千幻流光诀；但第二次，他们就没有这么幸运了。

\vspace{2em}

真正的危机，来自内部的背叛。

吴刚发展的地下网络，不可能永远铁板一块；在巨大的压力、或许还有守旧联盟暗中开出的价码诱惑下，人性中的脆弱与贪婪总会暴露。

一个代号“灰鼠”、负责传递情报和转运少量物资的下线，被当地一个依附丹鼎阁的赵家秘密策反了。

赵家老祖是真符境摹形期的修士，故而在本地势力颇大；早就对境内“不安定因素”深感恼火，得到了上峰“不惜代价肃清”的严令。

毕竟人都是有软肋的；有软肋，那就可以拿捏。

他们抓住了灰鼠一个至亲的把柄，并以重利许诺，逼其就范。

威逼利诱之下，灰鼠最终还是选择叛变。

或许是因为他想把队友分期卖个更高的价钱，也有可能是因为他良知尚存，灰鼠没有立刻暴露全部网络，而是选择向赵家提供了一条他认为是“大鱼”的情报：

三天后，游击队将秘密潜入镇外废弃的老矿洞区域，接收一批由吴刚网络从其他地方转运来的、较为珍贵的疗伤丹药和一批用于制造简易阵法的材料，并计划在那里休整一夜！

他赌的是用沈瑶这条“大鱼”，来换取自己和家人的安全，以及赵家许诺的荣华富贵。

情报被火速呈报。赵家老祖亲自出马，协调了当地紫阳宗驻军，抽调精锐，秘密在老矿洞周边布下了天罗地网。

他们甚至动用了能够屏蔽低阶传讯、干扰灵力感应的阵旗，只待沈瑶小队进入，便一网打尽。

沈瑶对此一无所知。

灰鼠这条线之前传递过几次无误的情报，且这次接头地点和时间，是通过吴刚网络内部相对安全的渠道确认的，应当无误。

直到他们深入矿洞区域腹地，准备发出接头的暗号时，异变陡生！

“轰！”

四周山壁上，预先埋设的爆裂符骤然激发，制造了相当的巨响和混乱，打断任何可能的遁术起手。

与此同时，至少超过五十道身影从废弃的矿坑、堆积的废石堆后、甚至是从地下破土而出！

居中那位白发老者，周身符光隐隐，身前悬浮着一枚不断变幻形状的土黄色符印，正是赵家那位真符境摹形期的老祖！

其左右，则是两名玄脉境后期的赵家长老；外围，是数十名紫阳宗驻军和赵家护卫组成的包围圈，刀剑出鞘，弓弩上弦，灵光锁定了中心的沈瑶小队。

“你就是沈瑶？好个云霞山余孽！束手就擒，否则今日此地，便是你的葬身之处！”

赵家老祖声如洪钟，带着志在必得的杀意。

中计了！

沈瑶瞬间明白过来，眼神冰冷如万载寒冰。她自认凭千幻流光领域以及堪比真符境的战斗力，这些人估计是留不下她本人的。但是问题是，如果赵家老祖故意放过她，选择强行留下其他游击队员……

她没有任何废话，清叱一声：“跟我冲！”

话音未落，她已身化一道变幻不定的流光，直扑包围圈西北角。

根据她瞬间的观察，那边地形复杂，有突围的可能。

千幻流光诀全力催动，她身影仿佛同时出现在数个方位，道道虚实难辨的刀光泼洒而出，瞬间将那片区域的守军搅得人仰马翻。

队员们也红了眼，知道已是绝境，唯有血战求生。

他们结成小型战阵，不顾自身防御，将攻击催发到极致，紧紧跟随着沈瑶撕开的缺口。

“想走？！”赵家老祖冷哼一声，身前土黄色符印光华大盛，瞬间化作一面巨大的、布满玄奥纹路的岩石屏障，轰然砸落在沈瑶的突围路径上，更有一股沉重的重力场弥漫开来，让沈瑶和队员们的动作骤然一沉。同时，他袖中飞出数道金光，竟是成套的飞剑符宝，凌厉无比地攒射向沈瑶。

另外两名玄脉境长老也联手攻来，术法光芒呼啸。

“靠，你们先走，真符我能拖住！”

沈瑶骂了一声，立即旋身下达指令。

但没有人听从命令。

“队长，我们跟您一起！”

“你们是笨蛋吗！”沈瑶急了。

但脑子长在别人身上。

很快，其中一名游击队员被飞剑符宝穿透胸膛，当场壮烈牺牲，眼见得救不活了。

另一名队员在引爆身上携带的所有爆炎符，暂时逼退一片敌人后，也被乱刀砍倒。

“要活的，你们干什么呢！”

赵家老祖狂吼道。

很快，鲜血染红了废弃的矿场。沈瑶自己手臂也被一道符宝金光擦过，受了些伤。

她咬紧牙关，知道再拖下去，全体都要葬送在此。

“散开！各自突围！老地方汇合！”她厉声喝道，同时拼着受了赵家老祖一记符印重击，虽被震伤，却借力向后急退。

忽然——

“弥焰斩！”

众多追兵忽然听得一声男子暴喝，却听不真切是哪里发出来的。

“什么，那个姓白的也来了？”

“那个弥焰斩可是擦着边就死、碰了些就亡啊！”

赵家老祖也忌惮起来，毕竟白言冰加上沈瑶，两个真符……

就在这众人因犹疑而迟滞的一瞬！

“轰隆——！！”

震耳欲聋的爆炸声中，电蛇乱窜，火光冲天，剧烈的黑紫剑气（？）风暴暂时遮蔽了视线和感知。

随即，是浓到化不开的烟尘。

沈瑶身影在爆炸的掩护下，左右手上各抓了一个受了伤、无法全力飞遁的家伙。随即，身形如鬼魅般连连闪烁，千幻流光诀催动到极致，气息变幻不定，甚至模仿出数道向不同方向逃窜的虚影。

赵家老祖第一个反应过来。

这黑紫剑气徒有虚表，毫无杀伤力啊！

还能这么打架的？

怒吼着驱散烟尘，神识疯狂扫荡，却只抓到几缕迅速消散的幻影和远处隐约的灵力波动，根本无法锁定沈瑶真身。

他脸色铁青，知道最关键的“大鱼”脱钩了。留下的两个游击队员也已经成了尸体，盘问算是盘问不出来了。

“娘西匹，给我追！”

——上哪去追啊！

\vspace{2em}

沈瑶全力飞遁数十里，才终于停下，算是勉强逃出了包围圈。

跟随她成功突围的，只有三名玄脉境巅峰的队员，都受了些轻伤；以及另外两名被沈瑶抓出来、身受重伤的队员，但算是彻底失去了战斗力。

“对了，白师兄哪里去了？”一名帮助沈瑶临时为重伤号包扎伤口的队员随口问道。

“白师兄？哪个白师兄？”沈瑶疑惑道。

“白言冰师兄啊！”

沈瑶白了他一眼，“那是我用幻术领域模仿的！你当他真能来啊，那山上防务谁负责？”

队员面色有些变了：没想到沈瑶的幻术竟然到了这种地步，甚至连白言冰的杀招——弥焰斩都能轻易模仿！

沈瑶看出他的惊奇，又补充道：“那黑紫剑气也是假的！”

队员哑然，笑也不是，不笑也不是。

沈瑶逃出生天、立即通过紧急渠道将“任务失败”的消息传递给吴刚。

吴刚暴怒，动用地下网络的力量查明事实后，在灰鼠试图领取赵家赏金的前夜，将其秘密处决。

经此一役，“沈瑶”这个名字及其大致特征，以及其余几个逃出生天的队员，被守旧联盟正式列为高度危险的通缉要犯，画像与悬赏令贴满了后方许多城镇的关口要道。

敌人已经深刻认识到她的威胁，并会调动更多资源、采取更狡猾的手段来对付她和她的队伍。

当然，这对沈瑶而言几乎毫无意义——千幻流光诀一开，她想换成什么模样，就是什么模样。

通缉令上那张脸，在内域大乱结束之前，估计永远也不会再次出现了。

单纯的通缉自然是无意义的，联盟随即在内域全域大索，大力镇压可能的反对派活动，搞得整个内域鸡犬不宁，民不聊生。

然而，高压往往带来更强烈的反弹。越是镇压，那些分到过粮食的贫民、那些受过盘剥的散修、那些偷偷修炼《养气篇》感受到体内微弱元气增长的年轻人，心中对守旧联盟的憎恶和对问道途的向往就越是强烈。

地下网络、游击队伍在一次次的打击中变得更加隐秘和坚韧，如同野草，火烧不尽。

\vspace{2em}

很快，消息通过秘密渠道断断续续传回云霞山。

尽管山上的防御战依然惨烈，物资依然紧缺，但当叶牧谦将沈瑶、吴刚他们在敌后发动群众、分发物资、传播理念的事迹选择性地告知守军时，一种难以言喻的振奋在山间弥漫开来。

他们不是在孤军奋战！

他们的抗争，正在遥远的地方激起回响，正在动摇敌人看似稳固的统治基础！

这种认知，极大地鼓舞了防守者的士气。

白言冰在指挥防御时，有时会吼出：“想想沈瑶师妹他们在敌后为民请命！我们守住这里，就是守住希望的火种！”

陈千羽在救治伤员时，也会轻声鼓励：“吴刚师兄他们找到了更多药材来源，大家坚持住，我们的药田又快丰收了。”

而叶牧谦收到密信时，更是久久不能平静。

他的思路，是正确的。

实践也检验了真理。

本想提笔回复一些鼓励的话，但最终千言万语，终究不如井冈山上那光耀千古的八个字：

“星星之火，可以燎原。”

山上山下，尽管被重兵隔断，却通过精神的共鸣和战略的联动，形成了一个坚韧的整体。

云霞山如同吸引雷霆的孤峰，承受着正面最大的压力；而敌后战场则如地下奔涌的暗流，不断侵蚀着敌人的根基。

当陈炎再次于中军大帐中接到后方接二连三关于“匪患”“邪说流传”的急报时，他的脸色终于不再是单纯的阴沉，而是带上了一丝不易察觉的焦躁。

他发现自己面对的，不再只是一个龟缩防守的军事据点，而是一个正在用他完全不能理解的方式，将战争扩展到整个社会层面的可怕对手。

“叶牧谦，你到底是什么人……？”

陈炎不禁冷汗淋漓起来……

\vspace{2em}

沈瑶在许多个“灰谷镇”斩开的仓门后纷飞的粮食，吴刚在暗室中与许多个“老黑石”们推心置腹勾勒出的利益与尊严的蓝图，成了守旧联盟大后方的星星之火。

看似微弱，却因其内核燃烧的是“公道”与“希望”的燃料，而拥有了顽强的生命力。它们开始在守旧联盟统治的厚重阴影下，倔强地闪烁、蔓延。

最直观的对比，发生在每一天，每一处。

守旧联盟的大军或其爪牙过境，无论面对的是名义上的附属城镇还是自家的后方腹地，“征发”是唯一的逻辑。

粮食、药材、矿石，乃至低阶修士的性命，都在“道争大业”的冠冕堂皇之名下，被理所当然地索取。

尚有良知、愿意通融的低阶军官固然存在，但执行这道命令的更多低阶军官、税吏、附属家族的狗腿子们，将这视为大发横财、作威作福的良机。

他们敲骨吸髓，中饱私囊，视散修与凡人为无血无泪、可以无限榨取的资源矿藏，其行径之暴虐贪婪，较之山野匪类有过之而无不及。

失去儿子的母亲眼泪未干，夺走口粮的差役鞭影又至。

与此形成刺眼到令人心颤对比的，是沈瑶、吴刚小队那秋毫无犯，开仓济民的作风。

许许多多个“灰谷镇”的故事，早已被口耳相传，添加上百姓最朴素的想象与期盼，演化成数个略有不同的版本，在周边的乡镇、矿场、散修聚集地悄悄流淌。

起初，人们只当是乱世中一个慰藉心灵的奇谈，带着深深的怀疑与本能的观望；然而，现实是最好的教员。

当守旧联盟的征粮队因“损失”而将亏空变本加厉摊派到邻近村镇时，当催缴“特别军税”的皮鞭再次抽打在早已家徒四壁的农户背上时，那份被强力压制的怨愤，如同被堵塞的岩浆，开始炽热地寻找每一个可能的裂隙。

第一个大胆的散修，在月黑风高的夜晚，怀揣着难以抑制的愤怒与一丝微茫的希望，循着模糊的传言，冒险找到了沈瑶小队曾经短暂歇脚、现已废弃的山洞。

他留下的不是金银，而是一张用炭笔画在粗麻布上的简陋草图，上面歪歪扭扭却清晰地标注了一支征粮队返回驻地的必经之路、人数、以及押运修士的大致修为。

这份冒着生命危险送出的情报，成了点燃更多火焰的火种。

沈瑶——这回谨慎得多了——根据情报，多次踩点之后精准设伏，不仅截获了粮队，更再次将粮食分发给沿途受尽盘剥的村庄。

观望，在这一刻，化为了实实在在、浸透着泪水的感激。

人心的天平，开始了无声却致命的倾斜。

支持的形式，从此如雨后山林里的菌菇，在仇恨与渴望的潮湿土壤中，悄然萌发，形态万千，生命力顽强。

有的支持，如同那位无名散修，是关键时刻递出的一张纸：一张写着某条补给线护卫换班间隙的纸条，塞进赶集农妇的菜篮底；几句关于某个矿点监工今日醉酒、守卫松懈的闲谈，“无意”飘进酒馆角落里看似打盹的旅人耳中。

这些零碎的情报，通过吴刚再次精心编织、已愈发细密坚韧的地下网络，被迅速筛选、拼接、传递。

最终，它们可能化为沈瑶“千幻流光”下的一次精准斩首，救出几名即将被折磨至死的矿奴；或是引发某个偏远税卡守卫的“哗变”，整队人马带着武器和税款消失在山林，据说是投奔了游击队。

\vspace{2em}

寒鸦岭，这名字听着就透着一股子穷苦与荒僻。

岭下散落着几十户人家，多是土坯垒墙、茅草覆顶，日子过得比岭上的石头还硬。张老栓家就在村西头最边上，两间矮趴趴的茅屋，像是随时会被山风吹散架。

约莫三天前的后半夜，张老栓被屋后柴垛一阵窸窸窣窣的动静惊醒，还夹杂着压抑的、粗重的喘息。他抄起门边的柴刀，壮着胆子摸黑出去，月光下，只见柴垛旁歪着个人，半边身子都被暗红色的血浸透了，脸色白得像死人，只有胸口还在微弱地起伏。

那人穿着看不出本色的粗布衣服，但手里还死死攥着一把卷了刃的刀，刀柄缠着的布条都被血粘住了。

张老栓的心提到了嗓子眼。

他认得这打扮，近来山外边不太平，传得沸沸扬扬，说有一伙专跟“老爷们”作对的好汉，领头的是个女侠，来去如风。

眼前这位，八成就是那伙人里的，看样子是在别处吃了亏，一路逃命躲到了这山旮旯。

救，还是不救？

张老栓蹲在黑影里，嘴里的旱烟吧嗒吧嗒，却半天没吸进一口。救了吧，万一被王老爷或者上面下来的兵爷知道，这就是包庇匪类，要杀头的。

不救吧……这后生眼看就剩一口气了，扔在外头，不是冻死就是被野狗啃了。

他想起了上个月，岭东头老赵家。赵家的独苗后生被王老爷的征粮队硬拉去当了脚夫，说是去前线运粮，结果人没回来，只捎回句话，说是“遭了流匪，不幸没了”。

而王老爷的管家来“抚慰”时，只丢下两斗谷子，那架势，跟打发叫花子没两样。

一个壮劳力的性命就顶两斗谷子！

老赵头当时就呕出口血，没挺过三天。

他又想起了大概半个月前，村里悄悄流传的消息。说百里外的“灰谷镇”，老爷家那个屯了满仓粮食、却年年加租的大粮仓，被一伙人给砸开了，里头的粮食全分给了镇上的穷户。

消息传得有鼻子有眼，还说那伙人只抢老爷的，不碰穷人家一粒米。

当时村里人听了，只当是故事，心里却都觉得解气，晚上睡觉都踏实了些。

柴垛边那伤员的呼吸更微弱了。

张老栓把烟锅子在鞋底磕了磕，站起身，像是下定了决心。

他左右瞅瞅，确认四下无人，这才费劲地把那人连拖带拽，弄进了屋。家里空荡，藏不住人，他一眼瞅见了屋角那堆平时做饭用的柴火。老伴去得早，儿子也被征走没了音讯，就他一个孤老头子，柴火用得少，堆得老高。

他挪开面上干燥的柴捆，露出后面墙壁。这墙是土坯的，早年防山匪，他在后面偷偷挖了个小窑洞，不大，也就刚好能蜷个人，洞口拿块破木板遮着，再堆上柴，外人绝看不出来。洞里阴冷潮湿，常年不见光，有一股子霉土味。

张老栓把伤员小心地塞进去，让他半靠着土壁。又摸黑从自己床铺底下，拿出小半块黑乎乎、硬得像石头的杂面饼——这是他半天的口粮。

再提起墙角一个破葫芦，里面装着半葫芦清水。

他蹲在洞口，把饼子和水递进去。伤员似乎恢复了一点意识，艰难地抬起眼皮，眼神里充满了警惕和茫然。

“吃点，喝口。”张老栓的声音干巴巴的，没什么情绪， “别出声。白天不能出来。”

那伤员盯着他看了几秒，似乎判断出眼前这个满脸沟壑的老农没有恶意，才用颤抖的手接过饼和水，费力地咬了一口，又灌了口水。

张老栓不再说话，把柴捆仔细挪回原处，遮严实了洞口，又洒了点灰土掩饰痕迹。

然后他像没事人一样，躺回自己那张吱呀作响的破木板床上，睁着眼，听着屋后山风吹过树梢的呜咽，和柴垛后那极其轻微的、啃食干粮的声音。

第二天，村里果然不太平。

王老爷府上的几个家丁，陪着一个穿皮甲、挎腰刀的士兵，凶神恶煞地挨家挨户盘问，有没有见过生人，特别是带伤的。

说是有一股流匪在这一带活动，上头有令，见者立即禀报，隐藏不报者杀头。

到了张老栓家，兵爷用刀鞘胡乱拨拉着空荡荡的屋角，眼睛像鹰一样四处扫视。家丁踢了踢屋角的柴堆，柴火哗啦响了一下。张老栓的心提到了嗓子眼，垂着头，腰弯得更厉害，脸上是惯有的木然和惶恐。

“老东西，见过生人没有？”士兵喝问。

张老栓连连摇头，舌头像是打了结：“没…没有，军爷……咱这地方，鬼都不来……”

兵爷嫌恶地看了他一眼，又瞅了瞅那一贫如洗、一眼就能望到头的屋子，料定这老棺材瓤子也没胆子藏人，更没油水可榨，骂骂咧咧地走了。

听着脚步声远去，张老栓才缓缓直起腰，走到灶台边，借着添柴的机会，轻轻拍了拍柴堆。里面什么回应也没有，但他知道人还在。

往后的两天，张老栓白天照常下地，侍弄他那几分贫瘠的山田；晚上回来，会多煮一把野菜汤，盛出满满一碗，放在柴堆旁边一个不起眼的破瓦罐里。第二天早上，碗总是空了，洗得干干净净放回原处。饼子也时会少一点。

两人没有对话，一种无言的默契在潮湿的柴垛后和寂静的茅屋里流淌。

直到第三天夜里，伤员似乎恢复了些力气，能够自己稍微挪动了。张老栓搬开柴火，看见他正试图包扎自己手臂上一道较深的伤口，但动作笨拙。

张老栓沉默地走过去，从他手里接过洗干净的破布条——那是从他唯一一件还算完整的旧褂子上撕下来的。老人粗糙得像树皮的手，动作却出奇地稳当，仔细地帮他重新捆扎好。昏黄的油灯光晕下，一老一少，谁也没看谁。

包扎完了，张老栓默默坐回床边，准备吹灯。一直没怎么开过口的伤员，忽然在黑暗中低声说了一句：“……多谢老丈救命之恩。连累您了。”

张老栓动作顿了一下，没回头，半晌，才用他那干涩的、没什么起伏的乡音，木木地回道：

“谢个啥。”

停顿片刻，像是在思索，又像是在陈述一个再简单不过的事实，他补充道：

“你们打了王老爷的粮仓，却没抢俺家的鸡。”

黑暗中，柴垛后的呼吸声似乎凝滞了一瞬，随即，传来一声极轻的、像是释然，又像是苦涩的叹息。

张老栓吹熄了油灯。月光从破窗棂漏进来，照在他布满皱纹的、木然的脸上。他不懂什么大道理，不知道什么“问道途”，也分不清那女侠是真是假。

他只知道，王老爷的人来，准没好事，不是要粮就是要命，连他最后一只下蛋的母鸡上个月都被抢走了。

而柴垛后面那个年轻人，和他们一伙的，抢的是王老爷的粮仓，分给了像他一样或许永远去不了灰谷镇的穷人。

这就够了。

又过了几日，张老栓睡下之后，只听得堂下窸窸窣窣。

次日，他扒开柴堆，柴堆后那人已经走了，留下了一锭银子，约莫有一两之数。

这大致相当于他辛辛苦苦劳作一年的收入。

\vspace{2em}

最大胆，也最根本的支持，则是对改变命运之可能的直接拥抱。

那些纸张被摩挲得毛边、字迹因反复传抄而甚至出现了些错误的《养气篇》入门法诀，开始在深夜的油灯下、在矿洞交接班的间隙、在远离人烟的破庙里，被秘密地实践。

疲惫到极点的矿工，在浑身酸痛中勉强按照法诀指引，凝神感知。

十天，半月，毫无所感，只有更深的疲惫与沮丧。但总有一些人，凭借稍好的天赋或钢铁般的毅力，在某个万籁俱寂的深夜，忽然于丹田深处，捕捉到一丝微弱却无比真实、缓缓流转的暖意！

这丝暖流无法让他们力能扛鼎，却切实地缓解了积年的劳损酸痛，带来一种久违的、对自身身体与精神的微弱掌控感。

这种体验，如同在无尽黑暗中瞥见的一缕微光，一旦看见，便再也无法忘怀。

它成了最具感染力的“福音”，在绝对信任的兄弟、挚友间心照不宣地分享、探讨。

虽然范围依然受限，保密极为严格，但它像渗入千年旱地的涓涓细流，坚定、缓慢而无孔不入地扩展着“问道途”的根系。

星星之火，渐成燎原之势，其影响开始触及更高、也更复杂的层面。

一些在守旧联盟庞大体系中处于边缘、备受排挤、或是看清了联盟涸泽而渔本质的中小家族、不入流的小门派，开始暗中活动。

他们或许尚不敢公然易帜，但通过曲折隐晦的渠道，向问道途递出了橄榄枝：在联盟征调物资的清单上做手脚，截留部分品质最好的“孝敬”给吴刚的渠道；对自己辖地内那些深夜练气的“异动”睁一只眼闭一只眼，甚至暗中提供些许庇护。

更为大胆的人，会将族中那些天赋平平、不受重视，却对“问道途”理念心向往之的边缘子弟，以“游历”、“访友”之名悄悄送走，最终这些人往往出现在沈瑶或吴刚的队伍名录中；甚至有的人心甘情愿地做起了问道途的间谍。

散修中愿意投入游击队者，也终于不再仅限于空有一腔热血的青年；许多饱经世故、看透联盟将散修纯粹视为“耗材”本质的老江湖，在吴刚精明的利益折算和对散修心态入木三分的把握下，也做出了选择。

他们相信，在“问道途”这边，至少能被当作一个“人”来对待，流的血、出的力，能看到回报，更能赢得些许尊严。

沈瑶的游击队里，多了擅长在山林间追踪索迹的老猎户；吴刚的地下党中，加入了曾经走南闯北、熟悉各地方言与人情世故的“路路通”。

得益于此，一张无形却无处不在的耳目之网，在守旧联盟自以为严密的统治区内，悄然织就。

\vspace{2em}

情报战的态势，发生了根本性的、对联盟而言堪称灾难的逆转。

过去，守旧联盟的探子绞尽脑汁试图潜入云霞山、却往往有去无回；如今，在广大的乡村野地、散修市镇，一张无形而致密的耳目之网已然织就。

“青芒山一带，近日有残党啸聚，滋扰乡里，劫掠税赋。着令你部，于三日后卯时正，自黑水河西进，兵分两路，扫荡青芒山北麓及东侧河谷地带，务求全歼，以儆效尤！”

紫阳宗外门长老王猛将盖好印信的军令递给身旁身着文士衫、面相精明的师爷，沉声道：“速去誊抄三份，一份存档，一份送达黑石集驻军李统领，另一份快马传至青石城备案。此番必要肃清匪患，以振军威！”

“长老放心，属下即刻去办。”师爷恭敬地接过军令，躬身退下。

他回到自己那间堆满卷宗、弥漫着墨香与些许霉味的值房，掩上门。

脸上那副恭敬谨慎的神情如潮水般褪去，取而代之的是一种深沉的疲惫与一丝不易察觉的讥诮。

他展开军令，目光迅速扫过关键信息：三日后，卯时，黑水河西进，分两路，北麓与东侧河谷。

他先提起笔，在一张用于起草文书的草稿纸上，随意地记录了几个数字和简略方位。

这些符号混杂在其他无关的账目数字与地名缩写之中，即便被人看见，也只会当做寻常的工作笔记。

然后，他将这张草稿纸折好，塞进袖中一个特制的夹层。

接着，他才开始正式誊抄军令。笔走龙蛇，一丝不苟。

他写得很快，因为心中的紧迫感在催促。第一份存档，放入标有“乙未年七月剿匪”的卷宗匣内。第二份，装入专用的硬皮令筒，用火漆封好，盖上印鉴。

他唤来一名在门外候命的年轻传令兵。“将此令送至黑石集驻军李统领处，不得有误。”

“是！”传令兵接过令筒，转身便要走。

“且慢。”师爷叫住他，从案头拿起一个水囊递过去，“天热路远，带上这个。王统领脾气急，送达后莫要多言，速回复命即可。”

“多谢师爷！”传令兵感激地接过。

望着传令兵远去的背影，师爷坐回椅中，闭上眼，揉了揉眉心。

他又想起那个神秘的“行商”与自己“偶遇”时说的话：“先生大才，却屈居于此，替盘剥乡里、视人命如草芥之辈刀笔刑名，不觉得明珠暗投么？问道途所求，不过是让有才者得其用，有力者得其酬，这世道，不该是某些人生来就高高在上，而大多数人连抬头看一眼天空的机会都没有。”

当时他嗤之以鼻。但自己随后特意关注的一些关于联盟内部倾轧、以及云霞山内部分配制度的零星信息，却与他暗中观察到的一些迹象隐隐吻合。

更重要的是，对方也没有要求他做什么，只是说：“先生不妨冷眼旁观，看看这天下大势，究竟流向何方。选择全在先生身上。我等可不会逼迫先生。”

\vspace{2em}

命令还在路上，情报却已通过多条彼此不知晓的平行渠道，开始向地下网络的节点流淌。

黑石集，那个总是一脸和气的酒馆老板，正在柜台后擦拭酒杯。

一个风尘仆仆、贩运兽皮的货郎走进来，要了碗最便宜的粗茶，坐在角落默默喝着。货郎离开时，指尖似乎无意地在油腻的木桌上划过。老板过来收拾，目光扫过桌面，那里有几个茶渍水痕形成的、极其短暂的凌乱符号。

抹布一抹，符号消失。

但在打烊后，他进入后厨，从灶台旁一块松动的砖石后，取出一枚小小的玉简，将脑海中记忆的符号形状，以微弱的精神力刻画进去。

这枚玉简，会在明天清晨，由他那个恰好要去邻镇采买鲜货的侄儿带走。

几乎同时，在黑石集通往山外的岔路口，一个卖山货的农妇正在整理她的篮子，底层是沾着泥土的草药，上面盖着些野果。

一个看起来像是收药材的贩子蹲下来，挑拣着草药，低声讨价还价。最后生意没谈成，贩子摇摇头走了。

农妇也嘟囔着收拾篮子，但在无人注意时，她快速地从篮子底部的夹层里，抽出一根中空的细芦苇杆，里面卷着一条窄窄的布条，上面用炭笔写着歪斜的记号。

她不识字，但她知道要把这玩意传给下一个人。

这些零碎的情报，如同汇入溪流的无数水滴，通过吴刚精心编织、已愈发细密坚韧的网络，被迅速汇集、筛选、交叉验证。那个师爷的草稿符号、酒馆老板的玉简、农妇的布条……

它们从不同的社会阶层、不同的地理位置发出，沿着不同的安全路线流动，最终在某些绝对安全的隐蔽节点完成拼接和相互印证。

不到一天时间，一份完整、准确的敌情通报，已经摆在了沈瑶和吴刚的面前。

不止有清剿的时间、兵力、路线，甚至还包括了带队军官的性格特点、部分部队近期的士气状况。

\vspace{2em}

三日后，卯时。

黑水河西，紫阳宗外门弟子李统领骑着高大的龙血马，看着麾下三百余名修士和辅兵组成的队伍，在晨雾中列队完毕，刀枪映着微凉的晨光。

他踌躇满志，这次若能一举剿灭青芒山的“残党”，可是大功一件。

“出发！按计划，一队随我走北麓，二队走东侧河谷！仔细搜索，不得放过任何可疑痕迹！”王统领挥手下令。

大军开拔，脚步声、马蹄声、铠甲碰撞声打破了山林的寂静。

他们按照命令，仔细地搜遍了北麓的每一片林子，东侧河谷的每一个山洞。然而，除了发现几处早已废弃的临时营地痕迹、以及一些疑似有人活动过但已无从追查的细微线索外，一无所获。

连个“残党”的影子都没看到！

王统领的脸色从最初的踌躇满志，逐渐变得阴沉，最后涨成了猪肝色。

“怎么可能？！难道他们能未卜先知？！”

他愤怒地咆哮，心中却升起一股寒意。

这种一拳打在空处的感觉，比惨烈的战斗更让人憋屈和不安。

而就在李统领的大军在青芒山徒劳搜捕的同时，距离青芒山一百五十里外，一处隶属于丹鼎阁某附属家族的小型物资中转站，迎来了灭顶之灾。

这里存放着近期从附近几个矿点征收上来、尚未运走的低阶矿石和一批药材。因为原本驻扎于此的部分护卫被抽调去参与青芒山的“大行动”，守备力量比平日薄弱了三成不止。

留守的管家和几名护卫修士还在抱怨自己被留下看仓库是“没油水的差事”。

午时刚过，正是人最困乏的时候。数道如同鬼魅般的身影从仓库围墙的阴影中、从堆叠的货箱后悄无声息地闪现。

刀光如冷电，瞬间割断了哨兵的喉咙；沈瑶的身影如同融入阳光的幻影，直接出现在那名正在打盹的灵枢境护卫头领面前，在他惊骇睁眼的刹那，冰冷的锋刃已洞穿其护体灵光。

战斗短促而激烈，完全是一场不对等的屠杀。

负隅顽抗者被迅速格杀，其余人见势不妙，立即丢下武器投降。

沈瑶也没有管这帮家伙，只是让游击队员看着这群人，不让他们乱动。

迅速打开仓库，将里面堆积的矿石和药材能带走的带走；带不走的，连同仓库建筑，付之一炬。

冲天而起的黑烟，方圆数十里都清晰可见。

在撤离前，让队员将几袋便于携带的粮食，故意“遗落”在附近一个以贫苦闻名的村子外。

\vspace{2em}

广泛的群众基础和重新构建起的活跃的情报网络，使得游击战术的精髓得以淋漓尽致地发挥。

沈瑶手下精干的队伍，可以随时拆分成三四个更小的作战单元，如同水银泻地，同时袭扰数条补给线上的多个节点，焚毁粮草，破坏关键桥梁或运输法器，令敌军顾此失彼。

而沈瑶本人，则择机聚合精锐，执行高价值目标的斩首；那些尤其酷烈、贪婪、“没有团结价值”（吴刚语）的中低层军官，往往在某个酒醉的夜晚或巡视的路上，便被一道不知从何处而来的流光夺去性命。

这些行动，单次看来或许无法歼灭敌军成建制的主力，无法夺取一座城池，但其累积效应是恐怖而致命的。

守旧联盟漫长的补给线变得千疮百孔，从后方运往前线的物资损耗率急剧飙升，护送任务成了人人避之不及的死亡差事。

军队的士气，再难用简单的劫掠许诺和“必胜”口号来维持，厌战与恐惧的情绪在底层士兵和低级军官中如瘟疫般悄无声息地蔓延。

更让坐镇中军的陈炎焦头烂额的是，为了扑灭后方的“烽烟”，保护这些维系前线存在的“血管”，他不得不一次又一次地从围攻云霞山的主力中，分派精锐部队回援。

这些部队疲于奔命，穿梭在陌生的乡村山野，要么找不到神出鬼没的游击队，空耗钱粮士气；要么一脚踏入精心设计的伏击圈，损兵折将，狼狈不堪。

云霞山正面战场所承受的压力，就在这敌后持续不断的、精准的放血战术下，得到了切实而显著的缓解。山上的将士们发现，敌军的攻势间隔变长了，强度也有所减弱；而山下敌营中，因伤兵增多、补给不继而引发的抱怨与混乱，时有耳闻。

陈炎的军队自然也不是铁板一块：屡不建功、山下干耗许久，使得王丹和刘长靖对陈炎也失去了信心。两人对陈炎也颇有微词起来，守旧联盟内部也第一次出现了裂痕。

正是在这种背景下，陈炎在又一次接到后方某重要粮仓被焚、押运队全军覆没的急报后，终于暴怒而无奈地做出了战略调整。

他脸色铁青地打断了帐中关于下一次试探性进攻方案——当然也包括“究竟要不要打”——的争吵，嘶声道：“云霞山，不打了！传令，主力分批回撤！”

他不得不承认，继续将主力钉死在云霞山下，已是徒劳。

他计划逐步抽离主力部队，只在山下保持足够的兵力维持包围态势，防止山上守军大规模突围。

撤回的兵力，一部分用于加固几条核心补给线的防护，打造成“铁桶”通道；另一部分则组成数支机动精锐“黑旗营”，派往那些报告显示“匪患最烈”、“民心最浮”的区域，进行严厉无情的清洗与绥靖，意图用铁血手段扑灭所有反抗火苗，重新稳固后方。

与此同时，他试图以战利品分配优先权、未来地盘划分等空头许诺，拉拢那些在地方盘根错节的豪强家族，希望借助他们的“本土智慧”和人脉，来分辨良莠，协助维持基层。

然而，这剂猛药非但没能治病，反而成了最烈的毒药。

\vspace{2em}

那些在地方上经营数代、与守旧联盟中央嫡系本就存在利益纠葛与矛盾的中小家族，对陈炎这迟来的、充满功利色彩的“拉拢”看得一清二楚。

在一座小城，城主刘家接到了配合军队清剿境内“问道途暗桩”的严令。

刘家家主，一个面白无须、眼神精明的中年修士，恭敬地送走了杀气腾腾的联盟特使，转身回到密室，脸上的笑容瞬间消退。

“让我刘家去当这个恶人，替他紫阳宗趟雷？”他冷笑一声，对几位族老道，“问道途的人是那么好抓的？他们神出鬼没，在穷鬼里名声还好。我们真下死力去抓，抓得到几个不好说，这十里八乡的怨气，可全得记在我刘家头上。到时候，他紫阳宗拍拍屁股走了，我们还得在这片地上过日子！”

另一位族老捻着胡须，阴恻恻地补充：“正是。况且，谁能保证这问道途将来成不了气候？我看那叶牧谦，不是池中之物。现在把事做绝，将来何以自处？”

于是，刘家开始阳奉阴违起来。

他们派出的协查队伍，多是些老弱惫懒的家丁，每日在几个固定的、早就空无一人的废弃矿洞或山林外围转悠一圈，便算交了差。

偶尔上面催得急了，他们便去佃户里寻两个平日最不服管教、或是家里无力缴纳足额租子，做事还不老实的刺头，扣上“通匪”的帽子，五花大绑送去交差。

至于那些真正可能藏有伤员的村落，或者与吴刚网络有零星贸易往来的散修集市，刘家的人马总是“恰好”错过，或者“收到风声太迟”。

在更北边的平昌镇，地方豪强郝氏的做法则更显“高明”。

他们一方面大张旗鼓地配合联盟军队设立关卡，盘查行人，做出尽心竭力的姿态，甚至主动“贡献”出家族控制的几处偏远山庄，作为联盟军队的临时驻地。

但是另一方面，郝家族长却秘密吩咐心腹管家，通过绝对可靠的渠道，向吴刚的地下网络传递了一份经过筛选、不涉及核心利益的“清洗计划时间表”，同时附上一句：“迫于形势，虚应故事，万望海涵。贵方活动，请避开地图上划线两处。”

他们是在两头下注，用最小的代价，敷衍最大的凶险，同时为自己留一条后路。

这种“出工不出力”，抓替罪羊敷衍了事的现象，在非联盟嫡系的地方势力中成了普遍潜规则。

陈炎试图依靠的“本土秩序维护者”，首先维护的是他们自身的存续与利益，而非联盟那遥不可及且充满不确定性的“大业”。

这也是非常正常的。

哪个年代都不缺两头下注的骑墙派！

但这种骑墙派确实给问道途降低了巨量的压力！

\vspace{2em}

而那些被陈炎寄予厚望、授予全权的联盟嫡系机动精锐，则彻头彻尾地化为了噩梦的代名词。

这些“机动精锐”大多由紫阳宗失意弟子、依附家族的亡命徒、以及此前在正面战场或因剿匪不力受过处分、急于戴罪立功的军官组成，凶残暴戾，且被赋予“便宜行事”、“毋枉毋纵”的尚方宝剑。

他们面对的，也是陈炎和后方指挥部根本无法理解的现实困境：在问道途群众路线深耕已久的区域，“匪”与“民”的界限早已模糊不清，甚至浑然一体。

几乎家家户户都曾受过问道途小队或明或暗的恩惠——或许是深夜悄悄放在门前的半袋粮食，或许是家人重病时“恰好”路过游方郎中留下的药方，又或许只是在联盟税吏鞭打老人时，有陌生人站出来说了一句公道话。

几乎每个村落，都有年轻人偷偷修炼那本广为流传的《养气篇》，在体内生出第一缕微弱的、属于自己的元气。

对于这些手持利刃、背负军令、且本身就怀着抢掠发财欲望的“清剿”部队而言，分辨谁是“问道途”，成了不可能完成的任务。

部分人倒是能反应过来、悬崖勒马；但更多红了眼的人，认为分不清就不分了。

“传令麾下四王子，破城不须封刀匕！”

“宁可错杀一百，不可放过一人！”

“清洗”迅速、且必然地滑向了最野蛮、最血腥的深渊。

“黑旗营”开进村落之日，便是地狱洞开之时。

士兵们以搜捕“残党”、“收缴违禁功法”为名，如狼似虎地砸开每一扇门。

他们不在乎证据，只凭“眼神闪烁”“家中存粮稍多”“搜出写有不明字迹的纸片”等莫须有的理由，便将户主拖出，当着全家老小的面严刑拷打，逼问“同党”。

惨叫与哀求声响彻村落，但这只是前奏。

真正的灾难随之降临。

抢掠成了公开的行动。士兵们哄抢一切看得上眼的财物——不多的铜钱银角、稍微像样的衣物、甚至锅碗瓢盆。稍有反抗，便是刀剑加身。

村里的几个姑娘被拖进柴房，哭喊声很快被淫邪的笑声淹没。

一个老人仅仅是因为护着家里唯一的一只下蛋母鸡，便被当胸一刀刺倒。

最后，为了“以儆效尤”，也为了掩盖暴行，带队的校官下令点燃了几处被指认为“可能藏匿匪类”的房屋。

冲天的火光吞噬了屋舍，也吞噬了里面来不及逃出或不愿离开祖宅的老人。浓烟滚滚，混合着焦糊与血腥的气味，笼罩了这个曾经宁静的村落。

同样的场景，在多个被划为“重点清洗区”的村镇反复上演。

联盟军队带来的不是他们承诺的“秩序”，而是比传闻中任何“匪患”都更加无常、更加恐怖的噩梦。

在村民眼中，这帮疯子根本不是来平息动乱的官军，而是比盗匪更凶残百倍！

所过之处，十室九空，哀鸿遍野。

这种无差别的恐怖高压，在短期内确实取得了残酷的“效果”：不加区分的开杀，虽然大多数是错杀了，但也肯定有能歪打正着抓到几个真为问道途效力的。

一批忠诚而富有经验的联络员被捕后公开处决，首级被悬挂在城头示众。

沈瑶的游击队被迫化整为零，潜入更深的深山老林，活动几乎停滞，数次在针对性的围追堵截中险死还生，队伍减员严重。

吴刚的地下网络也遭受重创，许多线人被迫沉默或转移，情报传递变得异常艰难迟缓。

从表面战果看，陈炎的“清洗”似乎一度扼住了问道途敌后活动的咽喉，扑灭了若干“火焰”。但这恰恰是历史辩证法最无情的展现：高压，往往催生更猛烈、也更隐蔽的反抗。

表面的火焰被扑灭，灰烬之下，却是被烧得更透、更炽热、蔓延更广的草根与土壤！

粗暴野蛮的清洗，完成了问道途宣传队永远无法做到的、最彻底最深刻的动员。它将问道途这颗“杂草”的种子，用血与火作为肥料，更深、更牢地埋进了几乎每一个亲历者、乃至听闻者心底。

对于受过“官兵”劫掠的幸存者而言，当联盟士兵的刀砍向护着母鸡的老人、火把扔向祖屋时，所有的犹豫、恐惧和对“官府”残存的、惯性般的敬畏，都在那一刻燃烧殆尽了。

那个曾经只在夜间偷偷流传的“叶先生”和“问道途”，不再是一个模糊的符号或值得同情的“反贼”，而成了血海深仇中唯一的指望，活下去并寻求复仇的唯一可能路径。

甚至那些原本只是沉默忍受、内心隐隐偏向“秩序”——哪怕是不公的秩序——只求安稳度日的富农、小地主们，也在风暴中被席卷。

联盟军队并不总能准确区分“匪”和“民”——甚至更多的所谓“黑旗营官兵”已经不屑于去分辨了。

在抢红了眼的状态下，“通匪”的帽子可以扣给任何人，只为夺取其积蓄的财物。

当“秩序”的化身亲自演示了何谓无法无天的豺狼行径时，打破这“秩序”的“匪”，在百姓眼中便有了截然不同、甚至带上了悲壮与正义色彩的含义。

五帝三皇神圣事，骗了无涯过客；有多少风流人物？

盗跖庄蹻流誉后，更陈王奋起挥黄钺！

\vspace{2em}

民心的转变，不再是悄然的流淌，而是在绝望与暴政的熔炉中，完成了彻底、决绝的淬炼与逆转；也正是在这血腥的镇压与更顽强、更智慧的隐蔽反抗中，一种全新的、守旧联盟高层思维完全无法理解的战略格局，瓜熟蒂落，悄然成形。

云霞山，作为问道途不可撼动的战略支点、精神灯塔与正面战场，依然巍峨矗立，牢牢吸引并消耗着联盟剩余的注意力和部分兵力。

而广袤无垠的乡村野地、散修聚集的市镇、远离核心统治的中小城池——这些守旧联盟统治链条末梢、力量相对薄弱的区域——却在血腥的洗礼后，于无声处完成了人心的皈依，变成了问道途的海洋，革命的根据地初见雏形。

这里没有问道途的大军驻扎，没有飘扬的显赫旗帜，却有着无数双自愿为问道途守望的眼睛，无数张自愿为问道途传递消息的嘴，无数颗在黑夜中默默祈祷云霞山胜利、期待“叶先生”理念成真的心。

诉苦大会在深夜里秘密召开，红黑账本在暗地里严谨记录，互助性的、脱离守旧联盟掌控的自治组织在废墟上萌芽。

尽管守旧联盟的正规军依然控制着主要城池和交通要道，但他们惊恐地发现，一旦离开城墙的庇护，进入广大的乡村，他们就像滴入大海的墨点，被无尽的沉默与潜在的敌意所包围。

征粮变得异常困难，巡逻队时常遭遇冷箭或陷阱，甚至整队人马在熟悉的道路上诡异地“迷路”失踪。

联盟的军队可以凭借武力占领或摧毁任何一座他们看得见的城池，可以剿灭任何一支他们能追踪到的游击队，但他们永远无法占领每一寸土地，无法控制每一颗人心。

他们仿佛陷入了一片无形的沼泽，看似站在坚实的城池据点之上，实则举步维艰，因为承载他们的整个大地，已然悄悄改变了颜色。

农村包围城市——这一超越纯粹军事范畴、深植于社会土壤与人民意志的战略格局，在血与火的残酷浇灌下，在希望与绝望的反复淬炼中，终于第二次脱离纸面的构想，第一次走出历史的记载，而成为了内域大地上活生生的、呼吸着的现实。

上一次走向现实，还是叶牧谦的前世。

尽管敌后战场因残酷清洗而满目疮痍，付出了鲜血与牺牲的代价，但最深层的、决定战争最终走向的力量对比，已经发生了不可逆转的根本性变化。

陈炎的退兵与清洗，非但不是挽回颓势的妙手，反而成了加速这一历史进程的最后、也是最有力的推手。

他所面对的，早已不再是一个单纯的军事对手叶牧谦，而是一个从旧世界痛苦母体中孕育诞生、已深深扎根于亿万民众渴望之中的全新体系与浩荡潮流。

\vspace{2em}

在守旧联盟统治区域的东北部，有一座重要的中等城池“抚远”。

此城连接数条重要商路，是向东北方向诸多前哨、矿点转运物资的关键枢纽，城高池深，阵法齐全。常驻于此的超过五百名守军，由紫阳宗、丹鼎阁及其附属宗门的修士混编而成，统兵城守更是紫阳宗一位资历不浅、修为达到玄脉境融贯期的修士，赵镇岳。

在坐镇中军、遥控四方的陈炎看来，叶牧谦及其主力被死死钉在云霞山，抚远城地处后方，守备充足，将领可靠，乃是绝无可能出问题的稳固支点，是供应前线血脉上不可或缺的一环。

他忽略了一个朴素至理：再高的城墙，再深的护城河，终究是由一块块砖石垒砌，由一个个活生生的人驻守。而砖石会风化，人心，更会变。

赵镇岳年近五旬，面容方正，眉宇间带着经年军旅生涯留下的风霜与刻痕。他的修为在玄脉境融贯期停滞多年，大道前路已渺，这才被宗门派来担任这安稳、职责分明但远离权力核心的城守之职。他性格里有旧式军人的一丝耿直乃至迁腐，对麾下兵卒和城中百姓谈不上多么体恤仁慈，但至少遵循着一套在他看来天经地义的规矩：该发的饷银尽力筹措，该收的税赋照章办理，不似某些同门那般，将手中权柄视为肆意盘剥、中饱私囊的工具。

正因这份尚未完全泯灭的底线，他对近年来宗门——尤其是陈炎一系——为推动战争而越发酷烈无情的举措，内心积郁着越来越多难以言说的微词。他亲眼见过附属宗门的弟子如何借搜查“问道途奸细”之名，行敲诈勒索散修之实；见过丹鼎阁的商队如何凭借垄断地位，肆意压低本地药农的收购价，逼得人家破人亡。更让他夜深人静时感到隐隐不安的，是麾下那些出身寒微的士卒眼中，日益浓重、几乎化不开的麻木，以及麻木之下偶尔闪过的、对不公命运的愤懑之火。

转变的契机，始于半年前一次例行公事般的剿匪。

有探报称，一伙问道途乱党流窜至抚远城西面的莽莽山林，劫掠了一支运送“慰劳物资”前往某处前哨的小型车队。赵镇岳奉命率两百士卒出城清剿。

战斗过程近乎一场闹剧，那伙“乱党”人数寥寥，战力不明，甫一接触，便远遁山林，只留下些无关紧要的杂物和早已冰凉的篝火痕迹。

就在赵镇岳心中暗骂情报不准、准备收兵回城时，麾下一名负责侧翼侦查的什长——一个名叫韩铁、同来自紫阳宗、平素沉默寡言但办事稳妥的青年——带着两名手下，押着一个受伤被擒的年轻人来到他面前。

“大人，抓到一个落单的，”韩铁的声音压得很低，带着一贯的沉稳，“看其服饰举止……不像寻常山匪流寇。”

赵镇岳目光扫向那被俘者。

对方很年轻，约莫二十出头，衣衫是寻常散修穿的粗布材质，却浆洗得干净整洁，身上有几处新添的皮肉伤，鲜血浸染了布料，但其神情却不见多少惊慌失措，尤其那双眼睛，清澈亮澈，在逆境中仍保持着一种奇异的沉静。

最让赵镇岳心下惊疑的是，以他玄脉境的感知，竟在此人身上察觉不到丝毫灵枢径修士特有的灵力波动，反而有一种绵长沉凝、迥异于此世主流却又生机勃勃的气息，在其体内隐隐流转——这与宗门卷宗里对问道途修士那模糊却特征鲜明的描述，隐隐吻合。

赵镇岳挥手屏退左右亲兵，只留下似乎完全可信的韩铁在场，沉声喝问：“你是问道途的人？”

年轻人抬手擦去嘴角渗出的血丝，动作不见狼狈，坦然迎上赵镇岳审视的目光：“是。在下林河，云隐别院记名弟子。”

“好大的胆子！”赵镇岳须发微张，玄脉境修士的威压自然流露，“区区记名弟子，也敢劫掠我紫阳宗的车队？！”

林河闻言，却轻轻笑了起来，那笑容里没有猖狂，反而带着一种让赵镇岳颇为刺目的了然与讽刺：“车队？

“大人说的是那几辆满载着‘醉仙酿’、‘灵犀脯’，专供前方大营某几位长老享用的车队么？”

他顿了顿，目光变得更加锐利，“赵城守，前方将士在云霞山下日日浴血，听闻连最基础的疗伤丹药都时常短缺，这些所谓的‘慰劳品’，究竟是去慰劳谁的？”

赵镇岳喉咙一哽，竟一时语塞。

车队里装的是什么，他岂会不知？

宗门上层某些人的奢靡做派与前线艰困形成的刺眼对比，早已不是秘密。

但他身处其位，有些事看破不能说破。

他只能把脸一板，厉声道：“巧舌如簧！你们这些乱党，四处煽风点火，破坏后方安定，才是致使生灵涂炭的罪魁祸首！”

林河缓缓摇头，语气依旧平静，却字字如有千钧：“我们从未煽动任何人。我们只是将血淋淋的事实，将另一种活法的可能，摆在了那些被压迫、被蒙蔽的人面前。

“赵城守，你坐镇抚远多年，耳目通达。你可曾听闻，我问道途义军所到之处，主动劫掠过一户平民，欺辱过一名无辜散修？可曾见过我们攻占某地后，如贵方某些军爷那般，纵兵抢掠三日，或是强征暴敛、竭泽而渔？”

赵镇岳再次陷入沉默。

他确实听到过一些截然不同的风声。从那些商人、游侠零散的交谈中，从一些难以完全封锁的渠道里，他隐约得知，问道途义军行事极有章法，只针对联盟官仓和死硬分子，对平民百姓甚至堪称秋毫无犯，有时还会将缴获的粮食物资分发给穷苦人。这与宗门宣传中那烧杀抢掠、无恶不作的邪魔外道形象，相差何止万里。

见赵镇岳脸色变幻，林河知道话语已触及对方心防，便说出了那直指本心的一句：“《问道戒律》开宗明义，‘仰不愧于天，俯不怍于人’。我等所求，无非是开辟一条道路，让无数无根无基的散修，让终年劳作却不得温饱的百姓，也能凭借自身的汗水与毅力，存续性命，活得有几分人的尊严。”

林河的目光深深看进赵镇岳眼中，“请您扪心自问，在这世道，在这抚远城内城外，还有多少人，能堂堂正正地说自己活得‘不愧不怍’？您自己呢？”

“仰不愧于天，俯不怍于人”！

这十个字，如同一声暮鼓晨钟，又似一根烧红的钢针，狠狠刺入赵镇岳灵魂深处某个早已被世俗尘埃覆盖、却从未真正死去的角落。

他想起了年少时初入山门，也曾怀揣过斩妖除魔、护卫一方的热血梦想；想起了这些年来目睹的越来越多无力改变的龌龊与不公；更想起了麾下那些朴实士卒，在收到家书得知赋税又涨时，眼中那无法掩饰的忧虑与对前途的茫然。

复杂的情绪在胸中翻涌，最终，赵镇岳没有回答林河的问题，只是有些疲惫地挥了挥手，对侍立一旁的韩铁道：“给他包扎伤口，押回城中……单独关押于僻静处，不准用刑。等我亲自发落。”

回到抚远城，赵镇岳的心绪再也无法平静。

几日后的一个夜晚，他屏退左右，再一次独自去见了林河。

这一次，林河没有慷慨激昂地陈述理念，反而像一位关心民生疾苦的访客，细致地问起了抚远城的现状：今年粮价几何？各类税赋名目有多少，底层可能承受？守城士卒的饷银能否按时足额发放？城中散修坊市的生意是否兴旺，可有受到盘剥？

这些问题，有的赵镇岳清楚，有的却只能含糊以对——他虽为城守，但对底层具体生态的了解，远不如一个真正深入其中的观察者。

林河便不急不缓，向他讲述起其他类似规模城池，在守旧联盟治下与在问道途影响下的真实对比。数据详实，事例具体，何人因何破产，何处因何得以喘息，言之凿凿，逻辑清晰，由不得赵镇岳不信。

他第一次如此直观而深刻地意识到，问道途和他们一样，同样有着清晰的理念、严明的纪律，甚至是对未来如何治理一方、改善民生的具体构想。相比之下，宗门上层某些人穷兵黩武、竭泽而渔的做法，显得如此短视、粗暴，且不得人心。

与此同时，赵镇岳察觉到，自己颇为信任的部下韩铁，在汇报看守情况或谈及城中琐事时，也常常会“无意”地透露一些信息：某个老兵的家眷因交不起新加的“防务捐”被税吏打了；某个相识的散修朋友最近精神头好了不少，私下透露在练一套“奇特的养气法门”；城中酒肆茶楼里，对宗门久攻云霞山不下、反而赋税越来越重的抱怨声，正在悄然滋长，虽不敢明目张胆，却如地底暗流，涌动不休。

赵镇岳不是愚钝之人。他渐渐明白，韩铁恐怕也早已被渗透，甚至可能就是对方早早埋下的一枚棋子。

但他没有点破。

一方面，韩铁行事依旧稳妥可靠，并未做出任何损害他本人或城池防务的举动；另一方面，那些通过韩铁之口传来的消息，虽然刺耳，却往往是真实存在的、被他有意无意忽略的城中疾苦。一种微妙而奇异的默契，在赵镇岳、林河、韩铁这三个身份迥异的人之间，悄然形成。

导火索在三个月后被点燃。陈炎为应对日益猖獗、防不胜防的“后方袭扰”，决心强化清剿，命令各后方城池加征一批紧急物资，并抽调部分守军补充前线损耗。命令传到抚远，要求赵镇岳抽调一百名修士即刻前往报到，同时上缴城库中灵石储备的三分之一、疗伤丹药的一半。

这道命令，成了压垮赵镇岳心中那架摇摆天平的最后一块巨石。

抽调百名修士，几乎等于抽走了抚远城三分之一的机动防御力量，城池安全顿时空虚；而上缴大量关键物资，不仅动摇了守城的物质基础，更意味着本已紧张的城内储备将急剧缩水，随之而来的必是物价飞涨，民不聊生，怨声载道。

赵镇岳在城守府的书房内独坐了一整夜。案头那纸调令冰冷而沉重，窗外是抚远城沉睡的轮廓。林河关于民心向背的论述，韩铁平日透露的点点滴滴，宗门上层的苛刻短视，与这座他驻守多年、虽无深情却已有责任的城池可能面临的混乱破败未来，交织在一起，反复灼烧着他的理智。

天光将亮未亮之际，书房内的烛火燃尽了最后一滴油脂，悄然熄灭。在黎明前最深的黑暗中，赵镇岳抬起了头，眼中最后的犹豫已然散去，取而代之的是一种孤注一掷的决绝。

他秘密召来了韩铁。没有迂回，没有试探，声音因一夜未眠而沙哑，眼神却亮得惊人：“联系你们的人。”

韩铁身躯微不可察地一震，瞬间明白了这句话的含义，眼中爆发出压抑已久的、惊人的神采。

但他依然强行克制住激动，确认般低声问：“您指的是……紫阳宗方面？”

“当然不是紫阳宗。”

赵镇岳斩钉截铁，随即清晰地道出了自己的条件，“抚远城……可以易帜。但我有三个条件：第一，不得滥杀任何投降的守军及城中无辜百姓；第二，城中所获库藏，必须优先用于安抚民生、救治伤患；第三，我赵镇岳个人之生死荣辱，无足轻重，但请务必放过我家中老小，保他们平安离去。”

韩铁再无犹豫，单膝跪地，抱拳行礼，声音因激动而微颤：“赵城守深明大义！您所提条件，正是问道途一贯秉持之道义所为！韩铁以性命担保，必即刻将城守大义与条件，如实禀报白先生、沈姑娘！”

联络通过韩铁掌握的、深埋在敌后网络中的隐秘渠道，以最快速度传递出去。五日后的深夜里，一只毫不起眼、却附有特殊灵力印记的纸鹤，穿过层层阵法防护，悄然落入赵镇岳手中。他展开纸鹤，上面只有一行蝇头小楷，墨迹沉稳：

“七日后，子时三刻，东门。白、沈。”

赵镇岳将这短短一行字反复看了数遍，仿佛要将每一个笔画都刻入脑海。然后，他指尖腾起一缕火焰，将纸鹤点燃，看着它化为灰烬，飘散在夜风中。他知道，退路已断，从这一刻起，他要么与这座城一起获得新生，要么便坠入万劫不复的深渊。

接下来至关重要的七天，赵镇岳仿佛在刀尖上行走。

他必须表现得一切如常，甚至更加勤勉。

面对宗门派来催促物资和兵员的使者，他做出了恰到好处的为难与抱怨，以“清点库藏需要时间”、“士卒调动涉及防务交接”等理由，成功地将时间拖延了下来。与此同时，他暗中以“整备防务”、“预防奸细”为名，进行了一系列至关重要的布置：将库房中部分关键物资转移到更隐秘或更易控制的地点；将几个要害岗位的守卫，尤其是东门附近的哨位，逐步换上了经过韩铁筛选、绝对可靠或已被争取过来的士卒；对少数几个紫阳宗、丹鼎阁的死忠分子及其掌控的小队，则拟定了详细的监控与制伏计划。

整个抚远城，表面上依旧秩序井然，运转如常，但在平静的水面之下，一股决定城池命运的暗流，已经蓄势待发。

第七日，夜幕早早降临，厚重如墨，无星无月，正是改天换地之时。

子时将近，抚远城仿佛沉入酣睡，只有城墙上的气死风灯在夜风中孤零零地摇曳，昏黄的光晕照亮了值守士卒们略带困倦的脸庞。东门城楼上，今夜的值守官，正是什长韩铁。他像往常一样，带着十余名亲信士卒，神情严肃地巡视着城防。

时间一点一滴流逝，指向那个约定的时刻。

子时三刻整。

韩铁脚步沉稳地走到城楼一角，那里看似随意地堆放了一些用于夜间照明和烽火传讯的干燥柴薪、油脂等引火之物。他警惕地扫视四周，确认无人异常关注此处，随即迅速从怀中掏出火折，迎风一晃，幽蓝的火苗窜起，被他果断地引向柴堆。

“嗤啦——！”

干燥的物料遇火即燃，火舌猛地蹿起，越烧越旺，橘红色的烈焰在漆黑如铁的夜幕映衬下，显得格外刺目而耀眼！

“走水了！东门楼走水了！”几乎是立刻，城头响起了“惊慌失措”的喊叫，警锣被奋力敲响，刺耳的“铛铛”声划破夜的寂静，迅速传遍附近城墙。城头守军一阵骚动，附近几座哨塔的视线和注意力，全被这突如其来的“意外”火灾吸引了过去。

然而，这烽火，并非意外，而是号令！

几乎就在东门楼火光亮起的同一刹那，在抚远城外数里处一片早已埋伏多日的密林中，数百双在黑暗中依旧清亮的眼睛，同时睁开！

以白言冰、沈瑶为首，吴刚、陈默等核心骨干协同，超过三百名历经血火淬炼的问道途精锐游击队战士，以及闻讯响应、自愿前来参战的附近区域散修、乃至部分被成功策反的小家族私兵，总数近五百人，如同蛰伏已久的猛兽，骤然跃出藏身之地！

急促而整齐的脚步声、兵刃与轻甲摩擦的沙沙声，汇成一股低沉却充满压迫感的洪流，向着洞开的城门方向（他们确信那里将会洞开），发起了无声而迅猛的冲锋！

队伍最前方，白言冰一袭朴素的青色劲装，手握长剑，身影如电，眼神沉静如古井寒潭，不见波澜，却蕴含着劈开一切阻碍的决心。

紧随其侧的沈瑶，黑衣如墨，身形在夜色中几乎融入背景，唯有那双眼眸，在疾驰中闪烁着千幻流光般冰冷而锐利的光芒，仿佛已锁定了城中所有需要拔除的顽抗节点。他们，以及他们身后每一个战士，为这一天，已经等待了太久，积蓄了太久。

城头，真正的混乱与更替，在火光与喊叫的完美掩饰下，高效而致命地展开了！

在火灾吸引注意力的瞬间，韩铁和他那十余名早已准备就绪的心腹，如同捕猎的豹群般突然发动！他们目标明确，配合默契，以迅雷不及掩耳之势，扑向东门附近少数尚未反应过来的、已知的死忠紫阳宗、丹鼎阁修士及其铁杆附庸。

刀光剑影在火光映照下短暂闪烁，闷哼与短促的惨叫声被更大的嘈杂掩盖，反抗在组织性的突袭面前迅速瓦解。

“控制门闸！开城门！”韩铁低吼一声，挥刀砍断粗重门闩上的铁锁，与几名膂力惊人的壮汉合力，推动那扇厚重的包铁城门。

“吱嘎嘎——轰隆！！”

城门洞开的沉重声响，如同惊雷，在寂静的夜晚传出去极远，也传入了城外汹涌而来的洪流耳中。

“城门已开！随我进城！按预定计划，夺取要点，镇压顽抗，保护百姓！”白言冰的命令穿透夜色，清晰地传入每一个冲锋战士的耳中。

随即，白言冰一马当先，身形化作一道青色流光，瞬息间掠过护城河，穿入洞开的城门。沈瑶如影随形，刀光在她手中仿佛拥有了生命，每一次闪烁，都精准地带走一名试图组织抵抗的士兵的性命，为后续大军扫清道路。

数百精锐如决堤之水，汹涌灌入抚远城，随即按照早已演练过无数次的方案，迅速分成数十支精干小队，扑向各处战略要害：城守府、中心武库、传送阵法枢纽、各大官仓、监狱以及通往其他城门的交通要道。

与此同时，城内的内应行动也全面爆发。赵镇岳亲自坐镇，以“救火”、“有大量匪徒趁乱潜入城内”为名，调动了自己绝对掌控的亲兵队，以及被韩铁等人长期渗透、拉拢的绝大部分中下层军官和士卒，向那些驻扎在军营、武库等地的死硬分子发动了内外夹击的突袭。

许多紫阳宗、丹鼎阁的修士在睡梦中被同僚的刀剑指住咽喉，或是在执勤时被身后“战友”突然反水，大多茫然失措，零星抵抗在早有准备的镇压下迅速熄灭。

夺取抚远城的战斗，激烈却短暂。反抗主要集中于几处据点，在问道途精锐与起义守军的联手打击下，未能形成有效抵抗。天色将明未明之际，城中主要区域已基本被控制。

白言冰与沈瑶亲自带队，直扑城守府。在那里，他们见到了早已卸去甲胄、只着一身整洁常服，独自立于庭院中央等待的赵镇岳。他脚下，横放着自己那柄伴随多年的佩剑。

见白言冰、沈瑶到来，赵镇岳抱拳，深深一礼，声音带着一夜未眠的疲惫，却异常清晰：“罪将赵镇岳，约束部下不力，致使城池易手，愿领一切责罚。只求……只求贵军信守前约。”

白言冰快步上前，郑重拱手还礼，语气诚恳而有力：“赵城守深明大义，于紧要关头弃暗投明，不仅使无数将士百姓免遭无谓死伤，更保全此城元气，功莫大焉！我白言冰，谨代表云隐别院叶牧谦先生，欢迎赵城守加入我问道途行列，共襄义举！此前约定诸事，我问道途必一字不易，全力履行！”

沈瑶亦微微颔首，清冷的容颜上罕见地露出一丝温和：“赵城守能以满城生灵为念，沈瑶感佩。眼下城中初定，百废待兴，诸多秩序维系、人心安抚事宜，还需赵城守鼎力相助。”

赵镇岳听闻此言，心中那块高悬已久、重于千钧的巨石，轰然落地。一股混合着如释重负、愧疚与新生希望的热流冲上眼眶，他强行忍住，再次重重抱拳，声音微哽：“镇岳……敢不从命！定竭尽所能，助两位稳定抚远，戴罪立功！”

而韩铁与白言冰两人相遇时，更是相互抱拳一礼。

两年半之前，青鸾秘境之中，救命之恩，谁敢忘记？谁又能忘记？！

当第一缕晨光刺破东方的云层，照耀在抚远城头时，一面崭新的旗帜，在万众瞩目下缓缓升起，取代了那面象征紫阳宗权威的烈焰战旗。

旗帜底色是洁净的素白，中央以遒劲的笔墨书写着十个大字——“仰不愧于天，俯不怍于人”。在这核心信念之下，又分三行列着更为具体、直指行止的律条：

“一切行动听指挥”

“不拿群众一针一线”

“一切缴获要归公”

这，便是叶牧谦根据前世宝贵经验，结合此世实际，为这支新生人民军队确立的“三大纪律”。它简单，直白，却蕴含着超越所有宗门戒律的、与最广大民众血肉相连的深刻力量。

城中百姓在胆战心惊中推开家门，预料中的烧杀抢掠并未发生。街道上虽然有列队巡逻或匆匆往来的士卒，其中许多还是他们熟悉的本地守军，但秩序井然。

更令人惊异的是，街头巷尾的显眼处，一夜之间贴满了墨迹未干的安民告示。告示以浅显易懂的文字宣布：即刻废除守旧联盟在抚远城推行的一切苛捐杂税；打开官仓，以极平价出售粮食物资，对鳏寡孤独及赤贫者进行赈济；城中所有伤患，无论军民，皆可前往指定地点，获得问道途医修的免费救治。

而最让城中数量庞大的散修群体心跳加速、热血上涌的，是告示最后一条：即日起，于城中原演武场设立“讲道坪”，每日由问道途修士公开讲授《养气篇》基础法诀与修行疑难，不问出身，不论来历，皆可前来听讲！

疑虑、恐惧、观望……在随后几个时辰里，迅速被一件件实实在在发生的事情碾碎。当第一个大胆的贫苦老人真的用几枚铜钱从官仓买到足以果腹的灵米，当第一个在昨夜零星冲突中受伤的孩童被随军的医修温柔包扎、喂下药散，当第一个散修怀着朝圣般的心情踏入“讲道坪”，听到那深入浅出、直指大道的讲解时……积蓄已久的民心，如同冰封的河面在春日暖阳下轰然迸裂，化作汹涌澎湃的春潮！

“问道途万岁！”

“叶先生万岁！”

“白将军！沈仙子！”

欢呼声起初在几个街区零星响起，随即如同野火燎原，迅速蔓延至全城。那不是被胁迫的呼喊，而是发自肺腑的、对新生与希望的狂热礼赞！

白言冰、沈瑶等人深知，夺取城池只是第一步，巩固胜利果实更为关键。他们迅速清点战果：缴获的灵石、各类矿石、药材、粮食堆积如山，数量之巨，不仅足以支撑云霞山前线数月之用，更能极大缓解己方各根据地的物资压力。另一面，吴刚的后勤网络很快跟进，抚远城从一个孤立的堡垒，瞬间转变为连接云霞山与广袤敌后农村根据地的战略支点，物资人员流转效率陡增，真正实现了农村与城市的联动，进可攻，退可守。

然而，抚远城易帜的最大意义，远不止于军事和物资层面。它成了一个极具冲击力和示范效应的宣传样板。白言冰命人将战役概况、入城后的安民措施、新旧治理的鲜明对比，通过早已枝繁叶茂的地下网络，多渠道、全方位地迅速散播出去。

消息所至，石破天惊！

那些仍在守旧联盟统治下苦苦支撑、观望犹豫的中小城池守军，那些饱受压迫却不知出路在何方的散修民众，第一次如此清晰地看到：原来反抗并非必死，易帜竟能带来更好的生活；原来那面“仰不愧天，俯不怍人”的旗帜，不仅是一句口号，更能实实在在地庇护一方，带来公平与希望。守旧联盟后方本就因“农村包围城市”而暗流汹涌的人心，被“抚远事件”这把烈炬彻底点燃，彻底沸腾！

更多的“韩铁”在黑暗中更加活跃地串联，更多的“赵镇岳”在内心的挣扎与权衡中，悄然倾向另一边。

抚远一夜，乾坤始定。

这座中型枢纽城池的易手，如同在守旧联盟看似铁板一块的统治版图上，用最猛烈的方式凿开了一个巨大的、鲜血淋漓的缺口。光芒与希望从这个缺口汹涌灌入，而恐惧与动摇则从此处向内疯狂蔓延。

这不再仅仅是一次战术层面的奇袭胜利，而是一次政治与战略上的决定性转折，标志着问道途从长期的战略防御与农村渗透，正式迈入了武装夺取政权、巩固扩大根据地的战略反攻阶段。

内域道争的天平，经此一役，倾斜的速度骤然加剧，再难逆转。

消息通过特殊渠道，以最快速度传回巍峨屹立的云霞山。当叶牧谦展开密信，看到“抚远已定，白、沈、赵等皆安，人心归附，三大纪律张布于城头”等字样时，他久久伫立，遥望东北方向，脸上露出了自战争爆发以来，最为舒展、如释重负的欣然微笑。

那笑容里，有对弟子们成长的欣慰，有对策略成功的印证，更有对那不可阻挡的历史车轮滚滚向前的深切感悟与感慨。

而与此同时，紫阳宗内，接到十万火急、损失详报的陈炎，在短暂的死寂之后，猛然暴起，将他平日最珍爱的一套千年暖玉茶具狠狠掼在地上，摔得粉碎！瓷片与茶水四溅，映照着他那张因极度震怒、挫败以及一丝连他自己都不愿承认的恐慌而彻底扭曲铁青的脸。

“叶牧谦……韩铁……赵镇岳……好，好得很！”他咬牙切齿，声音从喉咙深处挤压出来，带着血腥气，“你们以为，夺下一城，就能翻天覆地吗？！”

他猛地转向噤若寒蝉的传令修士，眼中闪烁着近乎疯狂的凶光，嘶声吼道：

“传我急令！前线围山部队，再给我抽调……不！急令周文彬，让他立刻放下手头一切事务，滚来见我！立刻！马上！”

风暴，将以抚远为中心，向着更广阔、更惨烈的领域，悍然席卷。真正的战略决战，已被这成功策反、武装夺取的第一座城池，推到了历史的门前。